% Times-like text font with real bold and italic shapes (the template's own Times shapes are undefined in this TeX engine)
\usepackage{fontspec}
\setmainfont{texgyretermes}[Extension=.otf,UprightFont=*-regular,BoldFont=*-bold,ItalicFont=*-italic,BoldItalicFont=*-bolditalic]
% Extra packages. wacv.sty already loads times, xcolor, graphicx, amsmath, amssymb, booktabs, natbib,
% caption, subcaption, cleveref, enumitem, lineno (review mode) and url.
\usepackage{multirow}
\usepackage{tikz}
\usetikzlibrary{arrows.meta,positioning,fit,calc}

\newcommand{\dacis}{DACIS\xspace}
\newcommand{\pmp}{PMP\xspace}
\newcommand{\Gs}{\mathcal{G}}
\newcommand{\Vs}{\mathcal{V}}
\newcommand{\Ds}{\mathcal{D}}
\newcommand{\pp}{\,pp\xspace}

% \num{key}: a number written by the experiment code (generated/numbers.tex). A missing key prints ??.
\newcommand{\num}[1]{\ifcsname num@#1\endcsname\csname num@#1\endcsname\else\textbf{??}\fi}
% \numif{key}{text}: print text only if the experiment code wrote the key (used for the optional second-machine replication)
\newcommand{\numif}[2]{\ifcsname num@#1\endcsname #2\fi}
\IfFileExists{generated/numbers.tex}{% generated by pmp_dacis.report.write_numbers; do not edit
\expandafter\gdef\csname num@resisc45:seeds\endcsname{5}
\expandafter\gdef\csname num@resisc45:A2_D_P_30:novel_1\endcsname{62.3}
\expandafter\gdef\csname num@resisc45:A2_D_P_30:novel_1:sd\endcsname{4.6}
\expandafter\gdef\csname num@resisc45:A2_D_P_30:novel_1:min\endcsname{54.5}
\expandafter\gdef\csname num@resisc45:A2_D_P_30:novel_1:max\endcsname{65.8}
\expandafter\gdef\csname num@resisc45:A2_D_P_30:novel_5\endcsname{79.3}
\expandafter\gdef\csname num@resisc45:A2_D_P_30:novel_5:sd\endcsname{4.1}
\expandafter\gdef\csname num@resisc45:A2_D_P_30:novel_5:min\endcsname{72.2}
\expandafter\gdef\csname num@resisc45:A2_D_P_30:novel_5:max\endcsname{82.6}
\expandafter\gdef\csname num@resisc45:A2_D_P_30:novel_10\endcsname{82.1}
\expandafter\gdef\csname num@resisc45:A2_D_P_30:novel_10:sd\endcsname{4.1}
\expandafter\gdef\csname num@resisc45:A2_D_P_30:novel_10:min\endcsname{75.1}
\expandafter\gdef\csname num@resisc45:A2_D_P_30:novel_10:max\endcsname{85.0}
\expandafter\gdef\csname num@resisc45:A2_D_P_30:pct\endcsname{30.0}
\expandafter\gdef\csname num@resisc45:A2_D_P_30:params\endcsname{3.35}
\expandafter\gdef\csname num@resisc45:A2_D_P_30:macs_128\endcsname{0.251}
\expandafter\gdef\csname num@resisc45:A2_D_P_30:macs_224\endcsname{0.770}
\expandafter\gdef\csname num@resisc45:A2_dacis_P_30:novel_1\endcsname{62.5}
\expandafter\gdef\csname num@resisc45:A2_dacis_P_30:novel_1:sd\endcsname{5.2}
\expandafter\gdef\csname num@resisc45:A2_dacis_P_30:novel_1:min\endcsname{53.7}
\expandafter\gdef\csname num@resisc45:A2_dacis_P_30:novel_1:max\endcsname{66.9}
\expandafter\gdef\csname num@resisc45:A2_dacis_P_30:novel_5\endcsname{79.7}
\expandafter\gdef\csname num@resisc45:A2_dacis_P_30:novel_5:sd\endcsname{4.1}
\expandafter\gdef\csname num@resisc45:A2_dacis_P_30:novel_5:min\endcsname{72.8}
\expandafter\gdef\csname num@resisc45:A2_dacis_P_30:novel_5:max\endcsname{82.8}
\expandafter\gdef\csname num@resisc45:A2_dacis_P_30:novel_10\endcsname{82.7}
\expandafter\gdef\csname num@resisc45:A2_dacis_P_30:novel_10:sd\endcsname{3.8}
\expandafter\gdef\csname num@resisc45:A2_dacis_P_30:novel_10:min\endcsname{76.1}
\expandafter\gdef\csname num@resisc45:A2_dacis_P_30:novel_10:max\endcsname{85.4}
\expandafter\gdef\csname num@resisc45:A2_dacis_P_30:pct\endcsname{30.0}
\expandafter\gdef\csname num@resisc45:A2_dacis_P_30:params\endcsname{3.35}
\expandafter\gdef\csname num@resisc45:A2_dacis_P_30:macs_128\endcsname{0.251}
\expandafter\gdef\csname num@resisc45:A2_dacis_P_30:macs_224\endcsname{0.770}
\expandafter\gdef\csname num@resisc45:A2_l1_P_30:novel_1\endcsname{64.5}
\expandafter\gdef\csname num@resisc45:A2_l1_P_30:novel_1:sd\endcsname{5.5}
\expandafter\gdef\csname num@resisc45:A2_l1_P_30:novel_1:min\endcsname{54.8}
\expandafter\gdef\csname num@resisc45:A2_l1_P_30:novel_1:max\endcsname{68.0}
\expandafter\gdef\csname num@resisc45:A2_l1_P_30:novel_5\endcsname{80.5}
\expandafter\gdef\csname num@resisc45:A2_l1_P_30:novel_5:sd\endcsname{4.5}
\expandafter\gdef\csname num@resisc45:A2_l1_P_30:novel_5:min\endcsname{72.7}
\expandafter\gdef\csname num@resisc45:A2_l1_P_30:novel_5:max\endcsname{83.8}
\expandafter\gdef\csname num@resisc45:A2_l1_P_30:novel_10\endcsname{83.3}
\expandafter\gdef\csname num@resisc45:A2_l1_P_30:novel_10:sd\endcsname{4.1}
\expandafter\gdef\csname num@resisc45:A2_l1_P_30:novel_10:min\endcsname{76.1}
\expandafter\gdef\csname num@resisc45:A2_l1_P_30:novel_10:max\endcsname{86.1}
\expandafter\gdef\csname num@resisc45:A2_l1_P_30:pct\endcsname{30.0}
\expandafter\gdef\csname num@resisc45:A2_l1_P_30:params\endcsname{3.35}
\expandafter\gdef\csname num@resisc45:A2_l1_P_30:macs_128\endcsname{0.251}
\expandafter\gdef\csname num@resisc45:A2_l1_P_30:macs_224\endcsname{0.770}
\expandafter\gdef\csname num@resisc45:A2_random_P_30:novel_1\endcsname{63.0}
\expandafter\gdef\csname num@resisc45:A2_random_P_30:novel_1:sd\endcsname{5.4}
\expandafter\gdef\csname num@resisc45:A2_random_P_30:novel_1:min\endcsname{54.1}
\expandafter\gdef\csname num@resisc45:A2_random_P_30:novel_1:max\endcsname{67.7}
\expandafter\gdef\csname num@resisc45:A2_random_P_30:novel_5\endcsname{79.7}
\expandafter\gdef\csname num@resisc45:A2_random_P_30:novel_5:sd\endcsname{4.9}
\expandafter\gdef\csname num@resisc45:A2_random_P_30:novel_5:min\endcsname{71.5}
\expandafter\gdef\csname num@resisc45:A2_random_P_30:novel_5:max\endcsname{83.7}
\expandafter\gdef\csname num@resisc45:A2_random_P_30:novel_10\endcsname{82.6}
\expandafter\gdef\csname num@resisc45:A2_random_P_30:novel_10:sd\endcsname{4.5}
\expandafter\gdef\csname num@resisc45:A2_random_P_30:novel_10:min\endcsname{74.9}
\expandafter\gdef\csname num@resisc45:A2_random_P_30:novel_10:max\endcsname{86.0}
\expandafter\gdef\csname num@resisc45:A2_random_P_30:pct\endcsname{30.0}
\expandafter\gdef\csname num@resisc45:A2_random_P_30:params\endcsname{3.35}
\expandafter\gdef\csname num@resisc45:A2_random_P_30:macs_128\endcsname{0.251}
\expandafter\gdef\csname num@resisc45:A2_random_P_30:macs_224\endcsname{0.770}
\expandafter\gdef\csname num@resisc45:A2_taylor_P_30:novel_1\endcsname{64.7}
\expandafter\gdef\csname num@resisc45:A2_taylor_P_30:novel_1:sd\endcsname{5.4}
\expandafter\gdef\csname num@resisc45:A2_taylor_P_30:novel_1:min\endcsname{55.3}
\expandafter\gdef\csname num@resisc45:A2_taylor_P_30:novel_1:max\endcsname{68.8}
\expandafter\gdef\csname num@resisc45:A2_taylor_P_30:novel_5\endcsname{82.2}
\expandafter\gdef\csname num@resisc45:A2_taylor_P_30:novel_5:sd\endcsname{3.9}
\expandafter\gdef\csname num@resisc45:A2_taylor_P_30:novel_5:min\endcsname{75.4}
\expandafter\gdef\csname num@resisc45:A2_taylor_P_30:novel_5:max\endcsname{84.8}
\expandafter\gdef\csname num@resisc45:A2_taylor_P_30:novel_10\endcsname{84.9}
\expandafter\gdef\csname num@resisc45:A2_taylor_P_30:novel_10:sd\endcsname{3.7}
\expandafter\gdef\csname num@resisc45:A2_taylor_P_30:novel_10:min\endcsname{78.5}
\expandafter\gdef\csname num@resisc45:A2_taylor_P_30:novel_10:max\endcsname{87.1}
\expandafter\gdef\csname num@resisc45:A2_taylor_P_30:pct\endcsname{30.0}
\expandafter\gdef\csname num@resisc45:A2_taylor_P_30:params\endcsname{3.35}
\expandafter\gdef\csname num@resisc45:A2_taylor_P_30:macs_128\endcsname{0.251}
\expandafter\gdef\csname num@resisc45:A2_taylor_P_30:macs_224\endcsname{0.770}
\expandafter\gdef\csname num@resisc45:A_D_PMP_30:novel_1\endcsname{62.3}
\expandafter\gdef\csname num@resisc45:A_D_PMP_30:novel_1:sd\endcsname{4.9}
\expandafter\gdef\csname num@resisc45:A_D_PMP_30:novel_1:min\endcsname{54.0}
\expandafter\gdef\csname num@resisc45:A_D_PMP_30:novel_1:max\endcsname{66.6}
\expandafter\gdef\csname num@resisc45:A_D_PMP_30:novel_5\endcsname{79.5}
\expandafter\gdef\csname num@resisc45:A_D_PMP_30:novel_5:sd\endcsname{3.9}
\expandafter\gdef\csname num@resisc45:A_D_PMP_30:novel_5:min\endcsname{72.9}
\expandafter\gdef\csname num@resisc45:A_D_PMP_30:novel_5:max\endcsname{82.9}
\expandafter\gdef\csname num@resisc45:A_D_PMP_30:novel_10\endcsname{82.4}
\expandafter\gdef\csname num@resisc45:A_D_PMP_30:novel_10:sd\endcsname{3.8}
\expandafter\gdef\csname num@resisc45:A_D_PMP_30:novel_10:min\endcsname{76.0}
\expandafter\gdef\csname num@resisc45:A_D_PMP_30:novel_10:max\endcsname{85.5}
\expandafter\gdef\csname num@resisc45:A_D_PMP_30:pct\endcsname{30.0}
\expandafter\gdef\csname num@resisc45:A_D_PMP_30:params\endcsname{3.35}
\expandafter\gdef\csname num@resisc45:A_D_PMP_30:macs_128\endcsname{0.251}
\expandafter\gdef\csname num@resisc45:A_D_PMP_30:macs_224\endcsname{0.770}
\expandafter\gdef\csname num@resisc45:A_G_PMP_30:novel_1\endcsname{63.2}
\expandafter\gdef\csname num@resisc45:A_G_PMP_30:novel_1:sd\endcsname{5.6}
\expandafter\gdef\csname num@resisc45:A_G_PMP_30:novel_1:min\endcsname{53.2}
\expandafter\gdef\csname num@resisc45:A_G_PMP_30:novel_1:max\endcsname{66.2}
\expandafter\gdef\csname num@resisc45:A_G_PMP_30:novel_5\endcsname{81.0}
\expandafter\gdef\csname num@resisc45:A_G_PMP_30:novel_5:sd\endcsname{4.5}
\expandafter\gdef\csname num@resisc45:A_G_PMP_30:novel_5:min\endcsname{73.0}
\expandafter\gdef\csname num@resisc45:A_G_PMP_30:novel_5:max\endcsname{83.6}
\expandafter\gdef\csname num@resisc45:A_G_PMP_30:novel_10\endcsname{83.9}
\expandafter\gdef\csname num@resisc45:A_G_PMP_30:novel_10:sd\endcsname{4.3}
\expandafter\gdef\csname num@resisc45:A_G_PMP_30:novel_10:min\endcsname{76.4}
\expandafter\gdef\csname num@resisc45:A_G_PMP_30:novel_10:max\endcsname{86.3}
\expandafter\gdef\csname num@resisc45:A_G_PMP_30:pct\endcsname{30.0}
\expandafter\gdef\csname num@resisc45:A_G_PMP_30:params\endcsname{3.35}
\expandafter\gdef\csname num@resisc45:A_G_PMP_30:macs_128\endcsname{0.251}
\expandafter\gdef\csname num@resisc45:A_G_PMP_30:macs_224\endcsname{0.770}
\expandafter\gdef\csname num@resisc45:A_V_PMP_30:novel_1\endcsname{63.5}
\expandafter\gdef\csname num@resisc45:A_V_PMP_30:novel_1:sd\endcsname{5.1}
\expandafter\gdef\csname num@resisc45:A_V_PMP_30:novel_1:min\endcsname{54.8}
\expandafter\gdef\csname num@resisc45:A_V_PMP_30:novel_1:max\endcsname{67.8}
\expandafter\gdef\csname num@resisc45:A_V_PMP_30:novel_5\endcsname{81.0}
\expandafter\gdef\csname num@resisc45:A_V_PMP_30:novel_5:sd\endcsname{4.0}
\expandafter\gdef\csname num@resisc45:A_V_PMP_30:novel_5:min\endcsname{74.3}
\expandafter\gdef\csname num@resisc45:A_V_PMP_30:novel_5:max\endcsname{84.5}
\expandafter\gdef\csname num@resisc45:A_V_PMP_30:novel_10\endcsname{83.9}
\expandafter\gdef\csname num@resisc45:A_V_PMP_30:novel_10:sd\endcsname{3.6}
\expandafter\gdef\csname num@resisc45:A_V_PMP_30:novel_10:min\endcsname{77.7}
\expandafter\gdef\csname num@resisc45:A_V_PMP_30:novel_10:max\endcsname{86.8}
\expandafter\gdef\csname num@resisc45:A_V_PMP_30:pct\endcsname{30.0}
\expandafter\gdef\csname num@resisc45:A_V_PMP_30:params\endcsname{3.35}
\expandafter\gdef\csname num@resisc45:A_V_PMP_30:macs_128\endcsname{0.251}
\expandafter\gdef\csname num@resisc45:A_V_PMP_30:macs_224\endcsname{0.770}
\expandafter\gdef\csname num@resisc45:A_bn_gamma_PMP_30:novel_1\endcsname{63.1}
\expandafter\gdef\csname num@resisc45:A_bn_gamma_PMP_30:novel_1:sd\endcsname{5.0}
\expandafter\gdef\csname num@resisc45:A_bn_gamma_PMP_30:novel_1:min\endcsname{54.2}
\expandafter\gdef\csname num@resisc45:A_bn_gamma_PMP_30:novel_1:max\endcsname{66.7}
\expandafter\gdef\csname num@resisc45:A_bn_gamma_PMP_30:novel_5\endcsname{79.7}
\expandafter\gdef\csname num@resisc45:A_bn_gamma_PMP_30:novel_5:sd\endcsname{4.2}
\expandafter\gdef\csname num@resisc45:A_bn_gamma_PMP_30:novel_5:min\endcsname{72.3}
\expandafter\gdef\csname num@resisc45:A_bn_gamma_PMP_30:novel_5:max\endcsname{83.1}
\expandafter\gdef\csname num@resisc45:A_bn_gamma_PMP_30:novel_10\endcsname{82.5}
\expandafter\gdef\csname num@resisc45:A_bn_gamma_PMP_30:novel_10:sd\endcsname{4.0}
\expandafter\gdef\csname num@resisc45:A_bn_gamma_PMP_30:novel_10:min\endcsname{75.5}
\expandafter\gdef\csname num@resisc45:A_bn_gamma_PMP_30:novel_10:max\endcsname{85.5}
\expandafter\gdef\csname num@resisc45:A_bn_gamma_PMP_30:pct\endcsname{30.0}
\expandafter\gdef\csname num@resisc45:A_bn_gamma_PMP_30:params\endcsname{3.35}
\expandafter\gdef\csname num@resisc45:A_bn_gamma_PMP_30:macs_128\endcsname{0.251}
\expandafter\gdef\csname num@resisc45:A_bn_gamma_PMP_30:macs_224\endcsname{0.770}
\expandafter\gdef\csname num@resisc45:A_dacis_PMP_30:novel_1\endcsname{62.9}
\expandafter\gdef\csname num@resisc45:A_dacis_PMP_30:novel_1:sd\endcsname{5.1}
\expandafter\gdef\csname num@resisc45:A_dacis_PMP_30:novel_1:min\endcsname{54.0}
\expandafter\gdef\csname num@resisc45:A_dacis_PMP_30:novel_1:max\endcsname{66.5}
\expandafter\gdef\csname num@resisc45:A_dacis_PMP_30:novel_5\endcsname{80.5}
\expandafter\gdef\csname num@resisc45:A_dacis_PMP_30:novel_5:sd\endcsname{4.2}
\expandafter\gdef\csname num@resisc45:A_dacis_PMP_30:novel_5:min\endcsname{73.3}
\expandafter\gdef\csname num@resisc45:A_dacis_PMP_30:novel_5:max\endcsname{83.3}
\expandafter\gdef\csname num@resisc45:A_dacis_PMP_30:novel_10\endcsname{83.4}
\expandafter\gdef\csname num@resisc45:A_dacis_PMP_30:novel_10:sd\endcsname{3.8}
\expandafter\gdef\csname num@resisc45:A_dacis_PMP_30:novel_10:min\endcsname{76.9}
\expandafter\gdef\csname num@resisc45:A_dacis_PMP_30:novel_10:max\endcsname{85.9}
\expandafter\gdef\csname num@resisc45:A_dacis_PMP_30:pct\endcsname{30.0}
\expandafter\gdef\csname num@resisc45:A_dacis_PMP_30:params\endcsname{3.35}
\expandafter\gdef\csname num@resisc45:A_dacis_PMP_30:macs_128\endcsname{0.251}
\expandafter\gdef\csname num@resisc45:A_dacis_PMP_30:macs_224\endcsname{0.770}
\expandafter\gdef\csname num@resisc45:A_dacis_PMP_30:cpu_1thread_b1_ms\endcsname{6.5}
\expandafter\gdef\csname num@resisc45:A_dacis_PMP_30:cpu_all_threads_b1_ms\endcsname{3.3}
\expandafter\gdef\csname num@resisc45:A_dacis_PMP_30:gpu_b1_ms\endcsname{1.0}
\expandafter\gdef\csname num@resisc45:A_dacis_PMP_30:gpu_ips\endcsname{40567}
\expandafter\gdef\csname num@resisc45:A_l1_PMP_30:novel_1\endcsname{65.1}
\expandafter\gdef\csname num@resisc45:A_l1_PMP_30:novel_1:sd\endcsname{5.0}
\expandafter\gdef\csname num@resisc45:A_l1_PMP_30:novel_1:min\endcsname{56.4}
\expandafter\gdef\csname num@resisc45:A_l1_PMP_30:novel_1:max\endcsname{68.6}
\expandafter\gdef\csname num@resisc45:A_l1_PMP_30:novel_5\endcsname{81.0}
\expandafter\gdef\csname num@resisc45:A_l1_PMP_30:novel_5:sd\endcsname{4.2}
\expandafter\gdef\csname num@resisc45:A_l1_PMP_30:novel_5:min\endcsname{73.9}
\expandafter\gdef\csname num@resisc45:A_l1_PMP_30:novel_5:max\endcsname{84.2}
\expandafter\gdef\csname num@resisc45:A_l1_PMP_30:novel_10\endcsname{83.8}
\expandafter\gdef\csname num@resisc45:A_l1_PMP_30:novel_10:sd\endcsname{3.9}
\expandafter\gdef\csname num@resisc45:A_l1_PMP_30:novel_10:min\endcsname{77.1}
\expandafter\gdef\csname num@resisc45:A_l1_PMP_30:novel_10:max\endcsname{86.7}
\expandafter\gdef\csname num@resisc45:A_l1_PMP_30:pct\endcsname{30.0}
\expandafter\gdef\csname num@resisc45:A_l1_PMP_30:params\endcsname{3.35}
\expandafter\gdef\csname num@resisc45:A_l1_PMP_30:macs_128\endcsname{0.251}
\expandafter\gdef\csname num@resisc45:A_l1_PMP_30:macs_224\endcsname{0.770}
\expandafter\gdef\csname num@resisc45:A_l1_PMP_30:cpu_1thread_b1_ms\endcsname{6.6}
\expandafter\gdef\csname num@resisc45:A_l1_PMP_30:cpu_all_threads_b1_ms\endcsname{3.2}
\expandafter\gdef\csname num@resisc45:A_l1_PMP_30:gpu_b1_ms\endcsname{1.0}
\expandafter\gdef\csname num@resisc45:A_l1_PMP_30:gpu_ips\endcsname{40431}
\expandafter\gdef\csname num@resisc45:A_random_PMP_30:novel_1\endcsname{64.4}
\expandafter\gdef\csname num@resisc45:A_random_PMP_30:novel_1:sd\endcsname{5.1}
\expandafter\gdef\csname num@resisc45:A_random_PMP_30:novel_1:min\endcsname{56.0}
\expandafter\gdef\csname num@resisc45:A_random_PMP_30:novel_1:max\endcsname{69.2}
\expandafter\gdef\csname num@resisc45:A_random_PMP_30:novel_5\endcsname{81.0}
\expandafter\gdef\csname num@resisc45:A_random_PMP_30:novel_5:sd\endcsname{4.2}
\expandafter\gdef\csname num@resisc45:A_random_PMP_30:novel_5:min\endcsname{73.8}
\expandafter\gdef\csname num@resisc45:A_random_PMP_30:novel_5:max\endcsname{84.5}
\expandafter\gdef\csname num@resisc45:A_random_PMP_30:novel_10\endcsname{83.9}
\expandafter\gdef\csname num@resisc45:A_random_PMP_30:novel_10:sd\endcsname{4.0}
\expandafter\gdef\csname num@resisc45:A_random_PMP_30:novel_10:min\endcsname{76.9}
\expandafter\gdef\csname num@resisc45:A_random_PMP_30:novel_10:max\endcsname{86.7}
\expandafter\gdef\csname num@resisc45:A_random_PMP_30:pct\endcsname{30.0}
\expandafter\gdef\csname num@resisc45:A_random_PMP_30:params\endcsname{3.35}
\expandafter\gdef\csname num@resisc45:A_random_PMP_30:macs_128\endcsname{0.251}
\expandafter\gdef\csname num@resisc45:A_random_PMP_30:macs_224\endcsname{0.770}
\expandafter\gdef\csname num@resisc45:A_taylor_PMP_30:novel_1\endcsname{64.6}
\expandafter\gdef\csname num@resisc45:A_taylor_PMP_30:novel_1:sd\endcsname{4.8}
\expandafter\gdef\csname num@resisc45:A_taylor_PMP_30:novel_1:min\endcsname{56.2}
\expandafter\gdef\csname num@resisc45:A_taylor_PMP_30:novel_1:max\endcsname{68.0}
\expandafter\gdef\csname num@resisc45:A_taylor_PMP_30:novel_5\endcsname{82.2}
\expandafter\gdef\csname num@resisc45:A_taylor_PMP_30:novel_5:sd\endcsname{3.9}
\expandafter\gdef\csname num@resisc45:A_taylor_PMP_30:novel_5:min\endcsname{75.5}
\expandafter\gdef\csname num@resisc45:A_taylor_PMP_30:novel_5:max\endcsname{84.7}
\expandafter\gdef\csname num@resisc45:A_taylor_PMP_30:novel_10\endcsname{84.9}
\expandafter\gdef\csname num@resisc45:A_taylor_PMP_30:novel_10:sd\endcsname{3.6}
\expandafter\gdef\csname num@resisc45:A_taylor_PMP_30:novel_10:min\endcsname{78.5}
\expandafter\gdef\csname num@resisc45:A_taylor_PMP_30:novel_10:max\endcsname{87.2}
\expandafter\gdef\csname num@resisc45:A_taylor_PMP_30:pct\endcsname{30.0}
\expandafter\gdef\csname num@resisc45:A_taylor_PMP_30:params\endcsname{3.35}
\expandafter\gdef\csname num@resisc45:A_taylor_PMP_30:macs_128\endcsname{0.251}
\expandafter\gdef\csname num@resisc45:A_taylor_PMP_30:macs_224\endcsname{0.770}
\expandafter\gdef\csname num@resisc45:B_dacis_MP_30:novel_1\endcsname{62.9}
\expandafter\gdef\csname num@resisc45:B_dacis_MP_30:novel_1:sd\endcsname{5.5}
\expandafter\gdef\csname num@resisc45:B_dacis_MP_30:novel_1:min\endcsname{53.5}
\expandafter\gdef\csname num@resisc45:B_dacis_MP_30:novel_1:max\endcsname{66.5}
\expandafter\gdef\csname num@resisc45:B_dacis_MP_30:novel_5\endcsname{80.5}
\expandafter\gdef\csname num@resisc45:B_dacis_MP_30:novel_5:sd\endcsname{4.5}
\expandafter\gdef\csname num@resisc45:B_dacis_MP_30:novel_5:min\endcsname{73.0}
\expandafter\gdef\csname num@resisc45:B_dacis_MP_30:novel_5:max\endcsname{83.5}
\expandafter\gdef\csname num@resisc45:B_dacis_MP_30:novel_10\endcsname{83.3}
\expandafter\gdef\csname num@resisc45:B_dacis_MP_30:novel_10:sd\endcsname{4.2}
\expandafter\gdef\csname num@resisc45:B_dacis_MP_30:novel_10:min\endcsname{76.3}
\expandafter\gdef\csname num@resisc45:B_dacis_MP_30:novel_10:max\endcsname{86.0}
\expandafter\gdef\csname num@resisc45:B_dacis_MP_30:pct\endcsname{30.0}
\expandafter\gdef\csname num@resisc45:B_dacis_MP_30:params\endcsname{3.35}
\expandafter\gdef\csname num@resisc45:B_dacis_MP_30:macs_128\endcsname{0.251}
\expandafter\gdef\csname num@resisc45:B_dacis_MP_30:macs_224\endcsname{0.770}
\expandafter\gdef\csname num@resisc45:B_dacis_PGP_30:novel_1\endcsname{63.4}
\expandafter\gdef\csname num@resisc45:B_dacis_PGP_30:novel_1:sd\endcsname{5.4}
\expandafter\gdef\csname num@resisc45:B_dacis_PGP_30:novel_1:min\endcsname{54.1}
\expandafter\gdef\csname num@resisc45:B_dacis_PGP_30:novel_1:max\endcsname{67.0}
\expandafter\gdef\csname num@resisc45:B_dacis_PGP_30:novel_5\endcsname{80.7}
\expandafter\gdef\csname num@resisc45:B_dacis_PGP_30:novel_5:sd\endcsname{4.3}
\expandafter\gdef\csname num@resisc45:B_dacis_PGP_30:novel_5:min\endcsname{73.3}
\expandafter\gdef\csname num@resisc45:B_dacis_PGP_30:novel_5:max\endcsname{83.8}
\expandafter\gdef\csname num@resisc45:B_dacis_PGP_30:novel_10\endcsname{83.6}
\expandafter\gdef\csname num@resisc45:B_dacis_PGP_30:novel_10:sd\endcsname{3.9}
\expandafter\gdef\csname num@resisc45:B_dacis_PGP_30:novel_10:min\endcsname{76.7}
\expandafter\gdef\csname num@resisc45:B_dacis_PGP_30:novel_10:max\endcsname{86.3}
\expandafter\gdef\csname num@resisc45:B_dacis_PGP_30:pct\endcsname{30.0}
\expandafter\gdef\csname num@resisc45:B_dacis_PGP_30:params\endcsname{3.35}
\expandafter\gdef\csname num@resisc45:B_dacis_PGP_30:macs_128\endcsname{0.251}
\expandafter\gdef\csname num@resisc45:B_dacis_PGP_30:macs_224\endcsname{0.770}
\expandafter\gdef\csname num@resisc45:B_dacis_PM_30:novel_1\endcsname{63.7}
\expandafter\gdef\csname num@resisc45:B_dacis_PM_30:novel_1:sd\endcsname{5.4}
\expandafter\gdef\csname num@resisc45:B_dacis_PM_30:novel_1:min\endcsname{54.9}
\expandafter\gdef\csname num@resisc45:B_dacis_PM_30:novel_1:max\endcsname{68.1}
\expandafter\gdef\csname num@resisc45:B_dacis_PM_30:novel_5\endcsname{79.7}
\expandafter\gdef\csname num@resisc45:B_dacis_PM_30:novel_5:sd\endcsname{4.6}
\expandafter\gdef\csname num@resisc45:B_dacis_PM_30:novel_5:min\endcsname{72.1}
\expandafter\gdef\csname num@resisc45:B_dacis_PM_30:novel_5:max\endcsname{83.0}
\expandafter\gdef\csname num@resisc45:B_dacis_PM_30:novel_10\endcsname{82.5}
\expandafter\gdef\csname num@resisc45:B_dacis_PM_30:novel_10:sd\endcsname{4.2}
\expandafter\gdef\csname num@resisc45:B_dacis_PM_30:novel_10:min\endcsname{75.4}
\expandafter\gdef\csname num@resisc45:B_dacis_PM_30:novel_10:max\endcsname{85.4}
\expandafter\gdef\csname num@resisc45:B_dacis_PM_30:pct\endcsname{30.0}
\expandafter\gdef\csname num@resisc45:B_dacis_PM_30:params\endcsname{3.35}
\expandafter\gdef\csname num@resisc45:B_dacis_PM_30:macs_128\endcsname{0.251}
\expandafter\gdef\csname num@resisc45:B_dacis_PM_30:macs_224\endcsname{0.770}
\expandafter\gdef\csname num@resisc45:C_dacis_PMP_10:novel_1\endcsname{54.9}
\expandafter\gdef\csname num@resisc45:C_dacis_PMP_10:novel_1:sd\endcsname{5.3}
\expandafter\gdef\csname num@resisc45:C_dacis_PMP_10:novel_1:min\endcsname{45.6}
\expandafter\gdef\csname num@resisc45:C_dacis_PMP_10:novel_1:max\endcsname{58.3}
\expandafter\gdef\csname num@resisc45:C_dacis_PMP_10:novel_5\endcsname{71.9}
\expandafter\gdef\csname num@resisc45:C_dacis_PMP_10:novel_5:sd\endcsname{4.8}
\expandafter\gdef\csname num@resisc45:C_dacis_PMP_10:novel_5:min\endcsname{63.6}
\expandafter\gdef\csname num@resisc45:C_dacis_PMP_10:novel_5:max\endcsname{75.5}
\expandafter\gdef\csname num@resisc45:C_dacis_PMP_10:novel_10\endcsname{75.1}
\expandafter\gdef\csname num@resisc45:C_dacis_PMP_10:novel_10:sd\endcsname{4.7}
\expandafter\gdef\csname num@resisc45:C_dacis_PMP_10:novel_10:min\endcsname{67.1}
\expandafter\gdef\csname num@resisc45:C_dacis_PMP_10:novel_10:max\endcsname{78.5}
\expandafter\gdef\csname num@resisc45:C_dacis_PMP_10:pct\endcsname{10.0}
\expandafter\gdef\csname num@resisc45:C_dacis_PMP_10:params\endcsname{1.11}
\expandafter\gdef\csname num@resisc45:C_dacis_PMP_10:macs_128\endcsname{0.105}
\expandafter\gdef\csname num@resisc45:C_dacis_PMP_10:macs_224\endcsname{0.322}
\expandafter\gdef\csname num@resisc45:C_dacis_PMP_10:cpu_1thread_b1_ms\endcsname{4.6}
\expandafter\gdef\csname num@resisc45:C_dacis_PMP_10:cpu_all_threads_b1_ms\endcsname{2.5}
\expandafter\gdef\csname num@resisc45:C_dacis_PMP_10:gpu_b1_ms\endcsname{1.0}
\expandafter\gdef\csname num@resisc45:C_dacis_PMP_10:gpu_ips\endcsname{45836}
\expandafter\gdef\csname num@resisc45:C_dacis_PMP_20:novel_1\endcsname{59.9}
\expandafter\gdef\csname num@resisc45:C_dacis_PMP_20:novel_1:sd\endcsname{5.6}
\expandafter\gdef\csname num@resisc45:C_dacis_PMP_20:novel_1:min\endcsname{50.8}
\expandafter\gdef\csname num@resisc45:C_dacis_PMP_20:novel_1:max\endcsname{64.9}
\expandafter\gdef\csname num@resisc45:C_dacis_PMP_20:novel_5\endcsname{77.4}
\expandafter\gdef\csname num@resisc45:C_dacis_PMP_20:novel_5:sd\endcsname{4.9}
\expandafter\gdef\csname num@resisc45:C_dacis_PMP_20:novel_5:min\endcsname{69.6}
\expandafter\gdef\csname num@resisc45:C_dacis_PMP_20:novel_5:max\endcsname{81.7}
\expandafter\gdef\csname num@resisc45:C_dacis_PMP_20:novel_10\endcsname{80.6}
\expandafter\gdef\csname num@resisc45:C_dacis_PMP_20:novel_10:sd\endcsname{4.6}
\expandafter\gdef\csname num@resisc45:C_dacis_PMP_20:novel_10:min\endcsname{73.3}
\expandafter\gdef\csname num@resisc45:C_dacis_PMP_20:novel_10:max\endcsname{84.6}
\expandafter\gdef\csname num@resisc45:C_dacis_PMP_20:pct\endcsname{20.0}
\expandafter\gdef\csname num@resisc45:C_dacis_PMP_20:params\endcsname{2.24}
\expandafter\gdef\csname num@resisc45:C_dacis_PMP_20:macs_128\endcsname{0.178}
\expandafter\gdef\csname num@resisc45:C_dacis_PMP_20:macs_224\endcsname{0.544}
\expandafter\gdef\csname num@resisc45:C_dacis_PMP_20:cpu_1thread_b1_ms\endcsname{5.7}
\expandafter\gdef\csname num@resisc45:C_dacis_PMP_20:cpu_all_threads_b1_ms\endcsname{3.0}
\expandafter\gdef\csname num@resisc45:C_dacis_PMP_20:gpu_b1_ms\endcsname{1.0}
\expandafter\gdef\csname num@resisc45:C_dacis_PMP_20:gpu_ips\endcsname{41542}
\expandafter\gdef\csname num@resisc45:C_dacis_PMP_50:novel_1\endcsname{66.8}
\expandafter\gdef\csname num@resisc45:C_dacis_PMP_50:novel_1:sd\endcsname{4.8}
\expandafter\gdef\csname num@resisc45:C_dacis_PMP_50:novel_1:min\endcsname{58.5}
\expandafter\gdef\csname num@resisc45:C_dacis_PMP_50:novel_1:max\endcsname{70.6}
\expandafter\gdef\csname num@resisc45:C_dacis_PMP_50:novel_5\endcsname{83.8}
\expandafter\gdef\csname num@resisc45:C_dacis_PMP_50:novel_5:sd\endcsname{3.7}
\expandafter\gdef\csname num@resisc45:C_dacis_PMP_50:novel_5:min\endcsname{77.3}
\expandafter\gdef\csname num@resisc45:C_dacis_PMP_50:novel_5:max\endcsname{86.5}
\expandafter\gdef\csname num@resisc45:C_dacis_PMP_50:novel_10\endcsname{86.5}
\expandafter\gdef\csname num@resisc45:C_dacis_PMP_50:novel_10:sd\endcsname{3.4}
\expandafter\gdef\csname num@resisc45:C_dacis_PMP_50:novel_10:min\endcsname{80.5}
\expandafter\gdef\csname num@resisc45:C_dacis_PMP_50:novel_10:max\endcsname{89.0}
\expandafter\gdef\csname num@resisc45:C_dacis_PMP_50:pct\endcsname{50.0}
\expandafter\gdef\csname num@resisc45:C_dacis_PMP_50:params\endcsname{5.59}
\expandafter\gdef\csname num@resisc45:C_dacis_PMP_50:macs_128\endcsname{0.398}
\expandafter\gdef\csname num@resisc45:C_dacis_PMP_50:macs_224\endcsname{1.218}
\expandafter\gdef\csname num@resisc45:C_dacis_PMP_50:cpu_1thread_b1_ms\endcsname{8.6}
\expandafter\gdef\csname num@resisc45:C_dacis_PMP_50:cpu_all_threads_b1_ms\endcsname{3.6}
\expandafter\gdef\csname num@resisc45:C_dacis_PMP_50:gpu_b1_ms\endcsname{1.0}
\expandafter\gdef\csname num@resisc45:C_dacis_PMP_50:gpu_ips\endcsname{38095}
\expandafter\gdef\csname num@resisc45:C_l1_PMP_10:novel_1\endcsname{55.8}
\expandafter\gdef\csname num@resisc45:C_l1_PMP_10:novel_1:sd\endcsname{5.9}
\expandafter\gdef\csname num@resisc45:C_l1_PMP_10:novel_1:min\endcsname{45.8}
\expandafter\gdef\csname num@resisc45:C_l1_PMP_10:novel_1:max\endcsname{59.9}
\expandafter\gdef\csname num@resisc45:C_l1_PMP_10:novel_5\endcsname{70.8}
\expandafter\gdef\csname num@resisc45:C_l1_PMP_10:novel_5:sd\endcsname{5.5}
\expandafter\gdef\csname num@resisc45:C_l1_PMP_10:novel_5:min\endcsname{61.3}
\expandafter\gdef\csname num@resisc45:C_l1_PMP_10:novel_5:max\endcsname{74.8}
\expandafter\gdef\csname num@resisc45:C_l1_PMP_10:novel_10\endcsname{73.9}
\expandafter\gdef\csname num@resisc45:C_l1_PMP_10:novel_10:sd\endcsname{5.2}
\expandafter\gdef\csname num@resisc45:C_l1_PMP_10:novel_10:min\endcsname{64.9}
\expandafter\gdef\csname num@resisc45:C_l1_PMP_10:novel_10:max\endcsname{77.6}
\expandafter\gdef\csname num@resisc45:C_l1_PMP_10:pct\endcsname{10.0}
\expandafter\gdef\csname num@resisc45:C_l1_PMP_10:params\endcsname{1.11}
\expandafter\gdef\csname num@resisc45:C_l1_PMP_10:macs_128\endcsname{0.105}
\expandafter\gdef\csname num@resisc45:C_l1_PMP_10:macs_224\endcsname{0.322}
\expandafter\gdef\csname num@resisc45:C_l1_PMP_10:cpu_1thread_b1_ms\endcsname{4.6}
\expandafter\gdef\csname num@resisc45:C_l1_PMP_10:cpu_all_threads_b1_ms\endcsname{2.4}
\expandafter\gdef\csname num@resisc45:C_l1_PMP_10:gpu_b1_ms\endcsname{1.0}
\expandafter\gdef\csname num@resisc45:C_l1_PMP_10:gpu_ips\endcsname{45865}
\expandafter\gdef\csname num@resisc45:C_l1_PMP_20:novel_1\endcsname{64.1}
\expandafter\gdef\csname num@resisc45:C_l1_PMP_20:novel_1:sd\endcsname{4.9}
\expandafter\gdef\csname num@resisc45:C_l1_PMP_20:novel_1:min\endcsname{55.8}
\expandafter\gdef\csname num@resisc45:C_l1_PMP_20:novel_1:max\endcsname{68.0}
\expandafter\gdef\csname num@resisc45:C_l1_PMP_20:novel_5\endcsname{79.4}
\expandafter\gdef\csname num@resisc45:C_l1_PMP_20:novel_5:sd\endcsname{4.4}
\expandafter\gdef\csname num@resisc45:C_l1_PMP_20:novel_5:min\endcsname{71.8}
\expandafter\gdef\csname num@resisc45:C_l1_PMP_20:novel_5:max\endcsname{82.8}
\expandafter\gdef\csname num@resisc45:C_l1_PMP_20:novel_10\endcsname{82.1}
\expandafter\gdef\csname num@resisc45:C_l1_PMP_20:novel_10:sd\endcsname{4.1}
\expandafter\gdef\csname num@resisc45:C_l1_PMP_20:novel_10:min\endcsname{74.9}
\expandafter\gdef\csname num@resisc45:C_l1_PMP_20:novel_10:max\endcsname{84.8}
\expandafter\gdef\csname num@resisc45:C_l1_PMP_20:pct\endcsname{20.0}
\expandafter\gdef\csname num@resisc45:C_l1_PMP_20:params\endcsname{2.24}
\expandafter\gdef\csname num@resisc45:C_l1_PMP_20:macs_128\endcsname{0.178}
\expandafter\gdef\csname num@resisc45:C_l1_PMP_20:macs_224\endcsname{0.544}
\expandafter\gdef\csname num@resisc45:C_l1_PMP_20:cpu_1thread_b1_ms\endcsname{5.7}
\expandafter\gdef\csname num@resisc45:C_l1_PMP_20:cpu_all_threads_b1_ms\endcsname{3.1}
\expandafter\gdef\csname num@resisc45:C_l1_PMP_20:gpu_b1_ms\endcsname{1.0}
\expandafter\gdef\csname num@resisc45:C_l1_PMP_20:gpu_ips\endcsname{41645}
\expandafter\gdef\csname num@resisc45:C_l1_PMP_50:novel_1\endcsname{68.1}
\expandafter\gdef\csname num@resisc45:C_l1_PMP_50:novel_1:sd\endcsname{4.8}
\expandafter\gdef\csname num@resisc45:C_l1_PMP_50:novel_1:min\endcsname{59.9}
\expandafter\gdef\csname num@resisc45:C_l1_PMP_50:novel_1:max\endcsname{72.2}
\expandafter\gdef\csname num@resisc45:C_l1_PMP_50:novel_5\endcsname{84.1}
\expandafter\gdef\csname num@resisc45:C_l1_PMP_50:novel_5:sd\endcsname{3.7}
\expandafter\gdef\csname num@resisc45:C_l1_PMP_50:novel_5:min\endcsname{77.8}
\expandafter\gdef\csname num@resisc45:C_l1_PMP_50:novel_5:max\endcsname{86.9}
\expandafter\gdef\csname num@resisc45:C_l1_PMP_50:novel_10\endcsname{86.7}
\expandafter\gdef\csname num@resisc45:C_l1_PMP_50:novel_10:sd\endcsname{3.5}
\expandafter\gdef\csname num@resisc45:C_l1_PMP_50:novel_10:min\endcsname{80.8}
\expandafter\gdef\csname num@resisc45:C_l1_PMP_50:novel_10:max\endcsname{89.2}
\expandafter\gdef\csname num@resisc45:C_l1_PMP_50:pct\endcsname{50.0}
\expandafter\gdef\csname num@resisc45:C_l1_PMP_50:params\endcsname{5.59}
\expandafter\gdef\csname num@resisc45:C_l1_PMP_50:macs_128\endcsname{0.398}
\expandafter\gdef\csname num@resisc45:C_l1_PMP_50:macs_224\endcsname{1.218}
\expandafter\gdef\csname num@resisc45:C_l1_PMP_50:cpu_1thread_b1_ms\endcsname{8.6}
\expandafter\gdef\csname num@resisc45:C_l1_PMP_50:cpu_all_threads_b1_ms\endcsname{3.5}
\expandafter\gdef\csname num@resisc45:C_l1_PMP_50:gpu_b1_ms\endcsname{1.0}
\expandafter\gdef\csname num@resisc45:C_l1_PMP_50:gpu_ips\endcsname{38034}
\expandafter\gdef\csname num@resisc45:C_random_PMP_10:novel_1\endcsname{54.0}
\expandafter\gdef\csname num@resisc45:C_random_PMP_10:novel_1:sd\endcsname{6.5}
\expandafter\gdef\csname num@resisc45:C_random_PMP_10:novel_1:min\endcsname{42.9}
\expandafter\gdef\csname num@resisc45:C_random_PMP_10:novel_1:max\endcsname{59.7}
\expandafter\gdef\csname num@resisc45:C_random_PMP_10:novel_5\endcsname{68.9}
\expandafter\gdef\csname num@resisc45:C_random_PMP_10:novel_5:sd\endcsname{7.6}
\expandafter\gdef\csname num@resisc45:C_random_PMP_10:novel_5:min\endcsname{56.1}
\expandafter\gdef\csname num@resisc45:C_random_PMP_10:novel_5:max\endcsname{75.5}
\expandafter\gdef\csname num@resisc45:C_random_PMP_10:novel_10\endcsname{71.8}
\expandafter\gdef\csname num@resisc45:C_random_PMP_10:novel_10:sd\endcsname{7.6}
\expandafter\gdef\csname num@resisc45:C_random_PMP_10:novel_10:min\endcsname{59.0}
\expandafter\gdef\csname num@resisc45:C_random_PMP_10:novel_10:max\endcsname{78.1}
\expandafter\gdef\csname num@resisc45:C_random_PMP_10:pct\endcsname{10.0}
\expandafter\gdef\csname num@resisc45:C_random_PMP_10:params\endcsname{1.11}
\expandafter\gdef\csname num@resisc45:C_random_PMP_10:macs_128\endcsname{0.105}
\expandafter\gdef\csname num@resisc45:C_random_PMP_10:macs_224\endcsname{0.322}
\expandafter\gdef\csname num@resisc45:C_random_PMP_20:novel_1\endcsname{60.5}
\expandafter\gdef\csname num@resisc45:C_random_PMP_20:novel_1:sd\endcsname{5.2}
\expandafter\gdef\csname num@resisc45:C_random_PMP_20:novel_1:min\endcsname{52.1}
\expandafter\gdef\csname num@resisc45:C_random_PMP_20:novel_1:max\endcsname{64.7}
\expandafter\gdef\csname num@resisc45:C_random_PMP_20:novel_5\endcsname{77.1}
\expandafter\gdef\csname num@resisc45:C_random_PMP_20:novel_5:sd\endcsname{4.7}
\expandafter\gdef\csname num@resisc45:C_random_PMP_20:novel_5:min\endcsname{69.3}
\expandafter\gdef\csname num@resisc45:C_random_PMP_20:novel_5:max\endcsname{80.8}
\expandafter\gdef\csname num@resisc45:C_random_PMP_20:novel_10\endcsname{80.1}
\expandafter\gdef\csname num@resisc45:C_random_PMP_20:novel_10:sd\endcsname{4.6}
\expandafter\gdef\csname num@resisc45:C_random_PMP_20:novel_10:min\endcsname{72.4}
\expandafter\gdef\csname num@resisc45:C_random_PMP_20:novel_10:max\endcsname{83.7}
\expandafter\gdef\csname num@resisc45:C_random_PMP_20:pct\endcsname{20.0}
\expandafter\gdef\csname num@resisc45:C_random_PMP_20:params\endcsname{2.24}
\expandafter\gdef\csname num@resisc45:C_random_PMP_20:macs_128\endcsname{0.178}
\expandafter\gdef\csname num@resisc45:C_random_PMP_20:macs_224\endcsname{0.544}
\expandafter\gdef\csname num@resisc45:C_random_PMP_50:novel_1\endcsname{67.2}
\expandafter\gdef\csname num@resisc45:C_random_PMP_50:novel_1:sd\endcsname{4.3}
\expandafter\gdef\csname num@resisc45:C_random_PMP_50:novel_1:min\endcsname{59.8}
\expandafter\gdef\csname num@resisc45:C_random_PMP_50:novel_1:max\endcsname{70.6}
\expandafter\gdef\csname num@resisc45:C_random_PMP_50:novel_5\endcsname{84.1}
\expandafter\gdef\csname num@resisc45:C_random_PMP_50:novel_5:sd\endcsname{3.6}
\expandafter\gdef\csname num@resisc45:C_random_PMP_50:novel_5:min\endcsname{78.0}
\expandafter\gdef\csname num@resisc45:C_random_PMP_50:novel_5:max\endcsname{86.8}
\expandafter\gdef\csname num@resisc45:C_random_PMP_50:novel_10\endcsname{86.8}
\expandafter\gdef\csname num@resisc45:C_random_PMP_50:novel_10:sd\endcsname{3.4}
\expandafter\gdef\csname num@resisc45:C_random_PMP_50:novel_10:min\endcsname{81.0}
\expandafter\gdef\csname num@resisc45:C_random_PMP_50:novel_10:max\endcsname{89.3}
\expandafter\gdef\csname num@resisc45:C_random_PMP_50:pct\endcsname{50.0}
\expandafter\gdef\csname num@resisc45:C_random_PMP_50:params\endcsname{5.59}
\expandafter\gdef\csname num@resisc45:C_random_PMP_50:macs_128\endcsname{0.398}
\expandafter\gdef\csname num@resisc45:C_random_PMP_50:macs_224\endcsname{1.218}
\expandafter\gdef\csname num@resisc45:C_taylor_PMP_10:novel_1\endcsname{58.1}
\expandafter\gdef\csname num@resisc45:C_taylor_PMP_10:novel_1:sd\endcsname{6.2}
\expandafter\gdef\csname num@resisc45:C_taylor_PMP_10:novel_1:min\endcsname{47.7}
\expandafter\gdef\csname num@resisc45:C_taylor_PMP_10:novel_1:max\endcsname{63.0}
\expandafter\gdef\csname num@resisc45:C_taylor_PMP_10:novel_5\endcsname{75.0}
\expandafter\gdef\csname num@resisc45:C_taylor_PMP_10:novel_5:sd\endcsname{5.5}
\expandafter\gdef\csname num@resisc45:C_taylor_PMP_10:novel_5:min\endcsname{65.6}
\expandafter\gdef\csname num@resisc45:C_taylor_PMP_10:novel_5:max\endcsname{79.1}
\expandafter\gdef\csname num@resisc45:C_taylor_PMP_10:novel_10\endcsname{78.1}
\expandafter\gdef\csname num@resisc45:C_taylor_PMP_10:novel_10:sd\endcsname{5.2}
\expandafter\gdef\csname num@resisc45:C_taylor_PMP_10:novel_10:min\endcsname{69.2}
\expandafter\gdef\csname num@resisc45:C_taylor_PMP_10:novel_10:max\endcsname{81.7}
\expandafter\gdef\csname num@resisc45:C_taylor_PMP_10:pct\endcsname{10.0}
\expandafter\gdef\csname num@resisc45:C_taylor_PMP_10:params\endcsname{1.11}
\expandafter\gdef\csname num@resisc45:C_taylor_PMP_10:macs_128\endcsname{0.105}
\expandafter\gdef\csname num@resisc45:C_taylor_PMP_10:macs_224\endcsname{0.322}
\expandafter\gdef\csname num@resisc45:C_taylor_PMP_20:novel_1\endcsname{62.1}
\expandafter\gdef\csname num@resisc45:C_taylor_PMP_20:novel_1:sd\endcsname{5.2}
\expandafter\gdef\csname num@resisc45:C_taylor_PMP_20:novel_1:min\endcsname{53.3}
\expandafter\gdef\csname num@resisc45:C_taylor_PMP_20:novel_1:max\endcsname{67.1}
\expandafter\gdef\csname num@resisc45:C_taylor_PMP_20:novel_5\endcsname{79.6}
\expandafter\gdef\csname num@resisc45:C_taylor_PMP_20:novel_5:sd\endcsname{4.2}
\expandafter\gdef\csname num@resisc45:C_taylor_PMP_20:novel_5:min\endcsname{72.5}
\expandafter\gdef\csname num@resisc45:C_taylor_PMP_20:novel_5:max\endcsname{83.1}
\expandafter\gdef\csname num@resisc45:C_taylor_PMP_20:novel_10\endcsname{82.5}
\expandafter\gdef\csname num@resisc45:C_taylor_PMP_20:novel_10:sd\endcsname{3.8}
\expandafter\gdef\csname num@resisc45:C_taylor_PMP_20:novel_10:min\endcsname{76.1}
\expandafter\gdef\csname num@resisc45:C_taylor_PMP_20:novel_10:max\endcsname{85.3}
\expandafter\gdef\csname num@resisc45:C_taylor_PMP_20:pct\endcsname{20.0}
\expandafter\gdef\csname num@resisc45:C_taylor_PMP_20:params\endcsname{2.24}
\expandafter\gdef\csname num@resisc45:C_taylor_PMP_20:macs_128\endcsname{0.178}
\expandafter\gdef\csname num@resisc45:C_taylor_PMP_20:macs_224\endcsname{0.544}
\expandafter\gdef\csname num@resisc45:C_taylor_PMP_50:novel_1\endcsname{67.6}
\expandafter\gdef\csname num@resisc45:C_taylor_PMP_50:novel_1:sd\endcsname{5.0}
\expandafter\gdef\csname num@resisc45:C_taylor_PMP_50:novel_1:min\endcsname{59.0}
\expandafter\gdef\csname num@resisc45:C_taylor_PMP_50:novel_1:max\endcsname{71.3}
\expandafter\gdef\csname num@resisc45:C_taylor_PMP_50:novel_5\endcsname{84.7}
\expandafter\gdef\csname num@resisc45:C_taylor_PMP_50:novel_5:sd\endcsname{3.7}
\expandafter\gdef\csname num@resisc45:C_taylor_PMP_50:novel_5:min\endcsname{78.2}
\expandafter\gdef\csname num@resisc45:C_taylor_PMP_50:novel_5:max\endcsname{87.2}
\expandafter\gdef\csname num@resisc45:C_taylor_PMP_50:novel_10\endcsname{87.2}
\expandafter\gdef\csname num@resisc45:C_taylor_PMP_50:novel_10:sd\endcsname{3.4}
\expandafter\gdef\csname num@resisc45:C_taylor_PMP_50:novel_10:min\endcsname{81.2}
\expandafter\gdef\csname num@resisc45:C_taylor_PMP_50:novel_10:max\endcsname{89.7}
\expandafter\gdef\csname num@resisc45:C_taylor_PMP_50:pct\endcsname{50.0}
\expandafter\gdef\csname num@resisc45:C_taylor_PMP_50:params\endcsname{5.59}
\expandafter\gdef\csname num@resisc45:C_taylor_PMP_50:macs_128\endcsname{0.398}
\expandafter\gdef\csname num@resisc45:C_taylor_PMP_50:macs_224\endcsname{1.218}
\expandafter\gdef\csname num@resisc45:K_efficientnet_b0:novel_1\endcsname{76.8}
\expandafter\gdef\csname num@resisc45:K_efficientnet_b0:novel_1:sd\endcsname{4.0}
\expandafter\gdef\csname num@resisc45:K_efficientnet_b0:novel_1:min\endcsname{70.1}
\expandafter\gdef\csname num@resisc45:K_efficientnet_b0:novel_1:max\endcsname{79.7}
\expandafter\gdef\csname num@resisc45:K_efficientnet_b0:novel_5\endcsname{90.5}
\expandafter\gdef\csname num@resisc45:K_efficientnet_b0:novel_5:sd\endcsname{2.4}
\expandafter\gdef\csname num@resisc45:K_efficientnet_b0:novel_5:min\endcsname{86.3}
\expandafter\gdef\csname num@resisc45:K_efficientnet_b0:novel_5:max\endcsname{92.3}
\expandafter\gdef\csname num@resisc45:K_efficientnet_b0:novel_10\endcsname{92.4}
\expandafter\gdef\csname num@resisc45:K_efficientnet_b0:novel_10:sd\endcsname{2.1}
\expandafter\gdef\csname num@resisc45:K_efficientnet_b0:novel_10:min\endcsname{88.7}
\expandafter\gdef\csname num@resisc45:K_efficientnet_b0:novel_10:max\endcsname{93.9}
\expandafter\gdef\csname num@resisc45:K_efficientnet_b0:pct\endcsname{35.9}
\expandafter\gdef\csname num@resisc45:K_efficientnet_b0:params\endcsname{4.01}
\expandafter\gdef\csname num@resisc45:K_efficientnet_b0:macs_128\endcsname{0.126}
\expandafter\gdef\csname num@resisc45:K_efficientnet_b0:macs_224\endcsname{0.385}
\expandafter\gdef\csname num@resisc45:K_efficientnet_b0:cpu_1thread_b1_ms\endcsname{6.3}
\expandafter\gdef\csname num@resisc45:K_efficientnet_b0:cpu_all_threads_b1_ms\endcsname{7.1}
\expandafter\gdef\csname num@resisc45:K_efficientnet_b0:gpu_b1_ms\endcsname{2.8}
\expandafter\gdef\csname num@resisc45:K_efficientnet_b0:gpu_ips\endcsname{18928}
\expandafter\gdef\csname num@resisc45:K_mobilenet_v2:novel_1\endcsname{71.3}
\expandafter\gdef\csname num@resisc45:K_mobilenet_v2:novel_1:sd\endcsname{3.4}
\expandafter\gdef\csname num@resisc45:K_mobilenet_v2:novel_1:min\endcsname{65.2}
\expandafter\gdef\csname num@resisc45:K_mobilenet_v2:novel_1:max\endcsname{73.2}
\expandafter\gdef\csname num@resisc45:K_mobilenet_v2:novel_5\endcsname{86.9}
\expandafter\gdef\csname num@resisc45:K_mobilenet_v2:novel_5:sd\endcsname{3.1}
\expandafter\gdef\csname num@resisc45:K_mobilenet_v2:novel_5:min\endcsname{81.3}
\expandafter\gdef\csname num@resisc45:K_mobilenet_v2:novel_5:max\endcsname{89.0}
\expandafter\gdef\csname num@resisc45:K_mobilenet_v2:novel_10\endcsname{88.9}
\expandafter\gdef\csname num@resisc45:K_mobilenet_v2:novel_10:sd\endcsname{3.2}
\expandafter\gdef\csname num@resisc45:K_mobilenet_v2:novel_10:min\endcsname{83.2}
\expandafter\gdef\csname num@resisc45:K_mobilenet_v2:novel_10:max\endcsname{90.9}
\expandafter\gdef\csname num@resisc45:K_mobilenet_v2:pct\endcsname{19.9}
\expandafter\gdef\csname num@resisc45:K_mobilenet_v2:params\endcsname{2.22}
\expandafter\gdef\csname num@resisc45:K_mobilenet_v2:macs_128\endcsname{0.098}
\expandafter\gdef\csname num@resisc45:K_mobilenet_v2:macs_224\endcsname{0.299}
\expandafter\gdef\csname num@resisc45:K_mobilenet_v2:cpu_1thread_b1_ms\endcsname{4.3}
\expandafter\gdef\csname num@resisc45:K_mobilenet_v2:cpu_all_threads_b1_ms\endcsname{5.5}
\expandafter\gdef\csname num@resisc45:K_mobilenet_v2:gpu_b1_ms\endcsname{1.9}
\expandafter\gdef\csname num@resisc45:K_mobilenet_v2:gpu_ips\endcsname{24416}
\expandafter\gdef\csname num@resisc45:K_mobilenet_v3_small:novel_1\endcsname{73.3}
\expandafter\gdef\csname num@resisc45:K_mobilenet_v3_small:novel_1:sd\endcsname{5.0}
\expandafter\gdef\csname num@resisc45:K_mobilenet_v3_small:novel_1:min\endcsname{64.4}
\expandafter\gdef\csname num@resisc45:K_mobilenet_v3_small:novel_1:max\endcsname{76.8}
\expandafter\gdef\csname num@resisc45:K_mobilenet_v3_small:novel_5\endcsname{88.1}
\expandafter\gdef\csname num@resisc45:K_mobilenet_v3_small:novel_5:sd\endcsname{3.3}
\expandafter\gdef\csname num@resisc45:K_mobilenet_v3_small:novel_5:min\endcsname{82.2}
\expandafter\gdef\csname num@resisc45:K_mobilenet_v3_small:novel_5:max\endcsname{90.1}
\expandafter\gdef\csname num@resisc45:K_mobilenet_v3_small:novel_10\endcsname{90.3}
\expandafter\gdef\csname num@resisc45:K_mobilenet_v3_small:novel_10:sd\endcsname{3.0}
\expandafter\gdef\csname num@resisc45:K_mobilenet_v3_small:novel_10:min\endcsname{84.9}
\expandafter\gdef\csname num@resisc45:K_mobilenet_v3_small:novel_10:max\endcsname{92.0}
\expandafter\gdef\csname num@resisc45:K_mobilenet_v3_small:pct\endcsname{8.3}
\expandafter\gdef\csname num@resisc45:K_mobilenet_v3_small:params\endcsname{0.93}
\expandafter\gdef\csname num@resisc45:K_mobilenet_v3_small:macs_128\endcsname{0.018}
\expandafter\gdef\csname num@resisc45:K_mobilenet_v3_small:macs_224\endcsname{0.055}
\expandafter\gdef\csname num@resisc45:K_mobilenet_v3_small:cpu_1thread_b1_ms\endcsname{2.3}
\expandafter\gdef\csname num@resisc45:K_mobilenet_v3_small:cpu_all_threads_b1_ms\endcsname{4.0}
\expandafter\gdef\csname num@resisc45:K_mobilenet_v3_small:gpu_b1_ms\endcsname{1.9}
\expandafter\gdef\csname num@resisc45:K_mobilenet_v3_small:gpu_ips\endcsname{63688}
\expandafter\gdef\csname num@resisc45:dense_ce:novel_1\endcsname{70.8}
\expandafter\gdef\csname num@resisc45:dense_ce:novel_1:sd\endcsname{4.9}
\expandafter\gdef\csname num@resisc45:dense_ce:novel_1:min\endcsname{62.2}
\expandafter\gdef\csname num@resisc45:dense_ce:novel_1:max\endcsname{74.2}
\expandafter\gdef\csname num@resisc45:dense_ce:novel_5\endcsname{86.9}
\expandafter\gdef\csname num@resisc45:dense_ce:novel_5:sd\endcsname{3.4}
\expandafter\gdef\csname num@resisc45:dense_ce:novel_5:min\endcsname{81.0}
\expandafter\gdef\csname num@resisc45:dense_ce:novel_5:max\endcsname{89.2}
\expandafter\gdef\csname num@resisc45:dense_ce:novel_10\endcsname{89.2}
\expandafter\gdef\csname num@resisc45:dense_ce:novel_10:sd\endcsname{3.1}
\expandafter\gdef\csname num@resisc45:dense_ce:novel_10:min\endcsname{83.8}
\expandafter\gdef\csname num@resisc45:dense_ce:novel_10:max\endcsname{91.1}
\expandafter\gdef\csname num@resisc45:dense_ce:pct\endcsname{100.0}
\expandafter\gdef\csname num@resisc45:dense_ce:params\endcsname{11.18}
\expandafter\gdef\csname num@resisc45:dense_ce:macs_128\endcsname{0.592}
\expandafter\gdef\csname num@resisc45:dense_ce:macs_224\endcsname{1.814}
\expandafter\gdef\csname num@resisc45:dense_ce:cpu_1thread_b1_ms\endcsname{10.8}
\expandafter\gdef\csname num@resisc45:dense_ce:cpu_all_threads_b1_ms\endcsname{3.8}
\expandafter\gdef\csname num@resisc45:dense_ce:gpu_b1_ms\endcsname{0.9}
\expandafter\gdef\csname num@resisc45:dense_ce:gpu_ips\endcsname{36520}
\expandafter\gdef\csname num@resisc45:dense_m:novel_1\endcsname{71.2}
\expandafter\gdef\csname num@resisc45:dense_m:novel_1:sd\endcsname{4.6}
\expandafter\gdef\csname num@resisc45:dense_m:novel_1:min\endcsname{63.3}
\expandafter\gdef\csname num@resisc45:dense_m:novel_1:max\endcsname{74.7}
\expandafter\gdef\csname num@resisc45:dense_m:novel_5\endcsname{86.6}
\expandafter\gdef\csname num@resisc45:dense_m:novel_5:sd\endcsname{3.3}
\expandafter\gdef\csname num@resisc45:dense_m:novel_5:min\endcsname{80.8}
\expandafter\gdef\csname num@resisc45:dense_m:novel_5:max\endcsname{88.9}
\expandafter\gdef\csname num@resisc45:dense_m:novel_10\endcsname{88.8}
\expandafter\gdef\csname num@resisc45:dense_m:novel_10:sd\endcsname{3.0}
\expandafter\gdef\csname num@resisc45:dense_m:novel_10:min\endcsname{83.5}
\expandafter\gdef\csname num@resisc45:dense_m:novel_10:max\endcsname{90.8}
\expandafter\gdef\csname num@resisc45:dense_m:pct\endcsname{100.0}
\expandafter\gdef\csname num@resisc45:dense_m:params\endcsname{11.18}
\expandafter\gdef\csname num@resisc45:dense_m:macs_128\endcsname{0.592}
\expandafter\gdef\csname num@resisc45:dense_m:macs_224\endcsname{1.814}
\expandafter\gdef\csname num@resisc45:dense_rmr:novel_1\endcsname{70.4}
\expandafter\gdef\csname num@resisc45:dense_rmr:novel_1:sd\endcsname{4.8}
\expandafter\gdef\csname num@resisc45:dense_rmr:novel_1:min\endcsname{62.3}
\expandafter\gdef\csname num@resisc45:dense_rmr:novel_1:max\endcsname{74.2}
\expandafter\gdef\csname num@resisc45:dense_rmr:novel_5\endcsname{86.7}
\expandafter\gdef\csname num@resisc45:dense_rmr:novel_5:sd\endcsname{3.3}
\expandafter\gdef\csname num@resisc45:dense_rmr:novel_5:min\endcsname{80.9}
\expandafter\gdef\csname num@resisc45:dense_rmr:novel_5:max\endcsname{89.1}
\expandafter\gdef\csname num@resisc45:dense_rmr:novel_10\endcsname{89.0}
\expandafter\gdef\csname num@resisc45:dense_rmr:novel_10:sd\endcsname{3.1}
\expandafter\gdef\csname num@resisc45:dense_rmr:novel_10:min\endcsname{83.6}
\expandafter\gdef\csname num@resisc45:dense_rmr:novel_10:max\endcsname{91.1}
\expandafter\gdef\csname num@resisc45:dense_rmr:pct\endcsname{100.0}
\expandafter\gdef\csname num@resisc45:dense_rmr:params\endcsname{11.18}
\expandafter\gdef\csname num@resisc45:dense_rmr:macs_128\endcsname{0.592}
\expandafter\gdef\csname num@resisc45:dense_rmr:macs_224\endcsname{1.814}
\expandafter\gdef\csname num@eurosat13:seeds\endcsname{5}
\expandafter\gdef\csname num@eurosat13:A2_D_P_30:novel_1\endcsname{45.3}
\expandafter\gdef\csname num@eurosat13:A2_D_P_30:novel_1:sd\endcsname{6.6}
\expandafter\gdef\csname num@eurosat13:A2_D_P_30:novel_1:min\endcsname{36.1}
\expandafter\gdef\csname num@eurosat13:A2_D_P_30:novel_1:max\endcsname{53.4}
\expandafter\gdef\csname num@eurosat13:A2_D_P_30:novel_5\endcsname{55.8}
\expandafter\gdef\csname num@eurosat13:A2_D_P_30:novel_5:sd\endcsname{6.9}
\expandafter\gdef\csname num@eurosat13:A2_D_P_30:novel_5:min\endcsname{45.5}
\expandafter\gdef\csname num@eurosat13:A2_D_P_30:novel_5:max\endcsname{64.3}
\expandafter\gdef\csname num@eurosat13:A2_D_P_30:novel_10\endcsname{58.3}
\expandafter\gdef\csname num@eurosat13:A2_D_P_30:novel_10:sd\endcsname{6.4}
\expandafter\gdef\csname num@eurosat13:A2_D_P_30:novel_10:min\endcsname{48.8}
\expandafter\gdef\csname num@eurosat13:A2_D_P_30:novel_10:max\endcsname{66.1}
\expandafter\gdef\csname num@eurosat13:A2_D_P_30:pct\endcsname{30.0}
\expandafter\gdef\csname num@eurosat13:A2_D_P_30:params\endcsname{3.36}
\expandafter\gdef\csname num@eurosat13:A2_D_P_30:macs_128\endcsname{0.378}
\expandafter\gdef\csname num@eurosat13:A2_D_P_30:macs_224\endcsname{1.157}
\expandafter\gdef\csname num@eurosat13:A2_dacis_P_30:novel_1\endcsname{45.4}
\expandafter\gdef\csname num@eurosat13:A2_dacis_P_30:novel_1:sd\endcsname{6.3}
\expandafter\gdef\csname num@eurosat13:A2_dacis_P_30:novel_1:min\endcsname{36.7}
\expandafter\gdef\csname num@eurosat13:A2_dacis_P_30:novel_1:max\endcsname{53.2}
\expandafter\gdef\csname num@eurosat13:A2_dacis_P_30:novel_5\endcsname{56.3}
\expandafter\gdef\csname num@eurosat13:A2_dacis_P_30:novel_5:sd\endcsname{6.7}
\expandafter\gdef\csname num@eurosat13:A2_dacis_P_30:novel_5:min\endcsname{46.6}
\expandafter\gdef\csname num@eurosat13:A2_dacis_P_30:novel_5:max\endcsname{64.3}
\expandafter\gdef\csname num@eurosat13:A2_dacis_P_30:novel_10\endcsname{58.6}
\expandafter\gdef\csname num@eurosat13:A2_dacis_P_30:novel_10:sd\endcsname{6.5}
\expandafter\gdef\csname num@eurosat13:A2_dacis_P_30:novel_10:min\endcsname{49.4}
\expandafter\gdef\csname num@eurosat13:A2_dacis_P_30:novel_10:max\endcsname{66.0}
\expandafter\gdef\csname num@eurosat13:A2_dacis_P_30:pct\endcsname{30.0}
\expandafter\gdef\csname num@eurosat13:A2_dacis_P_30:params\endcsname{3.36}
\expandafter\gdef\csname num@eurosat13:A2_dacis_P_30:macs_128\endcsname{0.378}
\expandafter\gdef\csname num@eurosat13:A2_dacis_P_30:macs_224\endcsname{1.157}
\expandafter\gdef\csname num@eurosat13:A2_l1_P_30:novel_1\endcsname{45.1}
\expandafter\gdef\csname num@eurosat13:A2_l1_P_30:novel_1:sd\endcsname{7.1}
\expandafter\gdef\csname num@eurosat13:A2_l1_P_30:novel_1:min\endcsname{34.5}
\expandafter\gdef\csname num@eurosat13:A2_l1_P_30:novel_1:max\endcsname{54.3}
\expandafter\gdef\csname num@eurosat13:A2_l1_P_30:novel_5\endcsname{56.2}
\expandafter\gdef\csname num@eurosat13:A2_l1_P_30:novel_5:sd\endcsname{7.9}
\expandafter\gdef\csname num@eurosat13:A2_l1_P_30:novel_5:min\endcsname{43.4}
\expandafter\gdef\csname num@eurosat13:A2_l1_P_30:novel_5:max\endcsname{64.9}
\expandafter\gdef\csname num@eurosat13:A2_l1_P_30:novel_10\endcsname{58.7}
\expandafter\gdef\csname num@eurosat13:A2_l1_P_30:novel_10:sd\endcsname{7.3}
\expandafter\gdef\csname num@eurosat13:A2_l1_P_30:novel_10:min\endcsname{46.9}
\expandafter\gdef\csname num@eurosat13:A2_l1_P_30:novel_10:max\endcsname{66.8}
\expandafter\gdef\csname num@eurosat13:A2_l1_P_30:pct\endcsname{30.0}
\expandafter\gdef\csname num@eurosat13:A2_l1_P_30:params\endcsname{3.36}
\expandafter\gdef\csname num@eurosat13:A2_l1_P_30:macs_128\endcsname{0.378}
\expandafter\gdef\csname num@eurosat13:A2_l1_P_30:macs_224\endcsname{1.157}
\expandafter\gdef\csname num@eurosat13:A2_random_P_30:novel_1\endcsname{44.0}
\expandafter\gdef\csname num@eurosat13:A2_random_P_30:novel_1:sd\endcsname{5.6}
\expandafter\gdef\csname num@eurosat13:A2_random_P_30:novel_1:min\endcsname{36.0}
\expandafter\gdef\csname num@eurosat13:A2_random_P_30:novel_1:max\endcsname{50.1}
\expandafter\gdef\csname num@eurosat13:A2_random_P_30:novel_5\endcsname{55.6}
\expandafter\gdef\csname num@eurosat13:A2_random_P_30:novel_5:sd\endcsname{6.4}
\expandafter\gdef\csname num@eurosat13:A2_random_P_30:novel_5:min\endcsname{46.5}
\expandafter\gdef\csname num@eurosat13:A2_random_P_30:novel_5:max\endcsname{61.9}
\expandafter\gdef\csname num@eurosat13:A2_random_P_30:novel_10\endcsname{58.3}
\expandafter\gdef\csname num@eurosat13:A2_random_P_30:novel_10:sd\endcsname{6.1}
\expandafter\gdef\csname num@eurosat13:A2_random_P_30:novel_10:min\endcsname{49.6}
\expandafter\gdef\csname num@eurosat13:A2_random_P_30:novel_10:max\endcsname{63.8}
\expandafter\gdef\csname num@eurosat13:A2_random_P_30:pct\endcsname{30.0}
\expandafter\gdef\csname num@eurosat13:A2_random_P_30:params\endcsname{3.36}
\expandafter\gdef\csname num@eurosat13:A2_random_P_30:macs_128\endcsname{0.378}
\expandafter\gdef\csname num@eurosat13:A2_random_P_30:macs_224\endcsname{1.157}
\expandafter\gdef\csname num@eurosat13:A2_taylor_P_30:novel_1\endcsname{46.8}
\expandafter\gdef\csname num@eurosat13:A2_taylor_P_30:novel_1:sd\endcsname{6.3}
\expandafter\gdef\csname num@eurosat13:A2_taylor_P_30:novel_1:min\endcsname{36.9}
\expandafter\gdef\csname num@eurosat13:A2_taylor_P_30:novel_1:max\endcsname{54.3}
\expandafter\gdef\csname num@eurosat13:A2_taylor_P_30:novel_5\endcsname{58.1}
\expandafter\gdef\csname num@eurosat13:A2_taylor_P_30:novel_5:sd\endcsname{6.6}
\expandafter\gdef\csname num@eurosat13:A2_taylor_P_30:novel_5:min\endcsname{47.0}
\expandafter\gdef\csname num@eurosat13:A2_taylor_P_30:novel_5:max\endcsname{64.5}
\expandafter\gdef\csname num@eurosat13:A2_taylor_P_30:novel_10\endcsname{60.4}
\expandafter\gdef\csname num@eurosat13:A2_taylor_P_30:novel_10:sd\endcsname{6.1}
\expandafter\gdef\csname num@eurosat13:A2_taylor_P_30:novel_10:min\endcsname{50.2}
\expandafter\gdef\csname num@eurosat13:A2_taylor_P_30:novel_10:max\endcsname{66.0}
\expandafter\gdef\csname num@eurosat13:A2_taylor_P_30:pct\endcsname{30.0}
\expandafter\gdef\csname num@eurosat13:A2_taylor_P_30:params\endcsname{3.36}
\expandafter\gdef\csname num@eurosat13:A2_taylor_P_30:macs_128\endcsname{0.378}
\expandafter\gdef\csname num@eurosat13:A2_taylor_P_30:macs_224\endcsname{1.157}
\expandafter\gdef\csname num@eurosat13:A_D_PMP_30:novel_1\endcsname{45.3}
\expandafter\gdef\csname num@eurosat13:A_D_PMP_30:novel_1:sd\endcsname{7.4}
\expandafter\gdef\csname num@eurosat13:A_D_PMP_30:novel_1:min\endcsname{35.4}
\expandafter\gdef\csname num@eurosat13:A_D_PMP_30:novel_1:max\endcsname{55.4}
\expandafter\gdef\csname num@eurosat13:A_D_PMP_30:novel_5\endcsname{56.0}
\expandafter\gdef\csname num@eurosat13:A_D_PMP_30:novel_5:sd\endcsname{7.5}
\expandafter\gdef\csname num@eurosat13:A_D_PMP_30:novel_5:min\endcsname{45.3}
\expandafter\gdef\csname num@eurosat13:A_D_PMP_30:novel_5:max\endcsname{65.6}
\expandafter\gdef\csname num@eurosat13:A_D_PMP_30:novel_10\endcsname{58.6}
\expandafter\gdef\csname num@eurosat13:A_D_PMP_30:novel_10:sd\endcsname{7.0}
\expandafter\gdef\csname num@eurosat13:A_D_PMP_30:novel_10:min\endcsname{48.7}
\expandafter\gdef\csname num@eurosat13:A_D_PMP_30:novel_10:max\endcsname{67.5}
\expandafter\gdef\csname num@eurosat13:A_D_PMP_30:pct\endcsname{30.0}
\expandafter\gdef\csname num@eurosat13:A_D_PMP_30:params\endcsname{3.36}
\expandafter\gdef\csname num@eurosat13:A_D_PMP_30:macs_128\endcsname{0.378}
\expandafter\gdef\csname num@eurosat13:A_D_PMP_30:macs_224\endcsname{1.157}
\expandafter\gdef\csname num@eurosat13:A_G_PMP_30:novel_1\endcsname{45.2}
\expandafter\gdef\csname num@eurosat13:A_G_PMP_30:novel_1:sd\endcsname{6.6}
\expandafter\gdef\csname num@eurosat13:A_G_PMP_30:novel_1:min\endcsname{36.2}
\expandafter\gdef\csname num@eurosat13:A_G_PMP_30:novel_1:max\endcsname{52.7}
\expandafter\gdef\csname num@eurosat13:A_G_PMP_30:novel_5\endcsname{56.4}
\expandafter\gdef\csname num@eurosat13:A_G_PMP_30:novel_5:sd\endcsname{6.5}
\expandafter\gdef\csname num@eurosat13:A_G_PMP_30:novel_5:min\endcsname{47.2}
\expandafter\gdef\csname num@eurosat13:A_G_PMP_30:novel_5:max\endcsname{64.1}
\expandafter\gdef\csname num@eurosat13:A_G_PMP_30:novel_10\endcsname{58.9}
\expandafter\gdef\csname num@eurosat13:A_G_PMP_30:novel_10:sd\endcsname{6.0}
\expandafter\gdef\csname num@eurosat13:A_G_PMP_30:novel_10:min\endcsname{50.4}
\expandafter\gdef\csname num@eurosat13:A_G_PMP_30:novel_10:max\endcsname{65.7}
\expandafter\gdef\csname num@eurosat13:A_G_PMP_30:pct\endcsname{30.0}
\expandafter\gdef\csname num@eurosat13:A_G_PMP_30:params\endcsname{3.36}
\expandafter\gdef\csname num@eurosat13:A_G_PMP_30:macs_128\endcsname{0.378}
\expandafter\gdef\csname num@eurosat13:A_G_PMP_30:macs_224\endcsname{1.157}
\expandafter\gdef\csname num@eurosat13:A_V_PMP_30:novel_1\endcsname{45.5}
\expandafter\gdef\csname num@eurosat13:A_V_PMP_30:novel_1:sd\endcsname{5.7}
\expandafter\gdef\csname num@eurosat13:A_V_PMP_30:novel_1:min\endcsname{38.9}
\expandafter\gdef\csname num@eurosat13:A_V_PMP_30:novel_1:max\endcsname{53.7}
\expandafter\gdef\csname num@eurosat13:A_V_PMP_30:novel_5\endcsname{56.5}
\expandafter\gdef\csname num@eurosat13:A_V_PMP_30:novel_5:sd\endcsname{5.2}
\expandafter\gdef\csname num@eurosat13:A_V_PMP_30:novel_5:min\endcsname{50.4}
\expandafter\gdef\csname num@eurosat13:A_V_PMP_30:novel_5:max\endcsname{64.0}
\expandafter\gdef\csname num@eurosat13:A_V_PMP_30:novel_10\endcsname{59.0}
\expandafter\gdef\csname num@eurosat13:A_V_PMP_30:novel_10:sd\endcsname{4.8}
\expandafter\gdef\csname num@eurosat13:A_V_PMP_30:novel_10:min\endcsname{53.7}
\expandafter\gdef\csname num@eurosat13:A_V_PMP_30:novel_10:max\endcsname{65.9}
\expandafter\gdef\csname num@eurosat13:A_V_PMP_30:pct\endcsname{30.0}
\expandafter\gdef\csname num@eurosat13:A_V_PMP_30:params\endcsname{3.36}
\expandafter\gdef\csname num@eurosat13:A_V_PMP_30:macs_128\endcsname{0.378}
\expandafter\gdef\csname num@eurosat13:A_V_PMP_30:macs_224\endcsname{1.157}
\expandafter\gdef\csname num@eurosat13:A_bn_gamma_PMP_30:novel_1\endcsname{44.0}
\expandafter\gdef\csname num@eurosat13:A_bn_gamma_PMP_30:novel_1:sd\endcsname{6.0}
\expandafter\gdef\csname num@eurosat13:A_bn_gamma_PMP_30:novel_1:min\endcsname{35.1}
\expandafter\gdef\csname num@eurosat13:A_bn_gamma_PMP_30:novel_1:max\endcsname{50.0}
\expandafter\gdef\csname num@eurosat13:A_bn_gamma_PMP_30:novel_5\endcsname{55.0}
\expandafter\gdef\csname num@eurosat13:A_bn_gamma_PMP_30:novel_5:sd\endcsname{6.1}
\expandafter\gdef\csname num@eurosat13:A_bn_gamma_PMP_30:novel_5:min\endcsname{45.3}
\expandafter\gdef\csname num@eurosat13:A_bn_gamma_PMP_30:novel_5:max\endcsname{60.7}
\expandafter\gdef\csname num@eurosat13:A_bn_gamma_PMP_30:novel_10\endcsname{57.3}
\expandafter\gdef\csname num@eurosat13:A_bn_gamma_PMP_30:novel_10:sd\endcsname{5.6}
\expandafter\gdef\csname num@eurosat13:A_bn_gamma_PMP_30:novel_10:min\endcsname{48.2}
\expandafter\gdef\csname num@eurosat13:A_bn_gamma_PMP_30:novel_10:max\endcsname{62.1}
\expandafter\gdef\csname num@eurosat13:A_bn_gamma_PMP_30:pct\endcsname{30.0}
\expandafter\gdef\csname num@eurosat13:A_bn_gamma_PMP_30:params\endcsname{3.36}
\expandafter\gdef\csname num@eurosat13:A_bn_gamma_PMP_30:macs_128\endcsname{0.378}
\expandafter\gdef\csname num@eurosat13:A_bn_gamma_PMP_30:macs_224\endcsname{1.157}
\expandafter\gdef\csname num@eurosat13:A_dacis_PMP_30:novel_1\endcsname{46.0}
\expandafter\gdef\csname num@eurosat13:A_dacis_PMP_30:novel_1:sd\endcsname{5.4}
\expandafter\gdef\csname num@eurosat13:A_dacis_PMP_30:novel_1:min\endcsname{38.3}
\expandafter\gdef\csname num@eurosat13:A_dacis_PMP_30:novel_1:max\endcsname{51.1}
\expandafter\gdef\csname num@eurosat13:A_dacis_PMP_30:novel_5\endcsname{57.1}
\expandafter\gdef\csname num@eurosat13:A_dacis_PMP_30:novel_5:sd\endcsname{4.9}
\expandafter\gdef\csname num@eurosat13:A_dacis_PMP_30:novel_5:min\endcsname{50.2}
\expandafter\gdef\csname num@eurosat13:A_dacis_PMP_30:novel_5:max\endcsname{61.9}
\expandafter\gdef\csname num@eurosat13:A_dacis_PMP_30:novel_10\endcsname{59.4}
\expandafter\gdef\csname num@eurosat13:A_dacis_PMP_30:novel_10:sd\endcsname{4.1}
\expandafter\gdef\csname num@eurosat13:A_dacis_PMP_30:novel_10:min\endcsname{53.7}
\expandafter\gdef\csname num@eurosat13:A_dacis_PMP_30:novel_10:max\endcsname{63.2}
\expandafter\gdef\csname num@eurosat13:A_dacis_PMP_30:pct\endcsname{30.0}
\expandafter\gdef\csname num@eurosat13:A_dacis_PMP_30:params\endcsname{3.36}
\expandafter\gdef\csname num@eurosat13:A_dacis_PMP_30:macs_128\endcsname{0.378}
\expandafter\gdef\csname num@eurosat13:A_dacis_PMP_30:macs_224\endcsname{1.157}
\expandafter\gdef\csname num@eurosat13:A_dacis_PMP_30:cpu_1thread_b1_ms\endcsname{3.2}
\expandafter\gdef\csname num@eurosat13:A_dacis_PMP_30:cpu_all_threads_b1_ms\endcsname{1.9}
\expandafter\gdef\csname num@eurosat13:A_dacis_PMP_30:gpu_b1_ms\endcsname{1.0}
\expandafter\gdef\csname num@eurosat13:A_dacis_PMP_30:gpu_ips\endcsname{77646}
\expandafter\gdef\csname num@eurosat13:A_l1_PMP_30:novel_1\endcsname{46.1}
\expandafter\gdef\csname num@eurosat13:A_l1_PMP_30:novel_1:sd\endcsname{6.5}
\expandafter\gdef\csname num@eurosat13:A_l1_PMP_30:novel_1:min\endcsname{37.5}
\expandafter\gdef\csname num@eurosat13:A_l1_PMP_30:novel_1:max\endcsname{55.5}
\expandafter\gdef\csname num@eurosat13:A_l1_PMP_30:novel_5\endcsname{57.0}
\expandafter\gdef\csname num@eurosat13:A_l1_PMP_30:novel_5:sd\endcsname{6.0}
\expandafter\gdef\csname num@eurosat13:A_l1_PMP_30:novel_5:min\endcsname{48.6}
\expandafter\gdef\csname num@eurosat13:A_l1_PMP_30:novel_5:max\endcsname{64.9}
\expandafter\gdef\csname num@eurosat13:A_l1_PMP_30:novel_10\endcsname{59.3}
\expandafter\gdef\csname num@eurosat13:A_l1_PMP_30:novel_10:sd\endcsname{5.2}
\expandafter\gdef\csname num@eurosat13:A_l1_PMP_30:novel_10:min\endcsname{52.4}
\expandafter\gdef\csname num@eurosat13:A_l1_PMP_30:novel_10:max\endcsname{66.3}
\expandafter\gdef\csname num@eurosat13:A_l1_PMP_30:pct\endcsname{30.0}
\expandafter\gdef\csname num@eurosat13:A_l1_PMP_30:params\endcsname{3.36}
\expandafter\gdef\csname num@eurosat13:A_l1_PMP_30:macs_128\endcsname{0.378}
\expandafter\gdef\csname num@eurosat13:A_l1_PMP_30:macs_224\endcsname{1.157}
\expandafter\gdef\csname num@eurosat13:A_l1_PMP_30:cpu_1thread_b1_ms\endcsname{3.2}
\expandafter\gdef\csname num@eurosat13:A_l1_PMP_30:cpu_all_threads_b1_ms\endcsname{1.8}
\expandafter\gdef\csname num@eurosat13:A_l1_PMP_30:gpu_b1_ms\endcsname{1.0}
\expandafter\gdef\csname num@eurosat13:A_l1_PMP_30:gpu_ips\endcsname{77599}
\expandafter\gdef\csname num@eurosat13:A_random_PMP_30:novel_1\endcsname{46.8}
\expandafter\gdef\csname num@eurosat13:A_random_PMP_30:novel_1:sd\endcsname{5.9}
\expandafter\gdef\csname num@eurosat13:A_random_PMP_30:novel_1:min\endcsname{38.8}
\expandafter\gdef\csname num@eurosat13:A_random_PMP_30:novel_1:max\endcsname{54.6}
\expandafter\gdef\csname num@eurosat13:A_random_PMP_30:novel_5\endcsname{57.6}
\expandafter\gdef\csname num@eurosat13:A_random_PMP_30:novel_5:sd\endcsname{6.0}
\expandafter\gdef\csname num@eurosat13:A_random_PMP_30:novel_5:min\endcsname{48.7}
\expandafter\gdef\csname num@eurosat13:A_random_PMP_30:novel_5:max\endcsname{65.2}
\expandafter\gdef\csname num@eurosat13:A_random_PMP_30:novel_10\endcsname{59.9}
\expandafter\gdef\csname num@eurosat13:A_random_PMP_30:novel_10:sd\endcsname{5.7}
\expandafter\gdef\csname num@eurosat13:A_random_PMP_30:novel_10:min\endcsname{51.4}
\expandafter\gdef\csname num@eurosat13:A_random_PMP_30:novel_10:max\endcsname{67.2}
\expandafter\gdef\csname num@eurosat13:A_random_PMP_30:pct\endcsname{30.0}
\expandafter\gdef\csname num@eurosat13:A_random_PMP_30:params\endcsname{3.36}
\expandafter\gdef\csname num@eurosat13:A_random_PMP_30:macs_128\endcsname{0.378}
\expandafter\gdef\csname num@eurosat13:A_random_PMP_30:macs_224\endcsname{1.157}
\expandafter\gdef\csname num@eurosat13:A_taylor_PMP_30:novel_1\endcsname{45.9}
\expandafter\gdef\csname num@eurosat13:A_taylor_PMP_30:novel_1:sd\endcsname{7.8}
\expandafter\gdef\csname num@eurosat13:A_taylor_PMP_30:novel_1:min\endcsname{34.8}
\expandafter\gdef\csname num@eurosat13:A_taylor_PMP_30:novel_1:max\endcsname{56.2}
\expandafter\gdef\csname num@eurosat13:A_taylor_PMP_30:novel_5\endcsname{57.0}
\expandafter\gdef\csname num@eurosat13:A_taylor_PMP_30:novel_5:sd\endcsname{8.0}
\expandafter\gdef\csname num@eurosat13:A_taylor_PMP_30:novel_5:min\endcsname{45.0}
\expandafter\gdef\csname num@eurosat13:A_taylor_PMP_30:novel_5:max\endcsname{66.6}
\expandafter\gdef\csname num@eurosat13:A_taylor_PMP_30:novel_10\endcsname{59.4}
\expandafter\gdef\csname num@eurosat13:A_taylor_PMP_30:novel_10:sd\endcsname{7.4}
\expandafter\gdef\csname num@eurosat13:A_taylor_PMP_30:novel_10:min\endcsname{48.6}
\expandafter\gdef\csname num@eurosat13:A_taylor_PMP_30:novel_10:max\endcsname{68.2}
\expandafter\gdef\csname num@eurosat13:A_taylor_PMP_30:pct\endcsname{30.0}
\expandafter\gdef\csname num@eurosat13:A_taylor_PMP_30:params\endcsname{3.36}
\expandafter\gdef\csname num@eurosat13:A_taylor_PMP_30:macs_128\endcsname{0.378}
\expandafter\gdef\csname num@eurosat13:A_taylor_PMP_30:macs_224\endcsname{1.157}
\expandafter\gdef\csname num@eurosat13:B_dacis_MP_30:novel_1\endcsname{45.8}
\expandafter\gdef\csname num@eurosat13:B_dacis_MP_30:novel_1:sd\endcsname{6.9}
\expandafter\gdef\csname num@eurosat13:B_dacis_MP_30:novel_1:min\endcsname{36.3}
\expandafter\gdef\csname num@eurosat13:B_dacis_MP_30:novel_1:max\endcsname{54.9}
\expandafter\gdef\csname num@eurosat13:B_dacis_MP_30:novel_5\endcsname{56.5}
\expandafter\gdef\csname num@eurosat13:B_dacis_MP_30:novel_5:sd\endcsname{6.8}
\expandafter\gdef\csname num@eurosat13:B_dacis_MP_30:novel_5:min\endcsname{46.4}
\expandafter\gdef\csname num@eurosat13:B_dacis_MP_30:novel_5:max\endcsname{64.6}
\expandafter\gdef\csname num@eurosat13:B_dacis_MP_30:novel_10\endcsname{58.8}
\expandafter\gdef\csname num@eurosat13:B_dacis_MP_30:novel_10:sd\endcsname{6.2}
\expandafter\gdef\csname num@eurosat13:B_dacis_MP_30:novel_10:min\endcsname{49.5}
\expandafter\gdef\csname num@eurosat13:B_dacis_MP_30:novel_10:max\endcsname{65.9}
\expandafter\gdef\csname num@eurosat13:B_dacis_MP_30:pct\endcsname{30.0}
\expandafter\gdef\csname num@eurosat13:B_dacis_MP_30:params\endcsname{3.36}
\expandafter\gdef\csname num@eurosat13:B_dacis_MP_30:macs_128\endcsname{0.378}
\expandafter\gdef\csname num@eurosat13:B_dacis_MP_30:macs_224\endcsname{1.157}
\expandafter\gdef\csname num@eurosat13:B_dacis_PGP_30:novel_1\endcsname{45.9}
\expandafter\gdef\csname num@eurosat13:B_dacis_PGP_30:novel_1:sd\endcsname{6.1}
\expandafter\gdef\csname num@eurosat13:B_dacis_PGP_30:novel_1:min\endcsname{36.7}
\expandafter\gdef\csname num@eurosat13:B_dacis_PGP_30:novel_1:max\endcsname{53.5}
\expandafter\gdef\csname num@eurosat13:B_dacis_PGP_30:novel_5\endcsname{56.6}
\expandafter\gdef\csname num@eurosat13:B_dacis_PGP_30:novel_5:sd\endcsname{6.1}
\expandafter\gdef\csname num@eurosat13:B_dacis_PGP_30:novel_5:min\endcsname{47.2}
\expandafter\gdef\csname num@eurosat13:B_dacis_PGP_30:novel_5:max\endcsname{64.0}
\expandafter\gdef\csname num@eurosat13:B_dacis_PGP_30:novel_10\endcsname{58.9}
\expandafter\gdef\csname num@eurosat13:B_dacis_PGP_30:novel_10:sd\endcsname{5.5}
\expandafter\gdef\csname num@eurosat13:B_dacis_PGP_30:novel_10:min\endcsname{50.7}
\expandafter\gdef\csname num@eurosat13:B_dacis_PGP_30:novel_10:max\endcsname{65.7}
\expandafter\gdef\csname num@eurosat13:B_dacis_PGP_30:pct\endcsname{30.0}
\expandafter\gdef\csname num@eurosat13:B_dacis_PGP_30:params\endcsname{3.36}
\expandafter\gdef\csname num@eurosat13:B_dacis_PGP_30:macs_128\endcsname{0.378}
\expandafter\gdef\csname num@eurosat13:B_dacis_PGP_30:macs_224\endcsname{1.157}
\expandafter\gdef\csname num@eurosat13:B_dacis_PM_30:novel_1\endcsname{45.2}
\expandafter\gdef\csname num@eurosat13:B_dacis_PM_30:novel_1:sd\endcsname{7.0}
\expandafter\gdef\csname num@eurosat13:B_dacis_PM_30:novel_1:min\endcsname{37.1}
\expandafter\gdef\csname num@eurosat13:B_dacis_PM_30:novel_1:max\endcsname{54.6}
\expandafter\gdef\csname num@eurosat13:B_dacis_PM_30:novel_5\endcsname{55.8}
\expandafter\gdef\csname num@eurosat13:B_dacis_PM_30:novel_5:sd\endcsname{6.9}
\expandafter\gdef\csname num@eurosat13:B_dacis_PM_30:novel_5:min\endcsname{47.0}
\expandafter\gdef\csname num@eurosat13:B_dacis_PM_30:novel_5:max\endcsname{64.5}
\expandafter\gdef\csname num@eurosat13:B_dacis_PM_30:novel_10\endcsname{58.2}
\expandafter\gdef\csname num@eurosat13:B_dacis_PM_30:novel_10:sd\endcsname{6.4}
\expandafter\gdef\csname num@eurosat13:B_dacis_PM_30:novel_10:min\endcsname{50.3}
\expandafter\gdef\csname num@eurosat13:B_dacis_PM_30:novel_10:max\endcsname{65.9}
\expandafter\gdef\csname num@eurosat13:B_dacis_PM_30:pct\endcsname{30.0}
\expandafter\gdef\csname num@eurosat13:B_dacis_PM_30:params\endcsname{3.36}
\expandafter\gdef\csname num@eurosat13:B_dacis_PM_30:macs_128\endcsname{0.378}
\expandafter\gdef\csname num@eurosat13:B_dacis_PM_30:macs_224\endcsname{1.157}
\expandafter\gdef\csname num@eurosat13:C_dacis_PMP_10:novel_1\endcsname{41.8}
\expandafter\gdef\csname num@eurosat13:C_dacis_PMP_10:novel_1:sd\endcsname{6.7}
\expandafter\gdef\csname num@eurosat13:C_dacis_PMP_10:novel_1:min\endcsname{33.6}
\expandafter\gdef\csname num@eurosat13:C_dacis_PMP_10:novel_1:max\endcsname{52.1}
\expandafter\gdef\csname num@eurosat13:C_dacis_PMP_10:novel_5\endcsname{51.4}
\expandafter\gdef\csname num@eurosat13:C_dacis_PMP_10:novel_5:sd\endcsname{7.5}
\expandafter\gdef\csname num@eurosat13:C_dacis_PMP_10:novel_5:min\endcsname{41.4}
\expandafter\gdef\csname num@eurosat13:C_dacis_PMP_10:novel_5:max\endcsname{62.4}
\expandafter\gdef\csname num@eurosat13:C_dacis_PMP_10:novel_10\endcsname{53.7}
\expandafter\gdef\csname num@eurosat13:C_dacis_PMP_10:novel_10:sd\endcsname{7.1}
\expandafter\gdef\csname num@eurosat13:C_dacis_PMP_10:novel_10:min\endcsname{44.2}
\expandafter\gdef\csname num@eurosat13:C_dacis_PMP_10:novel_10:max\endcsname{64.1}
\expandafter\gdef\csname num@eurosat13:C_dacis_PMP_10:pct\endcsname{10.0}
\expandafter\gdef\csname num@eurosat13:C_dacis_PMP_10:params\endcsname{1.12}
\expandafter\gdef\csname num@eurosat13:C_dacis_PMP_10:macs_128\endcsname{0.231}
\expandafter\gdef\csname num@eurosat13:C_dacis_PMP_10:macs_224\endcsname{0.708}
\expandafter\gdef\csname num@eurosat13:C_dacis_PMP_10:cpu_1thread_b1_ms\endcsname{2.4}
\expandafter\gdef\csname num@eurosat13:C_dacis_PMP_10:cpu_all_threads_b1_ms\endcsname{1.6}
\expandafter\gdef\csname num@eurosat13:C_dacis_PMP_10:gpu_b1_ms\endcsname{0.9}
\expandafter\gdef\csname num@eurosat13:C_dacis_PMP_10:gpu_ips\endcsname{86800}
\expandafter\gdef\csname num@eurosat13:C_dacis_PMP_20:novel_1\endcsname{44.7}
\expandafter\gdef\csname num@eurosat13:C_dacis_PMP_20:novel_1:sd\endcsname{7.4}
\expandafter\gdef\csname num@eurosat13:C_dacis_PMP_20:novel_1:min\endcsname{33.5}
\expandafter\gdef\csname num@eurosat13:C_dacis_PMP_20:novel_1:max\endcsname{53.5}
\expandafter\gdef\csname num@eurosat13:C_dacis_PMP_20:novel_5\endcsname{54.5}
\expandafter\gdef\csname num@eurosat13:C_dacis_PMP_20:novel_5:sd\endcsname{8.1}
\expandafter\gdef\csname num@eurosat13:C_dacis_PMP_20:novel_5:min\endcsname{41.7}
\expandafter\gdef\csname num@eurosat13:C_dacis_PMP_20:novel_5:max\endcsname{63.6}
\expandafter\gdef\csname num@eurosat13:C_dacis_PMP_20:novel_10\endcsname{56.6}
\expandafter\gdef\csname num@eurosat13:C_dacis_PMP_20:novel_10:sd\endcsname{7.6}
\expandafter\gdef\csname num@eurosat13:C_dacis_PMP_20:novel_10:min\endcsname{44.6}
\expandafter\gdef\csname num@eurosat13:C_dacis_PMP_20:novel_10:max\endcsname{64.8}
\expandafter\gdef\csname num@eurosat13:C_dacis_PMP_20:pct\endcsname{20.0}
\expandafter\gdef\csname num@eurosat13:C_dacis_PMP_20:params\endcsname{2.24}
\expandafter\gdef\csname num@eurosat13:C_dacis_PMP_20:macs_128\endcsname{0.304}
\expandafter\gdef\csname num@eurosat13:C_dacis_PMP_20:macs_224\endcsname{0.932}
\expandafter\gdef\csname num@eurosat13:C_dacis_PMP_20:cpu_1thread_b1_ms\endcsname{2.8}
\expandafter\gdef\csname num@eurosat13:C_dacis_PMP_20:cpu_all_threads_b1_ms\endcsname{1.6}
\expandafter\gdef\csname num@eurosat13:C_dacis_PMP_20:gpu_b1_ms\endcsname{1.0}
\expandafter\gdef\csname num@eurosat13:C_dacis_PMP_20:gpu_ips\endcsname{81435}
\expandafter\gdef\csname num@eurosat13:C_dacis_PMP_50:novel_1\endcsname{47.0}
\expandafter\gdef\csname num@eurosat13:C_dacis_PMP_50:novel_1:sd\endcsname{5.7}
\expandafter\gdef\csname num@eurosat13:C_dacis_PMP_50:novel_1:min\endcsname{38.1}
\expandafter\gdef\csname num@eurosat13:C_dacis_PMP_50:novel_1:max\endcsname{52.0}
\expandafter\gdef\csname num@eurosat13:C_dacis_PMP_50:novel_5\endcsname{58.6}
\expandafter\gdef\csname num@eurosat13:C_dacis_PMP_50:novel_5:sd\endcsname{5.7}
\expandafter\gdef\csname num@eurosat13:C_dacis_PMP_50:novel_5:min\endcsname{49.8}
\expandafter\gdef\csname num@eurosat13:C_dacis_PMP_50:novel_5:max\endcsname{64.2}
\expandafter\gdef\csname num@eurosat13:C_dacis_PMP_50:novel_10\endcsname{61.2}
\expandafter\gdef\csname num@eurosat13:C_dacis_PMP_50:novel_10:sd\endcsname{5.3}
\expandafter\gdef\csname num@eurosat13:C_dacis_PMP_50:novel_10:min\endcsname{53.5}
\expandafter\gdef\csname num@eurosat13:C_dacis_PMP_50:novel_10:max\endcsname{66.2}
\expandafter\gdef\csname num@eurosat13:C_dacis_PMP_50:pct\endcsname{50.0}
\expandafter\gdef\csname num@eurosat13:C_dacis_PMP_50:params\endcsname{5.61}
\expandafter\gdef\csname num@eurosat13:C_dacis_PMP_50:macs_128\endcsname{0.522}
\expandafter\gdef\csname num@eurosat13:C_dacis_PMP_50:macs_224\endcsname{1.600}
\expandafter\gdef\csname num@eurosat13:C_dacis_PMP_50:cpu_1thread_b1_ms\endcsname{4.0}
\expandafter\gdef\csname num@eurosat13:C_dacis_PMP_50:cpu_all_threads_b1_ms\endcsname{2.0}
\expandafter\gdef\csname num@eurosat13:C_dacis_PMP_50:gpu_b1_ms\endcsname{1.0}
\expandafter\gdef\csname num@eurosat13:C_dacis_PMP_50:gpu_ips\endcsname{73373}
\expandafter\gdef\csname num@eurosat13:C_l1_PMP_10:novel_1\endcsname{42.8}
\expandafter\gdef\csname num@eurosat13:C_l1_PMP_10:novel_1:sd\endcsname{9.0}
\expandafter\gdef\csname num@eurosat13:C_l1_PMP_10:novel_1:min\endcsname{29.7}
\expandafter\gdef\csname num@eurosat13:C_l1_PMP_10:novel_1:max\endcsname{55.0}
\expandafter\gdef\csname num@eurosat13:C_l1_PMP_10:novel_5\endcsname{50.9}
\expandafter\gdef\csname num@eurosat13:C_l1_PMP_10:novel_5:sd\endcsname{10.4}
\expandafter\gdef\csname num@eurosat13:C_l1_PMP_10:novel_5:min\endcsname{35.2}
\expandafter\gdef\csname num@eurosat13:C_l1_PMP_10:novel_5:max\endcsname{63.9}
\expandafter\gdef\csname num@eurosat13:C_l1_PMP_10:novel_10\endcsname{52.9}
\expandafter\gdef\csname num@eurosat13:C_l1_PMP_10:novel_10:sd\endcsname{9.9}
\expandafter\gdef\csname num@eurosat13:C_l1_PMP_10:novel_10:min\endcsname{38.0}
\expandafter\gdef\csname num@eurosat13:C_l1_PMP_10:novel_10:max\endcsname{65.2}
\expandafter\gdef\csname num@eurosat13:C_l1_PMP_10:pct\endcsname{10.0}
\expandafter\gdef\csname num@eurosat13:C_l1_PMP_10:params\endcsname{1.12}
\expandafter\gdef\csname num@eurosat13:C_l1_PMP_10:macs_128\endcsname{0.231}
\expandafter\gdef\csname num@eurosat13:C_l1_PMP_10:macs_224\endcsname{0.708}
\expandafter\gdef\csname num@eurosat13:C_l1_PMP_10:cpu_1thread_b1_ms\endcsname{2.4}
\expandafter\gdef\csname num@eurosat13:C_l1_PMP_10:cpu_all_threads_b1_ms\endcsname{1.5}
\expandafter\gdef\csname num@eurosat13:C_l1_PMP_10:gpu_b1_ms\endcsname{0.9}
\expandafter\gdef\csname num@eurosat13:C_l1_PMP_10:gpu_ips\endcsname{86721}
\expandafter\gdef\csname num@eurosat13:C_l1_PMP_20:novel_1\endcsname{44.0}
\expandafter\gdef\csname num@eurosat13:C_l1_PMP_20:novel_1:sd\endcsname{7.9}
\expandafter\gdef\csname num@eurosat13:C_l1_PMP_20:novel_1:min\endcsname{34.8}
\expandafter\gdef\csname num@eurosat13:C_l1_PMP_20:novel_1:max\endcsname{55.8}
\expandafter\gdef\csname num@eurosat13:C_l1_PMP_20:novel_5\endcsname{54.6}
\expandafter\gdef\csname num@eurosat13:C_l1_PMP_20:novel_5:sd\endcsname{8.7}
\expandafter\gdef\csname num@eurosat13:C_l1_PMP_20:novel_5:min\endcsname{43.2}
\expandafter\gdef\csname num@eurosat13:C_l1_PMP_20:novel_5:max\endcsname{66.6}
\expandafter\gdef\csname num@eurosat13:C_l1_PMP_20:novel_10\endcsname{56.9}
\expandafter\gdef\csname num@eurosat13:C_l1_PMP_20:novel_10:sd\endcsname{8.1}
\expandafter\gdef\csname num@eurosat13:C_l1_PMP_20:novel_10:min\endcsname{46.9}
\expandafter\gdef\csname num@eurosat13:C_l1_PMP_20:novel_10:max\endcsname{68.3}
\expandafter\gdef\csname num@eurosat13:C_l1_PMP_20:pct\endcsname{20.0}
\expandafter\gdef\csname num@eurosat13:C_l1_PMP_20:params\endcsname{2.24}
\expandafter\gdef\csname num@eurosat13:C_l1_PMP_20:macs_128\endcsname{0.304}
\expandafter\gdef\csname num@eurosat13:C_l1_PMP_20:macs_224\endcsname{0.932}
\expandafter\gdef\csname num@eurosat13:C_l1_PMP_20:cpu_1thread_b1_ms\endcsname{2.8}
\expandafter\gdef\csname num@eurosat13:C_l1_PMP_20:cpu_all_threads_b1_ms\endcsname{1.7}
\expandafter\gdef\csname num@eurosat13:C_l1_PMP_20:gpu_b1_ms\endcsname{1.0}
\expandafter\gdef\csname num@eurosat13:C_l1_PMP_20:gpu_ips\endcsname{82050}
\expandafter\gdef\csname num@eurosat13:C_l1_PMP_50:novel_1\endcsname{47.9}
\expandafter\gdef\csname num@eurosat13:C_l1_PMP_50:novel_1:sd\endcsname{5.1}
\expandafter\gdef\csname num@eurosat13:C_l1_PMP_50:novel_1:min\endcsname{40.9}
\expandafter\gdef\csname num@eurosat13:C_l1_PMP_50:novel_1:max\endcsname{54.2}
\expandafter\gdef\csname num@eurosat13:C_l1_PMP_50:novel_5\endcsname{60.1}
\expandafter\gdef\csname num@eurosat13:C_l1_PMP_50:novel_5:sd\endcsname{4.4}
\expandafter\gdef\csname num@eurosat13:C_l1_PMP_50:novel_5:min\endcsname{53.8}
\expandafter\gdef\csname num@eurosat13:C_l1_PMP_50:novel_5:max\endcsname{65.2}
\expandafter\gdef\csname num@eurosat13:C_l1_PMP_50:novel_10\endcsname{62.7}
\expandafter\gdef\csname num@eurosat13:C_l1_PMP_50:novel_10:sd\endcsname{4.0}
\expandafter\gdef\csname num@eurosat13:C_l1_PMP_50:novel_10:min\endcsname{57.0}
\expandafter\gdef\csname num@eurosat13:C_l1_PMP_50:novel_10:max\endcsname{67.3}
\expandafter\gdef\csname num@eurosat13:C_l1_PMP_50:pct\endcsname{50.0}
\expandafter\gdef\csname num@eurosat13:C_l1_PMP_50:params\endcsname{5.61}
\expandafter\gdef\csname num@eurosat13:C_l1_PMP_50:macs_128\endcsname{0.522}
\expandafter\gdef\csname num@eurosat13:C_l1_PMP_50:macs_224\endcsname{1.600}
\expandafter\gdef\csname num@eurosat13:C_l1_PMP_50:cpu_1thread_b1_ms\endcsname{4.1}
\expandafter\gdef\csname num@eurosat13:C_l1_PMP_50:cpu_all_threads_b1_ms\endcsname{2.1}
\expandafter\gdef\csname num@eurosat13:C_l1_PMP_50:gpu_b1_ms\endcsname{1.0}
\expandafter\gdef\csname num@eurosat13:C_l1_PMP_50:gpu_ips\endcsname{73256}
\expandafter\gdef\csname num@eurosat13:C_random_PMP_10:novel_1\endcsname{40.2}
\expandafter\gdef\csname num@eurosat13:C_random_PMP_10:novel_1:sd\endcsname{6.6}
\expandafter\gdef\csname num@eurosat13:C_random_PMP_10:novel_1:min\endcsname{33.8}
\expandafter\gdef\csname num@eurosat13:C_random_PMP_10:novel_1:max\endcsname{50.8}
\expandafter\gdef\csname num@eurosat13:C_random_PMP_10:novel_5\endcsname{49.2}
\expandafter\gdef\csname num@eurosat13:C_random_PMP_10:novel_5:sd\endcsname{7.4}
\expandafter\gdef\csname num@eurosat13:C_random_PMP_10:novel_5:min\endcsname{42.4}
\expandafter\gdef\csname num@eurosat13:C_random_PMP_10:novel_5:max\endcsname{61.1}
\expandafter\gdef\csname num@eurosat13:C_random_PMP_10:novel_10\endcsname{51.4}
\expandafter\gdef\csname num@eurosat13:C_random_PMP_10:novel_10:sd\endcsname{7.2}
\expandafter\gdef\csname num@eurosat13:C_random_PMP_10:novel_10:min\endcsname{45.3}
\expandafter\gdef\csname num@eurosat13:C_random_PMP_10:novel_10:max\endcsname{63.3}
\expandafter\gdef\csname num@eurosat13:C_random_PMP_10:pct\endcsname{10.0}
\expandafter\gdef\csname num@eurosat13:C_random_PMP_10:params\endcsname{1.12}
\expandafter\gdef\csname num@eurosat13:C_random_PMP_10:macs_128\endcsname{0.231}
\expandafter\gdef\csname num@eurosat13:C_random_PMP_10:macs_224\endcsname{0.708}
\expandafter\gdef\csname num@eurosat13:C_random_PMP_20:novel_1\endcsname{43.9}
\expandafter\gdef\csname num@eurosat13:C_random_PMP_20:novel_1:sd\endcsname{7.1}
\expandafter\gdef\csname num@eurosat13:C_random_PMP_20:novel_1:min\endcsname{34.3}
\expandafter\gdef\csname num@eurosat13:C_random_PMP_20:novel_1:max\endcsname{53.9}
\expandafter\gdef\csname num@eurosat13:C_random_PMP_20:novel_5\endcsname{54.5}
\expandafter\gdef\csname num@eurosat13:C_random_PMP_20:novel_5:sd\endcsname{7.4}
\expandafter\gdef\csname num@eurosat13:C_random_PMP_20:novel_5:min\endcsname{43.5}
\expandafter\gdef\csname num@eurosat13:C_random_PMP_20:novel_5:max\endcsname{63.8}
\expandafter\gdef\csname num@eurosat13:C_random_PMP_20:novel_10\endcsname{57.1}
\expandafter\gdef\csname num@eurosat13:C_random_PMP_20:novel_10:sd\endcsname{6.7}
\expandafter\gdef\csname num@eurosat13:C_random_PMP_20:novel_10:min\endcsname{47.0}
\expandafter\gdef\csname num@eurosat13:C_random_PMP_20:novel_10:max\endcsname{65.4}
\expandafter\gdef\csname num@eurosat13:C_random_PMP_20:pct\endcsname{20.0}
\expandafter\gdef\csname num@eurosat13:C_random_PMP_20:params\endcsname{2.24}
\expandafter\gdef\csname num@eurosat13:C_random_PMP_20:macs_128\endcsname{0.304}
\expandafter\gdef\csname num@eurosat13:C_random_PMP_20:macs_224\endcsname{0.932}
\expandafter\gdef\csname num@eurosat13:C_random_PMP_50:novel_1\endcsname{47.0}
\expandafter\gdef\csname num@eurosat13:C_random_PMP_50:novel_1:sd\endcsname{5.4}
\expandafter\gdef\csname num@eurosat13:C_random_PMP_50:novel_1:min\endcsname{38.1}
\expandafter\gdef\csname num@eurosat13:C_random_PMP_50:novel_1:max\endcsname{52.1}
\expandafter\gdef\csname num@eurosat13:C_random_PMP_50:novel_5\endcsname{59.4}
\expandafter\gdef\csname num@eurosat13:C_random_PMP_50:novel_5:sd\endcsname{5.5}
\expandafter\gdef\csname num@eurosat13:C_random_PMP_50:novel_5:min\endcsname{50.4}
\expandafter\gdef\csname num@eurosat13:C_random_PMP_50:novel_5:max\endcsname{64.2}
\expandafter\gdef\csname num@eurosat13:C_random_PMP_50:novel_10\endcsname{61.8}
\expandafter\gdef\csname num@eurosat13:C_random_PMP_50:novel_10:sd\endcsname{5.0}
\expandafter\gdef\csname num@eurosat13:C_random_PMP_50:novel_10:min\endcsname{54.0}
\expandafter\gdef\csname num@eurosat13:C_random_PMP_50:novel_10:max\endcsname{66.1}
\expandafter\gdef\csname num@eurosat13:C_random_PMP_50:pct\endcsname{50.0}
\expandafter\gdef\csname num@eurosat13:C_random_PMP_50:params\endcsname{5.61}
\expandafter\gdef\csname num@eurosat13:C_random_PMP_50:macs_128\endcsname{0.522}
\expandafter\gdef\csname num@eurosat13:C_random_PMP_50:macs_224\endcsname{1.600}
\expandafter\gdef\csname num@eurosat13:C_taylor_PMP_10:novel_1\endcsname{44.9}
\expandafter\gdef\csname num@eurosat13:C_taylor_PMP_10:novel_1:sd\endcsname{8.5}
\expandafter\gdef\csname num@eurosat13:C_taylor_PMP_10:novel_1:min\endcsname{36.2}
\expandafter\gdef\csname num@eurosat13:C_taylor_PMP_10:novel_1:max\endcsname{58.8}
\expandafter\gdef\csname num@eurosat13:C_taylor_PMP_10:novel_5\endcsname{54.3}
\expandafter\gdef\csname num@eurosat13:C_taylor_PMP_10:novel_5:sd\endcsname{8.3}
\expandafter\gdef\csname num@eurosat13:C_taylor_PMP_10:novel_5:min\endcsname{45.3}
\expandafter\gdef\csname num@eurosat13:C_taylor_PMP_10:novel_5:max\endcsname{67.4}
\expandafter\gdef\csname num@eurosat13:C_taylor_PMP_10:novel_10\endcsname{56.5}
\expandafter\gdef\csname num@eurosat13:C_taylor_PMP_10:novel_10:sd\endcsname{7.8}
\expandafter\gdef\csname num@eurosat13:C_taylor_PMP_10:novel_10:min\endcsname{48.4}
\expandafter\gdef\csname num@eurosat13:C_taylor_PMP_10:novel_10:max\endcsname{68.9}
\expandafter\gdef\csname num@eurosat13:C_taylor_PMP_10:pct\endcsname{10.0}
\expandafter\gdef\csname num@eurosat13:C_taylor_PMP_10:params\endcsname{1.12}
\expandafter\gdef\csname num@eurosat13:C_taylor_PMP_10:macs_128\endcsname{0.231}
\expandafter\gdef\csname num@eurosat13:C_taylor_PMP_10:macs_224\endcsname{0.708}
\expandafter\gdef\csname num@eurosat13:C_taylor_PMP_20:novel_1\endcsname{44.3}
\expandafter\gdef\csname num@eurosat13:C_taylor_PMP_20:novel_1:sd\endcsname{7.3}
\expandafter\gdef\csname num@eurosat13:C_taylor_PMP_20:novel_1:min\endcsname{35.3}
\expandafter\gdef\csname num@eurosat13:C_taylor_PMP_20:novel_1:max\endcsname{54.3}
\expandafter\gdef\csname num@eurosat13:C_taylor_PMP_20:novel_5\endcsname{54.8}
\expandafter\gdef\csname num@eurosat13:C_taylor_PMP_20:novel_5:sd\endcsname{7.3}
\expandafter\gdef\csname num@eurosat13:C_taylor_PMP_20:novel_5:min\endcsname{45.1}
\expandafter\gdef\csname num@eurosat13:C_taylor_PMP_20:novel_5:max\endcsname{64.3}
\expandafter\gdef\csname num@eurosat13:C_taylor_PMP_20:novel_10\endcsname{57.4}
\expandafter\gdef\csname num@eurosat13:C_taylor_PMP_20:novel_10:sd\endcsname{6.9}
\expandafter\gdef\csname num@eurosat13:C_taylor_PMP_20:novel_10:min\endcsname{48.4}
\expandafter\gdef\csname num@eurosat13:C_taylor_PMP_20:novel_10:max\endcsname{66.1}
\expandafter\gdef\csname num@eurosat13:C_taylor_PMP_20:pct\endcsname{20.0}
\expandafter\gdef\csname num@eurosat13:C_taylor_PMP_20:params\endcsname{2.24}
\expandafter\gdef\csname num@eurosat13:C_taylor_PMP_20:macs_128\endcsname{0.304}
\expandafter\gdef\csname num@eurosat13:C_taylor_PMP_20:macs_224\endcsname{0.932}
\expandafter\gdef\csname num@eurosat13:C_taylor_PMP_50:novel_1\endcsname{47.0}
\expandafter\gdef\csname num@eurosat13:C_taylor_PMP_50:novel_1:sd\endcsname{6.0}
\expandafter\gdef\csname num@eurosat13:C_taylor_PMP_50:novel_1:min\endcsname{37.9}
\expandafter\gdef\csname num@eurosat13:C_taylor_PMP_50:novel_1:max\endcsname{54.4}
\expandafter\gdef\csname num@eurosat13:C_taylor_PMP_50:novel_5\endcsname{59.3}
\expandafter\gdef\csname num@eurosat13:C_taylor_PMP_50:novel_5:sd\endcsname{6.1}
\expandafter\gdef\csname num@eurosat13:C_taylor_PMP_50:novel_5:min\endcsname{49.7}
\expandafter\gdef\csname num@eurosat13:C_taylor_PMP_50:novel_5:max\endcsname{66.3}
\expandafter\gdef\csname num@eurosat13:C_taylor_PMP_50:novel_10\endcsname{61.9}
\expandafter\gdef\csname num@eurosat13:C_taylor_PMP_50:novel_10:sd\endcsname{5.7}
\expandafter\gdef\csname num@eurosat13:C_taylor_PMP_50:novel_10:min\endcsname{53.2}
\expandafter\gdef\csname num@eurosat13:C_taylor_PMP_50:novel_10:max\endcsname{68.5}
\expandafter\gdef\csname num@eurosat13:C_taylor_PMP_50:pct\endcsname{50.0}
\expandafter\gdef\csname num@eurosat13:C_taylor_PMP_50:params\endcsname{5.61}
\expandafter\gdef\csname num@eurosat13:C_taylor_PMP_50:macs_128\endcsname{0.522}
\expandafter\gdef\csname num@eurosat13:C_taylor_PMP_50:macs_224\endcsname{1.600}
\expandafter\gdef\csname num@eurosat13:K_efficientnet_b0:novel_1\endcsname{54.4}
\expandafter\gdef\csname num@eurosat13:K_efficientnet_b0:novel_1:sd\endcsname{6.8}
\expandafter\gdef\csname num@eurosat13:K_efficientnet_b0:novel_1:min\endcsname{49.1}
\expandafter\gdef\csname num@eurosat13:K_efficientnet_b0:novel_1:max\endcsname{64.6}
\expandafter\gdef\csname num@eurosat13:K_efficientnet_b0:novel_5\endcsname{69.9}
\expandafter\gdef\csname num@eurosat13:K_efficientnet_b0:novel_5:sd\endcsname{5.1}
\expandafter\gdef\csname num@eurosat13:K_efficientnet_b0:novel_5:min\endcsname{64.2}
\expandafter\gdef\csname num@eurosat13:K_efficientnet_b0:novel_5:max\endcsname{77.4}
\expandafter\gdef\csname num@eurosat13:K_efficientnet_b0:novel_10\endcsname{74.1}
\expandafter\gdef\csname num@eurosat13:K_efficientnet_b0:novel_10:sd\endcsname{4.9}
\expandafter\gdef\csname num@eurosat13:K_efficientnet_b0:novel_10:min\endcsname{67.7}
\expandafter\gdef\csname num@eurosat13:K_efficientnet_b0:novel_10:max\endcsname{81.1}
\expandafter\gdef\csname num@eurosat13:K_efficientnet_b0:pct\endcsname{35.8}
\expandafter\gdef\csname num@eurosat13:K_efficientnet_b0:params\endcsname{4.01}
\expandafter\gdef\csname num@eurosat13:K_efficientnet_b0:macs_128\endcsname{0.138}
\expandafter\gdef\csname num@eurosat13:K_efficientnet_b0:macs_224\endcsname{0.421}
\expandafter\gdef\csname num@eurosat13:K_efficientnet_b0:cpu_1thread_b1_ms\endcsname{4.3}
\expandafter\gdef\csname num@eurosat13:K_efficientnet_b0:cpu_all_threads_b1_ms\endcsname{5.2}
\expandafter\gdef\csname num@eurosat13:K_efficientnet_b0:gpu_b1_ms\endcsname{2.9}
\expandafter\gdef\csname num@eurosat13:K_efficientnet_b0:gpu_ips\endcsname{31521}
\expandafter\gdef\csname num@eurosat13:K_mobilenet_v2:novel_1\endcsname{53.6}
\expandafter\gdef\csname num@eurosat13:K_mobilenet_v2:novel_1:sd\endcsname{4.5}
\expandafter\gdef\csname num@eurosat13:K_mobilenet_v2:novel_1:min\endcsname{48.2}
\expandafter\gdef\csname num@eurosat13:K_mobilenet_v2:novel_1:max\endcsname{59.5}
\expandafter\gdef\csname num@eurosat13:K_mobilenet_v2:novel_5\endcsname{65.4}
\expandafter\gdef\csname num@eurosat13:K_mobilenet_v2:novel_5:sd\endcsname{3.5}
\expandafter\gdef\csname num@eurosat13:K_mobilenet_v2:novel_5:min\endcsname{61.6}
\expandafter\gdef\csname num@eurosat13:K_mobilenet_v2:novel_5:max\endcsname{68.6}
\expandafter\gdef\csname num@eurosat13:K_mobilenet_v2:novel_10\endcsname{67.6}
\expandafter\gdef\csname num@eurosat13:K_mobilenet_v2:novel_10:sd\endcsname{3.3}
\expandafter\gdef\csname num@eurosat13:K_mobilenet_v2:novel_10:min\endcsname{63.5}
\expandafter\gdef\csname num@eurosat13:K_mobilenet_v2:novel_10:max\endcsname{71.3}
\expandafter\gdef\csname num@eurosat13:K_mobilenet_v2:pct\endcsname{19.9}
\expandafter\gdef\csname num@eurosat13:K_mobilenet_v2:params\endcsname{2.23}
\expandafter\gdef\csname num@eurosat13:K_mobilenet_v2:macs_128\endcsname{0.110}
\expandafter\gdef\csname num@eurosat13:K_mobilenet_v2:macs_224\endcsname{0.336}
\expandafter\gdef\csname num@eurosat13:K_mobilenet_v2:cpu_1thread_b1_ms\endcsname{2.9}
\expandafter\gdef\csname num@eurosat13:K_mobilenet_v2:cpu_all_threads_b1_ms\endcsname{4.3}
\expandafter\gdef\csname num@eurosat13:K_mobilenet_v2:gpu_b1_ms\endcsname{1.8}
\expandafter\gdef\csname num@eurosat13:K_mobilenet_v2:gpu_ips\endcsname{43023}
\expandafter\gdef\csname num@eurosat13:K_mobilenet_v3_small:novel_1\endcsname{57.7}
\expandafter\gdef\csname num@eurosat13:K_mobilenet_v3_small:novel_1:sd\endcsname{7.7}
\expandafter\gdef\csname num@eurosat13:K_mobilenet_v3_small:novel_1:min\endcsname{51.3}
\expandafter\gdef\csname num@eurosat13:K_mobilenet_v3_small:novel_1:max\endcsname{67.4}
\expandafter\gdef\csname num@eurosat13:K_mobilenet_v3_small:novel_5\endcsname{72.4}
\expandafter\gdef\csname num@eurosat13:K_mobilenet_v3_small:novel_5:sd\endcsname{4.9}
\expandafter\gdef\csname num@eurosat13:K_mobilenet_v3_small:novel_5:min\endcsname{66.8}
\expandafter\gdef\csname num@eurosat13:K_mobilenet_v3_small:novel_5:max\endcsname{78.3}
\expandafter\gdef\csname num@eurosat13:K_mobilenet_v3_small:novel_10\endcsname{75.7}
\expandafter\gdef\csname num@eurosat13:K_mobilenet_v3_small:novel_10:sd\endcsname{4.5}
\expandafter\gdef\csname num@eurosat13:K_mobilenet_v3_small:novel_10:min\endcsname{70.1}
\expandafter\gdef\csname num@eurosat13:K_mobilenet_v3_small:novel_10:max\endcsname{81.1}
\expandafter\gdef\csname num@eurosat13:K_mobilenet_v3_small:pct\endcsname{8.3}
\expandafter\gdef\csname num@eurosat13:K_mobilenet_v3_small:params\endcsname{0.93}
\expandafter\gdef\csname num@eurosat13:K_mobilenet_v3_small:macs_128\endcsname{0.024}
\expandafter\gdef\csname num@eurosat13:K_mobilenet_v3_small:macs_224\endcsname{0.073}
\expandafter\gdef\csname num@eurosat13:K_mobilenet_v3_small:cpu_1thread_b1_ms\endcsname{1.9}
\expandafter\gdef\csname num@eurosat13:K_mobilenet_v3_small:cpu_all_threads_b1_ms\endcsname{4.7}
\expandafter\gdef\csname num@eurosat13:K_mobilenet_v3_small:gpu_b1_ms\endcsname{1.9}
\expandafter\gdef\csname num@eurosat13:K_mobilenet_v3_small:gpu_ips\endcsname{65619}
\expandafter\gdef\csname num@eurosat13:dense_ce:novel_1\endcsname{49.4}
\expandafter\gdef\csname num@eurosat13:dense_ce:novel_1:sd\endcsname{5.5}
\expandafter\gdef\csname num@eurosat13:dense_ce:novel_1:min\endcsname{40.8}
\expandafter\gdef\csname num@eurosat13:dense_ce:novel_1:max\endcsname{55.2}
\expandafter\gdef\csname num@eurosat13:dense_ce:novel_5\endcsname{62.1}
\expandafter\gdef\csname num@eurosat13:dense_ce:novel_5:sd\endcsname{5.3}
\expandafter\gdef\csname num@eurosat13:dense_ce:novel_5:min\endcsname{54.0}
\expandafter\gdef\csname num@eurosat13:dense_ce:novel_5:max\endcsname{67.7}
\expandafter\gdef\csname num@eurosat13:dense_ce:novel_10\endcsname{64.9}
\expandafter\gdef\csname num@eurosat13:dense_ce:novel_10:sd\endcsname{5.0}
\expandafter\gdef\csname num@eurosat13:dense_ce:novel_10:min\endcsname{57.5}
\expandafter\gdef\csname num@eurosat13:dense_ce:novel_10:max\endcsname{70.5}
\expandafter\gdef\csname num@eurosat13:dense_ce:pct\endcsname{100.0}
\expandafter\gdef\csname num@eurosat13:dense_ce:params\endcsname{11.21}
\expandafter\gdef\csname num@eurosat13:dense_ce:macs_128\endcsname{0.721}
\expandafter\gdef\csname num@eurosat13:dense_ce:macs_224\endcsname{2.207}
\expandafter\gdef\csname num@eurosat13:dense_ce:cpu_1thread_b1_ms\endcsname{5.7}
\expandafter\gdef\csname num@eurosat13:dense_ce:cpu_all_threads_b1_ms\endcsname{2.2}
\expandafter\gdef\csname num@eurosat13:dense_ce:gpu_b1_ms\endcsname{1.0}
\expandafter\gdef\csname num@eurosat13:dense_ce:gpu_ips\endcsname{67644}
\expandafter\gdef\csname num@eurosat13:dense_m:novel_1\endcsname{48.9}
\expandafter\gdef\csname num@eurosat13:dense_m:novel_1:sd\endcsname{5.4}
\expandafter\gdef\csname num@eurosat13:dense_m:novel_1:min\endcsname{41.5}
\expandafter\gdef\csname num@eurosat13:dense_m:novel_1:max\endcsname{56.0}
\expandafter\gdef\csname num@eurosat13:dense_m:novel_5\endcsname{61.3}
\expandafter\gdef\csname num@eurosat13:dense_m:novel_5:sd\endcsname{4.4}
\expandafter\gdef\csname num@eurosat13:dense_m:novel_5:min\endcsname{55.3}
\expandafter\gdef\csname num@eurosat13:dense_m:novel_5:max\endcsname{66.7}
\expandafter\gdef\csname num@eurosat13:dense_m:novel_10\endcsname{63.9}
\expandafter\gdef\csname num@eurosat13:dense_m:novel_10:sd\endcsname{3.9}
\expandafter\gdef\csname num@eurosat13:dense_m:novel_10:min\endcsname{58.8}
\expandafter\gdef\csname num@eurosat13:dense_m:novel_10:max\endcsname{68.7}
\expandafter\gdef\csname num@eurosat13:dense_m:pct\endcsname{100.0}
\expandafter\gdef\csname num@eurosat13:dense_m:params\endcsname{11.21}
\expandafter\gdef\csname num@eurosat13:dense_m:macs_128\endcsname{0.721}
\expandafter\gdef\csname num@eurosat13:dense_m:macs_224\endcsname{2.207}
\expandafter\gdef\csname num@eurosat13:dense_rmr:novel_1\endcsname{49.0}
\expandafter\gdef\csname num@eurosat13:dense_rmr:novel_1:sd\endcsname{5.2}
\expandafter\gdef\csname num@eurosat13:dense_rmr:novel_1:min\endcsname{40.6}
\expandafter\gdef\csname num@eurosat13:dense_rmr:novel_1:max\endcsname{54.3}
\expandafter\gdef\csname num@eurosat13:dense_rmr:novel_5\endcsname{61.7}
\expandafter\gdef\csname num@eurosat13:dense_rmr:novel_5:sd\endcsname{5.1}
\expandafter\gdef\csname num@eurosat13:dense_rmr:novel_5:min\endcsname{53.5}
\expandafter\gdef\csname num@eurosat13:dense_rmr:novel_5:max\endcsname{67.0}
\expandafter\gdef\csname num@eurosat13:dense_rmr:novel_10\endcsname{64.3}
\expandafter\gdef\csname num@eurosat13:dense_rmr:novel_10:sd\endcsname{5.0}
\expandafter\gdef\csname num@eurosat13:dense_rmr:novel_10:min\endcsname{56.7}
\expandafter\gdef\csname num@eurosat13:dense_rmr:novel_10:max\endcsname{69.9}
\expandafter\gdef\csname num@eurosat13:dense_rmr:pct\endcsname{100.0}
\expandafter\gdef\csname num@eurosat13:dense_rmr:params\endcsname{11.21}
\expandafter\gdef\csname num@eurosat13:dense_rmr:macs_128\endcsname{0.721}
\expandafter\gdef\csname num@eurosat13:dense_rmr:macs_224\endcsname{2.207}
\expandafter\gdef\csname num@eurosat3:seeds\endcsname{5}
\expandafter\gdef\csname num@eurosat3:A_dacis_PMP_30:novel_1\endcsname{50.2}
\expandafter\gdef\csname num@eurosat3:A_dacis_PMP_30:novel_1:sd\endcsname{3.5}
\expandafter\gdef\csname num@eurosat3:A_dacis_PMP_30:novel_1:min\endcsname{46.1}
\expandafter\gdef\csname num@eurosat3:A_dacis_PMP_30:novel_1:max\endcsname{55.1}
\expandafter\gdef\csname num@eurosat3:A_dacis_PMP_30:novel_5\endcsname{63.6}
\expandafter\gdef\csname num@eurosat3:A_dacis_PMP_30:novel_5:sd\endcsname{3.9}
\expandafter\gdef\csname num@eurosat3:A_dacis_PMP_30:novel_5:min\endcsname{58.2}
\expandafter\gdef\csname num@eurosat3:A_dacis_PMP_30:novel_5:max\endcsname{68.0}
\expandafter\gdef\csname num@eurosat3:A_dacis_PMP_30:novel_10\endcsname{66.2}
\expandafter\gdef\csname num@eurosat3:A_dacis_PMP_30:novel_10:sd\endcsname{3.7}
\expandafter\gdef\csname num@eurosat3:A_dacis_PMP_30:novel_10:min\endcsname{60.7}
\expandafter\gdef\csname num@eurosat3:A_dacis_PMP_30:novel_10:max\endcsname{69.9}
\expandafter\gdef\csname num@eurosat3:A_dacis_PMP_30:pct\endcsname{30.0}
\expandafter\gdef\csname num@eurosat3:A_dacis_PMP_30:params\endcsname{3.35}
\expandafter\gdef\csname num@eurosat3:A_dacis_PMP_30:macs_128\endcsname{0.251}
\expandafter\gdef\csname num@eurosat3:A_dacis_PMP_30:macs_224\endcsname{0.770}
\expandafter\gdef\csname num@eurosat3:A_dacis_PMP_30:cpu_1thread_b1_ms\endcsname{2.8}
\expandafter\gdef\csname num@eurosat3:A_dacis_PMP_30:cpu_all_threads_b1_ms\endcsname{1.8}
\expandafter\gdef\csname num@eurosat3:A_dacis_PMP_30:gpu_b1_ms\endcsname{0.9}
\expandafter\gdef\csname num@eurosat3:A_dacis_PMP_30:gpu_ips\endcsname{95036}
\expandafter\gdef\csname num@eurosat3:A_l1_PMP_30:novel_1\endcsname{49.8}
\expandafter\gdef\csname num@eurosat3:A_l1_PMP_30:novel_1:sd\endcsname{2.4}
\expandafter\gdef\csname num@eurosat3:A_l1_PMP_30:novel_1:min\endcsname{48.4}
\expandafter\gdef\csname num@eurosat3:A_l1_PMP_30:novel_1:max\endcsname{54.0}
\expandafter\gdef\csname num@eurosat3:A_l1_PMP_30:novel_5\endcsname{63.9}
\expandafter\gdef\csname num@eurosat3:A_l1_PMP_30:novel_5:sd\endcsname{2.6}
\expandafter\gdef\csname num@eurosat3:A_l1_PMP_30:novel_5:min\endcsname{62.0}
\expandafter\gdef\csname num@eurosat3:A_l1_PMP_30:novel_5:max\endcsname{68.3}
\expandafter\gdef\csname num@eurosat3:A_l1_PMP_30:novel_10\endcsname{66.6}
\expandafter\gdef\csname num@eurosat3:A_l1_PMP_30:novel_10:sd\endcsname{2.3}
\expandafter\gdef\csname num@eurosat3:A_l1_PMP_30:novel_10:min\endcsname{64.3}
\expandafter\gdef\csname num@eurosat3:A_l1_PMP_30:novel_10:max\endcsname{70.3}
\expandafter\gdef\csname num@eurosat3:A_l1_PMP_30:pct\endcsname{30.0}
\expandafter\gdef\csname num@eurosat3:A_l1_PMP_30:params\endcsname{3.35}
\expandafter\gdef\csname num@eurosat3:A_l1_PMP_30:macs_128\endcsname{0.251}
\expandafter\gdef\csname num@eurosat3:A_l1_PMP_30:macs_224\endcsname{0.770}
\expandafter\gdef\csname num@eurosat3:A_l1_PMP_30:cpu_1thread_b1_ms\endcsname{2.8}
\expandafter\gdef\csname num@eurosat3:A_l1_PMP_30:cpu_all_threads_b1_ms\endcsname{1.8}
\expandafter\gdef\csname num@eurosat3:A_l1_PMP_30:gpu_b1_ms\endcsname{0.9}
\expandafter\gdef\csname num@eurosat3:A_l1_PMP_30:gpu_ips\endcsname{94937}
\expandafter\gdef\csname num@eurosat3:K_mobilenet_v3_small:novel_1\endcsname{58.3}
\expandafter\gdef\csname num@eurosat3:K_mobilenet_v3_small:novel_1:sd\endcsname{3.7}
\expandafter\gdef\csname num@eurosat3:K_mobilenet_v3_small:novel_1:min\endcsname{54.5}
\expandafter\gdef\csname num@eurosat3:K_mobilenet_v3_small:novel_1:max\endcsname{64.0}
\expandafter\gdef\csname num@eurosat3:K_mobilenet_v3_small:novel_5\endcsname{74.7}
\expandafter\gdef\csname num@eurosat3:K_mobilenet_v3_small:novel_5:sd\endcsname{4.1}
\expandafter\gdef\csname num@eurosat3:K_mobilenet_v3_small:novel_5:min\endcsname{71.0}
\expandafter\gdef\csname num@eurosat3:K_mobilenet_v3_small:novel_5:max\endcsname{81.0}
\expandafter\gdef\csname num@eurosat3:K_mobilenet_v3_small:novel_10\endcsname{78.0}
\expandafter\gdef\csname num@eurosat3:K_mobilenet_v3_small:novel_10:sd\endcsname{4.1}
\expandafter\gdef\csname num@eurosat3:K_mobilenet_v3_small:novel_10:min\endcsname{74.1}
\expandafter\gdef\csname num@eurosat3:K_mobilenet_v3_small:novel_10:max\endcsname{84.2}
\expandafter\gdef\csname num@eurosat3:K_mobilenet_v3_small:pct\endcsname{8.3}
\expandafter\gdef\csname num@eurosat3:K_mobilenet_v3_small:params\endcsname{0.93}
\expandafter\gdef\csname num@eurosat3:K_mobilenet_v3_small:macs_128\endcsname{0.018}
\expandafter\gdef\csname num@eurosat3:K_mobilenet_v3_small:macs_224\endcsname{0.055}
\expandafter\gdef\csname num@eurosat3:K_mobilenet_v3_small:cpu_1thread_b1_ms\endcsname{1.9}
\expandafter\gdef\csname num@eurosat3:K_mobilenet_v3_small:cpu_all_threads_b1_ms\endcsname{4.9}
\expandafter\gdef\csname num@eurosat3:K_mobilenet_v3_small:gpu_b1_ms\endcsname{1.9}
\expandafter\gdef\csname num@eurosat3:K_mobilenet_v3_small:gpu_ips\endcsname{65463}
\expandafter\gdef\csname num@eurosat3:dense_ce:novel_1\endcsname{53.7}
\expandafter\gdef\csname num@eurosat3:dense_ce:novel_1:sd\endcsname{4.3}
\expandafter\gdef\csname num@eurosat3:dense_ce:novel_1:min\endcsname{50.0}
\expandafter\gdef\csname num@eurosat3:dense_ce:novel_1:max\endcsname{60.1}
\expandafter\gdef\csname num@eurosat3:dense_ce:novel_5\endcsname{68.7}
\expandafter\gdef\csname num@eurosat3:dense_ce:novel_5:sd\endcsname{3.9}
\expandafter\gdef\csname num@eurosat3:dense_ce:novel_5:min\endcsname{65.0}
\expandafter\gdef\csname num@eurosat3:dense_ce:novel_5:max\endcsname{73.0}
\expandafter\gdef\csname num@eurosat3:dense_ce:novel_10\endcsname{71.4}
\expandafter\gdef\csname num@eurosat3:dense_ce:novel_10:sd\endcsname{3.6}
\expandafter\gdef\csname num@eurosat3:dense_ce:novel_10:min\endcsname{67.7}
\expandafter\gdef\csname num@eurosat3:dense_ce:novel_10:max\endcsname{75.7}
\expandafter\gdef\csname num@eurosat3:dense_ce:pct\endcsname{100.0}
\expandafter\gdef\csname num@eurosat3:dense_ce:params\endcsname{11.18}
\expandafter\gdef\csname num@eurosat3:dense_ce:macs_128\endcsname{0.592}
\expandafter\gdef\csname num@eurosat3:dense_ce:macs_224\endcsname{1.814}
\expandafter\gdef\csname num@eurosat3:dense_ce:cpu_1thread_b1_ms\endcsname{5.3}
\expandafter\gdef\csname num@eurosat3:dense_ce:cpu_all_threads_b1_ms\endcsname{2.2}
\expandafter\gdef\csname num@eurosat3:dense_ce:gpu_b1_ms\endcsname{1.0}
\expandafter\gdef\csname num@eurosat3:dense_ce:gpu_ips\endcsname{81631}
\expandafter\gdef\csname num@eurosat3:dense_m:novel_1\endcsname{54.2}
\expandafter\gdef\csname num@eurosat3:dense_m:novel_1:sd\endcsname{3.7}
\expandafter\gdef\csname num@eurosat3:dense_m:novel_1:min\endcsname{49.9}
\expandafter\gdef\csname num@eurosat3:dense_m:novel_1:max\endcsname{59.6}
\expandafter\gdef\csname num@eurosat3:dense_m:novel_5\endcsname{68.3}
\expandafter\gdef\csname num@eurosat3:dense_m:novel_5:sd\endcsname{3.5}
\expandafter\gdef\csname num@eurosat3:dense_m:novel_5:min\endcsname{65.2}
\expandafter\gdef\csname num@eurosat3:dense_m:novel_5:max\endcsname{72.6}
\expandafter\gdef\csname num@eurosat3:dense_m:novel_10\endcsname{70.8}
\expandafter\gdef\csname num@eurosat3:dense_m:novel_10:sd\endcsname{3.3}
\expandafter\gdef\csname num@eurosat3:dense_m:novel_10:min\endcsname{68.3}
\expandafter\gdef\csname num@eurosat3:dense_m:novel_10:max\endcsname{74.7}
\expandafter\gdef\csname num@eurosat3:dense_m:pct\endcsname{100.0}
\expandafter\gdef\csname num@eurosat3:dense_m:params\endcsname{11.18}
\expandafter\gdef\csname num@eurosat3:dense_m:macs_128\endcsname{0.592}
\expandafter\gdef\csname num@eurosat3:dense_m:macs_224\endcsname{1.814}
\expandafter\gdef\csname num@eurosat3:dense_rmr:novel_1\endcsname{53.1}
\expandafter\gdef\csname num@eurosat3:dense_rmr:novel_1:sd\endcsname{4.5}
\expandafter\gdef\csname num@eurosat3:dense_rmr:novel_1:min\endcsname{49.1}
\expandafter\gdef\csname num@eurosat3:dense_rmr:novel_1:max\endcsname{60.3}
\expandafter\gdef\csname num@eurosat3:dense_rmr:novel_5\endcsname{68.1}
\expandafter\gdef\csname num@eurosat3:dense_rmr:novel_5:sd\endcsname{3.8}
\expandafter\gdef\csname num@eurosat3:dense_rmr:novel_5:min\endcsname{64.3}
\expandafter\gdef\csname num@eurosat3:dense_rmr:novel_5:max\endcsname{73.0}
\expandafter\gdef\csname num@eurosat3:dense_rmr:novel_10\endcsname{70.9}
\expandafter\gdef\csname num@eurosat3:dense_rmr:novel_10:sd\endcsname{3.4}
\expandafter\gdef\csname num@eurosat3:dense_rmr:novel_10:min\endcsname{67.1}
\expandafter\gdef\csname num@eurosat3:dense_rmr:novel_10:max\endcsname{74.5}
\expandafter\gdef\csname num@eurosat3:dense_rmr:pct\endcsname{100.0}
\expandafter\gdef\csname num@eurosat3:dense_rmr:params\endcsname{11.18}
\expandafter\gdef\csname num@eurosat3:dense_rmr:macs_128\endcsname{0.592}
\expandafter\gdef\csname num@eurosat3:dense_rmr:macs_224\endcsname{1.814}
\expandafter\gdef\csname num@plantvillage:seeds\endcsname{5}
\expandafter\gdef\csname num@plantvillage:A2_D_P_30:novel_1\endcsname{71.9}
\expandafter\gdef\csname num@plantvillage:A2_D_P_30:novel_1:sd\endcsname{12.4}
\expandafter\gdef\csname num@plantvillage:A2_D_P_30:novel_1:min\endcsname{55.6}
\expandafter\gdef\csname num@plantvillage:A2_D_P_30:novel_1:max\endcsname{84.9}
\expandafter\gdef\csname num@plantvillage:A2_D_P_30:novel_5\endcsname{85.2}
\expandafter\gdef\csname num@plantvillage:A2_D_P_30:novel_5:sd\endcsname{8.9}
\expandafter\gdef\csname num@plantvillage:A2_D_P_30:novel_5:min\endcsname{73.1}
\expandafter\gdef\csname num@plantvillage:A2_D_P_30:novel_5:max\endcsname{94.1}
\expandafter\gdef\csname num@plantvillage:A2_D_P_30:novel_10\endcsname{87.7}
\expandafter\gdef\csname num@plantvillage:A2_D_P_30:novel_10:sd\endcsname{7.5}
\expandafter\gdef\csname num@plantvillage:A2_D_P_30:novel_10:min\endcsname{78.0}
\expandafter\gdef\csname num@plantvillage:A2_D_P_30:novel_10:max\endcsname{95.3}
\expandafter\gdef\csname num@plantvillage:A2_D_P_30:plantdoc_1\endcsname{34.8}
\expandafter\gdef\csname num@plantvillage:A2_D_P_30:plantdoc_1:sd\endcsname{0.6}
\expandafter\gdef\csname num@plantvillage:A2_D_P_30:plantdoc_1:min\endcsname{34.1}
\expandafter\gdef\csname num@plantvillage:A2_D_P_30:plantdoc_1:max\endcsname{35.7}
\expandafter\gdef\csname num@plantvillage:A2_D_P_30:plantdoc_5\endcsname{46.8}
\expandafter\gdef\csname num@plantvillage:A2_D_P_30:plantdoc_5:sd\endcsname{0.8}
\expandafter\gdef\csname num@plantvillage:A2_D_P_30:plantdoc_5:min\endcsname{46.1}
\expandafter\gdef\csname num@plantvillage:A2_D_P_30:plantdoc_5:max\endcsname{48.1}
\expandafter\gdef\csname num@plantvillage:A2_D_P_30:plantdoc_10\endcsname{51.2}
\expandafter\gdef\csname num@plantvillage:A2_D_P_30:plantdoc_10:sd\endcsname{1.0}
\expandafter\gdef\csname num@plantvillage:A2_D_P_30:plantdoc_10:min\endcsname{49.8}
\expandafter\gdef\csname num@plantvillage:A2_D_P_30:plantdoc_10:max\endcsname{52.6}
\expandafter\gdef\csname num@plantvillage:A2_D_P_30:pct\endcsname{30.0}
\expandafter\gdef\csname num@plantvillage:A2_D_P_30:params\endcsname{3.35}
\expandafter\gdef\csname num@plantvillage:A2_D_P_30:macs_128\endcsname{0.251}
\expandafter\gdef\csname num@plantvillage:A2_D_P_30:macs_224\endcsname{0.770}
\expandafter\gdef\csname num@plantvillage:A2_G_P_30:novel_1\endcsname{74.0}
\expandafter\gdef\csname num@plantvillage:A2_G_P_30:novel_1:sd\endcsname{11.0}
\expandafter\gdef\csname num@plantvillage:A2_G_P_30:novel_1:min\endcsname{59.4}
\expandafter\gdef\csname num@plantvillage:A2_G_P_30:novel_1:max\endcsname{86.2}
\expandafter\gdef\csname num@plantvillage:A2_G_P_30:novel_5\endcsname{87.3}
\expandafter\gdef\csname num@plantvillage:A2_G_P_30:novel_5:sd\endcsname{7.2}
\expandafter\gdef\csname num@plantvillage:A2_G_P_30:novel_5:min\endcsname{77.9}
\expandafter\gdef\csname num@plantvillage:A2_G_P_30:novel_5:max\endcsname{95.2}
\expandafter\gdef\csname num@plantvillage:A2_G_P_30:novel_10\endcsname{89.5}
\expandafter\gdef\csname num@plantvillage:A2_G_P_30:novel_10:sd\endcsname{5.9}
\expandafter\gdef\csname num@plantvillage:A2_G_P_30:novel_10:min\endcsname{81.9}
\expandafter\gdef\csname num@plantvillage:A2_G_P_30:novel_10:max\endcsname{96.1}
\expandafter\gdef\csname num@plantvillage:A2_G_P_30:plantdoc_1\endcsname{35.5}
\expandafter\gdef\csname num@plantvillage:A2_G_P_30:plantdoc_1:sd\endcsname{1.4}
\expandafter\gdef\csname num@plantvillage:A2_G_P_30:plantdoc_1:min\endcsname{33.5}
\expandafter\gdef\csname num@plantvillage:A2_G_P_30:plantdoc_1:max\endcsname{37.4}
\expandafter\gdef\csname num@plantvillage:A2_G_P_30:plantdoc_5\endcsname{47.6}
\expandafter\gdef\csname num@plantvillage:A2_G_P_30:plantdoc_5:sd\endcsname{1.5}
\expandafter\gdef\csname num@plantvillage:A2_G_P_30:plantdoc_5:min\endcsname{45.8}
\expandafter\gdef\csname num@plantvillage:A2_G_P_30:plantdoc_5:max\endcsname{49.9}
\expandafter\gdef\csname num@plantvillage:A2_G_P_30:plantdoc_10\endcsname{51.8}
\expandafter\gdef\csname num@plantvillage:A2_G_P_30:plantdoc_10:sd\endcsname{1.5}
\expandafter\gdef\csname num@plantvillage:A2_G_P_30:plantdoc_10:min\endcsname{50.3}
\expandafter\gdef\csname num@plantvillage:A2_G_P_30:plantdoc_10:max\endcsname{54.2}
\expandafter\gdef\csname num@plantvillage:A2_G_P_30:pct\endcsname{30.0}
\expandafter\gdef\csname num@plantvillage:A2_G_P_30:params\endcsname{3.35}
\expandafter\gdef\csname num@plantvillage:A2_G_P_30:macs_128\endcsname{0.251}
\expandafter\gdef\csname num@plantvillage:A2_G_P_30:macs_224\endcsname{0.770}
\expandafter\gdef\csname num@plantvillage:A2_V_P_30:novel_1\endcsname{72.9}
\expandafter\gdef\csname num@plantvillage:A2_V_P_30:novel_1:sd\endcsname{10.5}
\expandafter\gdef\csname num@plantvillage:A2_V_P_30:novel_1:min\endcsname{58.9}
\expandafter\gdef\csname num@plantvillage:A2_V_P_30:novel_1:max\endcsname{86.1}
\expandafter\gdef\csname num@plantvillage:A2_V_P_30:novel_5\endcsname{86.5}
\expandafter\gdef\csname num@plantvillage:A2_V_P_30:novel_5:sd\endcsname{6.9}
\expandafter\gdef\csname num@plantvillage:A2_V_P_30:novel_5:min\endcsname{77.3}
\expandafter\gdef\csname num@plantvillage:A2_V_P_30:novel_5:max\endcsname{94.8}
\expandafter\gdef\csname num@plantvillage:A2_V_P_30:novel_10\endcsname{88.8}
\expandafter\gdef\csname num@plantvillage:A2_V_P_30:novel_10:sd\endcsname{5.8}
\expandafter\gdef\csname num@plantvillage:A2_V_P_30:novel_10:min\endcsname{81.5}
\expandafter\gdef\csname num@plantvillage:A2_V_P_30:novel_10:max\endcsname{95.9}
\expandafter\gdef\csname num@plantvillage:A2_V_P_30:plantdoc_1\endcsname{35.5}
\expandafter\gdef\csname num@plantvillage:A2_V_P_30:plantdoc_1:sd\endcsname{1.3}
\expandafter\gdef\csname num@plantvillage:A2_V_P_30:plantdoc_1:min\endcsname{33.6}
\expandafter\gdef\csname num@plantvillage:A2_V_P_30:plantdoc_1:max\endcsname{36.7}
\expandafter\gdef\csname num@plantvillage:A2_V_P_30:plantdoc_5\endcsname{47.2}
\expandafter\gdef\csname num@plantvillage:A2_V_P_30:plantdoc_5:sd\endcsname{1.4}
\expandafter\gdef\csname num@plantvillage:A2_V_P_30:plantdoc_5:min\endcsname{44.9}
\expandafter\gdef\csname num@plantvillage:A2_V_P_30:plantdoc_5:max\endcsname{48.2}
\expandafter\gdef\csname num@plantvillage:A2_V_P_30:plantdoc_10\endcsname{51.4}
\expandafter\gdef\csname num@plantvillage:A2_V_P_30:plantdoc_10:sd\endcsname{1.5}
\expandafter\gdef\csname num@plantvillage:A2_V_P_30:plantdoc_10:min\endcsname{49.0}
\expandafter\gdef\csname num@plantvillage:A2_V_P_30:plantdoc_10:max\endcsname{52.5}
\expandafter\gdef\csname num@plantvillage:A2_V_P_30:pct\endcsname{30.0}
\expandafter\gdef\csname num@plantvillage:A2_V_P_30:params\endcsname{3.35}
\expandafter\gdef\csname num@plantvillage:A2_V_P_30:macs_128\endcsname{0.251}
\expandafter\gdef\csname num@plantvillage:A2_V_P_30:macs_224\endcsname{0.770}
\expandafter\gdef\csname num@plantvillage:A2_bn_gamma_P_30:novel_1\endcsname{71.7}
\expandafter\gdef\csname num@plantvillage:A2_bn_gamma_P_30:novel_1:sd\endcsname{12.9}
\expandafter\gdef\csname num@plantvillage:A2_bn_gamma_P_30:novel_1:min\endcsname{58.4}
\expandafter\gdef\csname num@plantvillage:A2_bn_gamma_P_30:novel_1:max\endcsname{84.8}
\expandafter\gdef\csname num@plantvillage:A2_bn_gamma_P_30:novel_5\endcsname{85.1}
\expandafter\gdef\csname num@plantvillage:A2_bn_gamma_P_30:novel_5:sd\endcsname{8.9}
\expandafter\gdef\csname num@plantvillage:A2_bn_gamma_P_30:novel_5:min\endcsname{74.7}
\expandafter\gdef\csname num@plantvillage:A2_bn_gamma_P_30:novel_5:max\endcsname{93.9}
\expandafter\gdef\csname num@plantvillage:A2_bn_gamma_P_30:novel_10\endcsname{87.5}
\expandafter\gdef\csname num@plantvillage:A2_bn_gamma_P_30:novel_10:sd\endcsname{7.3}
\expandafter\gdef\csname num@plantvillage:A2_bn_gamma_P_30:novel_10:min\endcsname{78.9}
\expandafter\gdef\csname num@plantvillage:A2_bn_gamma_P_30:novel_10:max\endcsname{95.0}
\expandafter\gdef\csname num@plantvillage:A2_bn_gamma_P_30:plantdoc_1\endcsname{33.2}
\expandafter\gdef\csname num@plantvillage:A2_bn_gamma_P_30:plantdoc_1:sd\endcsname{1.3}
\expandafter\gdef\csname num@plantvillage:A2_bn_gamma_P_30:plantdoc_1:min\endcsname{31.4}
\expandafter\gdef\csname num@plantvillage:A2_bn_gamma_P_30:plantdoc_1:max\endcsname{34.7}
\expandafter\gdef\csname num@plantvillage:A2_bn_gamma_P_30:plantdoc_5\endcsname{44.4}
\expandafter\gdef\csname num@plantvillage:A2_bn_gamma_P_30:plantdoc_5:sd\endcsname{1.7}
\expandafter\gdef\csname num@plantvillage:A2_bn_gamma_P_30:plantdoc_5:min\endcsname{41.6}
\expandafter\gdef\csname num@plantvillage:A2_bn_gamma_P_30:plantdoc_5:max\endcsname{45.5}
\expandafter\gdef\csname num@plantvillage:A2_bn_gamma_P_30:plantdoc_10\endcsname{48.5}
\expandafter\gdef\csname num@plantvillage:A2_bn_gamma_P_30:plantdoc_10:sd\endcsname{1.9}
\expandafter\gdef\csname num@plantvillage:A2_bn_gamma_P_30:plantdoc_10:min\endcsname{45.4}
\expandafter\gdef\csname num@plantvillage:A2_bn_gamma_P_30:plantdoc_10:max\endcsname{49.9}
\expandafter\gdef\csname num@plantvillage:A2_bn_gamma_P_30:pct\endcsname{30.0}
\expandafter\gdef\csname num@plantvillage:A2_bn_gamma_P_30:params\endcsname{3.35}
\expandafter\gdef\csname num@plantvillage:A2_bn_gamma_P_30:macs_128\endcsname{0.251}
\expandafter\gdef\csname num@plantvillage:A2_bn_gamma_P_30:macs_224\endcsname{0.770}
\expandafter\gdef\csname num@plantvillage:A2_dacis_P_30:novel_1\endcsname{72.7}
\expandafter\gdef\csname num@plantvillage:A2_dacis_P_30:novel_1:sd\endcsname{11.5}
\expandafter\gdef\csname num@plantvillage:A2_dacis_P_30:novel_1:min\endcsname{57.3}
\expandafter\gdef\csname num@plantvillage:A2_dacis_P_30:novel_1:max\endcsname{84.3}
\expandafter\gdef\csname num@plantvillage:A2_dacis_P_30:novel_5\endcsname{86.4}
\expandafter\gdef\csname num@plantvillage:A2_dacis_P_30:novel_5:sd\endcsname{7.7}
\expandafter\gdef\csname num@plantvillage:A2_dacis_P_30:novel_5:min\endcsname{75.7}
\expandafter\gdef\csname num@plantvillage:A2_dacis_P_30:novel_5:max\endcsname{93.9}
\expandafter\gdef\csname num@plantvillage:A2_dacis_P_30:novel_10\endcsname{88.6}
\expandafter\gdef\csname num@plantvillage:A2_dacis_P_30:novel_10:sd\endcsname{6.3}
\expandafter\gdef\csname num@plantvillage:A2_dacis_P_30:novel_10:min\endcsname{80.2}
\expandafter\gdef\csname num@plantvillage:A2_dacis_P_30:novel_10:max\endcsname{95.1}
\expandafter\gdef\csname num@plantvillage:A2_dacis_P_30:plantdoc_1\endcsname{35.2}
\expandafter\gdef\csname num@plantvillage:A2_dacis_P_30:plantdoc_1:sd\endcsname{1.2}
\expandafter\gdef\csname num@plantvillage:A2_dacis_P_30:plantdoc_1:min\endcsname{33.7}
\expandafter\gdef\csname num@plantvillage:A2_dacis_P_30:plantdoc_1:max\endcsname{37.0}
\expandafter\gdef\csname num@plantvillage:A2_dacis_P_30:plantdoc_5\endcsname{47.3}
\expandafter\gdef\csname num@plantvillage:A2_dacis_P_30:plantdoc_5:sd\endcsname{1.4}
\expandafter\gdef\csname num@plantvillage:A2_dacis_P_30:plantdoc_5:min\endcsname{45.9}
\expandafter\gdef\csname num@plantvillage:A2_dacis_P_30:plantdoc_5:max\endcsname{49.7}
\expandafter\gdef\csname num@plantvillage:A2_dacis_P_30:plantdoc_10\endcsname{51.4}
\expandafter\gdef\csname num@plantvillage:A2_dacis_P_30:plantdoc_10:sd\endcsname{1.4}
\expandafter\gdef\csname num@plantvillage:A2_dacis_P_30:plantdoc_10:min\endcsname{50.2}
\expandafter\gdef\csname num@plantvillage:A2_dacis_P_30:plantdoc_10:max\endcsname{53.6}
\expandafter\gdef\csname num@plantvillage:A2_dacis_P_30:pct\endcsname{30.0}
\expandafter\gdef\csname num@plantvillage:A2_dacis_P_30:params\endcsname{3.35}
\expandafter\gdef\csname num@plantvillage:A2_dacis_P_30:macs_128\endcsname{0.251}
\expandafter\gdef\csname num@plantvillage:A2_dacis_P_30:macs_224\endcsname{0.770}
\expandafter\gdef\csname num@plantvillage:A2_l1_P_30:novel_1\endcsname{72.9}
\expandafter\gdef\csname num@plantvillage:A2_l1_P_30:novel_1:sd\endcsname{11.5}
\expandafter\gdef\csname num@plantvillage:A2_l1_P_30:novel_1:min\endcsname{60.2}
\expandafter\gdef\csname num@plantvillage:A2_l1_P_30:novel_1:max\endcsname{85.5}
\expandafter\gdef\csname num@plantvillage:A2_l1_P_30:novel_5\endcsname{85.8}
\expandafter\gdef\csname num@plantvillage:A2_l1_P_30:novel_5:sd\endcsname{8.3}
\expandafter\gdef\csname num@plantvillage:A2_l1_P_30:novel_5:min\endcsname{76.4}
\expandafter\gdef\csname num@plantvillage:A2_l1_P_30:novel_5:max\endcsname{94.5}
\expandafter\gdef\csname num@plantvillage:A2_l1_P_30:novel_10\endcsname{88.2}
\expandafter\gdef\csname num@plantvillage:A2_l1_P_30:novel_10:sd\endcsname{7.2}
\expandafter\gdef\csname num@plantvillage:A2_l1_P_30:novel_10:min\endcsname{80.2}
\expandafter\gdef\csname num@plantvillage:A2_l1_P_30:novel_10:max\endcsname{95.8}
\expandafter\gdef\csname num@plantvillage:A2_l1_P_30:plantdoc_1\endcsname{37.0}
\expandafter\gdef\csname num@plantvillage:A2_l1_P_30:plantdoc_1:sd\endcsname{1.2}
\expandafter\gdef\csname num@plantvillage:A2_l1_P_30:plantdoc_1:min\endcsname{35.3}
\expandafter\gdef\csname num@plantvillage:A2_l1_P_30:plantdoc_1:max\endcsname{38.3}
\expandafter\gdef\csname num@plantvillage:A2_l1_P_30:plantdoc_5\endcsname{48.7}
\expandafter\gdef\csname num@plantvillage:A2_l1_P_30:plantdoc_5:sd\endcsname{0.9}
\expandafter\gdef\csname num@plantvillage:A2_l1_P_30:plantdoc_5:min\endcsname{47.5}
\expandafter\gdef\csname num@plantvillage:A2_l1_P_30:plantdoc_5:max\endcsname{50.0}
\expandafter\gdef\csname num@plantvillage:A2_l1_P_30:plantdoc_10\endcsname{52.7}
\expandafter\gdef\csname num@plantvillage:A2_l1_P_30:plantdoc_10:sd\endcsname{0.9}
\expandafter\gdef\csname num@plantvillage:A2_l1_P_30:plantdoc_10:min\endcsname{51.6}
\expandafter\gdef\csname num@plantvillage:A2_l1_P_30:plantdoc_10:max\endcsname{53.8}
\expandafter\gdef\csname num@plantvillage:A2_l1_P_30:pct\endcsname{30.0}
\expandafter\gdef\csname num@plantvillage:A2_l1_P_30:params\endcsname{3.35}
\expandafter\gdef\csname num@plantvillage:A2_l1_P_30:macs_128\endcsname{0.251}
\expandafter\gdef\csname num@plantvillage:A2_l1_P_30:macs_224\endcsname{0.770}
\expandafter\gdef\csname num@plantvillage:A2_random_P_30:novel_1\endcsname{71.7}
\expandafter\gdef\csname num@plantvillage:A2_random_P_30:novel_1:sd\endcsname{11.2}
\expandafter\gdef\csname num@plantvillage:A2_random_P_30:novel_1:min\endcsname{57.4}
\expandafter\gdef\csname num@plantvillage:A2_random_P_30:novel_1:max\endcsname{85.3}
\expandafter\gdef\csname num@plantvillage:A2_random_P_30:novel_5\endcsname{85.6}
\expandafter\gdef\csname num@plantvillage:A2_random_P_30:novel_5:sd\endcsname{8.1}
\expandafter\gdef\csname num@plantvillage:A2_random_P_30:novel_5:min\endcsname{74.9}
\expandafter\gdef\csname num@plantvillage:A2_random_P_30:novel_5:max\endcsname{94.9}
\expandafter\gdef\csname num@plantvillage:A2_random_P_30:novel_10\endcsname{88.1}
\expandafter\gdef\csname num@plantvillage:A2_random_P_30:novel_10:sd\endcsname{6.9}
\expandafter\gdef\csname num@plantvillage:A2_random_P_30:novel_10:min\endcsname{79.0}
\expandafter\gdef\csname num@plantvillage:A2_random_P_30:novel_10:max\endcsname{96.2}
\expandafter\gdef\csname num@plantvillage:A2_random_P_30:plantdoc_1\endcsname{35.0}
\expandafter\gdef\csname num@plantvillage:A2_random_P_30:plantdoc_1:sd\endcsname{1.3}
\expandafter\gdef\csname num@plantvillage:A2_random_P_30:plantdoc_1:min\endcsname{33.1}
\expandafter\gdef\csname num@plantvillage:A2_random_P_30:plantdoc_1:max\endcsname{36.1}
\expandafter\gdef\csname num@plantvillage:A2_random_P_30:plantdoc_5\endcsname{46.6}
\expandafter\gdef\csname num@plantvillage:A2_random_P_30:plantdoc_5:sd\endcsname{1.3}
\expandafter\gdef\csname num@plantvillage:A2_random_P_30:plantdoc_5:min\endcsname{44.7}
\expandafter\gdef\csname num@plantvillage:A2_random_P_30:plantdoc_5:max\endcsname{48.0}
\expandafter\gdef\csname num@plantvillage:A2_random_P_30:plantdoc_10\endcsname{50.9}
\expandafter\gdef\csname num@plantvillage:A2_random_P_30:plantdoc_10:sd\endcsname{1.3}
\expandafter\gdef\csname num@plantvillage:A2_random_P_30:plantdoc_10:min\endcsname{49.2}
\expandafter\gdef\csname num@plantvillage:A2_random_P_30:plantdoc_10:max\endcsname{52.4}
\expandafter\gdef\csname num@plantvillage:A2_random_P_30:pct\endcsname{30.0}
\expandafter\gdef\csname num@plantvillage:A2_random_P_30:params\endcsname{3.35}
\expandafter\gdef\csname num@plantvillage:A2_random_P_30:macs_128\endcsname{0.251}
\expandafter\gdef\csname num@plantvillage:A2_random_P_30:macs_224\endcsname{0.770}
\expandafter\gdef\csname num@plantvillage:A2_taylor_P_30:novel_1\endcsname{75.0}
\expandafter\gdef\csname num@plantvillage:A2_taylor_P_30:novel_1:sd\endcsname{10.3}
\expandafter\gdef\csname num@plantvillage:A2_taylor_P_30:novel_1:min\endcsname{61.1}
\expandafter\gdef\csname num@plantvillage:A2_taylor_P_30:novel_1:max\endcsname{85.9}
\expandafter\gdef\csname num@plantvillage:A2_taylor_P_30:novel_5\endcsname{88.3}
\expandafter\gdef\csname num@plantvillage:A2_taylor_P_30:novel_5:sd\endcsname{7.1}
\expandafter\gdef\csname num@plantvillage:A2_taylor_P_30:novel_5:min\endcsname{79.0}
\expandafter\gdef\csname num@plantvillage:A2_taylor_P_30:novel_5:max\endcsname{95.6}
\expandafter\gdef\csname num@plantvillage:A2_taylor_P_30:novel_10\endcsname{90.3}
\expandafter\gdef\csname num@plantvillage:A2_taylor_P_30:novel_10:sd\endcsname{5.8}
\expandafter\gdef\csname num@plantvillage:A2_taylor_P_30:novel_10:min\endcsname{82.9}
\expandafter\gdef\csname num@plantvillage:A2_taylor_P_30:novel_10:max\endcsname{96.4}
\expandafter\gdef\csname num@plantvillage:A2_taylor_P_30:plantdoc_1\endcsname{35.6}
\expandafter\gdef\csname num@plantvillage:A2_taylor_P_30:plantdoc_1:sd\endcsname{1.1}
\expandafter\gdef\csname num@plantvillage:A2_taylor_P_30:plantdoc_1:min\endcsname{33.7}
\expandafter\gdef\csname num@plantvillage:A2_taylor_P_30:plantdoc_1:max\endcsname{36.3}
\expandafter\gdef\csname num@plantvillage:A2_taylor_P_30:plantdoc_5\endcsname{47.9}
\expandafter\gdef\csname num@plantvillage:A2_taylor_P_30:plantdoc_5:sd\endcsname{1.2}
\expandafter\gdef\csname num@plantvillage:A2_taylor_P_30:plantdoc_5:min\endcsname{46.1}
\expandafter\gdef\csname num@plantvillage:A2_taylor_P_30:plantdoc_5:max\endcsname{49.2}
\expandafter\gdef\csname num@plantvillage:A2_taylor_P_30:plantdoc_10\endcsname{52.3}
\expandafter\gdef\csname num@plantvillage:A2_taylor_P_30:plantdoc_10:sd\endcsname{1.2}
\expandafter\gdef\csname num@plantvillage:A2_taylor_P_30:plantdoc_10:min\endcsname{50.7}
\expandafter\gdef\csname num@plantvillage:A2_taylor_P_30:plantdoc_10:max\endcsname{53.8}
\expandafter\gdef\csname num@plantvillage:A2_taylor_P_30:pct\endcsname{30.0}
\expandafter\gdef\csname num@plantvillage:A2_taylor_P_30:params\endcsname{3.35}
\expandafter\gdef\csname num@plantvillage:A2_taylor_P_30:macs_128\endcsname{0.251}
\expandafter\gdef\csname num@plantvillage:A2_taylor_P_30:macs_224\endcsname{0.770}
\expandafter\gdef\csname num@plantvillage:A_D_PMP_30:novel_1\endcsname{71.7}
\expandafter\gdef\csname num@plantvillage:A_D_PMP_30:novel_1:sd\endcsname{11.0}
\expandafter\gdef\csname num@plantvillage:A_D_PMP_30:novel_1:min\endcsname{55.9}
\expandafter\gdef\csname num@plantvillage:A_D_PMP_30:novel_1:max\endcsname{83.7}
\expandafter\gdef\csname num@plantvillage:A_D_PMP_30:novel_5\endcsname{85.3}
\expandafter\gdef\csname num@plantvillage:A_D_PMP_30:novel_5:sd\endcsname{8.0}
\expandafter\gdef\csname num@plantvillage:A_D_PMP_30:novel_5:min\endcsname{73.9}
\expandafter\gdef\csname num@plantvillage:A_D_PMP_30:novel_5:max\endcsname{93.8}
\expandafter\gdef\csname num@plantvillage:A_D_PMP_30:novel_10\endcsname{87.7}
\expandafter\gdef\csname num@plantvillage:A_D_PMP_30:novel_10:sd\endcsname{6.6}
\expandafter\gdef\csname num@plantvillage:A_D_PMP_30:novel_10:min\endcsname{78.6}
\expandafter\gdef\csname num@plantvillage:A_D_PMP_30:novel_10:max\endcsname{95.1}
\expandafter\gdef\csname num@plantvillage:A_D_PMP_30:plantdoc_1\endcsname{34.6}
\expandafter\gdef\csname num@plantvillage:A_D_PMP_30:plantdoc_1:sd\endcsname{0.9}
\expandafter\gdef\csname num@plantvillage:A_D_PMP_30:plantdoc_1:min\endcsname{33.7}
\expandafter\gdef\csname num@plantvillage:A_D_PMP_30:plantdoc_1:max\endcsname{35.9}
\expandafter\gdef\csname num@plantvillage:A_D_PMP_30:plantdoc_5\endcsname{46.4}
\expandafter\gdef\csname num@plantvillage:A_D_PMP_30:plantdoc_5:sd\endcsname{1.3}
\expandafter\gdef\csname num@plantvillage:A_D_PMP_30:plantdoc_5:min\endcsname{45.4}
\expandafter\gdef\csname num@plantvillage:A_D_PMP_30:plantdoc_5:max\endcsname{48.5}
\expandafter\gdef\csname num@plantvillage:A_D_PMP_30:plantdoc_10\endcsname{50.6}
\expandafter\gdef\csname num@plantvillage:A_D_PMP_30:plantdoc_10:sd\endcsname{1.5}
\expandafter\gdef\csname num@plantvillage:A_D_PMP_30:plantdoc_10:min\endcsname{49.1}
\expandafter\gdef\csname num@plantvillage:A_D_PMP_30:plantdoc_10:max\endcsname{53.0}
\expandafter\gdef\csname num@plantvillage:A_D_PMP_30:pct\endcsname{30.0}
\expandafter\gdef\csname num@plantvillage:A_D_PMP_30:params\endcsname{3.35}
\expandafter\gdef\csname num@plantvillage:A_D_PMP_30:macs_128\endcsname{0.251}
\expandafter\gdef\csname num@plantvillage:A_D_PMP_30:macs_224\endcsname{0.770}
\expandafter\gdef\csname num@plantvillage:A_GD_PMP_30:novel_1\endcsname{72.6}
\expandafter\gdef\csname num@plantvillage:A_GD_PMP_30:novel_1:sd\endcsname{10.5}
\expandafter\gdef\csname num@plantvillage:A_GD_PMP_30:novel_1:min\endcsname{58.0}
\expandafter\gdef\csname num@plantvillage:A_GD_PMP_30:novel_1:max\endcsname{83.5}
\expandafter\gdef\csname num@plantvillage:A_GD_PMP_30:novel_5\endcsname{86.5}
\expandafter\gdef\csname num@plantvillage:A_GD_PMP_30:novel_5:sd\endcsname{7.5}
\expandafter\gdef\csname num@plantvillage:A_GD_PMP_30:novel_5:min\endcsname{76.3}
\expandafter\gdef\csname num@plantvillage:A_GD_PMP_30:novel_5:max\endcsname{94.0}
\expandafter\gdef\csname num@plantvillage:A_GD_PMP_30:novel_10\endcsname{88.7}
\expandafter\gdef\csname num@plantvillage:A_GD_PMP_30:novel_10:sd\endcsname{6.3}
\expandafter\gdef\csname num@plantvillage:A_GD_PMP_30:novel_10:min\endcsname{80.3}
\expandafter\gdef\csname num@plantvillage:A_GD_PMP_30:novel_10:max\endcsname{95.3}
\expandafter\gdef\csname num@plantvillage:A_GD_PMP_30:plantdoc_1\endcsname{34.4}
\expandafter\gdef\csname num@plantvillage:A_GD_PMP_30:plantdoc_1:sd\endcsname{0.8}
\expandafter\gdef\csname num@plantvillage:A_GD_PMP_30:plantdoc_1:min\endcsname{32.9}
\expandafter\gdef\csname num@plantvillage:A_GD_PMP_30:plantdoc_1:max\endcsname{35.1}
\expandafter\gdef\csname num@plantvillage:A_GD_PMP_30:plantdoc_5\endcsname{46.5}
\expandafter\gdef\csname num@plantvillage:A_GD_PMP_30:plantdoc_5:sd\endcsname{0.9}
\expandafter\gdef\csname num@plantvillage:A_GD_PMP_30:plantdoc_5:min\endcsname{45.0}
\expandafter\gdef\csname num@plantvillage:A_GD_PMP_30:plantdoc_5:max\endcsname{47.2}
\expandafter\gdef\csname num@plantvillage:A_GD_PMP_30:plantdoc_10\endcsname{50.8}
\expandafter\gdef\csname num@plantvillage:A_GD_PMP_30:plantdoc_10:sd\endcsname{0.9}
\expandafter\gdef\csname num@plantvillage:A_GD_PMP_30:plantdoc_10:min\endcsname{49.5}
\expandafter\gdef\csname num@plantvillage:A_GD_PMP_30:plantdoc_10:max\endcsname{51.6}
\expandafter\gdef\csname num@plantvillage:A_GD_PMP_30:pct\endcsname{30.0}
\expandafter\gdef\csname num@plantvillage:A_GD_PMP_30:params\endcsname{3.35}
\expandafter\gdef\csname num@plantvillage:A_GD_PMP_30:macs_128\endcsname{0.251}
\expandafter\gdef\csname num@plantvillage:A_GD_PMP_30:macs_224\endcsname{0.770}
\expandafter\gdef\csname num@plantvillage:A_GV_PMP_30:novel_1\endcsname{73.0}
\expandafter\gdef\csname num@plantvillage:A_GV_PMP_30:novel_1:sd\endcsname{9.8}
\expandafter\gdef\csname num@plantvillage:A_GV_PMP_30:novel_1:min\endcsname{59.8}
\expandafter\gdef\csname num@plantvillage:A_GV_PMP_30:novel_1:max\endcsname{84.0}
\expandafter\gdef\csname num@plantvillage:A_GV_PMP_30:novel_5\endcsname{86.9}
\expandafter\gdef\csname num@plantvillage:A_GV_PMP_30:novel_5:sd\endcsname{7.0}
\expandafter\gdef\csname num@plantvillage:A_GV_PMP_30:novel_5:min\endcsname{77.7}
\expandafter\gdef\csname num@plantvillage:A_GV_PMP_30:novel_5:max\endcsname{94.4}
\expandafter\gdef\csname num@plantvillage:A_GV_PMP_30:novel_10\endcsname{89.2}
\expandafter\gdef\csname num@plantvillage:A_GV_PMP_30:novel_10:sd\endcsname{5.9}
\expandafter\gdef\csname num@plantvillage:A_GV_PMP_30:novel_10:min\endcsname{81.7}
\expandafter\gdef\csname num@plantvillage:A_GV_PMP_30:novel_10:max\endcsname{95.5}
\expandafter\gdef\csname num@plantvillage:A_GV_PMP_30:plantdoc_1\endcsname{35.1}
\expandafter\gdef\csname num@plantvillage:A_GV_PMP_30:plantdoc_1:sd\endcsname{1.2}
\expandafter\gdef\csname num@plantvillage:A_GV_PMP_30:plantdoc_1:min\endcsname{33.7}
\expandafter\gdef\csname num@plantvillage:A_GV_PMP_30:plantdoc_1:max\endcsname{36.5}
\expandafter\gdef\csname num@plantvillage:A_GV_PMP_30:plantdoc_5\endcsname{46.9}
\expandafter\gdef\csname num@plantvillage:A_GV_PMP_30:plantdoc_5:sd\endcsname{1.4}
\expandafter\gdef\csname num@plantvillage:A_GV_PMP_30:plantdoc_5:min\endcsname{45.2}
\expandafter\gdef\csname num@plantvillage:A_GV_PMP_30:plantdoc_5:max\endcsname{48.5}
\expandafter\gdef\csname num@plantvillage:A_GV_PMP_30:plantdoc_10\endcsname{51.3}
\expandafter\gdef\csname num@plantvillage:A_GV_PMP_30:plantdoc_10:sd\endcsname{1.3}
\expandafter\gdef\csname num@plantvillage:A_GV_PMP_30:plantdoc_10:min\endcsname{49.9}
\expandafter\gdef\csname num@plantvillage:A_GV_PMP_30:plantdoc_10:max\endcsname{53.0}
\expandafter\gdef\csname num@plantvillage:A_GV_PMP_30:pct\endcsname{30.0}
\expandafter\gdef\csname num@plantvillage:A_GV_PMP_30:params\endcsname{3.35}
\expandafter\gdef\csname num@plantvillage:A_GV_PMP_30:macs_128\endcsname{0.251}
\expandafter\gdef\csname num@plantvillage:A_GV_PMP_30:macs_224\endcsname{0.770}
\expandafter\gdef\csname num@plantvillage:A_G_PMP_30:novel_1\endcsname{72.8}
\expandafter\gdef\csname num@plantvillage:A_G_PMP_30:novel_1:sd\endcsname{11.4}
\expandafter\gdef\csname num@plantvillage:A_G_PMP_30:novel_1:min\endcsname{58.2}
\expandafter\gdef\csname num@plantvillage:A_G_PMP_30:novel_1:max\endcsname{84.9}
\expandafter\gdef\csname num@plantvillage:A_G_PMP_30:novel_5\endcsname{86.7}
\expandafter\gdef\csname num@plantvillage:A_G_PMP_30:novel_5:sd\endcsname{7.9}
\expandafter\gdef\csname num@plantvillage:A_G_PMP_30:novel_5:min\endcsname{76.3}
\expandafter\gdef\csname num@plantvillage:A_G_PMP_30:novel_5:max\endcsname{94.5}
\expandafter\gdef\csname num@plantvillage:A_G_PMP_30:novel_10\endcsname{88.9}
\expandafter\gdef\csname num@plantvillage:A_G_PMP_30:novel_10:sd\endcsname{6.5}
\expandafter\gdef\csname num@plantvillage:A_G_PMP_30:novel_10:min\endcsname{80.6}
\expandafter\gdef\csname num@plantvillage:A_G_PMP_30:novel_10:max\endcsname{95.3}
\expandafter\gdef\csname num@plantvillage:A_G_PMP_30:plantdoc_1\endcsname{34.5}
\expandafter\gdef\csname num@plantvillage:A_G_PMP_30:plantdoc_1:sd\endcsname{1.9}
\expandafter\gdef\csname num@plantvillage:A_G_PMP_30:plantdoc_1:min\endcsname{31.7}
\expandafter\gdef\csname num@plantvillage:A_G_PMP_30:plantdoc_1:max\endcsname{36.9}
\expandafter\gdef\csname num@plantvillage:A_G_PMP_30:plantdoc_5\endcsname{46.9}
\expandafter\gdef\csname num@plantvillage:A_G_PMP_30:plantdoc_5:sd\endcsname{2.5}
\expandafter\gdef\csname num@plantvillage:A_G_PMP_30:plantdoc_5:min\endcsname{43.2}
\expandafter\gdef\csname num@plantvillage:A_G_PMP_30:plantdoc_5:max\endcsname{50.3}
\expandafter\gdef\csname num@plantvillage:A_G_PMP_30:plantdoc_10\endcsname{51.3}
\expandafter\gdef\csname num@plantvillage:A_G_PMP_30:plantdoc_10:sd\endcsname{2.5}
\expandafter\gdef\csname num@plantvillage:A_G_PMP_30:plantdoc_10:min\endcsname{47.7}
\expandafter\gdef\csname num@plantvillage:A_G_PMP_30:plantdoc_10:max\endcsname{54.6}
\expandafter\gdef\csname num@plantvillage:A_G_PMP_30:pct\endcsname{30.0}
\expandafter\gdef\csname num@plantvillage:A_G_PMP_30:params\endcsname{3.35}
\expandafter\gdef\csname num@plantvillage:A_G_PMP_30:macs_128\endcsname{0.251}
\expandafter\gdef\csname num@plantvillage:A_G_PMP_30:macs_224\endcsname{0.770}
\expandafter\gdef\csname num@plantvillage:A_VD_PMP_30:novel_1\endcsname{72.7}
\expandafter\gdef\csname num@plantvillage:A_VD_PMP_30:novel_1:sd\endcsname{10.9}
\expandafter\gdef\csname num@plantvillage:A_VD_PMP_30:novel_1:min\endcsname{57.4}
\expandafter\gdef\csname num@plantvillage:A_VD_PMP_30:novel_1:max\endcsname{83.9}
\expandafter\gdef\csname num@plantvillage:A_VD_PMP_30:novel_5\endcsname{86.2}
\expandafter\gdef\csname num@plantvillage:A_VD_PMP_30:novel_5:sd\endcsname{7.6}
\expandafter\gdef\csname num@plantvillage:A_VD_PMP_30:novel_5:min\endcsname{75.4}
\expandafter\gdef\csname num@plantvillage:A_VD_PMP_30:novel_5:max\endcsname{94.0}
\expandafter\gdef\csname num@plantvillage:A_VD_PMP_30:novel_10\endcsname{88.4}
\expandafter\gdef\csname num@plantvillage:A_VD_PMP_30:novel_10:sd\endcsname{6.3}
\expandafter\gdef\csname num@plantvillage:A_VD_PMP_30:novel_10:min\endcsname{79.7}
\expandafter\gdef\csname num@plantvillage:A_VD_PMP_30:novel_10:max\endcsname{95.1}
\expandafter\gdef\csname num@plantvillage:A_VD_PMP_30:plantdoc_1\endcsname{34.9}
\expandafter\gdef\csname num@plantvillage:A_VD_PMP_30:plantdoc_1:sd\endcsname{1.2}
\expandafter\gdef\csname num@plantvillage:A_VD_PMP_30:plantdoc_1:min\endcsname{33.0}
\expandafter\gdef\csname num@plantvillage:A_VD_PMP_30:plantdoc_1:max\endcsname{36.1}
\expandafter\gdef\csname num@plantvillage:A_VD_PMP_30:plantdoc_5\endcsname{46.9}
\expandafter\gdef\csname num@plantvillage:A_VD_PMP_30:plantdoc_5:sd\endcsname{1.5}
\expandafter\gdef\csname num@plantvillage:A_VD_PMP_30:plantdoc_5:min\endcsname{44.7}
\expandafter\gdef\csname num@plantvillage:A_VD_PMP_30:plantdoc_5:max\endcsname{48.7}
\expandafter\gdef\csname num@plantvillage:A_VD_PMP_30:plantdoc_10\endcsname{51.2}
\expandafter\gdef\csname num@plantvillage:A_VD_PMP_30:plantdoc_10:sd\endcsname{1.6}
\expandafter\gdef\csname num@plantvillage:A_VD_PMP_30:plantdoc_10:min\endcsname{49.2}
\expandafter\gdef\csname num@plantvillage:A_VD_PMP_30:plantdoc_10:max\endcsname{53.1}
\expandafter\gdef\csname num@plantvillage:A_VD_PMP_30:pct\endcsname{30.0}
\expandafter\gdef\csname num@plantvillage:A_VD_PMP_30:params\endcsname{3.35}
\expandafter\gdef\csname num@plantvillage:A_VD_PMP_30:macs_128\endcsname{0.251}
\expandafter\gdef\csname num@plantvillage:A_VD_PMP_30:macs_224\endcsname{0.770}
\expandafter\gdef\csname num@plantvillage:A_V_PMP_30:novel_1\endcsname{72.5}
\expandafter\gdef\csname num@plantvillage:A_V_PMP_30:novel_1:sd\endcsname{10.3}
\expandafter\gdef\csname num@plantvillage:A_V_PMP_30:novel_1:min\endcsname{58.6}
\expandafter\gdef\csname num@plantvillage:A_V_PMP_30:novel_1:max\endcsname{83.9}
\expandafter\gdef\csname num@plantvillage:A_V_PMP_30:novel_5\endcsname{86.5}
\expandafter\gdef\csname num@plantvillage:A_V_PMP_30:novel_5:sd\endcsname{7.3}
\expandafter\gdef\csname num@plantvillage:A_V_PMP_30:novel_5:min\endcsname{76.9}
\expandafter\gdef\csname num@plantvillage:A_V_PMP_30:novel_5:max\endcsname{94.4}
\expandafter\gdef\csname num@plantvillage:A_V_PMP_30:novel_10\endcsname{88.8}
\expandafter\gdef\csname num@plantvillage:A_V_PMP_30:novel_10:sd\endcsname{6.1}
\expandafter\gdef\csname num@plantvillage:A_V_PMP_30:novel_10:min\endcsname{81.3}
\expandafter\gdef\csname num@plantvillage:A_V_PMP_30:novel_10:max\endcsname{95.7}
\expandafter\gdef\csname num@plantvillage:A_V_PMP_30:plantdoc_1\endcsname{35.4}
\expandafter\gdef\csname num@plantvillage:A_V_PMP_30:plantdoc_1:sd\endcsname{1.1}
\expandafter\gdef\csname num@plantvillage:A_V_PMP_30:plantdoc_1:min\endcsname{33.7}
\expandafter\gdef\csname num@plantvillage:A_V_PMP_30:plantdoc_1:max\endcsname{36.5}
\expandafter\gdef\csname num@plantvillage:A_V_PMP_30:plantdoc_5\endcsname{47.3}
\expandafter\gdef\csname num@plantvillage:A_V_PMP_30:plantdoc_5:sd\endcsname{1.3}
\expandafter\gdef\csname num@plantvillage:A_V_PMP_30:plantdoc_5:min\endcsname{45.3}
\expandafter\gdef\csname num@plantvillage:A_V_PMP_30:plantdoc_5:max\endcsname{48.5}
\expandafter\gdef\csname num@plantvillage:A_V_PMP_30:plantdoc_10\endcsname{51.5}
\expandafter\gdef\csname num@plantvillage:A_V_PMP_30:plantdoc_10:sd\endcsname{1.4}
\expandafter\gdef\csname num@plantvillage:A_V_PMP_30:plantdoc_10:min\endcsname{49.6}
\expandafter\gdef\csname num@plantvillage:A_V_PMP_30:plantdoc_10:max\endcsname{52.9}
\expandafter\gdef\csname num@plantvillage:A_V_PMP_30:pct\endcsname{30.0}
\expandafter\gdef\csname num@plantvillage:A_V_PMP_30:params\endcsname{3.35}
\expandafter\gdef\csname num@plantvillage:A_V_PMP_30:macs_128\endcsname{0.251}
\expandafter\gdef\csname num@plantvillage:A_V_PMP_30:macs_224\endcsname{0.770}
\expandafter\gdef\csname num@plantvillage:A_bn_gamma_PMP_30:novel_1\endcsname{71.6}
\expandafter\gdef\csname num@plantvillage:A_bn_gamma_PMP_30:novel_1:sd\endcsname{11.8}
\expandafter\gdef\csname num@plantvillage:A_bn_gamma_PMP_30:novel_1:min\endcsname{58.2}
\expandafter\gdef\csname num@plantvillage:A_bn_gamma_PMP_30:novel_1:max\endcsname{84.2}
\expandafter\gdef\csname num@plantvillage:A_bn_gamma_PMP_30:novel_5\endcsname{85.1}
\expandafter\gdef\csname num@plantvillage:A_bn_gamma_PMP_30:novel_5:sd\endcsname{8.7}
\expandafter\gdef\csname num@plantvillage:A_bn_gamma_PMP_30:novel_5:min\endcsname{75.4}
\expandafter\gdef\csname num@plantvillage:A_bn_gamma_PMP_30:novel_5:max\endcsname{93.7}
\expandafter\gdef\csname num@plantvillage:A_bn_gamma_PMP_30:novel_10\endcsname{87.5}
\expandafter\gdef\csname num@plantvillage:A_bn_gamma_PMP_30:novel_10:sd\endcsname{7.3}
\expandafter\gdef\csname num@plantvillage:A_bn_gamma_PMP_30:novel_10:min\endcsname{79.3}
\expandafter\gdef\csname num@plantvillage:A_bn_gamma_PMP_30:novel_10:max\endcsname{94.7}
\expandafter\gdef\csname num@plantvillage:A_bn_gamma_PMP_30:plantdoc_1\endcsname{33.9}
\expandafter\gdef\csname num@plantvillage:A_bn_gamma_PMP_30:plantdoc_1:sd\endcsname{0.8}
\expandafter\gdef\csname num@plantvillage:A_bn_gamma_PMP_30:plantdoc_1:min\endcsname{33.0}
\expandafter\gdef\csname num@plantvillage:A_bn_gamma_PMP_30:plantdoc_1:max\endcsname{34.8}
\expandafter\gdef\csname num@plantvillage:A_bn_gamma_PMP_30:plantdoc_5\endcsname{45.5}
\expandafter\gdef\csname num@plantvillage:A_bn_gamma_PMP_30:plantdoc_5:sd\endcsname{1.0}
\expandafter\gdef\csname num@plantvillage:A_bn_gamma_PMP_30:plantdoc_5:min\endcsname{44.3}
\expandafter\gdef\csname num@plantvillage:A_bn_gamma_PMP_30:plantdoc_5:max\endcsname{46.9}
\expandafter\gdef\csname num@plantvillage:A_bn_gamma_PMP_30:plantdoc_10\endcsname{49.8}
\expandafter\gdef\csname num@plantvillage:A_bn_gamma_PMP_30:plantdoc_10:sd\endcsname{1.2}
\expandafter\gdef\csname num@plantvillage:A_bn_gamma_PMP_30:plantdoc_10:min\endcsname{48.2}
\expandafter\gdef\csname num@plantvillage:A_bn_gamma_PMP_30:plantdoc_10:max\endcsname{51.2}
\expandafter\gdef\csname num@plantvillage:A_bn_gamma_PMP_30:pct\endcsname{30.0}
\expandafter\gdef\csname num@plantvillage:A_bn_gamma_PMP_30:params\endcsname{3.35}
\expandafter\gdef\csname num@plantvillage:A_bn_gamma_PMP_30:macs_128\endcsname{0.251}
\expandafter\gdef\csname num@plantvillage:A_bn_gamma_PMP_30:macs_224\endcsname{0.770}
\expandafter\gdef\csname num@plantvillage:A_dacis_PMP_30:novel_1\endcsname{72.7}
\expandafter\gdef\csname num@plantvillage:A_dacis_PMP_30:novel_1:sd\endcsname{10.8}
\expandafter\gdef\csname num@plantvillage:A_dacis_PMP_30:novel_1:min\endcsname{58.8}
\expandafter\gdef\csname num@plantvillage:A_dacis_PMP_30:novel_1:max\endcsname{83.7}
\expandafter\gdef\csname num@plantvillage:A_dacis_PMP_30:novel_5\endcsname{86.6}
\expandafter\gdef\csname num@plantvillage:A_dacis_PMP_30:novel_5:sd\endcsname{7.5}
\expandafter\gdef\csname num@plantvillage:A_dacis_PMP_30:novel_5:min\endcsname{76.7}
\expandafter\gdef\csname num@plantvillage:A_dacis_PMP_30:novel_5:max\endcsname{93.8}
\expandafter\gdef\csname num@plantvillage:A_dacis_PMP_30:novel_10\endcsname{88.9}
\expandafter\gdef\csname num@plantvillage:A_dacis_PMP_30:novel_10:sd\endcsname{6.3}
\expandafter\gdef\csname num@plantvillage:A_dacis_PMP_30:novel_10:min\endcsname{80.9}
\expandafter\gdef\csname num@plantvillage:A_dacis_PMP_30:novel_10:max\endcsname{95.1}
\expandafter\gdef\csname num@plantvillage:A_dacis_PMP_30:plantdoc_1\endcsname{35.2}
\expandafter\gdef\csname num@plantvillage:A_dacis_PMP_30:plantdoc_1:sd\endcsname{1.3}
\expandafter\gdef\csname num@plantvillage:A_dacis_PMP_30:plantdoc_1:min\endcsname{33.1}
\expandafter\gdef\csname num@plantvillage:A_dacis_PMP_30:plantdoc_1:max\endcsname{36.6}
\expandafter\gdef\csname num@plantvillage:A_dacis_PMP_30:plantdoc_5\endcsname{47.5}
\expandafter\gdef\csname num@plantvillage:A_dacis_PMP_30:plantdoc_5:sd\endcsname{1.6}
\expandafter\gdef\csname num@plantvillage:A_dacis_PMP_30:plantdoc_5:min\endcsname{45.2}
\expandafter\gdef\csname num@plantvillage:A_dacis_PMP_30:plantdoc_5:max\endcsname{49.5}
\expandafter\gdef\csname num@plantvillage:A_dacis_PMP_30:plantdoc_10\endcsname{51.8}
\expandafter\gdef\csname num@plantvillage:A_dacis_PMP_30:plantdoc_10:sd\endcsname{1.6}
\expandafter\gdef\csname num@plantvillage:A_dacis_PMP_30:plantdoc_10:min\endcsname{49.9}
\expandafter\gdef\csname num@plantvillage:A_dacis_PMP_30:plantdoc_10:max\endcsname{53.6}
\expandafter\gdef\csname num@plantvillage:A_dacis_PMP_30:pct\endcsname{30.0}
\expandafter\gdef\csname num@plantvillage:A_dacis_PMP_30:params\endcsname{3.35}
\expandafter\gdef\csname num@plantvillage:A_dacis_PMP_30:macs_128\endcsname{0.251}
\expandafter\gdef\csname num@plantvillage:A_dacis_PMP_30:macs_224\endcsname{0.770}
\expandafter\gdef\csname num@plantvillage:A_dacis_PMP_30:cpu_1thread_b1_ms\endcsname{6.6}
\expandafter\gdef\csname num@plantvillage:A_dacis_PMP_30:cpu_all_threads_b1_ms\endcsname{3.3}
\expandafter\gdef\csname num@plantvillage:A_dacis_PMP_30:gpu_b1_ms\endcsname{1.0}
\expandafter\gdef\csname num@plantvillage:A_dacis_PMP_30:gpu_ips\endcsname{40648}
\expandafter\gdef\csname num@plantvillage:A_l1_PMP_30:novel_1\endcsname{73.5}
\expandafter\gdef\csname num@plantvillage:A_l1_PMP_30:novel_1:sd\endcsname{11.1}
\expandafter\gdef\csname num@plantvillage:A_l1_PMP_30:novel_1:min\endcsname{59.8}
\expandafter\gdef\csname num@plantvillage:A_l1_PMP_30:novel_1:max\endcsname{86.5}
\expandafter\gdef\csname num@plantvillage:A_l1_PMP_30:novel_5\endcsname{86.4}
\expandafter\gdef\csname num@plantvillage:A_l1_PMP_30:novel_5:sd\endcsname{7.7}
\expandafter\gdef\csname num@plantvillage:A_l1_PMP_30:novel_5:min\endcsname{77.1}
\expandafter\gdef\csname num@plantvillage:A_l1_PMP_30:novel_5:max\endcsname{94.8}
\expandafter\gdef\csname num@plantvillage:A_l1_PMP_30:novel_10\endcsname{88.7}
\expandafter\gdef\csname num@plantvillage:A_l1_PMP_30:novel_10:sd\endcsname{6.4}
\expandafter\gdef\csname num@plantvillage:A_l1_PMP_30:novel_10:min\endcsname{81.4}
\expandafter\gdef\csname num@plantvillage:A_l1_PMP_30:novel_10:max\endcsname{95.9}
\expandafter\gdef\csname num@plantvillage:A_l1_PMP_30:plantdoc_1\endcsname{36.4}
\expandafter\gdef\csname num@plantvillage:A_l1_PMP_30:plantdoc_1:sd\endcsname{1.2}
\expandafter\gdef\csname num@plantvillage:A_l1_PMP_30:plantdoc_1:min\endcsname{34.9}
\expandafter\gdef\csname num@plantvillage:A_l1_PMP_30:plantdoc_1:max\endcsname{38.2}
\expandafter\gdef\csname num@plantvillage:A_l1_PMP_30:plantdoc_5\endcsname{48.1}
\expandafter\gdef\csname num@plantvillage:A_l1_PMP_30:plantdoc_5:sd\endcsname{0.8}
\expandafter\gdef\csname num@plantvillage:A_l1_PMP_30:plantdoc_5:min\endcsname{47.2}
\expandafter\gdef\csname num@plantvillage:A_l1_PMP_30:plantdoc_5:max\endcsname{49.3}
\expandafter\gdef\csname num@plantvillage:A_l1_PMP_30:plantdoc_10\endcsname{52.4}
\expandafter\gdef\csname num@plantvillage:A_l1_PMP_30:plantdoc_10:sd\endcsname{0.7}
\expandafter\gdef\csname num@plantvillage:A_l1_PMP_30:plantdoc_10:min\endcsname{51.6}
\expandafter\gdef\csname num@plantvillage:A_l1_PMP_30:plantdoc_10:max\endcsname{53.4}
\expandafter\gdef\csname num@plantvillage:A_l1_PMP_30:pct\endcsname{30.0}
\expandafter\gdef\csname num@plantvillage:A_l1_PMP_30:params\endcsname{3.35}
\expandafter\gdef\csname num@plantvillage:A_l1_PMP_30:macs_128\endcsname{0.251}
\expandafter\gdef\csname num@plantvillage:A_l1_PMP_30:macs_224\endcsname{0.770}
\expandafter\gdef\csname num@plantvillage:A_l1_PMP_30:cpu_1thread_b1_ms\endcsname{6.6}
\expandafter\gdef\csname num@plantvillage:A_l1_PMP_30:cpu_all_threads_b1_ms\endcsname{3.2}
\expandafter\gdef\csname num@plantvillage:A_l1_PMP_30:gpu_b1_ms\endcsname{1.0}
\expandafter\gdef\csname num@plantvillage:A_l1_PMP_30:gpu_ips\endcsname{40555}
\expandafter\gdef\csname num@plantvillage:A_random_PMP_30:novel_1\endcsname{73.6}
\expandafter\gdef\csname num@plantvillage:A_random_PMP_30:novel_1:sd\endcsname{10.9}
\expandafter\gdef\csname num@plantvillage:A_random_PMP_30:novel_1:min\endcsname{61.1}
\expandafter\gdef\csname num@plantvillage:A_random_PMP_30:novel_1:max\endcsname{85.1}
\expandafter\gdef\csname num@plantvillage:A_random_PMP_30:novel_5\endcsname{86.8}
\expandafter\gdef\csname num@plantvillage:A_random_PMP_30:novel_5:sd\endcsname{7.7}
\expandafter\gdef\csname num@plantvillage:A_random_PMP_30:novel_5:min\endcsname{78.6}
\expandafter\gdef\csname num@plantvillage:A_random_PMP_30:novel_5:max\endcsname{95.0}
\expandafter\gdef\csname num@plantvillage:A_random_PMP_30:novel_10\endcsname{88.9}
\expandafter\gdef\csname num@plantvillage:A_random_PMP_30:novel_10:sd\endcsname{6.5}
\expandafter\gdef\csname num@plantvillage:A_random_PMP_30:novel_10:min\endcsname{81.8}
\expandafter\gdef\csname num@plantvillage:A_random_PMP_30:novel_10:max\endcsname{96.2}
\expandafter\gdef\csname num@plantvillage:A_random_PMP_30:plantdoc_1\endcsname{34.9}
\expandafter\gdef\csname num@plantvillage:A_random_PMP_30:plantdoc_1:sd\endcsname{1.6}
\expandafter\gdef\csname num@plantvillage:A_random_PMP_30:plantdoc_1:min\endcsname{32.3}
\expandafter\gdef\csname num@plantvillage:A_random_PMP_30:plantdoc_1:max\endcsname{36.5}
\expandafter\gdef\csname num@plantvillage:A_random_PMP_30:plantdoc_5\endcsname{46.7}
\expandafter\gdef\csname num@plantvillage:A_random_PMP_30:plantdoc_5:sd\endcsname{1.7}
\expandafter\gdef\csname num@plantvillage:A_random_PMP_30:plantdoc_5:min\endcsname{43.8}
\expandafter\gdef\csname num@plantvillage:A_random_PMP_30:plantdoc_5:max\endcsname{47.8}
\expandafter\gdef\csname num@plantvillage:A_random_PMP_30:plantdoc_10\endcsname{50.9}
\expandafter\gdef\csname num@plantvillage:A_random_PMP_30:plantdoc_10:sd\endcsname{1.5}
\expandafter\gdef\csname num@plantvillage:A_random_PMP_30:plantdoc_10:min\endcsname{48.3}
\expandafter\gdef\csname num@plantvillage:A_random_PMP_30:plantdoc_10:max\endcsname{52.4}
\expandafter\gdef\csname num@plantvillage:A_random_PMP_30:pct\endcsname{30.0}
\expandafter\gdef\csname num@plantvillage:A_random_PMP_30:params\endcsname{3.35}
\expandafter\gdef\csname num@plantvillage:A_random_PMP_30:macs_128\endcsname{0.251}
\expandafter\gdef\csname num@plantvillage:A_random_PMP_30:macs_224\endcsname{0.770}
\expandafter\gdef\csname num@plantvillage:A_taylor_PMP_30:novel_1\endcsname{72.9}
\expandafter\gdef\csname num@plantvillage:A_taylor_PMP_30:novel_1:sd\endcsname{11.0}
\expandafter\gdef\csname num@plantvillage:A_taylor_PMP_30:novel_1:min\endcsname{59.8}
\expandafter\gdef\csname num@plantvillage:A_taylor_PMP_30:novel_1:max\endcsname{86.0}
\expandafter\gdef\csname num@plantvillage:A_taylor_PMP_30:novel_5\endcsname{87.1}
\expandafter\gdef\csname num@plantvillage:A_taylor_PMP_30:novel_5:sd\endcsname{7.1}
\expandafter\gdef\csname num@plantvillage:A_taylor_PMP_30:novel_5:min\endcsname{78.4}
\expandafter\gdef\csname num@plantvillage:A_taylor_PMP_30:novel_5:max\endcsname{95.3}
\expandafter\gdef\csname num@plantvillage:A_taylor_PMP_30:novel_10\endcsname{89.4}
\expandafter\gdef\csname num@plantvillage:A_taylor_PMP_30:novel_10:sd\endcsname{5.8}
\expandafter\gdef\csname num@plantvillage:A_taylor_PMP_30:novel_10:min\endcsname{82.3}
\expandafter\gdef\csname num@plantvillage:A_taylor_PMP_30:novel_10:max\endcsname{96.1}
\expandafter\gdef\csname num@plantvillage:A_taylor_PMP_30:plantdoc_1\endcsname{36.2}
\expandafter\gdef\csname num@plantvillage:A_taylor_PMP_30:plantdoc_1:sd\endcsname{1.2}
\expandafter\gdef\csname num@plantvillage:A_taylor_PMP_30:plantdoc_1:min\endcsname{34.1}
\expandafter\gdef\csname num@plantvillage:A_taylor_PMP_30:plantdoc_1:max\endcsname{37.1}
\expandafter\gdef\csname num@plantvillage:A_taylor_PMP_30:plantdoc_5\endcsname{48.5}
\expandafter\gdef\csname num@plantvillage:A_taylor_PMP_30:plantdoc_5:sd\endcsname{1.0}
\expandafter\gdef\csname num@plantvillage:A_taylor_PMP_30:plantdoc_5:min\endcsname{47.0}
\expandafter\gdef\csname num@plantvillage:A_taylor_PMP_30:plantdoc_5:max\endcsname{49.4}
\expandafter\gdef\csname num@plantvillage:A_taylor_PMP_30:plantdoc_10\endcsname{52.9}
\expandafter\gdef\csname num@plantvillage:A_taylor_PMP_30:plantdoc_10:sd\endcsname{0.7}
\expandafter\gdef\csname num@plantvillage:A_taylor_PMP_30:plantdoc_10:min\endcsname{51.8}
\expandafter\gdef\csname num@plantvillage:A_taylor_PMP_30:plantdoc_10:max\endcsname{53.7}
\expandafter\gdef\csname num@plantvillage:A_taylor_PMP_30:pct\endcsname{30.0}
\expandafter\gdef\csname num@plantvillage:A_taylor_PMP_30:params\endcsname{3.35}
\expandafter\gdef\csname num@plantvillage:A_taylor_PMP_30:macs_128\endcsname{0.251}
\expandafter\gdef\csname num@plantvillage:A_taylor_PMP_30:macs_224\endcsname{0.770}
\expandafter\gdef\csname num@plantvillage:B_dacis_MP_30:novel_1\endcsname{73.3}
\expandafter\gdef\csname num@plantvillage:B_dacis_MP_30:novel_1:sd\endcsname{10.6}
\expandafter\gdef\csname num@plantvillage:B_dacis_MP_30:novel_1:min\endcsname{59.6}
\expandafter\gdef\csname num@plantvillage:B_dacis_MP_30:novel_1:max\endcsname{84.0}
\expandafter\gdef\csname num@plantvillage:B_dacis_MP_30:novel_5\endcsname{86.7}
\expandafter\gdef\csname num@plantvillage:B_dacis_MP_30:novel_5:sd\endcsname{6.9}
\expandafter\gdef\csname num@plantvillage:B_dacis_MP_30:novel_5:min\endcsname{78.0}
\expandafter\gdef\csname num@plantvillage:B_dacis_MP_30:novel_5:max\endcsname{93.8}
\expandafter\gdef\csname num@plantvillage:B_dacis_MP_30:novel_10\endcsname{88.8}
\expandafter\gdef\csname num@plantvillage:B_dacis_MP_30:novel_10:sd\endcsname{5.6}
\expandafter\gdef\csname num@plantvillage:B_dacis_MP_30:novel_10:min\endcsname{82.4}
\expandafter\gdef\csname num@plantvillage:B_dacis_MP_30:novel_10:max\endcsname{95.1}
\expandafter\gdef\csname num@plantvillage:B_dacis_MP_30:plantdoc_1\endcsname{34.9}
\expandafter\gdef\csname num@plantvillage:B_dacis_MP_30:plantdoc_1:sd\endcsname{0.7}
\expandafter\gdef\csname num@plantvillage:B_dacis_MP_30:plantdoc_1:min\endcsname{33.9}
\expandafter\gdef\csname num@plantvillage:B_dacis_MP_30:plantdoc_1:max\endcsname{35.5}
\expandafter\gdef\csname num@plantvillage:B_dacis_MP_30:plantdoc_5\endcsname{46.7}
\expandafter\gdef\csname num@plantvillage:B_dacis_MP_30:plantdoc_5:sd\endcsname{0.9}
\expandafter\gdef\csname num@plantvillage:B_dacis_MP_30:plantdoc_5:min\endcsname{46.0}
\expandafter\gdef\csname num@plantvillage:B_dacis_MP_30:plantdoc_5:max\endcsname{48.2}
\expandafter\gdef\csname num@plantvillage:B_dacis_MP_30:plantdoc_10\endcsname{51.0}
\expandafter\gdef\csname num@plantvillage:B_dacis_MP_30:plantdoc_10:sd\endcsname{1.0}
\expandafter\gdef\csname num@plantvillage:B_dacis_MP_30:plantdoc_10:min\endcsname{50.4}
\expandafter\gdef\csname num@plantvillage:B_dacis_MP_30:plantdoc_10:max\endcsname{52.7}
\expandafter\gdef\csname num@plantvillage:B_dacis_MP_30:pct\endcsname{30.0}
\expandafter\gdef\csname num@plantvillage:B_dacis_MP_30:params\endcsname{3.35}
\expandafter\gdef\csname num@plantvillage:B_dacis_MP_30:macs_128\endcsname{0.251}
\expandafter\gdef\csname num@plantvillage:B_dacis_MP_30:macs_224\endcsname{0.770}
\expandafter\gdef\csname num@plantvillage:B_dacis_PGP_30:novel_1\endcsname{73.0}
\expandafter\gdef\csname num@plantvillage:B_dacis_PGP_30:novel_1:sd\endcsname{10.8}
\expandafter\gdef\csname num@plantvillage:B_dacis_PGP_30:novel_1:min\endcsname{58.7}
\expandafter\gdef\csname num@plantvillage:B_dacis_PGP_30:novel_1:max\endcsname{83.5}
\expandafter\gdef\csname num@plantvillage:B_dacis_PGP_30:novel_5\endcsname{86.8}
\expandafter\gdef\csname num@plantvillage:B_dacis_PGP_30:novel_5:sd\endcsname{7.4}
\expandafter\gdef\csname num@plantvillage:B_dacis_PGP_30:novel_5:min\endcsname{76.8}
\expandafter\gdef\csname num@plantvillage:B_dacis_PGP_30:novel_5:max\endcsname{93.6}
\expandafter\gdef\csname num@plantvillage:B_dacis_PGP_30:novel_10\endcsname{89.1}
\expandafter\gdef\csname num@plantvillage:B_dacis_PGP_30:novel_10:sd\endcsname{6.0}
\expandafter\gdef\csname num@plantvillage:B_dacis_PGP_30:novel_10:min\endcsname{81.3}
\expandafter\gdef\csname num@plantvillage:B_dacis_PGP_30:novel_10:max\endcsname{94.7}
\expandafter\gdef\csname num@plantvillage:B_dacis_PGP_30:plantdoc_1\endcsname{34.9}
\expandafter\gdef\csname num@plantvillage:B_dacis_PGP_30:plantdoc_1:sd\endcsname{1.3}
\expandafter\gdef\csname num@plantvillage:B_dacis_PGP_30:plantdoc_1:min\endcsname{33.0}
\expandafter\gdef\csname num@plantvillage:B_dacis_PGP_30:plantdoc_1:max\endcsname{36.1}
\expandafter\gdef\csname num@plantvillage:B_dacis_PGP_30:plantdoc_5\endcsname{47.1}
\expandafter\gdef\csname num@plantvillage:B_dacis_PGP_30:plantdoc_5:sd\endcsname{1.4}
\expandafter\gdef\csname num@plantvillage:B_dacis_PGP_30:plantdoc_5:min\endcsname{45.5}
\expandafter\gdef\csname num@plantvillage:B_dacis_PGP_30:plantdoc_5:max\endcsname{49.0}
\expandafter\gdef\csname num@plantvillage:B_dacis_PGP_30:plantdoc_10\endcsname{51.4}
\expandafter\gdef\csname num@plantvillage:B_dacis_PGP_30:plantdoc_10:sd\endcsname{1.5}
\expandafter\gdef\csname num@plantvillage:B_dacis_PGP_30:plantdoc_10:min\endcsname{49.7}
\expandafter\gdef\csname num@plantvillage:B_dacis_PGP_30:plantdoc_10:max\endcsname{53.4}
\expandafter\gdef\csname num@plantvillage:B_dacis_PGP_30:pct\endcsname{30.0}
\expandafter\gdef\csname num@plantvillage:B_dacis_PGP_30:params\endcsname{3.35}
\expandafter\gdef\csname num@plantvillage:B_dacis_PGP_30:macs_128\endcsname{0.251}
\expandafter\gdef\csname num@plantvillage:B_dacis_PGP_30:macs_224\endcsname{0.770}
\expandafter\gdef\csname num@plantvillage:B_dacis_PM2x_30:novel_1\endcsname{72.5}
\expandafter\gdef\csname num@plantvillage:B_dacis_PM2x_30:novel_1:sd\endcsname{11.2}
\expandafter\gdef\csname num@plantvillage:B_dacis_PM2x_30:novel_1:min\endcsname{57.8}
\expandafter\gdef\csname num@plantvillage:B_dacis_PM2x_30:novel_1:max\endcsname{84.4}
\expandafter\gdef\csname num@plantvillage:B_dacis_PM2x_30:novel_5\endcsname{85.6}
\expandafter\gdef\csname num@plantvillage:B_dacis_PM2x_30:novel_5:sd\endcsname{7.3}
\expandafter\gdef\csname num@plantvillage:B_dacis_PM2x_30:novel_5:min\endcsname{76.7}
\expandafter\gdef\csname num@plantvillage:B_dacis_PM2x_30:novel_5:max\endcsname{93.2}
\expandafter\gdef\csname num@plantvillage:B_dacis_PM2x_30:novel_10\endcsname{87.9}
\expandafter\gdef\csname num@plantvillage:B_dacis_PM2x_30:novel_10:sd\endcsname{6.1}
\expandafter\gdef\csname num@plantvillage:B_dacis_PM2x_30:novel_10:min\endcsname{81.1}
\expandafter\gdef\csname num@plantvillage:B_dacis_PM2x_30:novel_10:max\endcsname{94.4}
\expandafter\gdef\csname num@plantvillage:B_dacis_PM2x_30:plantdoc_1\endcsname{34.8}
\expandafter\gdef\csname num@plantvillage:B_dacis_PM2x_30:plantdoc_1:sd\endcsname{1.0}
\expandafter\gdef\csname num@plantvillage:B_dacis_PM2x_30:plantdoc_1:min\endcsname{33.7}
\expandafter\gdef\csname num@plantvillage:B_dacis_PM2x_30:plantdoc_1:max\endcsname{36.3}
\expandafter\gdef\csname num@plantvillage:B_dacis_PM2x_30:plantdoc_5\endcsname{46.4}
\expandafter\gdef\csname num@plantvillage:B_dacis_PM2x_30:plantdoc_5:sd\endcsname{1.2}
\expandafter\gdef\csname num@plantvillage:B_dacis_PM2x_30:plantdoc_5:min\endcsname{45.4}
\expandafter\gdef\csname num@plantvillage:B_dacis_PM2x_30:plantdoc_5:max\endcsname{48.5}
\expandafter\gdef\csname num@plantvillage:B_dacis_PM2x_30:plantdoc_10\endcsname{50.5}
\expandafter\gdef\csname num@plantvillage:B_dacis_PM2x_30:plantdoc_10:sd\endcsname{1.1}
\expandafter\gdef\csname num@plantvillage:B_dacis_PM2x_30:plantdoc_10:min\endcsname{49.7}
\expandafter\gdef\csname num@plantvillage:B_dacis_PM2x_30:plantdoc_10:max\endcsname{52.4}
\expandafter\gdef\csname num@plantvillage:B_dacis_PM2x_30:pct\endcsname{30.0}
\expandafter\gdef\csname num@plantvillage:B_dacis_PM2x_30:params\endcsname{3.35}
\expandafter\gdef\csname num@plantvillage:B_dacis_PM2x_30:macs_128\endcsname{0.251}
\expandafter\gdef\csname num@plantvillage:B_dacis_PM2x_30:macs_224\endcsname{0.770}
\expandafter\gdef\csname num@plantvillage:B_dacis_PMPM_30:novel_1\endcsname{73.2}
\expandafter\gdef\csname num@plantvillage:B_dacis_PMPM_30:novel_1:sd\endcsname{10.7}
\expandafter\gdef\csname num@plantvillage:B_dacis_PMPM_30:novel_1:min\endcsname{59.6}
\expandafter\gdef\csname num@plantvillage:B_dacis_PMPM_30:novel_1:max\endcsname{84.8}
\expandafter\gdef\csname num@plantvillage:B_dacis_PMPM_30:novel_5\endcsname{86.5}
\expandafter\gdef\csname num@plantvillage:B_dacis_PMPM_30:novel_5:sd\endcsname{7.4}
\expandafter\gdef\csname num@plantvillage:B_dacis_PMPM_30:novel_5:min\endcsname{77.7}
\expandafter\gdef\csname num@plantvillage:B_dacis_PMPM_30:novel_5:max\endcsname{94.4}
\expandafter\gdef\csname num@plantvillage:B_dacis_PMPM_30:novel_10\endcsname{88.6}
\expandafter\gdef\csname num@plantvillage:B_dacis_PMPM_30:novel_10:sd\endcsname{6.1}
\expandafter\gdef\csname num@plantvillage:B_dacis_PMPM_30:novel_10:min\endcsname{82.0}
\expandafter\gdef\csname num@plantvillage:B_dacis_PMPM_30:novel_10:max\endcsname{95.2}
\expandafter\gdef\csname num@plantvillage:B_dacis_PMPM_30:plantdoc_1\endcsname{35.6}
\expandafter\gdef\csname num@plantvillage:B_dacis_PMPM_30:plantdoc_1:sd\endcsname{1.3}
\expandafter\gdef\csname num@plantvillage:B_dacis_PMPM_30:plantdoc_1:min\endcsname{33.7}
\expandafter\gdef\csname num@plantvillage:B_dacis_PMPM_30:plantdoc_1:max\endcsname{37.4}
\expandafter\gdef\csname num@plantvillage:B_dacis_PMPM_30:plantdoc_5\endcsname{47.6}
\expandafter\gdef\csname num@plantvillage:B_dacis_PMPM_30:plantdoc_5:sd\endcsname{1.6}
\expandafter\gdef\csname num@plantvillage:B_dacis_PMPM_30:plantdoc_5:min\endcsname{45.5}
\expandafter\gdef\csname num@plantvillage:B_dacis_PMPM_30:plantdoc_5:max\endcsname{49.9}
\expandafter\gdef\csname num@plantvillage:B_dacis_PMPM_30:plantdoc_10\endcsname{51.9}
\expandafter\gdef\csname num@plantvillage:B_dacis_PMPM_30:plantdoc_10:sd\endcsname{1.5}
\expandafter\gdef\csname num@plantvillage:B_dacis_PMPM_30:plantdoc_10:min\endcsname{49.9}
\expandafter\gdef\csname num@plantvillage:B_dacis_PMPM_30:plantdoc_10:max\endcsname{54.1}
\expandafter\gdef\csname num@plantvillage:B_dacis_PMPM_30:pct\endcsname{30.0}
\expandafter\gdef\csname num@plantvillage:B_dacis_PMPM_30:params\endcsname{3.35}
\expandafter\gdef\csname num@plantvillage:B_dacis_PMPM_30:macs_128\endcsname{0.251}
\expandafter\gdef\csname num@plantvillage:B_dacis_PMPM_30:macs_224\endcsname{0.770}
\expandafter\gdef\csname num@plantvillage:B_dacis_PM_30:novel_1\endcsname{72.5}
\expandafter\gdef\csname num@plantvillage:B_dacis_PM_30:novel_1:sd\endcsname{11.3}
\expandafter\gdef\csname num@plantvillage:B_dacis_PM_30:novel_1:min\endcsname{58.6}
\expandafter\gdef\csname num@plantvillage:B_dacis_PM_30:novel_1:max\endcsname{85.2}
\expandafter\gdef\csname num@plantvillage:B_dacis_PM_30:novel_5\endcsname{85.8}
\expandafter\gdef\csname num@plantvillage:B_dacis_PM_30:novel_5:sd\endcsname{7.6}
\expandafter\gdef\csname num@plantvillage:B_dacis_PM_30:novel_5:min\endcsname{76.8}
\expandafter\gdef\csname num@plantvillage:B_dacis_PM_30:novel_5:max\endcsname{94.0}
\expandafter\gdef\csname num@plantvillage:B_dacis_PM_30:novel_10\endcsname{87.9}
\expandafter\gdef\csname num@plantvillage:B_dacis_PM_30:novel_10:sd\endcsname{6.4}
\expandafter\gdef\csname num@plantvillage:B_dacis_PM_30:novel_10:min\endcsname{80.9}
\expandafter\gdef\csname num@plantvillage:B_dacis_PM_30:novel_10:max\endcsname{95.0}
\expandafter\gdef\csname num@plantvillage:B_dacis_PM_30:plantdoc_1\endcsname{34.8}
\expandafter\gdef\csname num@plantvillage:B_dacis_PM_30:plantdoc_1:sd\endcsname{1.1}
\expandafter\gdef\csname num@plantvillage:B_dacis_PM_30:plantdoc_1:min\endcsname{33.7}
\expandafter\gdef\csname num@plantvillage:B_dacis_PM_30:plantdoc_1:max\endcsname{36.4}
\expandafter\gdef\csname num@plantvillage:B_dacis_PM_30:plantdoc_5\endcsname{46.5}
\expandafter\gdef\csname num@plantvillage:B_dacis_PM_30:plantdoc_5:sd\endcsname{1.3}
\expandafter\gdef\csname num@plantvillage:B_dacis_PM_30:plantdoc_5:min\endcsname{45.4}
\expandafter\gdef\csname num@plantvillage:B_dacis_PM_30:plantdoc_5:max\endcsname{48.4}
\expandafter\gdef\csname num@plantvillage:B_dacis_PM_30:plantdoc_10\endcsname{50.7}
\expandafter\gdef\csname num@plantvillage:B_dacis_PM_30:plantdoc_10:sd\endcsname{1.5}
\expandafter\gdef\csname num@plantvillage:B_dacis_PM_30:plantdoc_10:min\endcsname{49.1}
\expandafter\gdef\csname num@plantvillage:B_dacis_PM_30:plantdoc_10:max\endcsname{52.5}
\expandafter\gdef\csname num@plantvillage:B_dacis_PM_30:pct\endcsname{30.0}
\expandafter\gdef\csname num@plantvillage:B_dacis_PM_30:params\endcsname{3.35}
\expandafter\gdef\csname num@plantvillage:B_dacis_PM_30:macs_128\endcsname{0.251}
\expandafter\gdef\csname num@plantvillage:B_dacis_PM_30:macs_224\endcsname{0.770}
\expandafter\gdef\csname num@plantvillage:B_l1_PM_30:novel_1\endcsname{72.9}
\expandafter\gdef\csname num@plantvillage:B_l1_PM_30:novel_1:sd\endcsname{10.6}
\expandafter\gdef\csname num@plantvillage:B_l1_PM_30:novel_1:min\endcsname{60.4}
\expandafter\gdef\csname num@plantvillage:B_l1_PM_30:novel_1:max\endcsname{84.7}
\expandafter\gdef\csname num@plantvillage:B_l1_PM_30:novel_5\endcsname{85.6}
\expandafter\gdef\csname num@plantvillage:B_l1_PM_30:novel_5:sd\endcsname{7.8}
\expandafter\gdef\csname num@plantvillage:B_l1_PM_30:novel_5:min\endcsname{76.5}
\expandafter\gdef\csname num@plantvillage:B_l1_PM_30:novel_5:max\endcsname{94.4}
\expandafter\gdef\csname num@plantvillage:B_l1_PM_30:novel_10\endcsname{88.0}
\expandafter\gdef\csname num@plantvillage:B_l1_PM_30:novel_10:sd\endcsname{6.6}
\expandafter\gdef\csname num@plantvillage:B_l1_PM_30:novel_10:min\endcsname{80.4}
\expandafter\gdef\csname num@plantvillage:B_l1_PM_30:novel_10:max\endcsname{95.7}
\expandafter\gdef\csname num@plantvillage:B_l1_PM_30:plantdoc_1\endcsname{36.6}
\expandafter\gdef\csname num@plantvillage:B_l1_PM_30:plantdoc_1:sd\endcsname{1.3}
\expandafter\gdef\csname num@plantvillage:B_l1_PM_30:plantdoc_1:min\endcsname{34.7}
\expandafter\gdef\csname num@plantvillage:B_l1_PM_30:plantdoc_1:max\endcsname{38.0}
\expandafter\gdef\csname num@plantvillage:B_l1_PM_30:plantdoc_5\endcsname{48.3}
\expandafter\gdef\csname num@plantvillage:B_l1_PM_30:plantdoc_5:sd\endcsname{1.1}
\expandafter\gdef\csname num@plantvillage:B_l1_PM_30:plantdoc_5:min\endcsname{46.7}
\expandafter\gdef\csname num@plantvillage:B_l1_PM_30:plantdoc_5:max\endcsname{49.8}
\expandafter\gdef\csname num@plantvillage:B_l1_PM_30:plantdoc_10\endcsname{52.4}
\expandafter\gdef\csname num@plantvillage:B_l1_PM_30:plantdoc_10:sd\endcsname{1.1}
\expandafter\gdef\csname num@plantvillage:B_l1_PM_30:plantdoc_10:min\endcsname{50.9}
\expandafter\gdef\csname num@plantvillage:B_l1_PM_30:plantdoc_10:max\endcsname{53.8}
\expandafter\gdef\csname num@plantvillage:B_l1_PM_30:pct\endcsname{30.0}
\expandafter\gdef\csname num@plantvillage:B_l1_PM_30:params\endcsname{3.35}
\expandafter\gdef\csname num@plantvillage:B_l1_PM_30:macs_128\endcsname{0.251}
\expandafter\gdef\csname num@plantvillage:B_l1_PM_30:macs_224\endcsname{0.770}
\expandafter\gdef\csname num@plantvillage:C_dacis_PMP_10:novel_1\endcsname{63.7}
\expandafter\gdef\csname num@plantvillage:C_dacis_PMP_10:novel_1:sd\endcsname{10.9}
\expandafter\gdef\csname num@plantvillage:C_dacis_PMP_10:novel_1:min\endcsname{47.1}
\expandafter\gdef\csname num@plantvillage:C_dacis_PMP_10:novel_1:max\endcsname{73.0}
\expandafter\gdef\csname num@plantvillage:C_dacis_PMP_10:novel_5\endcsname{78.2}
\expandafter\gdef\csname num@plantvillage:C_dacis_PMP_10:novel_5:sd\endcsname{10.2}
\expandafter\gdef\csname num@plantvillage:C_dacis_PMP_10:novel_5:min\endcsname{63.0}
\expandafter\gdef\csname num@plantvillage:C_dacis_PMP_10:novel_5:max\endcsname{87.0}
\expandafter\gdef\csname num@plantvillage:C_dacis_PMP_10:novel_10\endcsname{81.0}
\expandafter\gdef\csname num@plantvillage:C_dacis_PMP_10:novel_10:sd\endcsname{9.3}
\expandafter\gdef\csname num@plantvillage:C_dacis_PMP_10:novel_10:min\endcsname{67.5}
\expandafter\gdef\csname num@plantvillage:C_dacis_PMP_10:novel_10:max\endcsname{89.1}
\expandafter\gdef\csname num@plantvillage:C_dacis_PMP_10:plantdoc_1\endcsname{31.0}
\expandafter\gdef\csname num@plantvillage:C_dacis_PMP_10:plantdoc_1:sd\endcsname{0.7}
\expandafter\gdef\csname num@plantvillage:C_dacis_PMP_10:plantdoc_1:min\endcsname{30.3}
\expandafter\gdef\csname num@plantvillage:C_dacis_PMP_10:plantdoc_1:max\endcsname{32.1}
\expandafter\gdef\csname num@plantvillage:C_dacis_PMP_10:plantdoc_5\endcsname{40.5}
\expandafter\gdef\csname num@plantvillage:C_dacis_PMP_10:plantdoc_5:sd\endcsname{1.1}
\expandafter\gdef\csname num@plantvillage:C_dacis_PMP_10:plantdoc_5:min\endcsname{39.6}
\expandafter\gdef\csname num@plantvillage:C_dacis_PMP_10:plantdoc_5:max\endcsname{42.3}
\expandafter\gdef\csname num@plantvillage:C_dacis_PMP_10:plantdoc_10\endcsname{44.2}
\expandafter\gdef\csname num@plantvillage:C_dacis_PMP_10:plantdoc_10:sd\endcsname{1.1}
\expandafter\gdef\csname num@plantvillage:C_dacis_PMP_10:plantdoc_10:min\endcsname{43.4}
\expandafter\gdef\csname num@plantvillage:C_dacis_PMP_10:plantdoc_10:max\endcsname{46.0}
\expandafter\gdef\csname num@plantvillage:C_dacis_PMP_10:pct\endcsname{10.0}
\expandafter\gdef\csname num@plantvillage:C_dacis_PMP_10:params\endcsname{1.11}
\expandafter\gdef\csname num@plantvillage:C_dacis_PMP_10:macs_128\endcsname{0.105}
\expandafter\gdef\csname num@plantvillage:C_dacis_PMP_10:macs_224\endcsname{0.322}
\expandafter\gdef\csname num@plantvillage:C_dacis_PMP_10:cpu_1thread_b1_ms\endcsname{4.6}
\expandafter\gdef\csname num@plantvillage:C_dacis_PMP_10:cpu_all_threads_b1_ms\endcsname{2.6}
\expandafter\gdef\csname num@plantvillage:C_dacis_PMP_10:gpu_b1_ms\endcsname{1.0}
\expandafter\gdef\csname num@plantvillage:C_dacis_PMP_10:gpu_ips\endcsname{45487}
\expandafter\gdef\csname num@plantvillage:C_dacis_PMP_20:novel_1\endcsname{69.5}
\expandafter\gdef\csname num@plantvillage:C_dacis_PMP_20:novel_1:sd\endcsname{10.8}
\expandafter\gdef\csname num@plantvillage:C_dacis_PMP_20:novel_1:min\endcsname{55.9}
\expandafter\gdef\csname num@plantvillage:C_dacis_PMP_20:novel_1:max\endcsname{81.2}
\expandafter\gdef\csname num@plantvillage:C_dacis_PMP_20:novel_5\endcsname{83.8}
\expandafter\gdef\csname num@plantvillage:C_dacis_PMP_20:novel_5:sd\endcsname{8.2}
\expandafter\gdef\csname num@plantvillage:C_dacis_PMP_20:novel_5:min\endcsname{72.9}
\expandafter\gdef\csname num@plantvillage:C_dacis_PMP_20:novel_5:max\endcsname{92.4}
\expandafter\gdef\csname num@plantvillage:C_dacis_PMP_20:novel_10\endcsname{86.3}
\expandafter\gdef\csname num@plantvillage:C_dacis_PMP_20:novel_10:sd\endcsname{7.1}
\expandafter\gdef\csname num@plantvillage:C_dacis_PMP_20:novel_10:min\endcsname{76.9}
\expandafter\gdef\csname num@plantvillage:C_dacis_PMP_20:novel_10:max\endcsname{93.5}
\expandafter\gdef\csname num@plantvillage:C_dacis_PMP_20:plantdoc_1\endcsname{33.0}
\expandafter\gdef\csname num@plantvillage:C_dacis_PMP_20:plantdoc_1:sd\endcsname{1.5}
\expandafter\gdef\csname num@plantvillage:C_dacis_PMP_20:plantdoc_1:min\endcsname{31.0}
\expandafter\gdef\csname num@plantvillage:C_dacis_PMP_20:plantdoc_1:max\endcsname{35.0}
\expandafter\gdef\csname num@plantvillage:C_dacis_PMP_20:plantdoc_5\endcsname{44.1}
\expandafter\gdef\csname num@plantvillage:C_dacis_PMP_20:plantdoc_5:sd\endcsname{2.0}
\expandafter\gdef\csname num@plantvillage:C_dacis_PMP_20:plantdoc_5:min\endcsname{41.5}
\expandafter\gdef\csname num@plantvillage:C_dacis_PMP_20:plantdoc_5:max\endcsname{47.0}
\expandafter\gdef\csname num@plantvillage:C_dacis_PMP_20:plantdoc_10\endcsname{48.0}
\expandafter\gdef\csname num@plantvillage:C_dacis_PMP_20:plantdoc_10:sd\endcsname{2.0}
\expandafter\gdef\csname num@plantvillage:C_dacis_PMP_20:plantdoc_10:min\endcsname{45.4}
\expandafter\gdef\csname num@plantvillage:C_dacis_PMP_20:plantdoc_10:max\endcsname{51.0}
\expandafter\gdef\csname num@plantvillage:C_dacis_PMP_20:pct\endcsname{20.0}
\expandafter\gdef\csname num@plantvillage:C_dacis_PMP_20:params\endcsname{2.24}
\expandafter\gdef\csname num@plantvillage:C_dacis_PMP_20:macs_128\endcsname{0.178}
\expandafter\gdef\csname num@plantvillage:C_dacis_PMP_20:macs_224\endcsname{0.544}
\expandafter\gdef\csname num@plantvillage:C_dacis_PMP_20:cpu_1thread_b1_ms\endcsname{5.7}
\expandafter\gdef\csname num@plantvillage:C_dacis_PMP_20:cpu_all_threads_b1_ms\endcsname{2.9}
\expandafter\gdef\csname num@plantvillage:C_dacis_PMP_20:gpu_b1_ms\endcsname{1.0}
\expandafter\gdef\csname num@plantvillage:C_dacis_PMP_20:gpu_ips\endcsname{41197}
\expandafter\gdef\csname num@plantvillage:C_dacis_PMP_50:novel_1\endcsname{75.5}
\expandafter\gdef\csname num@plantvillage:C_dacis_PMP_50:novel_1:sd\endcsname{10.6}
\expandafter\gdef\csname num@plantvillage:C_dacis_PMP_50:novel_1:min\endcsname{60.0}
\expandafter\gdef\csname num@plantvillage:C_dacis_PMP_50:novel_1:max\endcsname{86.0}
\expandafter\gdef\csname num@plantvillage:C_dacis_PMP_50:novel_5\endcsname{88.6}
\expandafter\gdef\csname num@plantvillage:C_dacis_PMP_50:novel_5:sd\endcsname{7.2}
\expandafter\gdef\csname num@plantvillage:C_dacis_PMP_50:novel_5:min\endcsname{78.5}
\expandafter\gdef\csname num@plantvillage:C_dacis_PMP_50:novel_5:max\endcsname{95.9}
\expandafter\gdef\csname num@plantvillage:C_dacis_PMP_50:novel_10\endcsname{90.6}
\expandafter\gdef\csname num@plantvillage:C_dacis_PMP_50:novel_10:sd\endcsname{5.9}
\expandafter\gdef\csname num@plantvillage:C_dacis_PMP_50:novel_10:min\endcsname{82.7}
\expandafter\gdef\csname num@plantvillage:C_dacis_PMP_50:novel_10:max\endcsname{96.8}
\expandafter\gdef\csname num@plantvillage:C_dacis_PMP_50:plantdoc_1\endcsname{38.1}
\expandafter\gdef\csname num@plantvillage:C_dacis_PMP_50:plantdoc_1:sd\endcsname{1.3}
\expandafter\gdef\csname num@plantvillage:C_dacis_PMP_50:plantdoc_1:min\endcsname{36.8}
\expandafter\gdef\csname num@plantvillage:C_dacis_PMP_50:plantdoc_1:max\endcsname{39.4}
\expandafter\gdef\csname num@plantvillage:C_dacis_PMP_50:plantdoc_5\endcsname{51.6}
\expandafter\gdef\csname num@plantvillage:C_dacis_PMP_50:plantdoc_5:sd\endcsname{1.4}
\expandafter\gdef\csname num@plantvillage:C_dacis_PMP_50:plantdoc_5:min\endcsname{50.3}
\expandafter\gdef\csname num@plantvillage:C_dacis_PMP_50:plantdoc_5:max\endcsname{53.5}
\expandafter\gdef\csname num@plantvillage:C_dacis_PMP_50:plantdoc_10\endcsname{56.1}
\expandafter\gdef\csname num@plantvillage:C_dacis_PMP_50:plantdoc_10:sd\endcsname{1.6}
\expandafter\gdef\csname num@plantvillage:C_dacis_PMP_50:plantdoc_10:min\endcsname{54.2}
\expandafter\gdef\csname num@plantvillage:C_dacis_PMP_50:plantdoc_10:max\endcsname{58.0}
\expandafter\gdef\csname num@plantvillage:C_dacis_PMP_50:pct\endcsname{50.0}
\expandafter\gdef\csname num@plantvillage:C_dacis_PMP_50:params\endcsname{5.59}
\expandafter\gdef\csname num@plantvillage:C_dacis_PMP_50:macs_128\endcsname{0.398}
\expandafter\gdef\csname num@plantvillage:C_dacis_PMP_50:macs_224\endcsname{1.218}
\expandafter\gdef\csname num@plantvillage:C_dacis_PMP_50:cpu_1thread_b1_ms\endcsname{8.5}
\expandafter\gdef\csname num@plantvillage:C_dacis_PMP_50:cpu_all_threads_b1_ms\endcsname{3.7}
\expandafter\gdef\csname num@plantvillage:C_dacis_PMP_50:gpu_b1_ms\endcsname{1.0}
\expandafter\gdef\csname num@plantvillage:C_dacis_PMP_50:gpu_ips\endcsname{37452}
\expandafter\gdef\csname num@plantvillage:C_l1_PMP_10:novel_1\endcsname{65.8}
\expandafter\gdef\csname num@plantvillage:C_l1_PMP_10:novel_1:sd\endcsname{12.3}
\expandafter\gdef\csname num@plantvillage:C_l1_PMP_10:novel_1:min\endcsname{50.6}
\expandafter\gdef\csname num@plantvillage:C_l1_PMP_10:novel_1:max\endcsname{81.5}
\expandafter\gdef\csname num@plantvillage:C_l1_PMP_10:novel_5\endcsname{79.3}
\expandafter\gdef\csname num@plantvillage:C_l1_PMP_10:novel_5:sd\endcsname{9.9}
\expandafter\gdef\csname num@plantvillage:C_l1_PMP_10:novel_5:min\endcsname{67.1}
\expandafter\gdef\csname num@plantvillage:C_l1_PMP_10:novel_5:max\endcsname{91.4}
\expandafter\gdef\csname num@plantvillage:C_l1_PMP_10:novel_10\endcsname{82.1}
\expandafter\gdef\csname num@plantvillage:C_l1_PMP_10:novel_10:sd\endcsname{8.7}
\expandafter\gdef\csname num@plantvillage:C_l1_PMP_10:novel_10:min\endcsname{70.9}
\expandafter\gdef\csname num@plantvillage:C_l1_PMP_10:novel_10:max\endcsname{92.6}
\expandafter\gdef\csname num@plantvillage:C_l1_PMP_10:plantdoc_1\endcsname{30.1}
\expandafter\gdef\csname num@plantvillage:C_l1_PMP_10:plantdoc_1:sd\endcsname{0.7}
\expandafter\gdef\csname num@plantvillage:C_l1_PMP_10:plantdoc_1:min\endcsname{29.0}
\expandafter\gdef\csname num@plantvillage:C_l1_PMP_10:plantdoc_1:max\endcsname{30.8}
\expandafter\gdef\csname num@plantvillage:C_l1_PMP_10:plantdoc_5\endcsname{39.1}
\expandafter\gdef\csname num@plantvillage:C_l1_PMP_10:plantdoc_5:sd\endcsname{1.1}
\expandafter\gdef\csname num@plantvillage:C_l1_PMP_10:plantdoc_5:min\endcsname{37.2}
\expandafter\gdef\csname num@plantvillage:C_l1_PMP_10:plantdoc_5:max\endcsname{39.9}
\expandafter\gdef\csname num@plantvillage:C_l1_PMP_10:plantdoc_10\endcsname{42.7}
\expandafter\gdef\csname num@plantvillage:C_l1_PMP_10:plantdoc_10:sd\endcsname{1.2}
\expandafter\gdef\csname num@plantvillage:C_l1_PMP_10:plantdoc_10:min\endcsname{40.6}
\expandafter\gdef\csname num@plantvillage:C_l1_PMP_10:plantdoc_10:max\endcsname{43.7}
\expandafter\gdef\csname num@plantvillage:C_l1_PMP_10:pct\endcsname{10.0}
\expandafter\gdef\csname num@plantvillage:C_l1_PMP_10:params\endcsname{1.11}
\expandafter\gdef\csname num@plantvillage:C_l1_PMP_10:macs_128\endcsname{0.105}
\expandafter\gdef\csname num@plantvillage:C_l1_PMP_10:macs_224\endcsname{0.322}
\expandafter\gdef\csname num@plantvillage:C_l1_PMP_10:cpu_1thread_b1_ms\endcsname{4.7}
\expandafter\gdef\csname num@plantvillage:C_l1_PMP_10:cpu_all_threads_b1_ms\endcsname{2.5}
\expandafter\gdef\csname num@plantvillage:C_l1_PMP_10:gpu_b1_ms\endcsname{1.0}
\expandafter\gdef\csname num@plantvillage:C_l1_PMP_10:gpu_ips\endcsname{45493}
\expandafter\gdef\csname num@plantvillage:C_l1_PMP_20:novel_1\endcsname{71.1}
\expandafter\gdef\csname num@plantvillage:C_l1_PMP_20:novel_1:sd\endcsname{11.3}
\expandafter\gdef\csname num@plantvillage:C_l1_PMP_20:novel_1:min\endcsname{57.9}
\expandafter\gdef\csname num@plantvillage:C_l1_PMP_20:novel_1:max\endcsname{84.3}
\expandafter\gdef\csname num@plantvillage:C_l1_PMP_20:novel_5\endcsname{84.2}
\expandafter\gdef\csname num@plantvillage:C_l1_PMP_20:novel_5:sd\endcsname{8.3}
\expandafter\gdef\csname num@plantvillage:C_l1_PMP_20:novel_5:min\endcsname{74.9}
\expandafter\gdef\csname num@plantvillage:C_l1_PMP_20:novel_5:max\endcsname{93.4}
\expandafter\gdef\csname num@plantvillage:C_l1_PMP_20:novel_10\endcsname{86.7}
\expandafter\gdef\csname num@plantvillage:C_l1_PMP_20:novel_10:sd\endcsname{7.2}
\expandafter\gdef\csname num@plantvillage:C_l1_PMP_20:novel_10:min\endcsname{78.6}
\expandafter\gdef\csname num@plantvillage:C_l1_PMP_20:novel_10:max\endcsname{94.6}
\expandafter\gdef\csname num@plantvillage:C_l1_PMP_20:plantdoc_1\endcsname{34.2}
\expandafter\gdef\csname num@plantvillage:C_l1_PMP_20:plantdoc_1:sd\endcsname{1.1}
\expandafter\gdef\csname num@plantvillage:C_l1_PMP_20:plantdoc_1:min\endcsname{32.4}
\expandafter\gdef\csname num@plantvillage:C_l1_PMP_20:plantdoc_1:max\endcsname{35.5}
\expandafter\gdef\csname num@plantvillage:C_l1_PMP_20:plantdoc_5\endcsname{45.3}
\expandafter\gdef\csname num@plantvillage:C_l1_PMP_20:plantdoc_5:sd\endcsname{1.0}
\expandafter\gdef\csname num@plantvillage:C_l1_PMP_20:plantdoc_5:min\endcsname{43.9}
\expandafter\gdef\csname num@plantvillage:C_l1_PMP_20:plantdoc_5:max\endcsname{46.9}
\expandafter\gdef\csname num@plantvillage:C_l1_PMP_20:plantdoc_10\endcsname{49.4}
\expandafter\gdef\csname num@plantvillage:C_l1_PMP_20:plantdoc_10:sd\endcsname{1.0}
\expandafter\gdef\csname num@plantvillage:C_l1_PMP_20:plantdoc_10:min\endcsname{48.1}
\expandafter\gdef\csname num@plantvillage:C_l1_PMP_20:plantdoc_10:max\endcsname{50.9}
\expandafter\gdef\csname num@plantvillage:C_l1_PMP_20:pct\endcsname{20.0}
\expandafter\gdef\csname num@plantvillage:C_l1_PMP_20:params\endcsname{2.24}
\expandafter\gdef\csname num@plantvillage:C_l1_PMP_20:macs_128\endcsname{0.178}
\expandafter\gdef\csname num@plantvillage:C_l1_PMP_20:macs_224\endcsname{0.544}
\expandafter\gdef\csname num@plantvillage:C_l1_PMP_20:cpu_1thread_b1_ms\endcsname{5.7}
\expandafter\gdef\csname num@plantvillage:C_l1_PMP_20:cpu_all_threads_b1_ms\endcsname{3.1}
\expandafter\gdef\csname num@plantvillage:C_l1_PMP_20:gpu_b1_ms\endcsname{1.0}
\expandafter\gdef\csname num@plantvillage:C_l1_PMP_20:gpu_ips\endcsname{41199}
\expandafter\gdef\csname num@plantvillage:C_l1_PMP_50:novel_1\endcsname{75.9}
\expandafter\gdef\csname num@plantvillage:C_l1_PMP_50:novel_1:sd\endcsname{11.2}
\expandafter\gdef\csname num@plantvillage:C_l1_PMP_50:novel_1:min\endcsname{62.4}
\expandafter\gdef\csname num@plantvillage:C_l1_PMP_50:novel_1:max\endcsname{88.4}
\expandafter\gdef\csname num@plantvillage:C_l1_PMP_50:novel_5\endcsname{89.0}
\expandafter\gdef\csname num@plantvillage:C_l1_PMP_50:novel_5:sd\endcsname{7.1}
\expandafter\gdef\csname num@plantvillage:C_l1_PMP_50:novel_5:min\endcsname{80.6}
\expandafter\gdef\csname num@plantvillage:C_l1_PMP_50:novel_5:max\endcsname{96.6}
\expandafter\gdef\csname num@plantvillage:C_l1_PMP_50:novel_10\endcsname{91.0}
\expandafter\gdef\csname num@plantvillage:C_l1_PMP_50:novel_10:sd\endcsname{5.9}
\expandafter\gdef\csname num@plantvillage:C_l1_PMP_50:novel_10:min\endcsname{84.2}
\expandafter\gdef\csname num@plantvillage:C_l1_PMP_50:novel_10:max\endcsname{97.4}
\expandafter\gdef\csname num@plantvillage:C_l1_PMP_50:plantdoc_1\endcsname{38.4}
\expandafter\gdef\csname num@plantvillage:C_l1_PMP_50:plantdoc_1:sd\endcsname{1.2}
\expandafter\gdef\csname num@plantvillage:C_l1_PMP_50:plantdoc_1:min\endcsname{36.7}
\expandafter\gdef\csname num@plantvillage:C_l1_PMP_50:plantdoc_1:max\endcsname{39.5}
\expandafter\gdef\csname num@plantvillage:C_l1_PMP_50:plantdoc_5\endcsname{51.3}
\expandafter\gdef\csname num@plantvillage:C_l1_PMP_50:plantdoc_5:sd\endcsname{1.1}
\expandafter\gdef\csname num@plantvillage:C_l1_PMP_50:plantdoc_5:min\endcsname{49.8}
\expandafter\gdef\csname num@plantvillage:C_l1_PMP_50:plantdoc_5:max\endcsname{52.3}
\expandafter\gdef\csname num@plantvillage:C_l1_PMP_50:plantdoc_10\endcsname{55.6}
\expandafter\gdef\csname num@plantvillage:C_l1_PMP_50:plantdoc_10:sd\endcsname{1.0}
\expandafter\gdef\csname num@plantvillage:C_l1_PMP_50:plantdoc_10:min\endcsname{54.1}
\expandafter\gdef\csname num@plantvillage:C_l1_PMP_50:plantdoc_10:max\endcsname{56.7}
\expandafter\gdef\csname num@plantvillage:C_l1_PMP_50:pct\endcsname{50.0}
\expandafter\gdef\csname num@plantvillage:C_l1_PMP_50:params\endcsname{5.59}
\expandafter\gdef\csname num@plantvillage:C_l1_PMP_50:macs_128\endcsname{0.398}
\expandafter\gdef\csname num@plantvillage:C_l1_PMP_50:macs_224\endcsname{1.218}
\expandafter\gdef\csname num@plantvillage:C_l1_PMP_50:cpu_1thread_b1_ms\endcsname{8.7}
\expandafter\gdef\csname num@plantvillage:C_l1_PMP_50:cpu_all_threads_b1_ms\endcsname{3.5}
\expandafter\gdef\csname num@plantvillage:C_l1_PMP_50:gpu_b1_ms\endcsname{1.0}
\expandafter\gdef\csname num@plantvillage:C_l1_PMP_50:gpu_ips\endcsname{37597}
\expandafter\gdef\csname num@plantvillage:E_dacis_PMP_30_minmax:novel_1\endcsname{72.2}
\expandafter\gdef\csname num@plantvillage:E_dacis_PMP_30_minmax:novel_1:sd\endcsname{11.0}
\expandafter\gdef\csname num@plantvillage:E_dacis_PMP_30_minmax:novel_1:min\endcsname{57.4}
\expandafter\gdef\csname num@plantvillage:E_dacis_PMP_30_minmax:novel_1:max\endcsname{83.9}
\expandafter\gdef\csname num@plantvillage:E_dacis_PMP_30_minmax:novel_5\endcsname{86.1}
\expandafter\gdef\csname num@plantvillage:E_dacis_PMP_30_minmax:novel_5:sd\endcsname{7.9}
\expandafter\gdef\csname num@plantvillage:E_dacis_PMP_30_minmax:novel_5:min\endcsname{75.4}
\expandafter\gdef\csname num@plantvillage:E_dacis_PMP_30_minmax:novel_5:max\endcsname{94.1}
\expandafter\gdef\csname num@plantvillage:E_dacis_PMP_30_minmax:novel_10\endcsname{88.4}
\expandafter\gdef\csname num@plantvillage:E_dacis_PMP_30_minmax:novel_10:sd\endcsname{6.6}
\expandafter\gdef\csname num@plantvillage:E_dacis_PMP_30_minmax:novel_10:min\endcsname{79.8}
\expandafter\gdef\csname num@plantvillage:E_dacis_PMP_30_minmax:novel_10:max\endcsname{95.2}
\expandafter\gdef\csname num@plantvillage:E_dacis_PMP_30_minmax:plantdoc_1\endcsname{34.8}
\expandafter\gdef\csname num@plantvillage:E_dacis_PMP_30_minmax:plantdoc_1:sd\endcsname{1.3}
\expandafter\gdef\csname num@plantvillage:E_dacis_PMP_30_minmax:plantdoc_1:min\endcsname{33.0}
\expandafter\gdef\csname num@plantvillage:E_dacis_PMP_30_minmax:plantdoc_1:max\endcsname{36.5}
\expandafter\gdef\csname num@plantvillage:E_dacis_PMP_30_minmax:plantdoc_5\endcsname{46.8}
\expandafter\gdef\csname num@plantvillage:E_dacis_PMP_30_minmax:plantdoc_5:sd\endcsname{1.5}
\expandafter\gdef\csname num@plantvillage:E_dacis_PMP_30_minmax:plantdoc_5:min\endcsname{45.5}
\expandafter\gdef\csname num@plantvillage:E_dacis_PMP_30_minmax:plantdoc_5:max\endcsname{49.3}
\expandafter\gdef\csname num@plantvillage:E_dacis_PMP_30_minmax:plantdoc_10\endcsname{51.2}
\expandafter\gdef\csname num@plantvillage:E_dacis_PMP_30_minmax:plantdoc_10:sd\endcsname{1.4}
\expandafter\gdef\csname num@plantvillage:E_dacis_PMP_30_minmax:plantdoc_10:min\endcsname{50.1}
\expandafter\gdef\csname num@plantvillage:E_dacis_PMP_30_minmax:plantdoc_10:max\endcsname{53.5}
\expandafter\gdef\csname num@plantvillage:E_dacis_PMP_30_minmax:pct\endcsname{30.0}
\expandafter\gdef\csname num@plantvillage:E_dacis_PMP_30_minmax:params\endcsname{3.35}
\expandafter\gdef\csname num@plantvillage:E_dacis_PMP_30_minmax:macs_128\endcsname{0.251}
\expandafter\gdef\csname num@plantvillage:E_dacis_PMP_30_minmax:macs_224\endcsname{0.770}
\expandafter\gdef\csname num@plantvillage:E_dacis_PMP_30_norefine:novel_1\endcsname{71.8}
\expandafter\gdef\csname num@plantvillage:E_dacis_PMP_30_norefine:novel_1:sd\endcsname{10.4}
\expandafter\gdef\csname num@plantvillage:E_dacis_PMP_30_norefine:novel_1:min\endcsname{58.3}
\expandafter\gdef\csname num@plantvillage:E_dacis_PMP_30_norefine:novel_1:max\endcsname{83.7}
\expandafter\gdef\csname num@plantvillage:E_dacis_PMP_30_norefine:novel_5\endcsname{85.6}
\expandafter\gdef\csname num@plantvillage:E_dacis_PMP_30_norefine:novel_5:sd\endcsname{7.6}
\expandafter\gdef\csname num@plantvillage:E_dacis_PMP_30_norefine:novel_5:min\endcsname{76.0}
\expandafter\gdef\csname num@plantvillage:E_dacis_PMP_30_norefine:novel_5:max\endcsname{93.9}
\expandafter\gdef\csname num@plantvillage:E_dacis_PMP_30_norefine:novel_10\endcsname{87.9}
\expandafter\gdef\csname num@plantvillage:E_dacis_PMP_30_norefine:novel_10:sd\endcsname{6.4}
\expandafter\gdef\csname num@plantvillage:E_dacis_PMP_30_norefine:novel_10:min\endcsname{80.5}
\expandafter\gdef\csname num@plantvillage:E_dacis_PMP_30_norefine:novel_10:max\endcsname{95.1}
\expandafter\gdef\csname num@plantvillage:E_dacis_PMP_30_norefine:plantdoc_1\endcsname{35.0}
\expandafter\gdef\csname num@plantvillage:E_dacis_PMP_30_norefine:plantdoc_1:sd\endcsname{0.7}
\expandafter\gdef\csname num@plantvillage:E_dacis_PMP_30_norefine:plantdoc_1:min\endcsname{33.8}
\expandafter\gdef\csname num@plantvillage:E_dacis_PMP_30_norefine:plantdoc_1:max\endcsname{35.5}
\expandafter\gdef\csname num@plantvillage:E_dacis_PMP_30_norefine:plantdoc_5\endcsname{47.0}
\expandafter\gdef\csname num@plantvillage:E_dacis_PMP_30_norefine:plantdoc_5:sd\endcsname{0.8}
\expandafter\gdef\csname num@plantvillage:E_dacis_PMP_30_norefine:plantdoc_5:min\endcsname{46.0}
\expandafter\gdef\csname num@plantvillage:E_dacis_PMP_30_norefine:plantdoc_5:max\endcsname{48.2}
\expandafter\gdef\csname num@plantvillage:E_dacis_PMP_30_norefine:plantdoc_10\endcsname{51.3}
\expandafter\gdef\csname num@plantvillage:E_dacis_PMP_30_norefine:plantdoc_10:sd\endcsname{0.9}
\expandafter\gdef\csname num@plantvillage:E_dacis_PMP_30_norefine:plantdoc_10:min\endcsname{50.3}
\expandafter\gdef\csname num@plantvillage:E_dacis_PMP_30_norefine:plantdoc_10:max\endcsname{52.5}
\expandafter\gdef\csname num@plantvillage:E_dacis_PMP_30_norefine:pct\endcsname{30.0}
\expandafter\gdef\csname num@plantvillage:E_dacis_PMP_30_norefine:params\endcsname{3.35}
\expandafter\gdef\csname num@plantvillage:E_dacis_PMP_30_norefine:macs_128\endcsname{0.251}
\expandafter\gdef\csname num@plantvillage:E_dacis_PMP_30_norefine:macs_224\endcsname{0.770}
\expandafter\gdef\csname num@plantvillage:E_dacis_PMP_30_uniform:novel_1\endcsname{71.8}
\expandafter\gdef\csname num@plantvillage:E_dacis_PMP_30_uniform:novel_1:sd\endcsname{10.5}
\expandafter\gdef\csname num@plantvillage:E_dacis_PMP_30_uniform:novel_1:min\endcsname{57.6}
\expandafter\gdef\csname num@plantvillage:E_dacis_PMP_30_uniform:novel_1:max\endcsname{82.8}
\expandafter\gdef\csname num@plantvillage:E_dacis_PMP_30_uniform:novel_5\endcsname{86.1}
\expandafter\gdef\csname num@plantvillage:E_dacis_PMP_30_uniform:novel_5:sd\endcsname{7.6}
\expandafter\gdef\csname num@plantvillage:E_dacis_PMP_30_uniform:novel_5:min\endcsname{76.1}
\expandafter\gdef\csname num@plantvillage:E_dacis_PMP_30_uniform:novel_5:max\endcsname{93.5}
\expandafter\gdef\csname num@plantvillage:E_dacis_PMP_30_uniform:novel_10\endcsname{88.4}
\expandafter\gdef\csname num@plantvillage:E_dacis_PMP_30_uniform:novel_10:sd\endcsname{6.5}
\expandafter\gdef\csname num@plantvillage:E_dacis_PMP_30_uniform:novel_10:min\endcsname{80.2}
\expandafter\gdef\csname num@plantvillage:E_dacis_PMP_30_uniform:novel_10:max\endcsname{95.1}
\expandafter\gdef\csname num@plantvillage:E_dacis_PMP_30_uniform:plantdoc_1\endcsname{34.2}
\expandafter\gdef\csname num@plantvillage:E_dacis_PMP_30_uniform:plantdoc_1:sd\endcsname{1.5}
\expandafter\gdef\csname num@plantvillage:E_dacis_PMP_30_uniform:plantdoc_1:min\endcsname{31.9}
\expandafter\gdef\csname num@plantvillage:E_dacis_PMP_30_uniform:plantdoc_1:max\endcsname{35.7}
\expandafter\gdef\csname num@plantvillage:E_dacis_PMP_30_uniform:plantdoc_5\endcsname{45.8}
\expandafter\gdef\csname num@plantvillage:E_dacis_PMP_30_uniform:plantdoc_5:sd\endcsname{2.0}
\expandafter\gdef\csname num@plantvillage:E_dacis_PMP_30_uniform:plantdoc_5:min\endcsname{42.8}
\expandafter\gdef\csname num@plantvillage:E_dacis_PMP_30_uniform:plantdoc_5:max\endcsname{48.2}
\expandafter\gdef\csname num@plantvillage:E_dacis_PMP_30_uniform:plantdoc_10\endcsname{50.0}
\expandafter\gdef\csname num@plantvillage:E_dacis_PMP_30_uniform:plantdoc_10:sd\endcsname{1.9}
\expandafter\gdef\csname num@plantvillage:E_dacis_PMP_30_uniform:plantdoc_10:min\endcsname{47.1}
\expandafter\gdef\csname num@plantvillage:E_dacis_PMP_30_uniform:plantdoc_10:max\endcsname{52.3}
\expandafter\gdef\csname num@plantvillage:E_dacis_PMP_30_uniform:pct\endcsname{30.1}
\expandafter\gdef\csname num@plantvillage:E_dacis_PMP_30_uniform:params\endcsname{3.36}
\expandafter\gdef\csname num@plantvillage:E_dacis_PMP_30_uniform:macs_128\endcsname{0.202}
\expandafter\gdef\csname num@plantvillage:E_dacis_PMP_30_uniform:macs_224\endcsname{0.618}
\expandafter\gdef\csname num@plantvillage:K_efficientnet_b0:novel_1\endcsname{82.7}
\expandafter\gdef\csname num@plantvillage:K_efficientnet_b0:novel_1:sd\endcsname{10.1}
\expandafter\gdef\csname num@plantvillage:K_efficientnet_b0:novel_1:min\endcsname{68.4}
\expandafter\gdef\csname num@plantvillage:K_efficientnet_b0:novel_1:max\endcsname{93.7}
\expandafter\gdef\csname num@plantvillage:K_efficientnet_b0:novel_5\endcsname{93.8}
\expandafter\gdef\csname num@plantvillage:K_efficientnet_b0:novel_5:sd\endcsname{5.2}
\expandafter\gdef\csname num@plantvillage:K_efficientnet_b0:novel_5:min\endcsname{86.7}
\expandafter\gdef\csname num@plantvillage:K_efficientnet_b0:novel_5:max\endcsname{98.7}
\expandafter\gdef\csname num@plantvillage:K_efficientnet_b0:novel_10\endcsname{95.4}
\expandafter\gdef\csname num@plantvillage:K_efficientnet_b0:novel_10:sd\endcsname{4.0}
\expandafter\gdef\csname num@plantvillage:K_efficientnet_b0:novel_10:min\endcsname{89.9}
\expandafter\gdef\csname num@plantvillage:K_efficientnet_b0:novel_10:max\endcsname{99.1}
\expandafter\gdef\csname num@plantvillage:K_efficientnet_b0:plantdoc_1\endcsname{43.4}
\expandafter\gdef\csname num@plantvillage:K_efficientnet_b0:plantdoc_1:sd\endcsname{0.6}
\expandafter\gdef\csname num@plantvillage:K_efficientnet_b0:plantdoc_1:min\endcsname{42.9}
\expandafter\gdef\csname num@plantvillage:K_efficientnet_b0:plantdoc_1:max\endcsname{44.3}
\expandafter\gdef\csname num@plantvillage:K_efficientnet_b0:plantdoc_5\endcsname{58.5}
\expandafter\gdef\csname num@plantvillage:K_efficientnet_b0:plantdoc_5:sd\endcsname{0.4}
\expandafter\gdef\csname num@plantvillage:K_efficientnet_b0:plantdoc_5:min\endcsname{57.9}
\expandafter\gdef\csname num@plantvillage:K_efficientnet_b0:plantdoc_5:max\endcsname{59.1}
\expandafter\gdef\csname num@plantvillage:K_efficientnet_b0:plantdoc_10\endcsname{63.5}
\expandafter\gdef\csname num@plantvillage:K_efficientnet_b0:plantdoc_10:sd\endcsname{0.6}
\expandafter\gdef\csname num@plantvillage:K_efficientnet_b0:plantdoc_10:min\endcsname{62.7}
\expandafter\gdef\csname num@plantvillage:K_efficientnet_b0:plantdoc_10:max\endcsname{64.3}
\expandafter\gdef\csname num@plantvillage:K_efficientnet_b0:pct\endcsname{35.9}
\expandafter\gdef\csname num@plantvillage:K_efficientnet_b0:params\endcsname{4.01}
\expandafter\gdef\csname num@plantvillage:K_efficientnet_b0:macs_128\endcsname{0.126}
\expandafter\gdef\csname num@plantvillage:K_efficientnet_b0:macs_224\endcsname{0.385}
\expandafter\gdef\csname num@plantvillage:K_efficientnet_b0:cpu_1thread_b1_ms\endcsname{6.2}
\expandafter\gdef\csname num@plantvillage:K_efficientnet_b0:cpu_all_threads_b1_ms\endcsname{7.2}
\expandafter\gdef\csname num@plantvillage:K_efficientnet_b0:gpu_b1_ms\endcsname{2.8}
\expandafter\gdef\csname num@plantvillage:K_efficientnet_b0:gpu_ips\endcsname{18728}
\expandafter\gdef\csname num@plantvillage:K_mobilenet_v2:novel_1\endcsname{79.2}
\expandafter\gdef\csname num@plantvillage:K_mobilenet_v2:novel_1:sd\endcsname{10.5}
\expandafter\gdef\csname num@plantvillage:K_mobilenet_v2:novel_1:min\endcsname{64.1}
\expandafter\gdef\csname num@plantvillage:K_mobilenet_v2:novel_1:max\endcsname{91.3}
\expandafter\gdef\csname num@plantvillage:K_mobilenet_v2:novel_5\endcsname{91.2}
\expandafter\gdef\csname num@plantvillage:K_mobilenet_v2:novel_5:sd\endcsname{6.3}
\expandafter\gdef\csname num@plantvillage:K_mobilenet_v2:novel_5:min\endcsname{82.3}
\expandafter\gdef\csname num@plantvillage:K_mobilenet_v2:novel_5:max\endcsname{97.8}
\expandafter\gdef\csname num@plantvillage:K_mobilenet_v2:novel_10\endcsname{93.1}
\expandafter\gdef\csname num@plantvillage:K_mobilenet_v2:novel_10:sd\endcsname{4.9}
\expandafter\gdef\csname num@plantvillage:K_mobilenet_v2:novel_10:min\endcsname{86.2}
\expandafter\gdef\csname num@plantvillage:K_mobilenet_v2:novel_10:max\endcsname{98.3}
\expandafter\gdef\csname num@plantvillage:K_mobilenet_v2:plantdoc_1\endcsname{40.9}
\expandafter\gdef\csname num@plantvillage:K_mobilenet_v2:plantdoc_1:sd\endcsname{0.8}
\expandafter\gdef\csname num@plantvillage:K_mobilenet_v2:plantdoc_1:min\endcsname{39.7}
\expandafter\gdef\csname num@plantvillage:K_mobilenet_v2:plantdoc_1:max\endcsname{41.6}
\expandafter\gdef\csname num@plantvillage:K_mobilenet_v2:plantdoc_5\endcsname{55.1}
\expandafter\gdef\csname num@plantvillage:K_mobilenet_v2:plantdoc_5:sd\endcsname{1.0}
\expandafter\gdef\csname num@plantvillage:K_mobilenet_v2:plantdoc_5:min\endcsname{53.8}
\expandafter\gdef\csname num@plantvillage:K_mobilenet_v2:plantdoc_5:max\endcsname{56.1}
\expandafter\gdef\csname num@plantvillage:K_mobilenet_v2:plantdoc_10\endcsname{59.8}
\expandafter\gdef\csname num@plantvillage:K_mobilenet_v2:plantdoc_10:sd\endcsname{1.3}
\expandafter\gdef\csname num@plantvillage:K_mobilenet_v2:plantdoc_10:min\endcsname{58.3}
\expandafter\gdef\csname num@plantvillage:K_mobilenet_v2:plantdoc_10:max\endcsname{61.3}
\expandafter\gdef\csname num@plantvillage:K_mobilenet_v2:pct\endcsname{19.9}
\expandafter\gdef\csname num@plantvillage:K_mobilenet_v2:params\endcsname{2.22}
\expandafter\gdef\csname num@plantvillage:K_mobilenet_v2:macs_128\endcsname{0.098}
\expandafter\gdef\csname num@plantvillage:K_mobilenet_v2:macs_224\endcsname{0.299}
\expandafter\gdef\csname num@plantvillage:K_mobilenet_v2:cpu_1thread_b1_ms\endcsname{4.3}
\expandafter\gdef\csname num@plantvillage:K_mobilenet_v2:cpu_all_threads_b1_ms\endcsname{5.7}
\expandafter\gdef\csname num@plantvillage:K_mobilenet_v2:gpu_b1_ms\endcsname{1.9}
\expandafter\gdef\csname num@plantvillage:K_mobilenet_v2:gpu_ips\endcsname{24222}
\expandafter\gdef\csname num@plantvillage:K_mobilenet_v3_small:novel_1\endcsname{78.0}
\expandafter\gdef\csname num@plantvillage:K_mobilenet_v3_small:novel_1:sd\endcsname{10.1}
\expandafter\gdef\csname num@plantvillage:K_mobilenet_v3_small:novel_1:min\endcsname{63.7}
\expandafter\gdef\csname num@plantvillage:K_mobilenet_v3_small:novel_1:max\endcsname{87.7}
\expandafter\gdef\csname num@plantvillage:K_mobilenet_v3_small:novel_5\endcsname{91.1}
\expandafter\gdef\csname num@plantvillage:K_mobilenet_v3_small:novel_5:sd\endcsname{5.9}
\expandafter\gdef\csname num@plantvillage:K_mobilenet_v3_small:novel_5:min\endcsname{83.2}
\expandafter\gdef\csname num@plantvillage:K_mobilenet_v3_small:novel_5:max\endcsname{96.4}
\expandafter\gdef\csname num@plantvillage:K_mobilenet_v3_small:novel_10\endcsname{93.1}
\expandafter\gdef\csname num@plantvillage:K_mobilenet_v3_small:novel_10:sd\endcsname{4.7}
\expandafter\gdef\csname num@plantvillage:K_mobilenet_v3_small:novel_10:min\endcsname{86.8}
\expandafter\gdef\csname num@plantvillage:K_mobilenet_v3_small:novel_10:max\endcsname{97.2}
\expandafter\gdef\csname num@plantvillage:K_mobilenet_v3_small:plantdoc_1\endcsname{36.9}
\expandafter\gdef\csname num@plantvillage:K_mobilenet_v3_small:plantdoc_1:sd\endcsname{0.6}
\expandafter\gdef\csname num@plantvillage:K_mobilenet_v3_small:plantdoc_1:min\endcsname{36.0}
\expandafter\gdef\csname num@plantvillage:K_mobilenet_v3_small:plantdoc_1:max\endcsname{37.8}
\expandafter\gdef\csname num@plantvillage:K_mobilenet_v3_small:plantdoc_5\endcsname{50.0}
\expandafter\gdef\csname num@plantvillage:K_mobilenet_v3_small:plantdoc_5:sd\endcsname{0.8}
\expandafter\gdef\csname num@plantvillage:K_mobilenet_v3_small:plantdoc_5:min\endcsname{49.1}
\expandafter\gdef\csname num@plantvillage:K_mobilenet_v3_small:plantdoc_5:max\endcsname{51.0}
\expandafter\gdef\csname num@plantvillage:K_mobilenet_v3_small:plantdoc_10\endcsname{54.7}
\expandafter\gdef\csname num@plantvillage:K_mobilenet_v3_small:plantdoc_10:sd\endcsname{0.5}
\expandafter\gdef\csname num@plantvillage:K_mobilenet_v3_small:plantdoc_10:min\endcsname{54.1}
\expandafter\gdef\csname num@plantvillage:K_mobilenet_v3_small:plantdoc_10:max\endcsname{55.4}
\expandafter\gdef\csname num@plantvillage:K_mobilenet_v3_small:pct\endcsname{8.3}
\expandafter\gdef\csname num@plantvillage:K_mobilenet_v3_small:params\endcsname{0.93}
\expandafter\gdef\csname num@plantvillage:K_mobilenet_v3_small:macs_128\endcsname{0.018}
\expandafter\gdef\csname num@plantvillage:K_mobilenet_v3_small:macs_224\endcsname{0.055}
\expandafter\gdef\csname num@plantvillage:K_mobilenet_v3_small:cpu_1thread_b1_ms\endcsname{2.2}
\expandafter\gdef\csname num@plantvillage:K_mobilenet_v3_small:cpu_all_threads_b1_ms\endcsname{4.1}
\expandafter\gdef\csname num@plantvillage:K_mobilenet_v3_small:gpu_b1_ms\endcsname{1.8}
\expandafter\gdef\csname num@plantvillage:K_mobilenet_v3_small:gpu_ips\endcsname{63459}
\expandafter\gdef\csname num@plantvillage:R_dacis_P_30_r3:novel_1\endcsname{73.2}
\expandafter\gdef\csname num@plantvillage:R_dacis_P_30_r3:novel_1:sd\endcsname{11.0}
\expandafter\gdef\csname num@plantvillage:R_dacis_P_30_r3:novel_1:min\endcsname{56.9}
\expandafter\gdef\csname num@plantvillage:R_dacis_P_30_r3:novel_1:max\endcsname{83.7}
\expandafter\gdef\csname num@plantvillage:R_dacis_P_30_r3:novel_5\endcsname{86.6}
\expandafter\gdef\csname num@plantvillage:R_dacis_P_30_r3:novel_5:sd\endcsname{7.7}
\expandafter\gdef\csname num@plantvillage:R_dacis_P_30_r3:novel_5:min\endcsname{75.3}
\expandafter\gdef\csname num@plantvillage:R_dacis_P_30_r3:novel_5:max\endcsname{93.8}
\expandafter\gdef\csname num@plantvillage:R_dacis_P_30_r3:novel_10\endcsname{88.7}
\expandafter\gdef\csname num@plantvillage:R_dacis_P_30_r3:novel_10:sd\endcsname{6.5}
\expandafter\gdef\csname num@plantvillage:R_dacis_P_30_r3:novel_10:min\endcsname{79.6}
\expandafter\gdef\csname num@plantvillage:R_dacis_P_30_r3:novel_10:max\endcsname{95.1}
\expandafter\gdef\csname num@plantvillage:R_dacis_P_30_r3:plantdoc_1\endcsname{35.0}
\expandafter\gdef\csname num@plantvillage:R_dacis_P_30_r3:plantdoc_1:sd\endcsname{0.7}
\expandafter\gdef\csname num@plantvillage:R_dacis_P_30_r3:plantdoc_1:min\endcsname{34.0}
\expandafter\gdef\csname num@plantvillage:R_dacis_P_30_r3:plantdoc_1:max\endcsname{36.0}
\expandafter\gdef\csname num@plantvillage:R_dacis_P_30_r3:plantdoc_5\endcsname{47.2}
\expandafter\gdef\csname num@plantvillage:R_dacis_P_30_r3:plantdoc_5:sd\endcsname{0.9}
\expandafter\gdef\csname num@plantvillage:R_dacis_P_30_r3:plantdoc_5:min\endcsname{46.5}
\expandafter\gdef\csname num@plantvillage:R_dacis_P_30_r3:plantdoc_5:max\endcsname{48.8}
\expandafter\gdef\csname num@plantvillage:R_dacis_P_30_r3:plantdoc_10\endcsname{51.4}
\expandafter\gdef\csname num@plantvillage:R_dacis_P_30_r3:plantdoc_10:sd\endcsname{1.0}
\expandafter\gdef\csname num@plantvillage:R_dacis_P_30_r3:plantdoc_10:min\endcsname{50.3}
\expandafter\gdef\csname num@plantvillage:R_dacis_P_30_r3:plantdoc_10:max\endcsname{53.0}
\expandafter\gdef\csname num@plantvillage:R_dacis_P_30_r3:pct\endcsname{30.0}
\expandafter\gdef\csname num@plantvillage:R_dacis_P_30_r3:params\endcsname{3.35}
\expandafter\gdef\csname num@plantvillage:R_dacis_P_30_r3:macs_128\endcsname{0.251}
\expandafter\gdef\csname num@plantvillage:R_dacis_P_30_r3:macs_224\endcsname{0.770}
\expandafter\gdef\csname num@plantvillage:R_dacis_P_30_r6:novel_1\endcsname{73.1}
\expandafter\gdef\csname num@plantvillage:R_dacis_P_30_r6:novel_1:sd\endcsname{10.8}
\expandafter\gdef\csname num@plantvillage:R_dacis_P_30_r6:novel_1:min\endcsname{57.2}
\expandafter\gdef\csname num@plantvillage:R_dacis_P_30_r6:novel_1:max\endcsname{83.1}
\expandafter\gdef\csname num@plantvillage:R_dacis_P_30_r6:novel_5\endcsname{86.5}
\expandafter\gdef\csname num@plantvillage:R_dacis_P_30_r6:novel_5:sd\endcsname{7.7}
\expandafter\gdef\csname num@plantvillage:R_dacis_P_30_r6:novel_5:min\endcsname{75.4}
\expandafter\gdef\csname num@plantvillage:R_dacis_P_30_r6:novel_5:max\endcsname{93.7}
\expandafter\gdef\csname num@plantvillage:R_dacis_P_30_r6:novel_10\endcsname{88.6}
\expandafter\gdef\csname num@plantvillage:R_dacis_P_30_r6:novel_10:sd\endcsname{6.4}
\expandafter\gdef\csname num@plantvillage:R_dacis_P_30_r6:novel_10:min\endcsname{79.8}
\expandafter\gdef\csname num@plantvillage:R_dacis_P_30_r6:novel_10:max\endcsname{94.8}
\expandafter\gdef\csname num@plantvillage:R_dacis_P_30_r6:plantdoc_1\endcsname{34.9}
\expandafter\gdef\csname num@plantvillage:R_dacis_P_30_r6:plantdoc_1:sd\endcsname{1.1}
\expandafter\gdef\csname num@plantvillage:R_dacis_P_30_r6:plantdoc_1:min\endcsname{33.1}
\expandafter\gdef\csname num@plantvillage:R_dacis_P_30_r6:plantdoc_1:max\endcsname{36.0}
\expandafter\gdef\csname num@plantvillage:R_dacis_P_30_r6:plantdoc_5\endcsname{46.9}
\expandafter\gdef\csname num@plantvillage:R_dacis_P_30_r6:plantdoc_5:sd\endcsname{1.5}
\expandafter\gdef\csname num@plantvillage:R_dacis_P_30_r6:plantdoc_5:min\endcsname{44.9}
\expandafter\gdef\csname num@plantvillage:R_dacis_P_30_r6:plantdoc_5:max\endcsname{49.0}
\expandafter\gdef\csname num@plantvillage:R_dacis_P_30_r6:plantdoc_10\endcsname{51.1}
\expandafter\gdef\csname num@plantvillage:R_dacis_P_30_r6:plantdoc_10:sd\endcsname{1.5}
\expandafter\gdef\csname num@plantvillage:R_dacis_P_30_r6:plantdoc_10:min\endcsname{49.4}
\expandafter\gdef\csname num@plantvillage:R_dacis_P_30_r6:plantdoc_10:max\endcsname{53.1}
\expandafter\gdef\csname num@plantvillage:R_dacis_P_30_r6:pct\endcsname{30.0}
\expandafter\gdef\csname num@plantvillage:R_dacis_P_30_r6:params\endcsname{3.35}
\expandafter\gdef\csname num@plantvillage:R_dacis_P_30_r6:macs_128\endcsname{0.251}
\expandafter\gdef\csname num@plantvillage:R_dacis_P_30_r6:macs_224\endcsname{0.770}
\expandafter\gdef\csname num@plantvillage:R_l1_P_30_r3:novel_1\endcsname{74.7}
\expandafter\gdef\csname num@plantvillage:R_l1_P_30_r3:novel_1:sd\endcsname{10.7}
\expandafter\gdef\csname num@plantvillage:R_l1_P_30_r3:novel_1:min\endcsname{63.2}
\expandafter\gdef\csname num@plantvillage:R_l1_P_30_r3:novel_1:max\endcsname{87.5}
\expandafter\gdef\csname num@plantvillage:R_l1_P_30_r3:novel_5\endcsname{87.4}
\expandafter\gdef\csname num@plantvillage:R_l1_P_30_r3:novel_5:sd\endcsname{7.3}
\expandafter\gdef\csname num@plantvillage:R_l1_P_30_r3:novel_5:min\endcsname{79.5}
\expandafter\gdef\csname num@plantvillage:R_l1_P_30_r3:novel_5:max\endcsname{95.4}
\expandafter\gdef\csname num@plantvillage:R_l1_P_30_r3:novel_10\endcsname{89.5}
\expandafter\gdef\csname num@plantvillage:R_l1_P_30_r3:novel_10:sd\endcsname{6.2}
\expandafter\gdef\csname num@plantvillage:R_l1_P_30_r3:novel_10:min\endcsname{82.3}
\expandafter\gdef\csname num@plantvillage:R_l1_P_30_r3:novel_10:max\endcsname{96.2}
\expandafter\gdef\csname num@plantvillage:R_l1_P_30_r3:plantdoc_1\endcsname{37.1}
\expandafter\gdef\csname num@plantvillage:R_l1_P_30_r3:plantdoc_1:sd\endcsname{1.2}
\expandafter\gdef\csname num@plantvillage:R_l1_P_30_r3:plantdoc_1:min\endcsname{35.4}
\expandafter\gdef\csname num@plantvillage:R_l1_P_30_r3:plantdoc_1:max\endcsname{38.5}
\expandafter\gdef\csname num@plantvillage:R_l1_P_30_r3:plantdoc_5\endcsname{49.2}
\expandafter\gdef\csname num@plantvillage:R_l1_P_30_r3:plantdoc_5:sd\endcsname{1.2}
\expandafter\gdef\csname num@plantvillage:R_l1_P_30_r3:plantdoc_5:min\endcsname{47.8}
\expandafter\gdef\csname num@plantvillage:R_l1_P_30_r3:plantdoc_5:max\endcsname{50.5}
\expandafter\gdef\csname num@plantvillage:R_l1_P_30_r3:plantdoc_10\endcsname{53.4}
\expandafter\gdef\csname num@plantvillage:R_l1_P_30_r3:plantdoc_10:sd\endcsname{1.2}
\expandafter\gdef\csname num@plantvillage:R_l1_P_30_r3:plantdoc_10:min\endcsname{52.1}
\expandafter\gdef\csname num@plantvillage:R_l1_P_30_r3:plantdoc_10:max\endcsname{54.7}
\expandafter\gdef\csname num@plantvillage:R_l1_P_30_r3:pct\endcsname{30.0}
\expandafter\gdef\csname num@plantvillage:R_l1_P_30_r3:params\endcsname{3.35}
\expandafter\gdef\csname num@plantvillage:R_l1_P_30_r3:macs_128\endcsname{0.251}
\expandafter\gdef\csname num@plantvillage:R_l1_P_30_r3:macs_224\endcsname{0.770}
\expandafter\gdef\csname num@plantvillage:R_l1_P_30_r6:novel_1\endcsname{74.4}
\expandafter\gdef\csname num@plantvillage:R_l1_P_30_r6:novel_1:sd\endcsname{10.6}
\expandafter\gdef\csname num@plantvillage:R_l1_P_30_r6:novel_1:min\endcsname{61.5}
\expandafter\gdef\csname num@plantvillage:R_l1_P_30_r6:novel_1:max\endcsname{86.3}
\expandafter\gdef\csname num@plantvillage:R_l1_P_30_r6:novel_5\endcsname{87.3}
\expandafter\gdef\csname num@plantvillage:R_l1_P_30_r6:novel_5:sd\endcsname{7.2}
\expandafter\gdef\csname num@plantvillage:R_l1_P_30_r6:novel_5:min\endcsname{78.8}
\expandafter\gdef\csname num@plantvillage:R_l1_P_30_r6:novel_5:max\endcsname{94.9}
\expandafter\gdef\csname num@plantvillage:R_l1_P_30_r6:novel_10\endcsname{89.5}
\expandafter\gdef\csname num@plantvillage:R_l1_P_30_r6:novel_10:sd\endcsname{6.0}
\expandafter\gdef\csname num@plantvillage:R_l1_P_30_r6:novel_10:min\endcsname{82.7}
\expandafter\gdef\csname num@plantvillage:R_l1_P_30_r6:novel_10:max\endcsname{95.9}
\expandafter\gdef\csname num@plantvillage:R_l1_P_30_r6:plantdoc_1\endcsname{37.5}
\expandafter\gdef\csname num@plantvillage:R_l1_P_30_r6:plantdoc_1:sd\endcsname{1.4}
\expandafter\gdef\csname num@plantvillage:R_l1_P_30_r6:plantdoc_1:min\endcsname{35.1}
\expandafter\gdef\csname num@plantvillage:R_l1_P_30_r6:plantdoc_1:max\endcsname{38.7}
\expandafter\gdef\csname num@plantvillage:R_l1_P_30_r6:plantdoc_5\endcsname{49.5}
\expandafter\gdef\csname num@plantvillage:R_l1_P_30_r6:plantdoc_5:sd\endcsname{1.2}
\expandafter\gdef\csname num@plantvillage:R_l1_P_30_r6:plantdoc_5:min\endcsname{47.8}
\expandafter\gdef\csname num@plantvillage:R_l1_P_30_r6:plantdoc_5:max\endcsname{50.8}
\expandafter\gdef\csname num@plantvillage:R_l1_P_30_r6:plantdoc_10\endcsname{53.7}
\expandafter\gdef\csname num@plantvillage:R_l1_P_30_r6:plantdoc_10:sd\endcsname{1.1}
\expandafter\gdef\csname num@plantvillage:R_l1_P_30_r6:plantdoc_10:min\endcsname{52.1}
\expandafter\gdef\csname num@plantvillage:R_l1_P_30_r6:plantdoc_10:max\endcsname{54.9}
\expandafter\gdef\csname num@plantvillage:R_l1_P_30_r6:pct\endcsname{30.0}
\expandafter\gdef\csname num@plantvillage:R_l1_P_30_r6:params\endcsname{3.35}
\expandafter\gdef\csname num@plantvillage:R_l1_P_30_r6:macs_128\endcsname{0.251}
\expandafter\gdef\csname num@plantvillage:R_l1_P_30_r6:macs_224\endcsname{0.770}
\expandafter\gdef\csname num@plantvillage:R_random_P_30_r3:novel_1\endcsname{73.5}
\expandafter\gdef\csname num@plantvillage:R_random_P_30_r3:novel_1:sd\endcsname{11.8}
\expandafter\gdef\csname num@plantvillage:R_random_P_30_r3:novel_1:min\endcsname{57.2}
\expandafter\gdef\csname num@plantvillage:R_random_P_30_r3:novel_1:max\endcsname{86.6}
\expandafter\gdef\csname num@plantvillage:R_random_P_30_r3:novel_5\endcsname{86.8}
\expandafter\gdef\csname num@plantvillage:R_random_P_30_r3:novel_5:sd\endcsname{8.1}
\expandafter\gdef\csname num@plantvillage:R_random_P_30_r3:novel_5:min\endcsname{75.5}
\expandafter\gdef\csname num@plantvillage:R_random_P_30_r3:novel_5:max\endcsname{95.2}
\expandafter\gdef\csname num@plantvillage:R_random_P_30_r3:novel_10\endcsname{89.1}
\expandafter\gdef\csname num@plantvillage:R_random_P_30_r3:novel_10:sd\endcsname{6.9}
\expandafter\gdef\csname num@plantvillage:R_random_P_30_r3:novel_10:min\endcsname{79.5}
\expandafter\gdef\csname num@plantvillage:R_random_P_30_r3:novel_10:max\endcsname{96.1}
\expandafter\gdef\csname num@plantvillage:R_random_P_30_r3:plantdoc_1\endcsname{35.6}
\expandafter\gdef\csname num@plantvillage:R_random_P_30_r3:plantdoc_1:sd\endcsname{1.5}
\expandafter\gdef\csname num@plantvillage:R_random_P_30_r3:plantdoc_1:min\endcsname{34.2}
\expandafter\gdef\csname num@plantvillage:R_random_P_30_r3:plantdoc_1:max\endcsname{37.8}
\expandafter\gdef\csname num@plantvillage:R_random_P_30_r3:plantdoc_5\endcsname{48.0}
\expandafter\gdef\csname num@plantvillage:R_random_P_30_r3:plantdoc_5:sd\endcsname{1.8}
\expandafter\gdef\csname num@plantvillage:R_random_P_30_r3:plantdoc_5:min\endcsname{46.4}
\expandafter\gdef\csname num@plantvillage:R_random_P_30_r3:plantdoc_5:max\endcsname{50.9}
\expandafter\gdef\csname num@plantvillage:R_random_P_30_r3:plantdoc_10\endcsname{52.2}
\expandafter\gdef\csname num@plantvillage:R_random_P_30_r3:plantdoc_10:sd\endcsname{1.6}
\expandafter\gdef\csname num@plantvillage:R_random_P_30_r3:plantdoc_10:min\endcsname{50.7}
\expandafter\gdef\csname num@plantvillage:R_random_P_30_r3:plantdoc_10:max\endcsname{54.7}
\expandafter\gdef\csname num@plantvillage:R_random_P_30_r3:pct\endcsname{30.0}
\expandafter\gdef\csname num@plantvillage:R_random_P_30_r3:params\endcsname{3.35}
\expandafter\gdef\csname num@plantvillage:R_random_P_30_r3:macs_128\endcsname{0.251}
\expandafter\gdef\csname num@plantvillage:R_random_P_30_r3:macs_224\endcsname{0.770}
\expandafter\gdef\csname num@plantvillage:R_random_P_30_r6:novel_1\endcsname{73.0}
\expandafter\gdef\csname num@plantvillage:R_random_P_30_r6:novel_1:sd\endcsname{11.1}
\expandafter\gdef\csname num@plantvillage:R_random_P_30_r6:novel_1:min\endcsname{58.7}
\expandafter\gdef\csname num@plantvillage:R_random_P_30_r6:novel_1:max\endcsname{85.9}
\expandafter\gdef\csname num@plantvillage:R_random_P_30_r6:novel_5\endcsname{86.3}
\expandafter\gdef\csname num@plantvillage:R_random_P_30_r6:novel_5:sd\endcsname{8.2}
\expandafter\gdef\csname num@plantvillage:R_random_P_30_r6:novel_5:min\endcsname{76.2}
\expandafter\gdef\csname num@plantvillage:R_random_P_30_r6:novel_5:max\endcsname{95.3}
\expandafter\gdef\csname num@plantvillage:R_random_P_30_r6:novel_10\endcsname{88.5}
\expandafter\gdef\csname num@plantvillage:R_random_P_30_r6:novel_10:sd\endcsname{7.1}
\expandafter\gdef\csname num@plantvillage:R_random_P_30_r6:novel_10:min\endcsname{80.1}
\expandafter\gdef\csname num@plantvillage:R_random_P_30_r6:novel_10:max\endcsname{96.3}
\expandafter\gdef\csname num@plantvillage:R_random_P_30_r6:plantdoc_1\endcsname{34.8}
\expandafter\gdef\csname num@plantvillage:R_random_P_30_r6:plantdoc_1:sd\endcsname{1.4}
\expandafter\gdef\csname num@plantvillage:R_random_P_30_r6:plantdoc_1:min\endcsname{32.6}
\expandafter\gdef\csname num@plantvillage:R_random_P_30_r6:plantdoc_1:max\endcsname{36.4}
\expandafter\gdef\csname num@plantvillage:R_random_P_30_r6:plantdoc_5\endcsname{46.6}
\expandafter\gdef\csname num@plantvillage:R_random_P_30_r6:plantdoc_5:sd\endcsname{1.8}
\expandafter\gdef\csname num@plantvillage:R_random_P_30_r6:plantdoc_5:min\endcsname{44.0}
\expandafter\gdef\csname num@plantvillage:R_random_P_30_r6:plantdoc_5:max\endcsname{48.6}
\expandafter\gdef\csname num@plantvillage:R_random_P_30_r6:plantdoc_10\endcsname{50.9}
\expandafter\gdef\csname num@plantvillage:R_random_P_30_r6:plantdoc_10:sd\endcsname{1.7}
\expandafter\gdef\csname num@plantvillage:R_random_P_30_r6:plantdoc_10:min\endcsname{48.5}
\expandafter\gdef\csname num@plantvillage:R_random_P_30_r6:plantdoc_10:max\endcsname{53.0}
\expandafter\gdef\csname num@plantvillage:R_random_P_30_r6:pct\endcsname{30.0}
\expandafter\gdef\csname num@plantvillage:R_random_P_30_r6:params\endcsname{3.35}
\expandafter\gdef\csname num@plantvillage:R_random_P_30_r6:macs_128\endcsname{0.251}
\expandafter\gdef\csname num@plantvillage:R_random_P_30_r6:macs_224\endcsname{0.770}
\expandafter\gdef\csname num@plantvillage:R_taylor_P_30_r3:novel_1\endcsname{74.7}
\expandafter\gdef\csname num@plantvillage:R_taylor_P_30_r3:novel_1:sd\endcsname{10.5}
\expandafter\gdef\csname num@plantvillage:R_taylor_P_30_r3:novel_1:min\endcsname{60.3}
\expandafter\gdef\csname num@plantvillage:R_taylor_P_30_r3:novel_1:max\endcsname{84.9}
\expandafter\gdef\csname num@plantvillage:R_taylor_P_30_r3:novel_5\endcsname{88.2}
\expandafter\gdef\csname num@plantvillage:R_taylor_P_30_r3:novel_5:sd\endcsname{7.3}
\expandafter\gdef\csname num@plantvillage:R_taylor_P_30_r3:novel_5:min\endcsname{78.5}
\expandafter\gdef\csname num@plantvillage:R_taylor_P_30_r3:novel_5:max\endcsname{95.4}
\expandafter\gdef\csname num@plantvillage:R_taylor_P_30_r3:novel_10\endcsname{90.2}
\expandafter\gdef\csname num@plantvillage:R_taylor_P_30_r3:novel_10:sd\endcsname{6.0}
\expandafter\gdef\csname num@plantvillage:R_taylor_P_30_r3:novel_10:min\endcsname{82.4}
\expandafter\gdef\csname num@plantvillage:R_taylor_P_30_r3:novel_10:max\endcsname{96.3}
\expandafter\gdef\csname num@plantvillage:R_taylor_P_30_r3:plantdoc_1\endcsname{35.9}
\expandafter\gdef\csname num@plantvillage:R_taylor_P_30_r3:plantdoc_1:sd\endcsname{1.4}
\expandafter\gdef\csname num@plantvillage:R_taylor_P_30_r3:plantdoc_1:min\endcsname{33.6}
\expandafter\gdef\csname num@plantvillage:R_taylor_P_30_r3:plantdoc_1:max\endcsname{37.4}
\expandafter\gdef\csname num@plantvillage:R_taylor_P_30_r3:plantdoc_5\endcsname{48.5}
\expandafter\gdef\csname num@plantvillage:R_taylor_P_30_r3:plantdoc_5:sd\endcsname{1.8}
\expandafter\gdef\csname num@plantvillage:R_taylor_P_30_r3:plantdoc_5:min\endcsname{46.4}
\expandafter\gdef\csname num@plantvillage:R_taylor_P_30_r3:plantdoc_5:max\endcsname{51.0}
\expandafter\gdef\csname num@plantvillage:R_taylor_P_30_r3:plantdoc_10\endcsname{52.8}
\expandafter\gdef\csname num@plantvillage:R_taylor_P_30_r3:plantdoc_10:sd\endcsname{1.7}
\expandafter\gdef\csname num@plantvillage:R_taylor_P_30_r3:plantdoc_10:min\endcsname{50.9}
\expandafter\gdef\csname num@plantvillage:R_taylor_P_30_r3:plantdoc_10:max\endcsname{55.4}
\expandafter\gdef\csname num@plantvillage:R_taylor_P_30_r3:pct\endcsname{30.0}
\expandafter\gdef\csname num@plantvillage:R_taylor_P_30_r3:params\endcsname{3.35}
\expandafter\gdef\csname num@plantvillage:R_taylor_P_30_r3:macs_128\endcsname{0.251}
\expandafter\gdef\csname num@plantvillage:R_taylor_P_30_r3:macs_224\endcsname{0.770}
\expandafter\gdef\csname num@plantvillage:R_taylor_P_30_r6:novel_1\endcsname{74.4}
\expandafter\gdef\csname num@plantvillage:R_taylor_P_30_r6:novel_1:sd\endcsname{10.7}
\expandafter\gdef\csname num@plantvillage:R_taylor_P_30_r6:novel_1:min\endcsname{59.7}
\expandafter\gdef\csname num@plantvillage:R_taylor_P_30_r6:novel_1:max\endcsname{84.5}
\expandafter\gdef\csname num@plantvillage:R_taylor_P_30_r6:novel_5\endcsname{87.8}
\expandafter\gdef\csname num@plantvillage:R_taylor_P_30_r6:novel_5:sd\endcsname{7.7}
\expandafter\gdef\csname num@plantvillage:R_taylor_P_30_r6:novel_5:min\endcsname{77.6}
\expandafter\gdef\csname num@plantvillage:R_taylor_P_30_r6:novel_5:max\endcsname{95.3}
\expandafter\gdef\csname num@plantvillage:R_taylor_P_30_r6:novel_10\endcsname{89.8}
\expandafter\gdef\csname num@plantvillage:R_taylor_P_30_r6:novel_10:sd\endcsname{6.4}
\expandafter\gdef\csname num@plantvillage:R_taylor_P_30_r6:novel_10:min\endcsname{81.7}
\expandafter\gdef\csname num@plantvillage:R_taylor_P_30_r6:novel_10:max\endcsname{96.3}
\expandafter\gdef\csname num@plantvillage:R_taylor_P_30_r6:plantdoc_1\endcsname{35.8}
\expandafter\gdef\csname num@plantvillage:R_taylor_P_30_r6:plantdoc_1:sd\endcsname{1.4}
\expandafter\gdef\csname num@plantvillage:R_taylor_P_30_r6:plantdoc_1:min\endcsname{33.4}
\expandafter\gdef\csname num@plantvillage:R_taylor_P_30_r6:plantdoc_1:max\endcsname{36.7}
\expandafter\gdef\csname num@plantvillage:R_taylor_P_30_r6:plantdoc_5\endcsname{48.2}
\expandafter\gdef\csname num@plantvillage:R_taylor_P_30_r6:plantdoc_5:sd\endcsname{1.4}
\expandafter\gdef\csname num@plantvillage:R_taylor_P_30_r6:plantdoc_5:min\endcsname{46.3}
\expandafter\gdef\csname num@plantvillage:R_taylor_P_30_r6:plantdoc_5:max\endcsname{50.1}
\expandafter\gdef\csname num@plantvillage:R_taylor_P_30_r6:plantdoc_10\endcsname{52.5}
\expandafter\gdef\csname num@plantvillage:R_taylor_P_30_r6:plantdoc_10:sd\endcsname{1.3}
\expandafter\gdef\csname num@plantvillage:R_taylor_P_30_r6:plantdoc_10:min\endcsname{50.8}
\expandafter\gdef\csname num@plantvillage:R_taylor_P_30_r6:plantdoc_10:max\endcsname{54.3}
\expandafter\gdef\csname num@plantvillage:R_taylor_P_30_r6:pct\endcsname{30.0}
\expandafter\gdef\csname num@plantvillage:R_taylor_P_30_r6:params\endcsname{3.35}
\expandafter\gdef\csname num@plantvillage:R_taylor_P_30_r6:macs_128\endcsname{0.251}
\expandafter\gdef\csname num@plantvillage:R_taylor_P_30_r6:macs_224\endcsname{0.770}
\expandafter\gdef\csname num@plantvillage:S_dacis_PMP_30_lam0:novel_1\endcsname{73.3}
\expandafter\gdef\csname num@plantvillage:S_dacis_PMP_30_lam0:novel_1:sd\endcsname{10.4}
\expandafter\gdef\csname num@plantvillage:S_dacis_PMP_30_lam0:novel_1:min\endcsname{58.5}
\expandafter\gdef\csname num@plantvillage:S_dacis_PMP_30_lam0:novel_1:max\endcsname{83.9}
\expandafter\gdef\csname num@plantvillage:S_dacis_PMP_30_lam0:novel_5\endcsname{86.9}
\expandafter\gdef\csname num@plantvillage:S_dacis_PMP_30_lam0:novel_5:sd\endcsname{7.2}
\expandafter\gdef\csname num@plantvillage:S_dacis_PMP_30_lam0:novel_5:min\endcsname{76.6}
\expandafter\gdef\csname num@plantvillage:S_dacis_PMP_30_lam0:novel_5:max\endcsname{93.8}
\expandafter\gdef\csname num@plantvillage:S_dacis_PMP_30_lam0:novel_10\endcsname{89.1}
\expandafter\gdef\csname num@plantvillage:S_dacis_PMP_30_lam0:novel_10:sd\endcsname{5.8}
\expandafter\gdef\csname num@plantvillage:S_dacis_PMP_30_lam0:novel_10:min\endcsname{81.0}
\expandafter\gdef\csname num@plantvillage:S_dacis_PMP_30_lam0:novel_10:max\endcsname{94.9}
\expandafter\gdef\csname num@plantvillage:S_dacis_PMP_30_lam0:plantdoc_1\endcsname{34.7}
\expandafter\gdef\csname num@plantvillage:S_dacis_PMP_30_lam0:plantdoc_1:sd\endcsname{1.4}
\expandafter\gdef\csname num@plantvillage:S_dacis_PMP_30_lam0:plantdoc_1:min\endcsname{32.5}
\expandafter\gdef\csname num@plantvillage:S_dacis_PMP_30_lam0:plantdoc_1:max\endcsname{35.9}
\expandafter\gdef\csname num@plantvillage:S_dacis_PMP_30_lam0:plantdoc_5\endcsname{46.9}
\expandafter\gdef\csname num@plantvillage:S_dacis_PMP_30_lam0:plantdoc_5:sd\endcsname{1.6}
\expandafter\gdef\csname num@plantvillage:S_dacis_PMP_30_lam0:plantdoc_5:min\endcsname{44.3}
\expandafter\gdef\csname num@plantvillage:S_dacis_PMP_30_lam0:plantdoc_5:max\endcsname{48.4}
\expandafter\gdef\csname num@plantvillage:S_dacis_PMP_30_lam0:plantdoc_10\endcsname{51.2}
\expandafter\gdef\csname num@plantvillage:S_dacis_PMP_30_lam0:plantdoc_10:sd\endcsname{1.5}
\expandafter\gdef\csname num@plantvillage:S_dacis_PMP_30_lam0:plantdoc_10:min\endcsname{49.0}
\expandafter\gdef\csname num@plantvillage:S_dacis_PMP_30_lam0:plantdoc_10:max\endcsname{52.6}
\expandafter\gdef\csname num@plantvillage:S_dacis_PMP_30_lam0:val_1\endcsname{74.3}
\expandafter\gdef\csname num@plantvillage:S_dacis_PMP_30_lam0:val_1:sd\endcsname{6.8}
\expandafter\gdef\csname num@plantvillage:S_dacis_PMP_30_lam0:val_1:min\endcsname{63.9}
\expandafter\gdef\csname num@plantvillage:S_dacis_PMP_30_lam0:val_1:max\endcsname{81.9}
\expandafter\gdef\csname num@plantvillage:S_dacis_PMP_30_lam0:val_5\endcsname{88.5}
\expandafter\gdef\csname num@plantvillage:S_dacis_PMP_30_lam0:val_5:sd\endcsname{4.2}
\expandafter\gdef\csname num@plantvillage:S_dacis_PMP_30_lam0:val_5:min\endcsname{81.6}
\expandafter\gdef\csname num@plantvillage:S_dacis_PMP_30_lam0:val_5:max\endcsname{91.6}
\expandafter\gdef\csname num@plantvillage:S_dacis_PMP_30_lam0:val_10\endcsname{90.7}
\expandafter\gdef\csname num@plantvillage:S_dacis_PMP_30_lam0:val_10:sd\endcsname{3.4}
\expandafter\gdef\csname num@plantvillage:S_dacis_PMP_30_lam0:val_10:min\endcsname{85.1}
\expandafter\gdef\csname num@plantvillage:S_dacis_PMP_30_lam0:val_10:max\endcsname{93.0}
\expandafter\gdef\csname num@plantvillage:S_dacis_PMP_30_lam0:pct\endcsname{30.0}
\expandafter\gdef\csname num@plantvillage:S_dacis_PMP_30_lam0:params\endcsname{3.35}
\expandafter\gdef\csname num@plantvillage:S_dacis_PMP_30_lam0:macs_128\endcsname{0.251}
\expandafter\gdef\csname num@plantvillage:S_dacis_PMP_30_lam0:macs_224\endcsname{0.770}
\expandafter\gdef\csname num@plantvillage:S_dacis_PMP_30_lam1:novel_1\endcsname{73.2}
\expandafter\gdef\csname num@plantvillage:S_dacis_PMP_30_lam1:novel_1:sd\endcsname{10.7}
\expandafter\gdef\csname num@plantvillage:S_dacis_PMP_30_lam1:novel_1:min\endcsname{59.0}
\expandafter\gdef\csname num@plantvillage:S_dacis_PMP_30_lam1:novel_1:max\endcsname{84.9}
\expandafter\gdef\csname num@plantvillage:S_dacis_PMP_30_lam1:novel_5\endcsname{86.9}
\expandafter\gdef\csname num@plantvillage:S_dacis_PMP_30_lam1:novel_5:sd\endcsname{7.5}
\expandafter\gdef\csname num@plantvillage:S_dacis_PMP_30_lam1:novel_5:min\endcsname{77.0}
\expandafter\gdef\csname num@plantvillage:S_dacis_PMP_30_lam1:novel_5:max\endcsname{94.9}
\expandafter\gdef\csname num@plantvillage:S_dacis_PMP_30_lam1:novel_10\endcsname{89.2}
\expandafter\gdef\csname num@plantvillage:S_dacis_PMP_30_lam1:novel_10:sd\endcsname{6.2}
\expandafter\gdef\csname num@plantvillage:S_dacis_PMP_30_lam1:novel_10:min\endcsname{81.2}
\expandafter\gdef\csname num@plantvillage:S_dacis_PMP_30_lam1:novel_10:max\endcsname{95.8}
\expandafter\gdef\csname num@plantvillage:S_dacis_PMP_30_lam1:plantdoc_1\endcsname{35.4}
\expandafter\gdef\csname num@plantvillage:S_dacis_PMP_30_lam1:plantdoc_1:sd\endcsname{1.3}
\expandafter\gdef\csname num@plantvillage:S_dacis_PMP_30_lam1:plantdoc_1:min\endcsname{33.9}
\expandafter\gdef\csname num@plantvillage:S_dacis_PMP_30_lam1:plantdoc_1:max\endcsname{37.5}
\expandafter\gdef\csname num@plantvillage:S_dacis_PMP_30_lam1:plantdoc_5\endcsname{47.8}
\expandafter\gdef\csname num@plantvillage:S_dacis_PMP_30_lam1:plantdoc_5:sd\endcsname{1.5}
\expandafter\gdef\csname num@plantvillage:S_dacis_PMP_30_lam1:plantdoc_5:min\endcsname{46.6}
\expandafter\gdef\csname num@plantvillage:S_dacis_PMP_30_lam1:plantdoc_5:max\endcsname{50.4}
\expandafter\gdef\csname num@plantvillage:S_dacis_PMP_30_lam1:plantdoc_10\endcsname{52.1}
\expandafter\gdef\csname num@plantvillage:S_dacis_PMP_30_lam1:plantdoc_10:sd\endcsname{1.4}
\expandafter\gdef\csname num@plantvillage:S_dacis_PMP_30_lam1:plantdoc_10:min\endcsname{51.2}
\expandafter\gdef\csname num@plantvillage:S_dacis_PMP_30_lam1:plantdoc_10:max\endcsname{54.5}
\expandafter\gdef\csname num@plantvillage:S_dacis_PMP_30_lam1:val_1\endcsname{73.7}
\expandafter\gdef\csname num@plantvillage:S_dacis_PMP_30_lam1:val_1:sd\endcsname{6.2}
\expandafter\gdef\csname num@plantvillage:S_dacis_PMP_30_lam1:val_1:min\endcsname{64.4}
\expandafter\gdef\csname num@plantvillage:S_dacis_PMP_30_lam1:val_1:max\endcsname{81.0}
\expandafter\gdef\csname num@plantvillage:S_dacis_PMP_30_lam1:val_5\endcsname{88.1}
\expandafter\gdef\csname num@plantvillage:S_dacis_PMP_30_lam1:val_5:sd\endcsname{4.0}
\expandafter\gdef\csname num@plantvillage:S_dacis_PMP_30_lam1:val_5:min\endcsname{81.6}
\expandafter\gdef\csname num@plantvillage:S_dacis_PMP_30_lam1:val_5:max\endcsname{91.1}
\expandafter\gdef\csname num@plantvillage:S_dacis_PMP_30_lam1:val_10\endcsname{90.3}
\expandafter\gdef\csname num@plantvillage:S_dacis_PMP_30_lam1:val_10:sd\endcsname{3.2}
\expandafter\gdef\csname num@plantvillage:S_dacis_PMP_30_lam1:val_10:min\endcsname{85.0}
\expandafter\gdef\csname num@plantvillage:S_dacis_PMP_30_lam1:val_10:max\endcsname{92.5}
\expandafter\gdef\csname num@plantvillage:S_dacis_PMP_30_lam1:pct\endcsname{30.0}
\expandafter\gdef\csname num@plantvillage:S_dacis_PMP_30_lam1:params\endcsname{3.35}
\expandafter\gdef\csname num@plantvillage:S_dacis_PMP_30_lam1:macs_128\endcsname{0.251}
\expandafter\gdef\csname num@plantvillage:S_dacis_PMP_30_lam1:macs_224\endcsname{0.770}
\expandafter\gdef\csname num@plantvillage:S_dacis_PMP_30_lam2:novel_1\endcsname{72.3}
\expandafter\gdef\csname num@plantvillage:S_dacis_PMP_30_lam2:novel_1:sd\endcsname{10.8}
\expandafter\gdef\csname num@plantvillage:S_dacis_PMP_30_lam2:novel_1:min\endcsname{58.5}
\expandafter\gdef\csname num@plantvillage:S_dacis_PMP_30_lam2:novel_1:max\endcsname{84.5}
\expandafter\gdef\csname num@plantvillage:S_dacis_PMP_30_lam2:novel_5\endcsname{86.2}
\expandafter\gdef\csname num@plantvillage:S_dacis_PMP_30_lam2:novel_5:sd\endcsname{7.7}
\expandafter\gdef\csname num@plantvillage:S_dacis_PMP_30_lam2:novel_5:min\endcsname{76.5}
\expandafter\gdef\csname num@plantvillage:S_dacis_PMP_30_lam2:novel_5:max\endcsname{94.6}
\expandafter\gdef\csname num@plantvillage:S_dacis_PMP_30_lam2:novel_10\endcsname{88.6}
\expandafter\gdef\csname num@plantvillage:S_dacis_PMP_30_lam2:novel_10:sd\endcsname{6.4}
\expandafter\gdef\csname num@plantvillage:S_dacis_PMP_30_lam2:novel_10:min\endcsname{81.0}
\expandafter\gdef\csname num@plantvillage:S_dacis_PMP_30_lam2:novel_10:max\endcsname{95.7}
\expandafter\gdef\csname num@plantvillage:S_dacis_PMP_30_lam2:plantdoc_1\endcsname{34.8}
\expandafter\gdef\csname num@plantvillage:S_dacis_PMP_30_lam2:plantdoc_1:sd\endcsname{1.5}
\expandafter\gdef\csname num@plantvillage:S_dacis_PMP_30_lam2:plantdoc_1:min\endcsname{32.5}
\expandafter\gdef\csname num@plantvillage:S_dacis_PMP_30_lam2:plantdoc_1:max\endcsname{36.4}
\expandafter\gdef\csname num@plantvillage:S_dacis_PMP_30_lam2:plantdoc_5\endcsname{47.0}
\expandafter\gdef\csname num@plantvillage:S_dacis_PMP_30_lam2:plantdoc_5:sd\endcsname{2.1}
\expandafter\gdef\csname num@plantvillage:S_dacis_PMP_30_lam2:plantdoc_5:min\endcsname{44.0}
\expandafter\gdef\csname num@plantvillage:S_dacis_PMP_30_lam2:plantdoc_5:max\endcsname{49.8}
\expandafter\gdef\csname num@plantvillage:S_dacis_PMP_30_lam2:plantdoc_10\endcsname{51.4}
\expandafter\gdef\csname num@plantvillage:S_dacis_PMP_30_lam2:plantdoc_10:sd\endcsname{2.0}
\expandafter\gdef\csname num@plantvillage:S_dacis_PMP_30_lam2:plantdoc_10:min\endcsname{48.3}
\expandafter\gdef\csname num@plantvillage:S_dacis_PMP_30_lam2:plantdoc_10:max\endcsname{54.0}
\expandafter\gdef\csname num@plantvillage:S_dacis_PMP_30_lam2:val_1\endcsname{73.6}
\expandafter\gdef\csname num@plantvillage:S_dacis_PMP_30_lam2:val_1:sd\endcsname{6.5}
\expandafter\gdef\csname num@plantvillage:S_dacis_PMP_30_lam2:val_1:min\endcsname{64.4}
\expandafter\gdef\csname num@plantvillage:S_dacis_PMP_30_lam2:val_1:max\endcsname{82.5}
\expandafter\gdef\csname num@plantvillage:S_dacis_PMP_30_lam2:val_5\endcsname{87.9}
\expandafter\gdef\csname num@plantvillage:S_dacis_PMP_30_lam2:val_5:sd\endcsname{4.5}
\expandafter\gdef\csname num@plantvillage:S_dacis_PMP_30_lam2:val_5:min\endcsname{80.7}
\expandafter\gdef\csname num@plantvillage:S_dacis_PMP_30_lam2:val_5:max\endcsname{92.6}
\expandafter\gdef\csname num@plantvillage:S_dacis_PMP_30_lam2:val_10\endcsname{90.1}
\expandafter\gdef\csname num@plantvillage:S_dacis_PMP_30_lam2:val_10:sd\endcsname{3.8}
\expandafter\gdef\csname num@plantvillage:S_dacis_PMP_30_lam2:val_10:min\endcsname{83.9}
\expandafter\gdef\csname num@plantvillage:S_dacis_PMP_30_lam2:val_10:max\endcsname{93.6}
\expandafter\gdef\csname num@plantvillage:S_dacis_PMP_30_lam2:pct\endcsname{30.0}
\expandafter\gdef\csname num@plantvillage:S_dacis_PMP_30_lam2:params\endcsname{3.35}
\expandafter\gdef\csname num@plantvillage:S_dacis_PMP_30_lam2:macs_128\endcsname{0.251}
\expandafter\gdef\csname num@plantvillage:S_dacis_PMP_30_lam2:macs_224\endcsname{0.770}
\expandafter\gdef\csname num@plantvillage:S_dacis_PMP_30_lam3:novel_1\endcsname{72.7}
\expandafter\gdef\csname num@plantvillage:S_dacis_PMP_30_lam3:novel_1:sd\endcsname{10.9}
\expandafter\gdef\csname num@plantvillage:S_dacis_PMP_30_lam3:novel_1:min\endcsname{58.5}
\expandafter\gdef\csname num@plantvillage:S_dacis_PMP_30_lam3:novel_1:max\endcsname{84.8}
\expandafter\gdef\csname num@plantvillage:S_dacis_PMP_30_lam3:novel_5\endcsname{86.5}
\expandafter\gdef\csname num@plantvillage:S_dacis_PMP_30_lam3:novel_5:sd\endcsname{7.3}
\expandafter\gdef\csname num@plantvillage:S_dacis_PMP_30_lam3:novel_5:min\endcsname{77.2}
\expandafter\gdef\csname num@plantvillage:S_dacis_PMP_30_lam3:novel_5:max\endcsname{94.3}
\expandafter\gdef\csname num@plantvillage:S_dacis_PMP_30_lam3:novel_10\endcsname{88.7}
\expandafter\gdef\csname num@plantvillage:S_dacis_PMP_30_lam3:novel_10:sd\endcsname{6.0}
\expandafter\gdef\csname num@plantvillage:S_dacis_PMP_30_lam3:novel_10:min\endcsname{81.5}
\expandafter\gdef\csname num@plantvillage:S_dacis_PMP_30_lam3:novel_10:max\endcsname{95.3}
\expandafter\gdef\csname num@plantvillage:S_dacis_PMP_30_lam3:plantdoc_1\endcsname{35.3}
\expandafter\gdef\csname num@plantvillage:S_dacis_PMP_30_lam3:plantdoc_1:sd\endcsname{1.2}
\expandafter\gdef\csname num@plantvillage:S_dacis_PMP_30_lam3:plantdoc_1:min\endcsname{33.3}
\expandafter\gdef\csname num@plantvillage:S_dacis_PMP_30_lam3:plantdoc_1:max\endcsname{36.3}
\expandafter\gdef\csname num@plantvillage:S_dacis_PMP_30_lam3:plantdoc_5\endcsname{47.3}
\expandafter\gdef\csname num@plantvillage:S_dacis_PMP_30_lam3:plantdoc_5:sd\endcsname{1.6}
\expandafter\gdef\csname num@plantvillage:S_dacis_PMP_30_lam3:plantdoc_5:min\endcsname{44.8}
\expandafter\gdef\csname num@plantvillage:S_dacis_PMP_30_lam3:plantdoc_5:max\endcsname{48.9}
\expandafter\gdef\csname num@plantvillage:S_dacis_PMP_30_lam3:plantdoc_10\endcsname{51.7}
\expandafter\gdef\csname num@plantvillage:S_dacis_PMP_30_lam3:plantdoc_10:sd\endcsname{1.5}
\expandafter\gdef\csname num@plantvillage:S_dacis_PMP_30_lam3:plantdoc_10:min\endcsname{49.5}
\expandafter\gdef\csname num@plantvillage:S_dacis_PMP_30_lam3:plantdoc_10:max\endcsname{53.2}
\expandafter\gdef\csname num@plantvillage:S_dacis_PMP_30_lam3:val_1\endcsname{74.3}
\expandafter\gdef\csname num@plantvillage:S_dacis_PMP_30_lam3:val_1:sd\endcsname{6.0}
\expandafter\gdef\csname num@plantvillage:S_dacis_PMP_30_lam3:val_1:min\endcsname{65.4}
\expandafter\gdef\csname num@plantvillage:S_dacis_PMP_30_lam3:val_1:max\endcsname{81.3}
\expandafter\gdef\csname num@plantvillage:S_dacis_PMP_30_lam3:val_5\endcsname{88.5}
\expandafter\gdef\csname num@plantvillage:S_dacis_PMP_30_lam3:val_5:sd\endcsname{4.0}
\expandafter\gdef\csname num@plantvillage:S_dacis_PMP_30_lam3:val_5:min\endcsname{82.4}
\expandafter\gdef\csname num@plantvillage:S_dacis_PMP_30_lam3:val_5:max\endcsname{91.9}
\expandafter\gdef\csname num@plantvillage:S_dacis_PMP_30_lam3:val_10\endcsname{90.6}
\expandafter\gdef\csname num@plantvillage:S_dacis_PMP_30_lam3:val_10:sd\endcsname{3.2}
\expandafter\gdef\csname num@plantvillage:S_dacis_PMP_30_lam3:val_10:min\endcsname{85.9}
\expandafter\gdef\csname num@plantvillage:S_dacis_PMP_30_lam3:val_10:max\endcsname{93.2}
\expandafter\gdef\csname num@plantvillage:S_dacis_PMP_30_lam3:pct\endcsname{30.0}
\expandafter\gdef\csname num@plantvillage:S_dacis_PMP_30_lam3:params\endcsname{3.35}
\expandafter\gdef\csname num@plantvillage:S_dacis_PMP_30_lam3:macs_128\endcsname{0.251}
\expandafter\gdef\csname num@plantvillage:S_dacis_PMP_30_lam3:macs_224\endcsname{0.770}
\expandafter\gdef\csname num@plantvillage:S_dacis_PMP_30_lam4:novel_1\endcsname{72.4}
\expandafter\gdef\csname num@plantvillage:S_dacis_PMP_30_lam4:novel_1:sd\endcsname{10.7}
\expandafter\gdef\csname num@plantvillage:S_dacis_PMP_30_lam4:novel_1:min\endcsname{57.5}
\expandafter\gdef\csname num@plantvillage:S_dacis_PMP_30_lam4:novel_1:max\endcsname{83.0}
\expandafter\gdef\csname num@plantvillage:S_dacis_PMP_30_lam4:novel_5\endcsname{86.2}
\expandafter\gdef\csname num@plantvillage:S_dacis_PMP_30_lam4:novel_5:sd\endcsname{7.8}
\expandafter\gdef\csname num@plantvillage:S_dacis_PMP_30_lam4:novel_5:min\endcsname{74.9}
\expandafter\gdef\csname num@plantvillage:S_dacis_PMP_30_lam4:novel_5:max\endcsname{93.6}
\expandafter\gdef\csname num@plantvillage:S_dacis_PMP_30_lam4:novel_10\endcsname{88.5}
\expandafter\gdef\csname num@plantvillage:S_dacis_PMP_30_lam4:novel_10:sd\endcsname{6.5}
\expandafter\gdef\csname num@plantvillage:S_dacis_PMP_30_lam4:novel_10:min\endcsname{79.4}
\expandafter\gdef\csname num@plantvillage:S_dacis_PMP_30_lam4:novel_10:max\endcsname{95.0}
\expandafter\gdef\csname num@plantvillage:S_dacis_PMP_30_lam4:plantdoc_1\endcsname{35.2}
\expandafter\gdef\csname num@plantvillage:S_dacis_PMP_30_lam4:plantdoc_1:sd\endcsname{1.5}
\expandafter\gdef\csname num@plantvillage:S_dacis_PMP_30_lam4:plantdoc_1:min\endcsname{33.3}
\expandafter\gdef\csname num@plantvillage:S_dacis_PMP_30_lam4:plantdoc_1:max\endcsname{37.4}
\expandafter\gdef\csname num@plantvillage:S_dacis_PMP_30_lam4:plantdoc_5\endcsname{47.3}
\expandafter\gdef\csname num@plantvillage:S_dacis_PMP_30_lam4:plantdoc_5:sd\endcsname{2.0}
\expandafter\gdef\csname num@plantvillage:S_dacis_PMP_30_lam4:plantdoc_5:min\endcsname{45.2}
\expandafter\gdef\csname num@plantvillage:S_dacis_PMP_30_lam4:plantdoc_5:max\endcsname{50.4}
\expandafter\gdef\csname num@plantvillage:S_dacis_PMP_30_lam4:plantdoc_10\endcsname{51.7}
\expandafter\gdef\csname num@plantvillage:S_dacis_PMP_30_lam4:plantdoc_10:sd\endcsname{2.0}
\expandafter\gdef\csname num@plantvillage:S_dacis_PMP_30_lam4:plantdoc_10:min\endcsname{49.8}
\expandafter\gdef\csname num@plantvillage:S_dacis_PMP_30_lam4:plantdoc_10:max\endcsname{54.8}
\expandafter\gdef\csname num@plantvillage:S_dacis_PMP_30_lam4:val_1\endcsname{74.3}
\expandafter\gdef\csname num@plantvillage:S_dacis_PMP_30_lam4:val_1:sd\endcsname{6.9}
\expandafter\gdef\csname num@plantvillage:S_dacis_PMP_30_lam4:val_1:min\endcsname{64.3}
\expandafter\gdef\csname num@plantvillage:S_dacis_PMP_30_lam4:val_1:max\endcsname{82.6}
\expandafter\gdef\csname num@plantvillage:S_dacis_PMP_30_lam4:val_5\endcsname{88.9}
\expandafter\gdef\csname num@plantvillage:S_dacis_PMP_30_lam4:val_5:sd\endcsname{4.0}
\expandafter\gdef\csname num@plantvillage:S_dacis_PMP_30_lam4:val_5:min\endcsname{83.0}
\expandafter\gdef\csname num@plantvillage:S_dacis_PMP_30_lam4:val_5:max\endcsname{92.4}
\expandafter\gdef\csname num@plantvillage:S_dacis_PMP_30_lam4:val_10\endcsname{91.2}
\expandafter\gdef\csname num@plantvillage:S_dacis_PMP_30_lam4:val_10:sd\endcsname{3.0}
\expandafter\gdef\csname num@plantvillage:S_dacis_PMP_30_lam4:val_10:min\endcsname{86.7}
\expandafter\gdef\csname num@plantvillage:S_dacis_PMP_30_lam4:val_10:max\endcsname{94.1}
\expandafter\gdef\csname num@plantvillage:S_dacis_PMP_30_lam4:pct\endcsname{30.0}
\expandafter\gdef\csname num@plantvillage:S_dacis_PMP_30_lam4:params\endcsname{3.35}
\expandafter\gdef\csname num@plantvillage:S_dacis_PMP_30_lam4:macs_128\endcsname{0.251}
\expandafter\gdef\csname num@plantvillage:S_dacis_PMP_30_lam4:macs_224\endcsname{0.770}
\expandafter\gdef\csname num@plantvillage:S_dacis_PMP_30_lam5:novel_1\endcsname{72.3}
\expandafter\gdef\csname num@plantvillage:S_dacis_PMP_30_lam5:novel_1:sd\endcsname{10.9}
\expandafter\gdef\csname num@plantvillage:S_dacis_PMP_30_lam5:novel_1:min\endcsname{56.8}
\expandafter\gdef\csname num@plantvillage:S_dacis_PMP_30_lam5:novel_1:max\endcsname{83.6}
\expandafter\gdef\csname num@plantvillage:S_dacis_PMP_30_lam5:novel_5\endcsname{86.1}
\expandafter\gdef\csname num@plantvillage:S_dacis_PMP_30_lam5:novel_5:sd\endcsname{7.8}
\expandafter\gdef\csname num@plantvillage:S_dacis_PMP_30_lam5:novel_5:min\endcsname{74.8}
\expandafter\gdef\csname num@plantvillage:S_dacis_PMP_30_lam5:novel_5:max\endcsname{94.2}
\expandafter\gdef\csname num@plantvillage:S_dacis_PMP_30_lam5:novel_10\endcsname{88.4}
\expandafter\gdef\csname num@plantvillage:S_dacis_PMP_30_lam5:novel_10:sd\endcsname{6.6}
\expandafter\gdef\csname num@plantvillage:S_dacis_PMP_30_lam5:novel_10:min\endcsname{79.1}
\expandafter\gdef\csname num@plantvillage:S_dacis_PMP_30_lam5:novel_10:max\endcsname{95.4}
\expandafter\gdef\csname num@plantvillage:S_dacis_PMP_30_lam5:plantdoc_1\endcsname{34.7}
\expandafter\gdef\csname num@plantvillage:S_dacis_PMP_30_lam5:plantdoc_1:sd\endcsname{1.3}
\expandafter\gdef\csname num@plantvillage:S_dacis_PMP_30_lam5:plantdoc_1:min\endcsname{32.8}
\expandafter\gdef\csname num@plantvillage:S_dacis_PMP_30_lam5:plantdoc_1:max\endcsname{36.3}
\expandafter\gdef\csname num@plantvillage:S_dacis_PMP_30_lam5:plantdoc_5\endcsname{46.8}
\expandafter\gdef\csname num@plantvillage:S_dacis_PMP_30_lam5:plantdoc_5:sd\endcsname{1.3}
\expandafter\gdef\csname num@plantvillage:S_dacis_PMP_30_lam5:plantdoc_5:min\endcsname{45.0}
\expandafter\gdef\csname num@plantvillage:S_dacis_PMP_30_lam5:plantdoc_5:max\endcsname{48.6}
\expandafter\gdef\csname num@plantvillage:S_dacis_PMP_30_lam5:plantdoc_10\endcsname{51.0}
\expandafter\gdef\csname num@plantvillage:S_dacis_PMP_30_lam5:plantdoc_10:sd\endcsname{1.3}
\expandafter\gdef\csname num@plantvillage:S_dacis_PMP_30_lam5:plantdoc_10:min\endcsname{49.4}
\expandafter\gdef\csname num@plantvillage:S_dacis_PMP_30_lam5:plantdoc_10:max\endcsname{52.8}
\expandafter\gdef\csname num@plantvillage:S_dacis_PMP_30_lam5:val_1\endcsname{74.3}
\expandafter\gdef\csname num@plantvillage:S_dacis_PMP_30_lam5:val_1:sd\endcsname{6.3}
\expandafter\gdef\csname num@plantvillage:S_dacis_PMP_30_lam5:val_1:min\endcsname{66.2}
\expandafter\gdef\csname num@plantvillage:S_dacis_PMP_30_lam5:val_1:max\endcsname{83.6}
\expandafter\gdef\csname num@plantvillage:S_dacis_PMP_30_lam5:val_5\endcsname{88.2}
\expandafter\gdef\csname num@plantvillage:S_dacis_PMP_30_lam5:val_5:sd\endcsname{4.2}
\expandafter\gdef\csname num@plantvillage:S_dacis_PMP_30_lam5:val_5:min\endcsname{81.7}
\expandafter\gdef\csname num@plantvillage:S_dacis_PMP_30_lam5:val_5:max\endcsname{92.5}
\expandafter\gdef\csname num@plantvillage:S_dacis_PMP_30_lam5:val_10\endcsname{90.4}
\expandafter\gdef\csname num@plantvillage:S_dacis_PMP_30_lam5:val_10:sd\endcsname{3.5}
\expandafter\gdef\csname num@plantvillage:S_dacis_PMP_30_lam5:val_10:min\endcsname{84.6}
\expandafter\gdef\csname num@plantvillage:S_dacis_PMP_30_lam5:val_10:max\endcsname{93.4}
\expandafter\gdef\csname num@plantvillage:S_dacis_PMP_30_lam5:pct\endcsname{30.0}
\expandafter\gdef\csname num@plantvillage:S_dacis_PMP_30_lam5:params\endcsname{3.35}
\expandafter\gdef\csname num@plantvillage:S_dacis_PMP_30_lam5:macs_128\endcsname{0.251}
\expandafter\gdef\csname num@plantvillage:S_dacis_PMP_30_lam5:macs_224\endcsname{0.770}
\expandafter\gdef\csname num@plantvillage:S_dacis_PMP_30_lam6:novel_1\endcsname{71.8}
\expandafter\gdef\csname num@plantvillage:S_dacis_PMP_30_lam6:novel_1:sd\endcsname{10.2}
\expandafter\gdef\csname num@plantvillage:S_dacis_PMP_30_lam6:novel_1:min\endcsname{56.8}
\expandafter\gdef\csname num@plantvillage:S_dacis_PMP_30_lam6:novel_1:max\endcsname{82.6}
\expandafter\gdef\csname num@plantvillage:S_dacis_PMP_30_lam6:novel_5\endcsname{85.9}
\expandafter\gdef\csname num@plantvillage:S_dacis_PMP_30_lam6:novel_5:sd\endcsname{7.8}
\expandafter\gdef\csname num@plantvillage:S_dacis_PMP_30_lam6:novel_5:min\endcsname{74.2}
\expandafter\gdef\csname num@plantvillage:S_dacis_PMP_30_lam6:novel_5:max\endcsname{93.6}
\expandafter\gdef\csname num@plantvillage:S_dacis_PMP_30_lam6:novel_10\endcsname{88.2}
\expandafter\gdef\csname num@plantvillage:S_dacis_PMP_30_lam6:novel_10:sd\endcsname{6.6}
\expandafter\gdef\csname num@plantvillage:S_dacis_PMP_30_lam6:novel_10:min\endcsname{78.6}
\expandafter\gdef\csname num@plantvillage:S_dacis_PMP_30_lam6:novel_10:max\endcsname{95.0}
\expandafter\gdef\csname num@plantvillage:S_dacis_PMP_30_lam6:plantdoc_1\endcsname{34.5}
\expandafter\gdef\csname num@plantvillage:S_dacis_PMP_30_lam6:plantdoc_1:sd\endcsname{0.8}
\expandafter\gdef\csname num@plantvillage:S_dacis_PMP_30_lam6:plantdoc_1:min\endcsname{33.4}
\expandafter\gdef\csname num@plantvillage:S_dacis_PMP_30_lam6:plantdoc_1:max\endcsname{35.5}
\expandafter\gdef\csname num@plantvillage:S_dacis_PMP_30_lam6:plantdoc_5\endcsname{46.6}
\expandafter\gdef\csname num@plantvillage:S_dacis_PMP_30_lam6:plantdoc_5:sd\endcsname{0.6}
\expandafter\gdef\csname num@plantvillage:S_dacis_PMP_30_lam6:plantdoc_5:min\endcsname{45.9}
\expandafter\gdef\csname num@plantvillage:S_dacis_PMP_30_lam6:plantdoc_5:max\endcsname{47.6}
\expandafter\gdef\csname num@plantvillage:S_dacis_PMP_30_lam6:plantdoc_10\endcsname{50.9}
\expandafter\gdef\csname num@plantvillage:S_dacis_PMP_30_lam6:plantdoc_10:sd\endcsname{0.6}
\expandafter\gdef\csname num@plantvillage:S_dacis_PMP_30_lam6:plantdoc_10:min\endcsname{50.3}
\expandafter\gdef\csname num@plantvillage:S_dacis_PMP_30_lam6:plantdoc_10:max\endcsname{51.8}
\expandafter\gdef\csname num@plantvillage:S_dacis_PMP_30_lam6:val_1\endcsname{74.1}
\expandafter\gdef\csname num@plantvillage:S_dacis_PMP_30_lam6:val_1:sd\endcsname{7.4}
\expandafter\gdef\csname num@plantvillage:S_dacis_PMP_30_lam6:val_1:min\endcsname{65.4}
\expandafter\gdef\csname num@plantvillage:S_dacis_PMP_30_lam6:val_1:max\endcsname{83.6}
\expandafter\gdef\csname num@plantvillage:S_dacis_PMP_30_lam6:val_5\endcsname{88.0}
\expandafter\gdef\csname num@plantvillage:S_dacis_PMP_30_lam6:val_5:sd\endcsname{4.6}
\expandafter\gdef\csname num@plantvillage:S_dacis_PMP_30_lam6:val_5:min\endcsname{82.0}
\expandafter\gdef\csname num@plantvillage:S_dacis_PMP_30_lam6:val_5:max\endcsname{92.5}
\expandafter\gdef\csname num@plantvillage:S_dacis_PMP_30_lam6:val_10\endcsname{90.1}
\expandafter\gdef\csname num@plantvillage:S_dacis_PMP_30_lam6:val_10:sd\endcsname{3.6}
\expandafter\gdef\csname num@plantvillage:S_dacis_PMP_30_lam6:val_10:min\endcsname{85.2}
\expandafter\gdef\csname num@plantvillage:S_dacis_PMP_30_lam6:val_10:max\endcsname{93.3}
\expandafter\gdef\csname num@plantvillage:S_dacis_PMP_30_lam6:pct\endcsname{30.0}
\expandafter\gdef\csname num@plantvillage:S_dacis_PMP_30_lam6:params\endcsname{3.35}
\expandafter\gdef\csname num@plantvillage:S_dacis_PMP_30_lam6:macs_128\endcsname{0.251}
\expandafter\gdef\csname num@plantvillage:S_dacis_PMP_30_lam6:macs_224\endcsname{0.770}
\expandafter\gdef\csname num@plantvillage:S_dacis_PMP_30_lam7:novel_1\endcsname{73.0}
\expandafter\gdef\csname num@plantvillage:S_dacis_PMP_30_lam7:novel_1:sd\endcsname{10.8}
\expandafter\gdef\csname num@plantvillage:S_dacis_PMP_30_lam7:novel_1:min\endcsname{59.0}
\expandafter\gdef\csname num@plantvillage:S_dacis_PMP_30_lam7:novel_1:max\endcsname{85.4}
\expandafter\gdef\csname num@plantvillage:S_dacis_PMP_30_lam7:novel_5\endcsname{86.9}
\expandafter\gdef\csname num@plantvillage:S_dacis_PMP_30_lam7:novel_5:sd\endcsname{7.0}
\expandafter\gdef\csname num@plantvillage:S_dacis_PMP_30_lam7:novel_5:min\endcsname{78.1}
\expandafter\gdef\csname num@plantvillage:S_dacis_PMP_30_lam7:novel_5:max\endcsname{94.9}
\expandafter\gdef\csname num@plantvillage:S_dacis_PMP_30_lam7:novel_10\endcsname{89.2}
\expandafter\gdef\csname num@plantvillage:S_dacis_PMP_30_lam7:novel_10:sd\endcsname{5.7}
\expandafter\gdef\csname num@plantvillage:S_dacis_PMP_30_lam7:novel_10:min\endcsname{82.4}
\expandafter\gdef\csname num@plantvillage:S_dacis_PMP_30_lam7:novel_10:max\endcsname{95.9}
\expandafter\gdef\csname num@plantvillage:S_dacis_PMP_30_lam7:plantdoc_1\endcsname{35.4}
\expandafter\gdef\csname num@plantvillage:S_dacis_PMP_30_lam7:plantdoc_1:sd\endcsname{1.3}
\expandafter\gdef\csname num@plantvillage:S_dacis_PMP_30_lam7:plantdoc_1:min\endcsname{33.6}
\expandafter\gdef\csname num@plantvillage:S_dacis_PMP_30_lam7:plantdoc_1:max\endcsname{36.4}
\expandafter\gdef\csname num@plantvillage:S_dacis_PMP_30_lam7:plantdoc_5\endcsname{47.7}
\expandafter\gdef\csname num@plantvillage:S_dacis_PMP_30_lam7:plantdoc_5:sd\endcsname{1.5}
\expandafter\gdef\csname num@plantvillage:S_dacis_PMP_30_lam7:plantdoc_5:min\endcsname{45.6}
\expandafter\gdef\csname num@plantvillage:S_dacis_PMP_30_lam7:plantdoc_5:max\endcsname{49.5}
\expandafter\gdef\csname num@plantvillage:S_dacis_PMP_30_lam7:plantdoc_10\endcsname{52.0}
\expandafter\gdef\csname num@plantvillage:S_dacis_PMP_30_lam7:plantdoc_10:sd\endcsname{1.3}
\expandafter\gdef\csname num@plantvillage:S_dacis_PMP_30_lam7:plantdoc_10:min\endcsname{50.2}
\expandafter\gdef\csname num@plantvillage:S_dacis_PMP_30_lam7:plantdoc_10:max\endcsname{53.6}
\expandafter\gdef\csname num@plantvillage:S_dacis_PMP_30_lam7:val_1\endcsname{74.1}
\expandafter\gdef\csname num@plantvillage:S_dacis_PMP_30_lam7:val_1:sd\endcsname{5.7}
\expandafter\gdef\csname num@plantvillage:S_dacis_PMP_30_lam7:val_1:min\endcsname{66.2}
\expandafter\gdef\csname num@plantvillage:S_dacis_PMP_30_lam7:val_1:max\endcsname{81.2}
\expandafter\gdef\csname num@plantvillage:S_dacis_PMP_30_lam7:val_5\endcsname{88.8}
\expandafter\gdef\csname num@plantvillage:S_dacis_PMP_30_lam7:val_5:sd\endcsname{3.4}
\expandafter\gdef\csname num@plantvillage:S_dacis_PMP_30_lam7:val_5:min\endcsname{84.0}
\expandafter\gdef\csname num@plantvillage:S_dacis_PMP_30_lam7:val_5:max\endcsname{92.2}
\expandafter\gdef\csname num@plantvillage:S_dacis_PMP_30_lam7:val_10\endcsname{91.0}
\expandafter\gdef\csname num@plantvillage:S_dacis_PMP_30_lam7:val_10:sd\endcsname{2.7}
\expandafter\gdef\csname num@plantvillage:S_dacis_PMP_30_lam7:val_10:min\endcsname{87.3}
\expandafter\gdef\csname num@plantvillage:S_dacis_PMP_30_lam7:val_10:max\endcsname{93.5}
\expandafter\gdef\csname num@plantvillage:S_dacis_PMP_30_lam7:pct\endcsname{30.0}
\expandafter\gdef\csname num@plantvillage:S_dacis_PMP_30_lam7:params\endcsname{3.35}
\expandafter\gdef\csname num@plantvillage:S_dacis_PMP_30_lam7:macs_128\endcsname{0.251}
\expandafter\gdef\csname num@plantvillage:S_dacis_PMP_30_lam7:macs_224\endcsname{0.770}
\expandafter\gdef\csname num@plantvillage:dense_ce:novel_1\endcsname{79.5}
\expandafter\gdef\csname num@plantvillage:dense_ce:novel_1:sd\endcsname{9.6}
\expandafter\gdef\csname num@plantvillage:dense_ce:novel_1:min\endcsname{65.8}
\expandafter\gdef\csname num@plantvillage:dense_ce:novel_1:max\endcsname{90.0}
\expandafter\gdef\csname num@plantvillage:dense_ce:novel_5\endcsname{91.4}
\expandafter\gdef\csname num@plantvillage:dense_ce:novel_5:sd\endcsname{5.6}
\expandafter\gdef\csname num@plantvillage:dense_ce:novel_5:min\endcsname{83.7}
\expandafter\gdef\csname num@plantvillage:dense_ce:novel_5:max\endcsname{97.3}
\expandafter\gdef\csname num@plantvillage:dense_ce:novel_10\endcsname{93.1}
\expandafter\gdef\csname num@plantvillage:dense_ce:novel_10:sd\endcsname{4.4}
\expandafter\gdef\csname num@plantvillage:dense_ce:novel_10:min\endcsname{87.1}
\expandafter\gdef\csname num@plantvillage:dense_ce:novel_10:max\endcsname{97.9}
\expandafter\gdef\csname num@plantvillage:dense_ce:plantdoc_1\endcsname{41.1}
\expandafter\gdef\csname num@plantvillage:dense_ce:plantdoc_1:sd\endcsname{1.2}
\expandafter\gdef\csname num@plantvillage:dense_ce:plantdoc_1:min\endcsname{39.8}
\expandafter\gdef\csname num@plantvillage:dense_ce:plantdoc_1:max\endcsname{42.4}
\expandafter\gdef\csname num@plantvillage:dense_ce:plantdoc_5\endcsname{55.4}
\expandafter\gdef\csname num@plantvillage:dense_ce:plantdoc_5:sd\endcsname{1.3}
\expandafter\gdef\csname num@plantvillage:dense_ce:plantdoc_5:min\endcsname{54.4}
\expandafter\gdef\csname num@plantvillage:dense_ce:plantdoc_5:max\endcsname{57.4}
\expandafter\gdef\csname num@plantvillage:dense_ce:plantdoc_10\endcsname{60.1}
\expandafter\gdef\csname num@plantvillage:dense_ce:plantdoc_10:sd\endcsname{1.2}
\expandafter\gdef\csname num@plantvillage:dense_ce:plantdoc_10:min\endcsname{58.6}
\expandafter\gdef\csname num@plantvillage:dense_ce:plantdoc_10:max\endcsname{61.8}
\expandafter\gdef\csname num@plantvillage:dense_ce:pct\endcsname{100.0}
\expandafter\gdef\csname num@plantvillage:dense_ce:params\endcsname{11.18}
\expandafter\gdef\csname num@plantvillage:dense_ce:macs_128\endcsname{0.592}
\expandafter\gdef\csname num@plantvillage:dense_ce:macs_224\endcsname{1.814}
\expandafter\gdef\csname num@plantvillage:dense_ce:cpu_1thread_b1_ms\endcsname{10.9}
\expandafter\gdef\csname num@plantvillage:dense_ce:cpu_all_threads_b1_ms\endcsname{3.9}
\expandafter\gdef\csname num@plantvillage:dense_ce:gpu_b1_ms\endcsname{0.9}
\expandafter\gdef\csname num@plantvillage:dense_ce:gpu_ips\endcsname{36430}
\expandafter\gdef\csname num@plantvillage:dense_m:novel_1\endcsname{79.5}
\expandafter\gdef\csname num@plantvillage:dense_m:novel_1:sd\endcsname{9.8}
\expandafter\gdef\csname num@plantvillage:dense_m:novel_1:min\endcsname{65.8}
\expandafter\gdef\csname num@plantvillage:dense_m:novel_1:max\endcsname{90.0}
\expandafter\gdef\csname num@plantvillage:dense_m:novel_5\endcsname{91.2}
\expandafter\gdef\csname num@plantvillage:dense_m:novel_5:sd\endcsname{5.7}
\expandafter\gdef\csname num@plantvillage:dense_m:novel_5:min\endcsname{83.6}
\expandafter\gdef\csname num@plantvillage:dense_m:novel_5:max\endcsname{97.1}
\expandafter\gdef\csname num@plantvillage:dense_m:novel_10\endcsname{92.8}
\expandafter\gdef\csname num@plantvillage:dense_m:novel_10:sd\endcsname{4.5}
\expandafter\gdef\csname num@plantvillage:dense_m:novel_10:min\endcsname{86.9}
\expandafter\gdef\csname num@plantvillage:dense_m:novel_10:max\endcsname{97.6}
\expandafter\gdef\csname num@plantvillage:dense_m:plantdoc_1\endcsname{41.2}
\expandafter\gdef\csname num@plantvillage:dense_m:plantdoc_1:sd\endcsname{1.2}
\expandafter\gdef\csname num@plantvillage:dense_m:plantdoc_1:min\endcsname{39.9}
\expandafter\gdef\csname num@plantvillage:dense_m:plantdoc_1:max\endcsname{42.7}
\expandafter\gdef\csname num@plantvillage:dense_m:plantdoc_5\endcsname{55.4}
\expandafter\gdef\csname num@plantvillage:dense_m:plantdoc_5:sd\endcsname{1.2}
\expandafter\gdef\csname num@plantvillage:dense_m:plantdoc_5:min\endcsname{54.4}
\expandafter\gdef\csname num@plantvillage:dense_m:plantdoc_5:max\endcsname{57.1}
\expandafter\gdef\csname num@plantvillage:dense_m:plantdoc_10\endcsname{59.9}
\expandafter\gdef\csname num@plantvillage:dense_m:plantdoc_10:sd\endcsname{1.2}
\expandafter\gdef\csname num@plantvillage:dense_m:plantdoc_10:min\endcsname{58.5}
\expandafter\gdef\csname num@plantvillage:dense_m:plantdoc_10:max\endcsname{61.5}
\expandafter\gdef\csname num@plantvillage:dense_m:pct\endcsname{100.0}
\expandafter\gdef\csname num@plantvillage:dense_m:params\endcsname{11.18}
\expandafter\gdef\csname num@plantvillage:dense_m:macs_128\endcsname{0.592}
\expandafter\gdef\csname num@plantvillage:dense_m:macs_224\endcsname{1.814}
\expandafter\gdef\csname num@plantvillage:dense_rmr:novel_1\endcsname{79.2}
\expandafter\gdef\csname num@plantvillage:dense_rmr:novel_1:sd\endcsname{9.5}
\expandafter\gdef\csname num@plantvillage:dense_rmr:novel_1:min\endcsname{65.2}
\expandafter\gdef\csname num@plantvillage:dense_rmr:novel_1:max\endcsname{89.3}
\expandafter\gdef\csname num@plantvillage:dense_rmr:novel_5\endcsname{91.3}
\expandafter\gdef\csname num@plantvillage:dense_rmr:novel_5:sd\endcsname{5.6}
\expandafter\gdef\csname num@plantvillage:dense_rmr:novel_5:min\endcsname{83.4}
\expandafter\gdef\csname num@plantvillage:dense_rmr:novel_5:max\endcsname{97.0}
\expandafter\gdef\csname num@plantvillage:dense_rmr:novel_10\endcsname{93.0}
\expandafter\gdef\csname num@plantvillage:dense_rmr:novel_10:sd\endcsname{4.3}
\expandafter\gdef\csname num@plantvillage:dense_rmr:novel_10:min\endcsname{87.2}
\expandafter\gdef\csname num@plantvillage:dense_rmr:novel_10:max\endcsname{97.7}
\expandafter\gdef\csname num@plantvillage:dense_rmr:plantdoc_1\endcsname{41.1}
\expandafter\gdef\csname num@plantvillage:dense_rmr:plantdoc_1:sd\endcsname{1.0}
\expandafter\gdef\csname num@plantvillage:dense_rmr:plantdoc_1:min\endcsname{40.1}
\expandafter\gdef\csname num@plantvillage:dense_rmr:plantdoc_1:max\endcsname{42.4}
\expandafter\gdef\csname num@plantvillage:dense_rmr:plantdoc_5\endcsname{55.5}
\expandafter\gdef\csname num@plantvillage:dense_rmr:plantdoc_5:sd\endcsname{1.1}
\expandafter\gdef\csname num@plantvillage:dense_rmr:plantdoc_5:min\endcsname{54.5}
\expandafter\gdef\csname num@plantvillage:dense_rmr:plantdoc_5:max\endcsname{57.0}
\expandafter\gdef\csname num@plantvillage:dense_rmr:plantdoc_10\endcsname{60.2}
\expandafter\gdef\csname num@plantvillage:dense_rmr:plantdoc_10:sd\endcsname{1.1}
\expandafter\gdef\csname num@plantvillage:dense_rmr:plantdoc_10:min\endcsname{58.7}
\expandafter\gdef\csname num@plantvillage:dense_rmr:plantdoc_10:max\endcsname{61.4}
\expandafter\gdef\csname num@plantvillage:dense_rmr:pct\endcsname{100.0}
\expandafter\gdef\csname num@plantvillage:dense_rmr:params\endcsname{11.18}
\expandafter\gdef\csname num@plantvillage:dense_rmr:macs_128\endcsname{0.592}
\expandafter\gdef\csname num@plantvillage:dense_rmr:macs_224\endcsname{1.814}
\expandafter\gdef\csname num@resisc45:P_random:novel_1\endcsname{63.0}
\expandafter\gdef\csname num@resisc45:P_random:novel_1:sd\endcsname{5.4}
\expandafter\gdef\csname num@resisc45:P_random:novel_1:min\endcsname{54.1}
\expandafter\gdef\csname num@resisc45:P_random:novel_1:max\endcsname{67.7}
\expandafter\gdef\csname num@resisc45:P_random:novel_5\endcsname{79.7}
\expandafter\gdef\csname num@resisc45:P_random:novel_5:sd\endcsname{4.9}
\expandafter\gdef\csname num@resisc45:P_random:novel_5:min\endcsname{71.5}
\expandafter\gdef\csname num@resisc45:P_random:novel_5:max\endcsname{83.7}
\expandafter\gdef\csname num@resisc45:P_random:novel_10\endcsname{82.6}
\expandafter\gdef\csname num@resisc45:P_random:novel_10:sd\endcsname{4.5}
\expandafter\gdef\csname num@resisc45:P_random:novel_10:min\endcsname{74.9}
\expandafter\gdef\csname num@resisc45:P_random:novel_10:max\endcsname{86.0}
\expandafter\gdef\csname num@resisc45:P_random:pct\endcsname{30.0}
\expandafter\gdef\csname num@resisc45:P_random:params\endcsname{3.35}
\expandafter\gdef\csname num@resisc45:P_random:macs_128\endcsname{0.251}
\expandafter\gdef\csname num@resisc45:P_random:macs_224\endcsname{0.770}
\expandafter\gdef\csname num@resisc45:P_l1:novel_1\endcsname{64.5}
\expandafter\gdef\csname num@resisc45:P_l1:novel_1:sd\endcsname{5.5}
\expandafter\gdef\csname num@resisc45:P_l1:novel_1:min\endcsname{54.8}
\expandafter\gdef\csname num@resisc45:P_l1:novel_1:max\endcsname{68.0}
\expandafter\gdef\csname num@resisc45:P_l1:novel_5\endcsname{80.5}
\expandafter\gdef\csname num@resisc45:P_l1:novel_5:sd\endcsname{4.5}
\expandafter\gdef\csname num@resisc45:P_l1:novel_5:min\endcsname{72.7}
\expandafter\gdef\csname num@resisc45:P_l1:novel_5:max\endcsname{83.8}
\expandafter\gdef\csname num@resisc45:P_l1:novel_10\endcsname{83.3}
\expandafter\gdef\csname num@resisc45:P_l1:novel_10:sd\endcsname{4.1}
\expandafter\gdef\csname num@resisc45:P_l1:novel_10:min\endcsname{76.1}
\expandafter\gdef\csname num@resisc45:P_l1:novel_10:max\endcsname{86.1}
\expandafter\gdef\csname num@resisc45:P_l1:pct\endcsname{30.0}
\expandafter\gdef\csname num@resisc45:P_l1:params\endcsname{3.35}
\expandafter\gdef\csname num@resisc45:P_l1:macs_128\endcsname{0.251}
\expandafter\gdef\csname num@resisc45:P_l1:macs_224\endcsname{0.770}
\expandafter\gdef\csname num@resisc45:P_taylor:novel_1\endcsname{64.7}
\expandafter\gdef\csname num@resisc45:P_taylor:novel_1:sd\endcsname{5.4}
\expandafter\gdef\csname num@resisc45:P_taylor:novel_1:min\endcsname{55.3}
\expandafter\gdef\csname num@resisc45:P_taylor:novel_1:max\endcsname{68.8}
\expandafter\gdef\csname num@resisc45:P_taylor:novel_5\endcsname{82.2}
\expandafter\gdef\csname num@resisc45:P_taylor:novel_5:sd\endcsname{3.9}
\expandafter\gdef\csname num@resisc45:P_taylor:novel_5:min\endcsname{75.4}
\expandafter\gdef\csname num@resisc45:P_taylor:novel_5:max\endcsname{84.8}
\expandafter\gdef\csname num@resisc45:P_taylor:novel_10\endcsname{84.9}
\expandafter\gdef\csname num@resisc45:P_taylor:novel_10:sd\endcsname{3.7}
\expandafter\gdef\csname num@resisc45:P_taylor:novel_10:min\endcsname{78.5}
\expandafter\gdef\csname num@resisc45:P_taylor:novel_10:max\endcsname{87.1}
\expandafter\gdef\csname num@resisc45:P_taylor:pct\endcsname{30.0}
\expandafter\gdef\csname num@resisc45:P_taylor:params\endcsname{3.35}
\expandafter\gdef\csname num@resisc45:P_taylor:macs_128\endcsname{0.251}
\expandafter\gdef\csname num@resisc45:P_taylor:macs_224\endcsname{0.770}
\expandafter\gdef\csname num@resisc45:P_D:novel_1\endcsname{62.3}
\expandafter\gdef\csname num@resisc45:P_D:novel_1:sd\endcsname{4.6}
\expandafter\gdef\csname num@resisc45:P_D:novel_1:min\endcsname{54.5}
\expandafter\gdef\csname num@resisc45:P_D:novel_1:max\endcsname{65.8}
\expandafter\gdef\csname num@resisc45:P_D:novel_5\endcsname{79.3}
\expandafter\gdef\csname num@resisc45:P_D:novel_5:sd\endcsname{4.1}
\expandafter\gdef\csname num@resisc45:P_D:novel_5:min\endcsname{72.2}
\expandafter\gdef\csname num@resisc45:P_D:novel_5:max\endcsname{82.6}
\expandafter\gdef\csname num@resisc45:P_D:novel_10\endcsname{82.1}
\expandafter\gdef\csname num@resisc45:P_D:novel_10:sd\endcsname{4.1}
\expandafter\gdef\csname num@resisc45:P_D:novel_10:min\endcsname{75.1}
\expandafter\gdef\csname num@resisc45:P_D:novel_10:max\endcsname{85.0}
\expandafter\gdef\csname num@resisc45:P_D:pct\endcsname{30.0}
\expandafter\gdef\csname num@resisc45:P_D:params\endcsname{3.35}
\expandafter\gdef\csname num@resisc45:P_D:macs_128\endcsname{0.251}
\expandafter\gdef\csname num@resisc45:P_D:macs_224\endcsname{0.770}
\expandafter\gdef\csname num@resisc45:P_dacis:novel_1\endcsname{62.5}
\expandafter\gdef\csname num@resisc45:P_dacis:novel_1:sd\endcsname{5.2}
\expandafter\gdef\csname num@resisc45:P_dacis:novel_1:min\endcsname{53.7}
\expandafter\gdef\csname num@resisc45:P_dacis:novel_1:max\endcsname{66.9}
\expandafter\gdef\csname num@resisc45:P_dacis:novel_5\endcsname{79.7}
\expandafter\gdef\csname num@resisc45:P_dacis:novel_5:sd\endcsname{4.1}
\expandafter\gdef\csname num@resisc45:P_dacis:novel_5:min\endcsname{72.8}
\expandafter\gdef\csname num@resisc45:P_dacis:novel_5:max\endcsname{82.8}
\expandafter\gdef\csname num@resisc45:P_dacis:novel_10\endcsname{82.7}
\expandafter\gdef\csname num@resisc45:P_dacis:novel_10:sd\endcsname{3.8}
\expandafter\gdef\csname num@resisc45:P_dacis:novel_10:min\endcsname{76.1}
\expandafter\gdef\csname num@resisc45:P_dacis:novel_10:max\endcsname{85.4}
\expandafter\gdef\csname num@resisc45:P_dacis:pct\endcsname{30.0}
\expandafter\gdef\csname num@resisc45:P_dacis:params\endcsname{3.35}
\expandafter\gdef\csname num@resisc45:P_dacis:macs_128\endcsname{0.251}
\expandafter\gdef\csname num@resisc45:P_dacis:macs_224\endcsname{0.770}
\expandafter\gdef\csname num@eurosat13:P_random:novel_1\endcsname{44.0}
\expandafter\gdef\csname num@eurosat13:P_random:novel_1:sd\endcsname{5.6}
\expandafter\gdef\csname num@eurosat13:P_random:novel_1:min\endcsname{36.0}
\expandafter\gdef\csname num@eurosat13:P_random:novel_1:max\endcsname{50.1}
\expandafter\gdef\csname num@eurosat13:P_random:novel_5\endcsname{55.6}
\expandafter\gdef\csname num@eurosat13:P_random:novel_5:sd\endcsname{6.4}
\expandafter\gdef\csname num@eurosat13:P_random:novel_5:min\endcsname{46.5}
\expandafter\gdef\csname num@eurosat13:P_random:novel_5:max\endcsname{61.9}
\expandafter\gdef\csname num@eurosat13:P_random:novel_10\endcsname{58.3}
\expandafter\gdef\csname num@eurosat13:P_random:novel_10:sd\endcsname{6.1}
\expandafter\gdef\csname num@eurosat13:P_random:novel_10:min\endcsname{49.6}
\expandafter\gdef\csname num@eurosat13:P_random:novel_10:max\endcsname{63.8}
\expandafter\gdef\csname num@eurosat13:P_random:pct\endcsname{30.0}
\expandafter\gdef\csname num@eurosat13:P_random:params\endcsname{3.36}
\expandafter\gdef\csname num@eurosat13:P_random:macs_128\endcsname{0.378}
\expandafter\gdef\csname num@eurosat13:P_random:macs_224\endcsname{1.157}
\expandafter\gdef\csname num@eurosat13:P_l1:novel_1\endcsname{45.1}
\expandafter\gdef\csname num@eurosat13:P_l1:novel_1:sd\endcsname{7.1}
\expandafter\gdef\csname num@eurosat13:P_l1:novel_1:min\endcsname{34.5}
\expandafter\gdef\csname num@eurosat13:P_l1:novel_1:max\endcsname{54.3}
\expandafter\gdef\csname num@eurosat13:P_l1:novel_5\endcsname{56.2}
\expandafter\gdef\csname num@eurosat13:P_l1:novel_5:sd\endcsname{7.9}
\expandafter\gdef\csname num@eurosat13:P_l1:novel_5:min\endcsname{43.4}
\expandafter\gdef\csname num@eurosat13:P_l1:novel_5:max\endcsname{64.9}
\expandafter\gdef\csname num@eurosat13:P_l1:novel_10\endcsname{58.7}
\expandafter\gdef\csname num@eurosat13:P_l1:novel_10:sd\endcsname{7.3}
\expandafter\gdef\csname num@eurosat13:P_l1:novel_10:min\endcsname{46.9}
\expandafter\gdef\csname num@eurosat13:P_l1:novel_10:max\endcsname{66.8}
\expandafter\gdef\csname num@eurosat13:P_l1:pct\endcsname{30.0}
\expandafter\gdef\csname num@eurosat13:P_l1:params\endcsname{3.36}
\expandafter\gdef\csname num@eurosat13:P_l1:macs_128\endcsname{0.378}
\expandafter\gdef\csname num@eurosat13:P_l1:macs_224\endcsname{1.157}
\expandafter\gdef\csname num@eurosat13:P_taylor:novel_1\endcsname{46.8}
\expandafter\gdef\csname num@eurosat13:P_taylor:novel_1:sd\endcsname{6.3}
\expandafter\gdef\csname num@eurosat13:P_taylor:novel_1:min\endcsname{36.9}
\expandafter\gdef\csname num@eurosat13:P_taylor:novel_1:max\endcsname{54.3}
\expandafter\gdef\csname num@eurosat13:P_taylor:novel_5\endcsname{58.1}
\expandafter\gdef\csname num@eurosat13:P_taylor:novel_5:sd\endcsname{6.6}
\expandafter\gdef\csname num@eurosat13:P_taylor:novel_5:min\endcsname{47.0}
\expandafter\gdef\csname num@eurosat13:P_taylor:novel_5:max\endcsname{64.5}
\expandafter\gdef\csname num@eurosat13:P_taylor:novel_10\endcsname{60.4}
\expandafter\gdef\csname num@eurosat13:P_taylor:novel_10:sd\endcsname{6.1}
\expandafter\gdef\csname num@eurosat13:P_taylor:novel_10:min\endcsname{50.2}
\expandafter\gdef\csname num@eurosat13:P_taylor:novel_10:max\endcsname{66.0}
\expandafter\gdef\csname num@eurosat13:P_taylor:pct\endcsname{30.0}
\expandafter\gdef\csname num@eurosat13:P_taylor:params\endcsname{3.36}
\expandafter\gdef\csname num@eurosat13:P_taylor:macs_128\endcsname{0.378}
\expandafter\gdef\csname num@eurosat13:P_taylor:macs_224\endcsname{1.157}
\expandafter\gdef\csname num@eurosat13:P_D:novel_1\endcsname{45.3}
\expandafter\gdef\csname num@eurosat13:P_D:novel_1:sd\endcsname{6.6}
\expandafter\gdef\csname num@eurosat13:P_D:novel_1:min\endcsname{36.1}
\expandafter\gdef\csname num@eurosat13:P_D:novel_1:max\endcsname{53.4}
\expandafter\gdef\csname num@eurosat13:P_D:novel_5\endcsname{55.8}
\expandafter\gdef\csname num@eurosat13:P_D:novel_5:sd\endcsname{6.9}
\expandafter\gdef\csname num@eurosat13:P_D:novel_5:min\endcsname{45.5}
\expandafter\gdef\csname num@eurosat13:P_D:novel_5:max\endcsname{64.3}
\expandafter\gdef\csname num@eurosat13:P_D:novel_10\endcsname{58.3}
\expandafter\gdef\csname num@eurosat13:P_D:novel_10:sd\endcsname{6.4}
\expandafter\gdef\csname num@eurosat13:P_D:novel_10:min\endcsname{48.8}
\expandafter\gdef\csname num@eurosat13:P_D:novel_10:max\endcsname{66.1}
\expandafter\gdef\csname num@eurosat13:P_D:pct\endcsname{30.0}
\expandafter\gdef\csname num@eurosat13:P_D:params\endcsname{3.36}
\expandafter\gdef\csname num@eurosat13:P_D:macs_128\endcsname{0.378}
\expandafter\gdef\csname num@eurosat13:P_D:macs_224\endcsname{1.157}
\expandafter\gdef\csname num@eurosat13:P_dacis:novel_1\endcsname{45.4}
\expandafter\gdef\csname num@eurosat13:P_dacis:novel_1:sd\endcsname{6.3}
\expandafter\gdef\csname num@eurosat13:P_dacis:novel_1:min\endcsname{36.7}
\expandafter\gdef\csname num@eurosat13:P_dacis:novel_1:max\endcsname{53.2}
\expandafter\gdef\csname num@eurosat13:P_dacis:novel_5\endcsname{56.3}
\expandafter\gdef\csname num@eurosat13:P_dacis:novel_5:sd\endcsname{6.7}
\expandafter\gdef\csname num@eurosat13:P_dacis:novel_5:min\endcsname{46.6}
\expandafter\gdef\csname num@eurosat13:P_dacis:novel_5:max\endcsname{64.3}
\expandafter\gdef\csname num@eurosat13:P_dacis:novel_10\endcsname{58.6}
\expandafter\gdef\csname num@eurosat13:P_dacis:novel_10:sd\endcsname{6.5}
\expandafter\gdef\csname num@eurosat13:P_dacis:novel_10:min\endcsname{49.4}
\expandafter\gdef\csname num@eurosat13:P_dacis:novel_10:max\endcsname{66.0}
\expandafter\gdef\csname num@eurosat13:P_dacis:pct\endcsname{30.0}
\expandafter\gdef\csname num@eurosat13:P_dacis:params\endcsname{3.36}
\expandafter\gdef\csname num@eurosat13:P_dacis:macs_128\endcsname{0.378}
\expandafter\gdef\csname num@eurosat13:P_dacis:macs_224\endcsname{1.157}
\expandafter\gdef\csname num@plantvillage:P_random:novel_1\endcsname{71.7}
\expandafter\gdef\csname num@plantvillage:P_random:novel_1:sd\endcsname{11.2}
\expandafter\gdef\csname num@plantvillage:P_random:novel_1:min\endcsname{57.4}
\expandafter\gdef\csname num@plantvillage:P_random:novel_1:max\endcsname{85.3}
\expandafter\gdef\csname num@plantvillage:P_random:novel_5\endcsname{85.6}
\expandafter\gdef\csname num@plantvillage:P_random:novel_5:sd\endcsname{8.1}
\expandafter\gdef\csname num@plantvillage:P_random:novel_5:min\endcsname{74.9}
\expandafter\gdef\csname num@plantvillage:P_random:novel_5:max\endcsname{94.9}
\expandafter\gdef\csname num@plantvillage:P_random:novel_10\endcsname{88.1}
\expandafter\gdef\csname num@plantvillage:P_random:novel_10:sd\endcsname{6.9}
\expandafter\gdef\csname num@plantvillage:P_random:novel_10:min\endcsname{79.0}
\expandafter\gdef\csname num@plantvillage:P_random:novel_10:max\endcsname{96.2}
\expandafter\gdef\csname num@plantvillage:P_random:plantdoc_1\endcsname{35.0}
\expandafter\gdef\csname num@plantvillage:P_random:plantdoc_1:sd\endcsname{1.3}
\expandafter\gdef\csname num@plantvillage:P_random:plantdoc_1:min\endcsname{33.1}
\expandafter\gdef\csname num@plantvillage:P_random:plantdoc_1:max\endcsname{36.1}
\expandafter\gdef\csname num@plantvillage:P_random:plantdoc_5\endcsname{46.6}
\expandafter\gdef\csname num@plantvillage:P_random:plantdoc_5:sd\endcsname{1.3}
\expandafter\gdef\csname num@plantvillage:P_random:plantdoc_5:min\endcsname{44.7}
\expandafter\gdef\csname num@plantvillage:P_random:plantdoc_5:max\endcsname{48.0}
\expandafter\gdef\csname num@plantvillage:P_random:plantdoc_10\endcsname{50.9}
\expandafter\gdef\csname num@plantvillage:P_random:plantdoc_10:sd\endcsname{1.3}
\expandafter\gdef\csname num@plantvillage:P_random:plantdoc_10:min\endcsname{49.2}
\expandafter\gdef\csname num@plantvillage:P_random:plantdoc_10:max\endcsname{52.4}
\expandafter\gdef\csname num@plantvillage:P_random:pct\endcsname{30.0}
\expandafter\gdef\csname num@plantvillage:P_random:params\endcsname{3.35}
\expandafter\gdef\csname num@plantvillage:P_random:macs_128\endcsname{0.251}
\expandafter\gdef\csname num@plantvillage:P_random:macs_224\endcsname{0.770}
\expandafter\gdef\csname num@plantvillage:P_l1:novel_1\endcsname{72.9}
\expandafter\gdef\csname num@plantvillage:P_l1:novel_1:sd\endcsname{11.5}
\expandafter\gdef\csname num@plantvillage:P_l1:novel_1:min\endcsname{60.2}
\expandafter\gdef\csname num@plantvillage:P_l1:novel_1:max\endcsname{85.5}
\expandafter\gdef\csname num@plantvillage:P_l1:novel_5\endcsname{85.8}
\expandafter\gdef\csname num@plantvillage:P_l1:novel_5:sd\endcsname{8.3}
\expandafter\gdef\csname num@plantvillage:P_l1:novel_5:min\endcsname{76.4}
\expandafter\gdef\csname num@plantvillage:P_l1:novel_5:max\endcsname{94.5}
\expandafter\gdef\csname num@plantvillage:P_l1:novel_10\endcsname{88.2}
\expandafter\gdef\csname num@plantvillage:P_l1:novel_10:sd\endcsname{7.2}
\expandafter\gdef\csname num@plantvillage:P_l1:novel_10:min\endcsname{80.2}
\expandafter\gdef\csname num@plantvillage:P_l1:novel_10:max\endcsname{95.8}
\expandafter\gdef\csname num@plantvillage:P_l1:plantdoc_1\endcsname{37.0}
\expandafter\gdef\csname num@plantvillage:P_l1:plantdoc_1:sd\endcsname{1.2}
\expandafter\gdef\csname num@plantvillage:P_l1:plantdoc_1:min\endcsname{35.3}
\expandafter\gdef\csname num@plantvillage:P_l1:plantdoc_1:max\endcsname{38.3}
\expandafter\gdef\csname num@plantvillage:P_l1:plantdoc_5\endcsname{48.7}
\expandafter\gdef\csname num@plantvillage:P_l1:plantdoc_5:sd\endcsname{0.9}
\expandafter\gdef\csname num@plantvillage:P_l1:plantdoc_5:min\endcsname{47.5}
\expandafter\gdef\csname num@plantvillage:P_l1:plantdoc_5:max\endcsname{50.0}
\expandafter\gdef\csname num@plantvillage:P_l1:plantdoc_10\endcsname{52.7}
\expandafter\gdef\csname num@plantvillage:P_l1:plantdoc_10:sd\endcsname{0.9}
\expandafter\gdef\csname num@plantvillage:P_l1:plantdoc_10:min\endcsname{51.6}
\expandafter\gdef\csname num@plantvillage:P_l1:plantdoc_10:max\endcsname{53.8}
\expandafter\gdef\csname num@plantvillage:P_l1:pct\endcsname{30.0}
\expandafter\gdef\csname num@plantvillage:P_l1:params\endcsname{3.35}
\expandafter\gdef\csname num@plantvillage:P_l1:macs_128\endcsname{0.251}
\expandafter\gdef\csname num@plantvillage:P_l1:macs_224\endcsname{0.770}
\expandafter\gdef\csname num@plantvillage:P_bn_gamma:novel_1\endcsname{71.7}
\expandafter\gdef\csname num@plantvillage:P_bn_gamma:novel_1:sd\endcsname{12.9}
\expandafter\gdef\csname num@plantvillage:P_bn_gamma:novel_1:min\endcsname{58.4}
\expandafter\gdef\csname num@plantvillage:P_bn_gamma:novel_1:max\endcsname{84.8}
\expandafter\gdef\csname num@plantvillage:P_bn_gamma:novel_5\endcsname{85.1}
\expandafter\gdef\csname num@plantvillage:P_bn_gamma:novel_5:sd\endcsname{8.9}
\expandafter\gdef\csname num@plantvillage:P_bn_gamma:novel_5:min\endcsname{74.7}
\expandafter\gdef\csname num@plantvillage:P_bn_gamma:novel_5:max\endcsname{93.9}
\expandafter\gdef\csname num@plantvillage:P_bn_gamma:novel_10\endcsname{87.5}
\expandafter\gdef\csname num@plantvillage:P_bn_gamma:novel_10:sd\endcsname{7.3}
\expandafter\gdef\csname num@plantvillage:P_bn_gamma:novel_10:min\endcsname{78.9}
\expandafter\gdef\csname num@plantvillage:P_bn_gamma:novel_10:max\endcsname{95.0}
\expandafter\gdef\csname num@plantvillage:P_bn_gamma:plantdoc_1\endcsname{33.2}
\expandafter\gdef\csname num@plantvillage:P_bn_gamma:plantdoc_1:sd\endcsname{1.3}
\expandafter\gdef\csname num@plantvillage:P_bn_gamma:plantdoc_1:min\endcsname{31.4}
\expandafter\gdef\csname num@plantvillage:P_bn_gamma:plantdoc_1:max\endcsname{34.7}
\expandafter\gdef\csname num@plantvillage:P_bn_gamma:plantdoc_5\endcsname{44.4}
\expandafter\gdef\csname num@plantvillage:P_bn_gamma:plantdoc_5:sd\endcsname{1.7}
\expandafter\gdef\csname num@plantvillage:P_bn_gamma:plantdoc_5:min\endcsname{41.6}
\expandafter\gdef\csname num@plantvillage:P_bn_gamma:plantdoc_5:max\endcsname{45.5}
\expandafter\gdef\csname num@plantvillage:P_bn_gamma:plantdoc_10\endcsname{48.5}
\expandafter\gdef\csname num@plantvillage:P_bn_gamma:plantdoc_10:sd\endcsname{1.9}
\expandafter\gdef\csname num@plantvillage:P_bn_gamma:plantdoc_10:min\endcsname{45.4}
\expandafter\gdef\csname num@plantvillage:P_bn_gamma:plantdoc_10:max\endcsname{49.9}
\expandafter\gdef\csname num@plantvillage:P_bn_gamma:pct\endcsname{30.0}
\expandafter\gdef\csname num@plantvillage:P_bn_gamma:params\endcsname{3.35}
\expandafter\gdef\csname num@plantvillage:P_bn_gamma:macs_128\endcsname{0.251}
\expandafter\gdef\csname num@plantvillage:P_bn_gamma:macs_224\endcsname{0.770}
\expandafter\gdef\csname num@plantvillage:P_taylor:novel_1\endcsname{75.0}
\expandafter\gdef\csname num@plantvillage:P_taylor:novel_1:sd\endcsname{10.3}
\expandafter\gdef\csname num@plantvillage:P_taylor:novel_1:min\endcsname{61.1}
\expandafter\gdef\csname num@plantvillage:P_taylor:novel_1:max\endcsname{85.9}
\expandafter\gdef\csname num@plantvillage:P_taylor:novel_5\endcsname{88.3}
\expandafter\gdef\csname num@plantvillage:P_taylor:novel_5:sd\endcsname{7.1}
\expandafter\gdef\csname num@plantvillage:P_taylor:novel_5:min\endcsname{79.0}
\expandafter\gdef\csname num@plantvillage:P_taylor:novel_5:max\endcsname{95.6}
\expandafter\gdef\csname num@plantvillage:P_taylor:novel_10\endcsname{90.3}
\expandafter\gdef\csname num@plantvillage:P_taylor:novel_10:sd\endcsname{5.8}
\expandafter\gdef\csname num@plantvillage:P_taylor:novel_10:min\endcsname{82.9}
\expandafter\gdef\csname num@plantvillage:P_taylor:novel_10:max\endcsname{96.4}
\expandafter\gdef\csname num@plantvillage:P_taylor:plantdoc_1\endcsname{35.6}
\expandafter\gdef\csname num@plantvillage:P_taylor:plantdoc_1:sd\endcsname{1.1}
\expandafter\gdef\csname num@plantvillage:P_taylor:plantdoc_1:min\endcsname{33.7}
\expandafter\gdef\csname num@plantvillage:P_taylor:plantdoc_1:max\endcsname{36.3}
\expandafter\gdef\csname num@plantvillage:P_taylor:plantdoc_5\endcsname{47.9}
\expandafter\gdef\csname num@plantvillage:P_taylor:plantdoc_5:sd\endcsname{1.2}
\expandafter\gdef\csname num@plantvillage:P_taylor:plantdoc_5:min\endcsname{46.1}
\expandafter\gdef\csname num@plantvillage:P_taylor:plantdoc_5:max\endcsname{49.2}
\expandafter\gdef\csname num@plantvillage:P_taylor:plantdoc_10\endcsname{52.3}
\expandafter\gdef\csname num@plantvillage:P_taylor:plantdoc_10:sd\endcsname{1.2}
\expandafter\gdef\csname num@plantvillage:P_taylor:plantdoc_10:min\endcsname{50.7}
\expandafter\gdef\csname num@plantvillage:P_taylor:plantdoc_10:max\endcsname{53.8}
\expandafter\gdef\csname num@plantvillage:P_taylor:pct\endcsname{30.0}
\expandafter\gdef\csname num@plantvillage:P_taylor:params\endcsname{3.35}
\expandafter\gdef\csname num@plantvillage:P_taylor:macs_128\endcsname{0.251}
\expandafter\gdef\csname num@plantvillage:P_taylor:macs_224\endcsname{0.770}
\expandafter\gdef\csname num@plantvillage:P_G:novel_1\endcsname{74.0}
\expandafter\gdef\csname num@plantvillage:P_G:novel_1:sd\endcsname{11.0}
\expandafter\gdef\csname num@plantvillage:P_G:novel_1:min\endcsname{59.4}
\expandafter\gdef\csname num@plantvillage:P_G:novel_1:max\endcsname{86.2}
\expandafter\gdef\csname num@plantvillage:P_G:novel_5\endcsname{87.3}
\expandafter\gdef\csname num@plantvillage:P_G:novel_5:sd\endcsname{7.2}
\expandafter\gdef\csname num@plantvillage:P_G:novel_5:min\endcsname{77.9}
\expandafter\gdef\csname num@plantvillage:P_G:novel_5:max\endcsname{95.2}
\expandafter\gdef\csname num@plantvillage:P_G:novel_10\endcsname{89.5}
\expandafter\gdef\csname num@plantvillage:P_G:novel_10:sd\endcsname{5.9}
\expandafter\gdef\csname num@plantvillage:P_G:novel_10:min\endcsname{81.9}
\expandafter\gdef\csname num@plantvillage:P_G:novel_10:max\endcsname{96.1}
\expandafter\gdef\csname num@plantvillage:P_G:plantdoc_1\endcsname{35.5}
\expandafter\gdef\csname num@plantvillage:P_G:plantdoc_1:sd\endcsname{1.4}
\expandafter\gdef\csname num@plantvillage:P_G:plantdoc_1:min\endcsname{33.5}
\expandafter\gdef\csname num@plantvillage:P_G:plantdoc_1:max\endcsname{37.4}
\expandafter\gdef\csname num@plantvillage:P_G:plantdoc_5\endcsname{47.6}
\expandafter\gdef\csname num@plantvillage:P_G:plantdoc_5:sd\endcsname{1.5}
\expandafter\gdef\csname num@plantvillage:P_G:plantdoc_5:min\endcsname{45.8}
\expandafter\gdef\csname num@plantvillage:P_G:plantdoc_5:max\endcsname{49.9}
\expandafter\gdef\csname num@plantvillage:P_G:plantdoc_10\endcsname{51.8}
\expandafter\gdef\csname num@plantvillage:P_G:plantdoc_10:sd\endcsname{1.5}
\expandafter\gdef\csname num@plantvillage:P_G:plantdoc_10:min\endcsname{50.3}
\expandafter\gdef\csname num@plantvillage:P_G:plantdoc_10:max\endcsname{54.2}
\expandafter\gdef\csname num@plantvillage:P_G:pct\endcsname{30.0}
\expandafter\gdef\csname num@plantvillage:P_G:params\endcsname{3.35}
\expandafter\gdef\csname num@plantvillage:P_G:macs_128\endcsname{0.251}
\expandafter\gdef\csname num@plantvillage:P_G:macs_224\endcsname{0.770}
\expandafter\gdef\csname num@plantvillage:P_V:novel_1\endcsname{72.9}
\expandafter\gdef\csname num@plantvillage:P_V:novel_1:sd\endcsname{10.5}
\expandafter\gdef\csname num@plantvillage:P_V:novel_1:min\endcsname{58.9}
\expandafter\gdef\csname num@plantvillage:P_V:novel_1:max\endcsname{86.1}
\expandafter\gdef\csname num@plantvillage:P_V:novel_5\endcsname{86.5}
\expandafter\gdef\csname num@plantvillage:P_V:novel_5:sd\endcsname{6.9}
\expandafter\gdef\csname num@plantvillage:P_V:novel_5:min\endcsname{77.3}
\expandafter\gdef\csname num@plantvillage:P_V:novel_5:max\endcsname{94.8}
\expandafter\gdef\csname num@plantvillage:P_V:novel_10\endcsname{88.8}
\expandafter\gdef\csname num@plantvillage:P_V:novel_10:sd\endcsname{5.8}
\expandafter\gdef\csname num@plantvillage:P_V:novel_10:min\endcsname{81.5}
\expandafter\gdef\csname num@plantvillage:P_V:novel_10:max\endcsname{95.9}
\expandafter\gdef\csname num@plantvillage:P_V:plantdoc_1\endcsname{35.5}
\expandafter\gdef\csname num@plantvillage:P_V:plantdoc_1:sd\endcsname{1.3}
\expandafter\gdef\csname num@plantvillage:P_V:plantdoc_1:min\endcsname{33.6}
\expandafter\gdef\csname num@plantvillage:P_V:plantdoc_1:max\endcsname{36.7}
\expandafter\gdef\csname num@plantvillage:P_V:plantdoc_5\endcsname{47.2}
\expandafter\gdef\csname num@plantvillage:P_V:plantdoc_5:sd\endcsname{1.4}
\expandafter\gdef\csname num@plantvillage:P_V:plantdoc_5:min\endcsname{44.9}
\expandafter\gdef\csname num@plantvillage:P_V:plantdoc_5:max\endcsname{48.2}
\expandafter\gdef\csname num@plantvillage:P_V:plantdoc_10\endcsname{51.4}
\expandafter\gdef\csname num@plantvillage:P_V:plantdoc_10:sd\endcsname{1.5}
\expandafter\gdef\csname num@plantvillage:P_V:plantdoc_10:min\endcsname{49.0}
\expandafter\gdef\csname num@plantvillage:P_V:plantdoc_10:max\endcsname{52.5}
\expandafter\gdef\csname num@plantvillage:P_V:pct\endcsname{30.0}
\expandafter\gdef\csname num@plantvillage:P_V:params\endcsname{3.35}
\expandafter\gdef\csname num@plantvillage:P_V:macs_128\endcsname{0.251}
\expandafter\gdef\csname num@plantvillage:P_V:macs_224\endcsname{0.770}
\expandafter\gdef\csname num@plantvillage:P_D:novel_1\endcsname{71.9}
\expandafter\gdef\csname num@plantvillage:P_D:novel_1:sd\endcsname{12.4}
\expandafter\gdef\csname num@plantvillage:P_D:novel_1:min\endcsname{55.6}
\expandafter\gdef\csname num@plantvillage:P_D:novel_1:max\endcsname{84.9}
\expandafter\gdef\csname num@plantvillage:P_D:novel_5\endcsname{85.2}
\expandafter\gdef\csname num@plantvillage:P_D:novel_5:sd\endcsname{8.9}
\expandafter\gdef\csname num@plantvillage:P_D:novel_5:min\endcsname{73.1}
\expandafter\gdef\csname num@plantvillage:P_D:novel_5:max\endcsname{94.1}
\expandafter\gdef\csname num@plantvillage:P_D:novel_10\endcsname{87.7}
\expandafter\gdef\csname num@plantvillage:P_D:novel_10:sd\endcsname{7.5}
\expandafter\gdef\csname num@plantvillage:P_D:novel_10:min\endcsname{78.0}
\expandafter\gdef\csname num@plantvillage:P_D:novel_10:max\endcsname{95.3}
\expandafter\gdef\csname num@plantvillage:P_D:plantdoc_1\endcsname{34.8}
\expandafter\gdef\csname num@plantvillage:P_D:plantdoc_1:sd\endcsname{0.6}
\expandafter\gdef\csname num@plantvillage:P_D:plantdoc_1:min\endcsname{34.1}
\expandafter\gdef\csname num@plantvillage:P_D:plantdoc_1:max\endcsname{35.7}
\expandafter\gdef\csname num@plantvillage:P_D:plantdoc_5\endcsname{46.8}
\expandafter\gdef\csname num@plantvillage:P_D:plantdoc_5:sd\endcsname{0.8}
\expandafter\gdef\csname num@plantvillage:P_D:plantdoc_5:min\endcsname{46.1}
\expandafter\gdef\csname num@plantvillage:P_D:plantdoc_5:max\endcsname{48.1}
\expandafter\gdef\csname num@plantvillage:P_D:plantdoc_10\endcsname{51.2}
\expandafter\gdef\csname num@plantvillage:P_D:plantdoc_10:sd\endcsname{1.0}
\expandafter\gdef\csname num@plantvillage:P_D:plantdoc_10:min\endcsname{49.8}
\expandafter\gdef\csname num@plantvillage:P_D:plantdoc_10:max\endcsname{52.6}
\expandafter\gdef\csname num@plantvillage:P_D:pct\endcsname{30.0}
\expandafter\gdef\csname num@plantvillage:P_D:params\endcsname{3.35}
\expandafter\gdef\csname num@plantvillage:P_D:macs_128\endcsname{0.251}
\expandafter\gdef\csname num@plantvillage:P_D:macs_224\endcsname{0.770}
\expandafter\gdef\csname num@plantvillage:P_dacis:novel_1\endcsname{72.7}
\expandafter\gdef\csname num@plantvillage:P_dacis:novel_1:sd\endcsname{11.5}
\expandafter\gdef\csname num@plantvillage:P_dacis:novel_1:min\endcsname{57.3}
\expandafter\gdef\csname num@plantvillage:P_dacis:novel_1:max\endcsname{84.3}
\expandafter\gdef\csname num@plantvillage:P_dacis:novel_5\endcsname{86.4}
\expandafter\gdef\csname num@plantvillage:P_dacis:novel_5:sd\endcsname{7.7}
\expandafter\gdef\csname num@plantvillage:P_dacis:novel_5:min\endcsname{75.7}
\expandafter\gdef\csname num@plantvillage:P_dacis:novel_5:max\endcsname{93.9}
\expandafter\gdef\csname num@plantvillage:P_dacis:novel_10\endcsname{88.6}
\expandafter\gdef\csname num@plantvillage:P_dacis:novel_10:sd\endcsname{6.3}
\expandafter\gdef\csname num@plantvillage:P_dacis:novel_10:min\endcsname{80.2}
\expandafter\gdef\csname num@plantvillage:P_dacis:novel_10:max\endcsname{95.1}
\expandafter\gdef\csname num@plantvillage:P_dacis:plantdoc_1\endcsname{35.2}
\expandafter\gdef\csname num@plantvillage:P_dacis:plantdoc_1:sd\endcsname{1.2}
\expandafter\gdef\csname num@plantvillage:P_dacis:plantdoc_1:min\endcsname{33.7}
\expandafter\gdef\csname num@plantvillage:P_dacis:plantdoc_1:max\endcsname{37.0}
\expandafter\gdef\csname num@plantvillage:P_dacis:plantdoc_5\endcsname{47.3}
\expandafter\gdef\csname num@plantvillage:P_dacis:plantdoc_5:sd\endcsname{1.4}
\expandafter\gdef\csname num@plantvillage:P_dacis:plantdoc_5:min\endcsname{45.9}
\expandafter\gdef\csname num@plantvillage:P_dacis:plantdoc_5:max\endcsname{49.7}
\expandafter\gdef\csname num@plantvillage:P_dacis:plantdoc_10\endcsname{51.4}
\expandafter\gdef\csname num@plantvillage:P_dacis:plantdoc_10:sd\endcsname{1.4}
\expandafter\gdef\csname num@plantvillage:P_dacis:plantdoc_10:min\endcsname{50.2}
\expandafter\gdef\csname num@plantvillage:P_dacis:plantdoc_10:max\endcsname{53.6}
\expandafter\gdef\csname num@plantvillage:P_dacis:pct\endcsname{30.0}
\expandafter\gdef\csname num@plantvillage:P_dacis:params\endcsname{3.35}
\expandafter\gdef\csname num@plantvillage:P_dacis:macs_128\endcsname{0.251}
\expandafter\gdef\csname num@plantvillage:P_dacis:macs_224\endcsname{0.770}
\expandafter\gdef\csname num@d:rs:novel_5:dacis-random:PMP\endcsname{\ensuremath{-}0.5}
\expandafter\gdef\csname num@d:rs:novel_5:dacis-random:PMP:lo\endcsname{\ensuremath{-}1.3}
\expandafter\gdef\csname num@d:rs:novel_5:dacis-random:PMP:hi\endcsname{+0.2}
\expandafter\gdef\csname num@d:rs:novel_5:dacis-random:PMP:abs\endcsname{0.5}
\expandafter\gdef\csname num@d:rs:novel_5:dacis-random:PMP:wins\endcsname{1}
\expandafter\gdef\csname num@d:rs:novel_5:dacis-random:PMP:n\endcsname{5}
\expandafter\gdef\csname num@d:rs:novel_5:dacis-random:PMP:worst\endcsname{1.3}
\expandafter\gdef\csname num@d:rs:novel_5:dacis-random:P\endcsname{+0.1}
\expandafter\gdef\csname num@d:rs:novel_5:dacis-random:P:lo\endcsname{\ensuremath{-}1.4}
\expandafter\gdef\csname num@d:rs:novel_5:dacis-random:P:hi\endcsname{+1.5}
\expandafter\gdef\csname num@d:rs:novel_5:dacis-random:P:abs\endcsname{0.1}
\expandafter\gdef\csname num@d:rs:novel_5:dacis-random:P:wins\endcsname{3}
\expandafter\gdef\csname num@d:rs:novel_5:dacis-random:P:n\endcsname{5}
\expandafter\gdef\csname num@d:rs:novel_5:dacis-random:P:worst\endcsname{1.5}
\expandafter\gdef\csname num@d:rs:novel_5:dacis-l1:PMP\endcsname{\ensuremath{-}0.6}
\expandafter\gdef\csname num@d:rs:novel_5:dacis-l1:PMP:lo\endcsname{\ensuremath{-}1.2}
\expandafter\gdef\csname num@d:rs:novel_5:dacis-l1:PMP:hi\endcsname{+0.0}
\expandafter\gdef\csname num@d:rs:novel_5:dacis-l1:PMP:abs\endcsname{0.6}
\expandafter\gdef\csname num@d:rs:novel_5:dacis-l1:PMP:wins\endcsname{1}
\expandafter\gdef\csname num@d:rs:novel_5:dacis-l1:PMP:n\endcsname{5}
\expandafter\gdef\csname num@d:rs:novel_5:dacis-l1:PMP:worst\endcsname{1.2}
\expandafter\gdef\csname num@d:rs:novel_5:dacis-l1:P\endcsname{\ensuremath{-}0.8}
\expandafter\gdef\csname num@d:rs:novel_5:dacis-l1:P:lo\endcsname{\ensuremath{-}2.0}
\expandafter\gdef\csname num@d:rs:novel_5:dacis-l1:P:hi\endcsname{+0.4}
\expandafter\gdef\csname num@d:rs:novel_5:dacis-l1:P:abs\endcsname{0.8}
\expandafter\gdef\csname num@d:rs:novel_5:dacis-l1:P:wins\endcsname{2}
\expandafter\gdef\csname num@d:rs:novel_5:dacis-l1:P:n\endcsname{5}
\expandafter\gdef\csname num@d:rs:novel_5:dacis-l1:P:worst\endcsname{2.0}
\expandafter\gdef\csname num@d:rs:novel_5:dacis-taylor:PMP\endcsname{\ensuremath{-}1.8}
\expandafter\gdef\csname num@d:rs:novel_5:dacis-taylor:PMP:lo\endcsname{\ensuremath{-}2.3}
\expandafter\gdef\csname num@d:rs:novel_5:dacis-taylor:PMP:hi\endcsname{\ensuremath{-}1.2}
\expandafter\gdef\csname num@d:rs:novel_5:dacis-taylor:PMP:abs\endcsname{1.8}
\expandafter\gdef\csname num@d:rs:novel_5:dacis-taylor:PMP:wins\endcsname{0}
\expandafter\gdef\csname num@d:rs:novel_5:dacis-taylor:PMP:n\endcsname{5}
\expandafter\gdef\csname num@d:rs:novel_5:dacis-taylor:PMP:worst\endcsname{2.3}
\expandafter\gdef\csname num@d:rs:novel_5:dacis-taylor:P\endcsname{\ensuremath{-}2.4}
\expandafter\gdef\csname num@d:rs:novel_5:dacis-taylor:P:lo\endcsname{\ensuremath{-}3.1}
\expandafter\gdef\csname num@d:rs:novel_5:dacis-taylor:P:hi\endcsname{\ensuremath{-}1.7}
\expandafter\gdef\csname num@d:rs:novel_5:dacis-taylor:P:abs\endcsname{2.4}
\expandafter\gdef\csname num@d:rs:novel_5:dacis-taylor:P:wins\endcsname{0}
\expandafter\gdef\csname num@d:rs:novel_5:dacis-taylor:P:n\endcsname{5}
\expandafter\gdef\csname num@d:rs:novel_5:dacis-taylor:P:worst\endcsname{3.1}
\expandafter\gdef\csname num@d:rs:novel_5:dacis-bn_gamma:PMP\endcsname{+0.8}
\expandafter\gdef\csname num@d:rs:novel_5:dacis-bn_gamma:PMP:lo\endcsname{\ensuremath{-}0.2}
\expandafter\gdef\csname num@d:rs:novel_5:dacis-bn_gamma:PMP:hi\endcsname{+1.7}
\expandafter\gdef\csname num@d:rs:novel_5:dacis-bn_gamma:PMP:abs\endcsname{0.8}
\expandafter\gdef\csname num@d:rs:novel_5:dacis-bn_gamma:PMP:wins\endcsname{4}
\expandafter\gdef\csname num@d:rs:novel_5:dacis-bn_gamma:PMP:n\endcsname{5}
\expandafter\gdef\csname num@d:rs:novel_5:dacis-bn_gamma:PMP:worst\endcsname{1.7}
\expandafter\gdef\csname num@d:rs:novel_5:dacis-D:PMP\endcsname{+1.0}
\expandafter\gdef\csname num@d:rs:novel_5:dacis-D:PMP:lo\endcsname{+0.0}
\expandafter\gdef\csname num@d:rs:novel_5:dacis-D:PMP:hi\endcsname{+1.9}
\expandafter\gdef\csname num@d:rs:novel_5:dacis-D:PMP:abs\endcsname{1.0}
\expandafter\gdef\csname num@d:rs:novel_5:dacis-D:PMP:wins\endcsname{5}
\expandafter\gdef\csname num@d:rs:novel_5:dacis-D:PMP:n\endcsname{5}
\expandafter\gdef\csname num@d:rs:novel_5:dacis-D:PMP:worst\endcsname{1.9}
\expandafter\gdef\csname num@d:rs:novel_5:dacis-D:P\endcsname{+0.5}
\expandafter\gdef\csname num@d:rs:novel_5:dacis-D:P:lo\endcsname{\ensuremath{-}0.4}
\expandafter\gdef\csname num@d:rs:novel_5:dacis-D:P:hi\endcsname{+1.3}
\expandafter\gdef\csname num@d:rs:novel_5:dacis-D:P:abs\endcsname{0.5}
\expandafter\gdef\csname num@d:rs:novel_5:dacis-D:P:wins\endcsname{4}
\expandafter\gdef\csname num@d:rs:novel_5:dacis-D:P:n\endcsname{5}
\expandafter\gdef\csname num@d:rs:novel_5:dacis-D:P:worst\endcsname{1.3}
\expandafter\gdef\csname num@d:rs:novel_5:dacis-G:PMP\endcsname{\ensuremath{-}0.5}
\expandafter\gdef\csname num@d:rs:novel_5:dacis-G:PMP:lo\endcsname{\ensuremath{-}1.4}
\expandafter\gdef\csname num@d:rs:novel_5:dacis-G:PMP:hi\endcsname{+0.4}
\expandafter\gdef\csname num@d:rs:novel_5:dacis-G:PMP:abs\endcsname{0.5}
\expandafter\gdef\csname num@d:rs:novel_5:dacis-G:PMP:wins\endcsname{1}
\expandafter\gdef\csname num@d:rs:novel_5:dacis-G:PMP:n\endcsname{5}
\expandafter\gdef\csname num@d:rs:novel_5:dacis-G:PMP:worst\endcsname{1.4}
\expandafter\gdef\csname num@d:rs:novel_5:dacis-V:PMP\endcsname{\ensuremath{-}0.5}
\expandafter\gdef\csname num@d:rs:novel_5:dacis-V:PMP:lo\endcsname{\ensuremath{-}1.5}
\expandafter\gdef\csname num@d:rs:novel_5:dacis-V:PMP:hi\endcsname{+0.5}
\expandafter\gdef\csname num@d:rs:novel_5:dacis-V:PMP:abs\endcsname{0.5}
\expandafter\gdef\csname num@d:rs:novel_5:dacis-V:PMP:wins\endcsname{1}
\expandafter\gdef\csname num@d:rs:novel_5:dacis-V:PMP:n\endcsname{5}
\expandafter\gdef\csname num@d:rs:novel_5:dacis-V:PMP:worst\endcsname{1.5}
\expandafter\gdef\csname num@d:rs:novel_5:l1-random:PMP\endcsname{+0.0}
\expandafter\gdef\csname num@d:rs:novel_5:l1-random:PMP:lo\endcsname{\ensuremath{-}0.5}
\expandafter\gdef\csname num@d:rs:novel_5:l1-random:PMP:hi\endcsname{+0.5}
\expandafter\gdef\csname num@d:rs:novel_5:l1-random:PMP:abs\endcsname{0.0}
\expandafter\gdef\csname num@d:rs:novel_5:l1-random:PMP:wins\endcsname{3}
\expandafter\gdef\csname num@d:rs:novel_5:l1-random:PMP:n\endcsname{5}
\expandafter\gdef\csname num@d:rs:novel_5:l1-random:PMP:worst\endcsname{0.5}
\expandafter\gdef\csname num@d:rs:novel_5:l1-random:P\endcsname{+0.9}
\expandafter\gdef\csname num@d:rs:novel_5:l1-random:P:lo\endcsname{\ensuremath{-}0.0}
\expandafter\gdef\csname num@d:rs:novel_5:l1-random:P:hi\endcsname{+1.7}
\expandafter\gdef\csname num@d:rs:novel_5:l1-random:P:abs\endcsname{0.9}
\expandafter\gdef\csname num@d:rs:novel_5:l1-random:P:wins\endcsname{5}
\expandafter\gdef\csname num@d:rs:novel_5:l1-random:P:n\endcsname{5}
\expandafter\gdef\csname num@d:rs:novel_5:l1-random:P:worst\endcsname{1.7}
\expandafter\gdef\csname num@d:rs:novel_5:taylor-random:PMP\endcsname{+1.2}
\expandafter\gdef\csname num@d:rs:novel_5:taylor-random:PMP:lo\endcsname{+0.4}
\expandafter\gdef\csname num@d:rs:novel_5:taylor-random:PMP:hi\endcsname{+2.0}
\expandafter\gdef\csname num@d:rs:novel_5:taylor-random:PMP:abs\endcsname{1.2}
\expandafter\gdef\csname num@d:rs:novel_5:taylor-random:PMP:wins\endcsname{5}
\expandafter\gdef\csname num@d:rs:novel_5:taylor-random:PMP:n\endcsname{5}
\expandafter\gdef\csname num@d:rs:novel_5:taylor-random:PMP:worst\endcsname{2.0}
\expandafter\gdef\csname num@d:rs:novel_5:taylor-random:P\endcsname{+2.5}
\expandafter\gdef\csname num@d:rs:novel_5:taylor-random:P:lo\endcsname{+1.2}
\expandafter\gdef\csname num@d:rs:novel_5:taylor-random:P:hi\endcsname{+3.8}
\expandafter\gdef\csname num@d:rs:novel_5:taylor-random:P:abs\endcsname{2.5}
\expandafter\gdef\csname num@d:rs:novel_5:taylor-random:P:wins\endcsname{5}
\expandafter\gdef\csname num@d:rs:novel_5:taylor-random:P:n\endcsname{5}
\expandafter\gdef\csname num@d:rs:novel_5:taylor-random:P:worst\endcsname{3.8}
\expandafter\gdef\csname num@d:rs:novel_5:bn_gamma-random:PMP\endcsname{\ensuremath{-}1.3}
\expandafter\gdef\csname num@d:rs:novel_5:bn_gamma-random:PMP:lo\endcsname{\ensuremath{-}1.6}
\expandafter\gdef\csname num@d:rs:novel_5:bn_gamma-random:PMP:hi\endcsname{\ensuremath{-}1.0}
\expandafter\gdef\csname num@d:rs:novel_5:bn_gamma-random:PMP:abs\endcsname{1.3}
\expandafter\gdef\csname num@d:rs:novel_5:bn_gamma-random:PMP:wins\endcsname{0}
\expandafter\gdef\csname num@d:rs:novel_5:bn_gamma-random:PMP:n\endcsname{5}
\expandafter\gdef\csname num@d:rs:novel_5:bn_gamma-random:PMP:worst\endcsname{1.6}
\expandafter\gdef\csname num@d:rs:novel_5:D-random:PMP\endcsname{\ensuremath{-}1.5}
\expandafter\gdef\csname num@d:rs:novel_5:D-random:PMP:lo\endcsname{\ensuremath{-}2.9}
\expandafter\gdef\csname num@d:rs:novel_5:D-random:PMP:hi\endcsname{\ensuremath{-}0.1}
\expandafter\gdef\csname num@d:rs:novel_5:D-random:PMP:abs\endcsname{1.5}
\expandafter\gdef\csname num@d:rs:novel_5:D-random:PMP:wins\endcsname{1}
\expandafter\gdef\csname num@d:rs:novel_5:D-random:PMP:n\endcsname{5}
\expandafter\gdef\csname num@d:rs:novel_5:D-random:PMP:worst\endcsname{2.9}
\expandafter\gdef\csname num@d:rs:novel_5:D-random:P\endcsname{\ensuremath{-}0.4}
\expandafter\gdef\csname num@d:rs:novel_5:D-random:P:lo\endcsname{\ensuremath{-}1.5}
\expandafter\gdef\csname num@d:rs:novel_5:D-random:P:hi\endcsname{+0.7}
\expandafter\gdef\csname num@d:rs:novel_5:D-random:P:abs\endcsname{0.4}
\expandafter\gdef\csname num@d:rs:novel_5:D-random:P:wins\endcsname{2}
\expandafter\gdef\csname num@d:rs:novel_5:D-random:P:n\endcsname{5}
\expandafter\gdef\csname num@d:rs:novel_5:D-random:P:worst\endcsname{1.5}
\expandafter\gdef\csname num@d:rs:novel_5:G-random:PMP\endcsname{\ensuremath{-}0.0}
\expandafter\gdef\csname num@d:rs:novel_5:G-random:PMP:lo\endcsname{\ensuremath{-}1.0}
\expandafter\gdef\csname num@d:rs:novel_5:G-random:PMP:hi\endcsname{+1.0}
\expandafter\gdef\csname num@d:rs:novel_5:G-random:PMP:abs\endcsname{0.0}
\expandafter\gdef\csname num@d:rs:novel_5:G-random:PMP:wins\endcsname{3}
\expandafter\gdef\csname num@d:rs:novel_5:G-random:PMP:n\endcsname{5}
\expandafter\gdef\csname num@d:rs:novel_5:G-random:PMP:worst\endcsname{1.0}
\expandafter\gdef\csname num@d:rs:novel_5:V-random:PMP\endcsname{\ensuremath{-}0.0}
\expandafter\gdef\csname num@d:rs:novel_5:V-random:PMP:lo\endcsname{\ensuremath{-}0.6}
\expandafter\gdef\csname num@d:rs:novel_5:V-random:PMP:hi\endcsname{+0.5}
\expandafter\gdef\csname num@d:rs:novel_5:V-random:PMP:abs\endcsname{0.0}
\expandafter\gdef\csname num@d:rs:novel_5:V-random:PMP:wins\endcsname{2}
\expandafter\gdef\csname num@d:rs:novel_5:V-random:PMP:n\endcsname{5}
\expandafter\gdef\csname num@d:rs:novel_5:V-random:PMP:worst\endcsname{0.6}
\expandafter\gdef\csname num@d:rs:novel_5:PM-P\endcsname{\ensuremath{-}0.0}
\expandafter\gdef\csname num@d:rs:novel_5:PM-P:lo\endcsname{\ensuremath{-}0.9}
\expandafter\gdef\csname num@d:rs:novel_5:PM-P:hi\endcsname{+0.8}
\expandafter\gdef\csname num@d:rs:novel_5:PM-P:abs\endcsname{0.0}
\expandafter\gdef\csname num@d:rs:novel_5:PM-P:wins\endcsname{3}
\expandafter\gdef\csname num@d:rs:novel_5:PM-P:n\endcsname{5}
\expandafter\gdef\csname num@d:rs:novel_5:PM-P:worst\endcsname{0.9}
\expandafter\gdef\csname num@d:rs:novel_5:MP-P\endcsname{+0.7}
\expandafter\gdef\csname num@d:rs:novel_5:MP-P:lo\endcsname{\ensuremath{-}0.1}
\expandafter\gdef\csname num@d:rs:novel_5:MP-P:hi\endcsname{+1.6}
\expandafter\gdef\csname num@d:rs:novel_5:MP-P:abs\endcsname{0.7}
\expandafter\gdef\csname num@d:rs:novel_5:MP-P:wins\endcsname{5}
\expandafter\gdef\csname num@d:rs:novel_5:MP-P:n\endcsname{5}
\expandafter\gdef\csname num@d:rs:novel_5:MP-P:worst\endcsname{1.6}
\expandafter\gdef\csname num@d:rs:novel_5:PGP-P\endcsname{+1.0}
\expandafter\gdef\csname num@d:rs:novel_5:PGP-P:lo\endcsname{+0.6}
\expandafter\gdef\csname num@d:rs:novel_5:PGP-P:hi\endcsname{+1.3}
\expandafter\gdef\csname num@d:rs:novel_5:PGP-P:abs\endcsname{1.0}
\expandafter\gdef\csname num@d:rs:novel_5:PGP-P:wins\endcsname{5}
\expandafter\gdef\csname num@d:rs:novel_5:PGP-P:n\endcsname{5}
\expandafter\gdef\csname num@d:rs:novel_5:PGP-P:worst\endcsname{1.3}
\expandafter\gdef\csname num@d:rs:novel_5:PMP-P\endcsname{+0.7}
\expandafter\gdef\csname num@d:rs:novel_5:PMP-P:lo\endcsname{+0.2}
\expandafter\gdef\csname num@d:rs:novel_5:PMP-P:hi\endcsname{+1.2}
\expandafter\gdef\csname num@d:rs:novel_5:PMP-P:abs\endcsname{0.7}
\expandafter\gdef\csname num@d:rs:novel_5:PMP-P:wins\endcsname{5}
\expandafter\gdef\csname num@d:rs:novel_5:PMP-P:n\endcsname{5}
\expandafter\gdef\csname num@d:rs:novel_5:PMP-P:worst\endcsname{1.2}
\expandafter\gdef\csname num@d:rs:novel_5:densem-dense\endcsname{\ensuremath{-}0.3}
\expandafter\gdef\csname num@d:rs:novel_5:densem-dense:lo\endcsname{\ensuremath{-}0.5}
\expandafter\gdef\csname num@d:rs:novel_5:densem-dense:hi\endcsname{\ensuremath{-}0.1}
\expandafter\gdef\csname num@d:rs:novel_5:densem-dense:abs\endcsname{0.3}
\expandafter\gdef\csname num@d:rs:novel_5:densem-dense:wins\endcsname{0}
\expandafter\gdef\csname num@d:rs:novel_5:densem-dense:n\endcsname{5}
\expandafter\gdef\csname num@d:rs:novel_5:densem-dense:worst\endcsname{0.5}
\expandafter\gdef\csname num@d:rs:novel_5:densermr-dense\endcsname{\ensuremath{-}0.2}
\expandafter\gdef\csname num@d:rs:novel_5:densermr-dense:lo\endcsname{\ensuremath{-}0.4}
\expandafter\gdef\csname num@d:rs:novel_5:densermr-dense:hi\endcsname{\ensuremath{-}0.0}
\expandafter\gdef\csname num@d:rs:novel_5:densermr-dense:abs\endcsname{0.2}
\expandafter\gdef\csname num@d:rs:novel_5:densermr-dense:wins\endcsname{0}
\expandafter\gdef\csname num@d:rs:novel_5:densermr-dense:n\endcsname{5}
\expandafter\gdef\csname num@d:rs:novel_5:densermr-dense:worst\endcsname{0.4}
\expandafter\gdef\csname num@d:rs:novel_5:dense-dacisP\endcsname{+7.1}
\expandafter\gdef\csname num@d:rs:novel_5:dense-dacisP:lo\endcsname{+5.9}
\expandafter\gdef\csname num@d:rs:novel_5:dense-dacisP:hi\endcsname{+8.4}
\expandafter\gdef\csname num@d:rs:novel_5:dense-dacisP:abs\endcsname{7.1}
\expandafter\gdef\csname num@d:rs:novel_5:dense-dacisP:wins\endcsname{5}
\expandafter\gdef\csname num@d:rs:novel_5:dense-dacisP:n\endcsname{5}
\expandafter\gdef\csname num@d:rs:novel_5:dense-dacisP:worst\endcsname{8.4}
\expandafter\gdef\csname num@d:rs:novel_5:dense-dacisPMP\endcsname{+6.4}
\expandafter\gdef\csname num@d:rs:novel_5:dense-dacisPMP:lo\endcsname{+5.0}
\expandafter\gdef\csname num@d:rs:novel_5:dense-dacisPMP:hi\endcsname{+7.8}
\expandafter\gdef\csname num@d:rs:novel_5:dense-dacisPMP:abs\endcsname{6.4}
\expandafter\gdef\csname num@d:rs:novel_5:dense-dacisPMP:wins\endcsname{5}
\expandafter\gdef\csname num@d:rs:novel_5:dense-dacisPMP:n\endcsname{5}
\expandafter\gdef\csname num@d:rs:novel_5:dense-dacisPMP:worst\endcsname{7.8}
\expandafter\gdef\csname num@d:rs:novel_5:dense-randomP\endcsname{+7.2}
\expandafter\gdef\csname num@d:rs:novel_5:dense-randomP:lo\endcsname{+4.9}
\expandafter\gdef\csname num@d:rs:novel_5:dense-randomP:hi\endcsname{+9.5}
\expandafter\gdef\csname num@d:rs:novel_5:dense-randomP:abs\endcsname{7.2}
\expandafter\gdef\csname num@d:rs:novel_5:dense-randomP:wins\endcsname{5}
\expandafter\gdef\csname num@d:rs:novel_5:dense-randomP:n\endcsname{5}
\expandafter\gdef\csname num@d:rs:novel_5:dense-randomP:worst\endcsname{9.5}
\expandafter\gdef\csname num@d:rs:novel_5:dense-randomPMP\endcsname{+5.9}
\expandafter\gdef\csname num@d:rs:novel_5:dense-randomPMP:lo\endcsname{+4.5}
\expandafter\gdef\csname num@d:rs:novel_5:dense-randomPMP:hi\endcsname{+7.3}
\expandafter\gdef\csname num@d:rs:novel_5:dense-randomPMP:abs\endcsname{5.9}
\expandafter\gdef\csname num@d:rs:novel_5:dense-randomPMP:wins\endcsname{5}
\expandafter\gdef\csname num@d:rs:novel_5:dense-randomPMP:n\endcsname{5}
\expandafter\gdef\csname num@d:rs:novel_5:dense-randomPMP:worst\endcsname{7.3}
\expandafter\gdef\csname num@d:rs:novel_5:mnv3-dacisPMP\endcsname{+7.7}
\expandafter\gdef\csname num@d:rs:novel_5:mnv3-dacisPMP:lo\endcsname{+6.2}
\expandafter\gdef\csname num@d:rs:novel_5:mnv3-dacisPMP:hi\endcsname{+9.1}
\expandafter\gdef\csname num@d:rs:novel_5:mnv3-dacisPMP:abs\endcsname{7.7}
\expandafter\gdef\csname num@d:rs:novel_5:mnv3-dacisPMP:wins\endcsname{5}
\expandafter\gdef\csname num@d:rs:novel_5:mnv3-dacisPMP:n\endcsname{5}
\expandafter\gdef\csname num@d:rs:novel_5:mnv3-dacisPMP:worst\endcsname{9.1}
\expandafter\gdef\csname num@d:rs:novel_5:effb0-dacisPMP\endcsname{+10.1}
\expandafter\gdef\csname num@d:rs:novel_5:effb0-dacisPMP:lo\endcsname{+7.7}
\expandafter\gdef\csname num@d:rs:novel_5:effb0-dacisPMP:hi\endcsname{+12.5}
\expandafter\gdef\csname num@d:rs:novel_5:effb0-dacisPMP:abs\endcsname{10.1}
\expandafter\gdef\csname num@d:rs:novel_5:effb0-dacisPMP:wins\endcsname{5}
\expandafter\gdef\csname num@d:rs:novel_5:effb0-dacisPMP:n\endcsname{5}
\expandafter\gdef\csname num@d:rs:novel_5:effb0-dacisPMP:worst\endcsname{12.5}
\expandafter\gdef\csname num@d:rs:novel_5:mnv2-dacisPMP\endcsname{+6.4}
\expandafter\gdef\csname num@d:rs:novel_5:mnv2-dacisPMP:lo\endcsname{+4.5}
\expandafter\gdef\csname num@d:rs:novel_5:mnv2-dacisPMP:hi\endcsname{+8.3}
\expandafter\gdef\csname num@d:rs:novel_5:mnv2-dacisPMP:abs\endcsname{6.4}
\expandafter\gdef\csname num@d:rs:novel_5:mnv2-dacisPMP:wins\endcsname{5}
\expandafter\gdef\csname num@d:rs:novel_5:mnv2-dacisPMP:n\endcsname{5}
\expandafter\gdef\csname num@d:rs:novel_5:mnv2-dacisPMP:worst\endcsname{8.3}
\expandafter\gdef\csname num@d:rs:novel_5:dacis-l1:PMP50\endcsname{\ensuremath{-}0.2}
\expandafter\gdef\csname num@d:rs:novel_5:dacis-l1:PMP50:lo\endcsname{\ensuremath{-}1.1}
\expandafter\gdef\csname num@d:rs:novel_5:dacis-l1:PMP50:hi\endcsname{+0.6}
\expandafter\gdef\csname num@d:rs:novel_5:dacis-l1:PMP50:abs\endcsname{0.2}
\expandafter\gdef\csname num@d:rs:novel_5:dacis-l1:PMP50:wins\endcsname{1}
\expandafter\gdef\csname num@d:rs:novel_5:dacis-l1:PMP50:n\endcsname{5}
\expandafter\gdef\csname num@d:rs:novel_5:dacis-l1:PMP50:worst\endcsname{1.1}
\expandafter\gdef\csname num@d:rs:novel_5:dacis-random:PMP50\endcsname{\ensuremath{-}0.3}
\expandafter\gdef\csname num@d:rs:novel_5:dacis-random:PMP50:lo\endcsname{\ensuremath{-}1.3}
\expandafter\gdef\csname num@d:rs:novel_5:dacis-random:PMP50:hi\endcsname{+0.7}
\expandafter\gdef\csname num@d:rs:novel_5:dacis-random:PMP50:abs\endcsname{0.3}
\expandafter\gdef\csname num@d:rs:novel_5:dacis-random:PMP50:wins\endcsname{2}
\expandafter\gdef\csname num@d:rs:novel_5:dacis-random:PMP50:n\endcsname{5}
\expandafter\gdef\csname num@d:rs:novel_5:dacis-random:PMP50:worst\endcsname{1.3}
\expandafter\gdef\csname num@d:rs:novel_5:dense-dacisPMP50\endcsname{+3.0}
\expandafter\gdef\csname num@d:rs:novel_5:dense-dacisPMP50:lo\endcsname{+2.5}
\expandafter\gdef\csname num@d:rs:novel_5:dense-dacisPMP50:hi\endcsname{+3.6}
\expandafter\gdef\csname num@d:rs:novel_5:dense-dacisPMP50:abs\endcsname{3.0}
\expandafter\gdef\csname num@d:rs:novel_5:dense-dacisPMP50:wins\endcsname{5}
\expandafter\gdef\csname num@d:rs:novel_5:dense-dacisPMP50:n\endcsname{5}
\expandafter\gdef\csname num@d:rs:novel_5:dense-dacisPMP50:worst\endcsname{3.6}
\expandafter\gdef\csname num@d:rs:novel_5:dacis-l1:PMP20\endcsname{\ensuremath{-}2.0}
\expandafter\gdef\csname num@d:rs:novel_5:dacis-l1:PMP20:lo\endcsname{\ensuremath{-}3.8}
\expandafter\gdef\csname num@d:rs:novel_5:dacis-l1:PMP20:hi\endcsname{\ensuremath{-}0.2}
\expandafter\gdef\csname num@d:rs:novel_5:dacis-l1:PMP20:abs\endcsname{2.0}
\expandafter\gdef\csname num@d:rs:novel_5:dacis-l1:PMP20:wins\endcsname{0}
\expandafter\gdef\csname num@d:rs:novel_5:dacis-l1:PMP20:n\endcsname{5}
\expandafter\gdef\csname num@d:rs:novel_5:dacis-l1:PMP20:worst\endcsname{3.8}
\expandafter\gdef\csname num@d:rs:novel_5:dacis-random:PMP20\endcsname{+0.3}
\expandafter\gdef\csname num@d:rs:novel_5:dacis-random:PMP20:lo\endcsname{\ensuremath{-}0.9}
\expandafter\gdef\csname num@d:rs:novel_5:dacis-random:PMP20:hi\endcsname{+1.5}
\expandafter\gdef\csname num@d:rs:novel_5:dacis-random:PMP20:abs\endcsname{0.3}
\expandafter\gdef\csname num@d:rs:novel_5:dacis-random:PMP20:wins\endcsname{3}
\expandafter\gdef\csname num@d:rs:novel_5:dacis-random:PMP20:n\endcsname{5}
\expandafter\gdef\csname num@d:rs:novel_5:dacis-random:PMP20:worst\endcsname{1.5}
\expandafter\gdef\csname num@d:rs:novel_5:dense-dacisPMP20\endcsname{+9.5}
\expandafter\gdef\csname num@d:rs:novel_5:dense-dacisPMP20:lo\endcsname{+6.6}
\expandafter\gdef\csname num@d:rs:novel_5:dense-dacisPMP20:hi\endcsname{+12.4}
\expandafter\gdef\csname num@d:rs:novel_5:dense-dacisPMP20:abs\endcsname{9.5}
\expandafter\gdef\csname num@d:rs:novel_5:dense-dacisPMP20:wins\endcsname{5}
\expandafter\gdef\csname num@d:rs:novel_5:dense-dacisPMP20:n\endcsname{5}
\expandafter\gdef\csname num@d:rs:novel_5:dense-dacisPMP20:worst\endcsname{12.4}
\expandafter\gdef\csname num@d:rs:novel_5:dacis-l1:PMP10\endcsname{+1.1}
\expandafter\gdef\csname num@d:rs:novel_5:dacis-l1:PMP10:lo\endcsname{+0.1}
\expandafter\gdef\csname num@d:rs:novel_5:dacis-l1:PMP10:hi\endcsname{+2.0}
\expandafter\gdef\csname num@d:rs:novel_5:dacis-l1:PMP10:abs\endcsname{1.1}
\expandafter\gdef\csname num@d:rs:novel_5:dacis-l1:PMP10:wins\endcsname{5}
\expandafter\gdef\csname num@d:rs:novel_5:dacis-l1:PMP10:n\endcsname{5}
\expandafter\gdef\csname num@d:rs:novel_5:dacis-l1:PMP10:worst\endcsname{2.0}
\expandafter\gdef\csname num@d:rs:novel_5:dacis-random:PMP10\endcsname{+2.9}
\expandafter\gdef\csname num@d:rs:novel_5:dacis-random:PMP10:lo\endcsname{\ensuremath{-}0.8}
\expandafter\gdef\csname num@d:rs:novel_5:dacis-random:PMP10:hi\endcsname{+6.7}
\expandafter\gdef\csname num@d:rs:novel_5:dacis-random:PMP10:abs\endcsname{2.9}
\expandafter\gdef\csname num@d:rs:novel_5:dacis-random:PMP10:wins\endcsname{4}
\expandafter\gdef\csname num@d:rs:novel_5:dacis-random:PMP10:n\endcsname{5}
\expandafter\gdef\csname num@d:rs:novel_5:dacis-random:PMP10:worst\endcsname{6.7}
\expandafter\gdef\csname num@d:rs:novel_5:dense-dacisPMP10\endcsname{+15.0}
\expandafter\gdef\csname num@d:rs:novel_5:dense-dacisPMP10:lo\endcsname{+12.8}
\expandafter\gdef\csname num@d:rs:novel_5:dense-dacisPMP10:hi\endcsname{+17.2}
\expandafter\gdef\csname num@d:rs:novel_5:dense-dacisPMP10:abs\endcsname{15.0}
\expandafter\gdef\csname num@d:rs:novel_5:dense-dacisPMP10:wins\endcsname{5}
\expandafter\gdef\csname num@d:rs:novel_5:dense-dacisPMP10:n\endcsname{5}
\expandafter\gdef\csname num@d:rs:novel_5:dense-dacisPMP10:worst\endcsname{17.2}
\expandafter\gdef\csname num@d:es:novel_5:dacis-random:PMP\endcsname{\ensuremath{-}0.4}
\expandafter\gdef\csname num@d:es:novel_5:dacis-random:PMP:lo\endcsname{\ensuremath{-}3.5}
\expandafter\gdef\csname num@d:es:novel_5:dacis-random:PMP:hi\endcsname{+2.7}
\expandafter\gdef\csname num@d:es:novel_5:dacis-random:PMP:abs\endcsname{0.4}
\expandafter\gdef\csname num@d:es:novel_5:dacis-random:PMP:wins\endcsname{3}
\expandafter\gdef\csname num@d:es:novel_5:dacis-random:PMP:n\endcsname{5}
\expandafter\gdef\csname num@d:es:novel_5:dacis-random:PMP:worst\endcsname{3.5}
\expandafter\gdef\csname num@d:es:novel_5:dacis-random:P\endcsname{+0.7}
\expandafter\gdef\csname num@d:es:novel_5:dacis-random:P:lo\endcsname{\ensuremath{-}1.6}
\expandafter\gdef\csname num@d:es:novel_5:dacis-random:P:hi\endcsname{+2.9}
\expandafter\gdef\csname num@d:es:novel_5:dacis-random:P:abs\endcsname{0.7}
\expandafter\gdef\csname num@d:es:novel_5:dacis-random:P:wins\endcsname{4}
\expandafter\gdef\csname num@d:es:novel_5:dacis-random:P:n\endcsname{5}
\expandafter\gdef\csname num@d:es:novel_5:dacis-random:P:worst\endcsname{2.9}
\expandafter\gdef\csname num@d:es:novel_5:dacis-l1:PMP\endcsname{+0.1}
\expandafter\gdef\csname num@d:es:novel_5:dacis-l1:PMP:lo\endcsname{\ensuremath{-}2.6}
\expandafter\gdef\csname num@d:es:novel_5:dacis-l1:PMP:hi\endcsname{+2.9}
\expandafter\gdef\csname num@d:es:novel_5:dacis-l1:PMP:abs\endcsname{0.1}
\expandafter\gdef\csname num@d:es:novel_5:dacis-l1:PMP:wins\endcsname{3}
\expandafter\gdef\csname num@d:es:novel_5:dacis-l1:PMP:n\endcsname{5}
\expandafter\gdef\csname num@d:es:novel_5:dacis-l1:PMP:worst\endcsname{2.9}
\expandafter\gdef\csname num@d:es:novel_5:dacis-l1:P\endcsname{+0.1}
\expandafter\gdef\csname num@d:es:novel_5:dacis-l1:P:lo\endcsname{\ensuremath{-}3.6}
\expandafter\gdef\csname num@d:es:novel_5:dacis-l1:P:hi\endcsname{+3.9}
\expandafter\gdef\csname num@d:es:novel_5:dacis-l1:P:abs\endcsname{0.1}
\expandafter\gdef\csname num@d:es:novel_5:dacis-l1:P:wins\endcsname{2}
\expandafter\gdef\csname num@d:es:novel_5:dacis-l1:P:n\endcsname{5}
\expandafter\gdef\csname num@d:es:novel_5:dacis-l1:P:worst\endcsname{3.9}
\expandafter\gdef\csname num@d:es:novel_5:dacis-taylor:PMP\endcsname{+0.2}
\expandafter\gdef\csname num@d:es:novel_5:dacis-taylor:PMP:lo\endcsname{\ensuremath{-}5.1}
\expandafter\gdef\csname num@d:es:novel_5:dacis-taylor:PMP:hi\endcsname{+5.4}
\expandafter\gdef\csname num@d:es:novel_5:dacis-taylor:PMP:abs\endcsname{0.2}
\expandafter\gdef\csname num@d:es:novel_5:dacis-taylor:PMP:wins\endcsname{3}
\expandafter\gdef\csname num@d:es:novel_5:dacis-taylor:PMP:n\endcsname{5}
\expandafter\gdef\csname num@d:es:novel_5:dacis-taylor:PMP:worst\endcsname{5.4}
\expandafter\gdef\csname num@d:es:novel_5:dacis-taylor:P\endcsname{\ensuremath{-}1.8}
\expandafter\gdef\csname num@d:es:novel_5:dacis-taylor:P:lo\endcsname{\ensuremath{-}5.2}
\expandafter\gdef\csname num@d:es:novel_5:dacis-taylor:P:hi\endcsname{+1.7}
\expandafter\gdef\csname num@d:es:novel_5:dacis-taylor:P:abs\endcsname{1.8}
\expandafter\gdef\csname num@d:es:novel_5:dacis-taylor:P:wins\endcsname{1}
\expandafter\gdef\csname num@d:es:novel_5:dacis-taylor:P:n\endcsname{5}
\expandafter\gdef\csname num@d:es:novel_5:dacis-taylor:P:worst\endcsname{5.2}
\expandafter\gdef\csname num@d:es:novel_5:dacis-bn_gamma:PMP\endcsname{+2.2}
\expandafter\gdef\csname num@d:es:novel_5:dacis-bn_gamma:PMP:lo\endcsname{+0.2}
\expandafter\gdef\csname num@d:es:novel_5:dacis-bn_gamma:PMP:hi\endcsname{+4.2}
\expandafter\gdef\csname num@d:es:novel_5:dacis-bn_gamma:PMP:abs\endcsname{2.2}
\expandafter\gdef\csname num@d:es:novel_5:dacis-bn_gamma:PMP:wins\endcsname{5}
\expandafter\gdef\csname num@d:es:novel_5:dacis-bn_gamma:PMP:n\endcsname{5}
\expandafter\gdef\csname num@d:es:novel_5:dacis-bn_gamma:PMP:worst\endcsname{4.2}
\expandafter\gdef\csname num@d:es:novel_5:dacis-D:PMP\endcsname{+1.1}
\expandafter\gdef\csname num@d:es:novel_5:dacis-D:PMP:lo\endcsname{\ensuremath{-}2.8}
\expandafter\gdef\csname num@d:es:novel_5:dacis-D:PMP:hi\endcsname{+5.0}
\expandafter\gdef\csname num@d:es:novel_5:dacis-D:PMP:abs\endcsname{1.1}
\expandafter\gdef\csname num@d:es:novel_5:dacis-D:PMP:wins\endcsname{4}
\expandafter\gdef\csname num@d:es:novel_5:dacis-D:PMP:n\endcsname{5}
\expandafter\gdef\csname num@d:es:novel_5:dacis-D:PMP:worst\endcsname{5.0}
\expandafter\gdef\csname num@d:es:novel_5:dacis-D:P\endcsname{+0.5}
\expandafter\gdef\csname num@d:es:novel_5:dacis-D:P:lo\endcsname{\ensuremath{-}0.7}
\expandafter\gdef\csname num@d:es:novel_5:dacis-D:P:hi\endcsname{+1.6}
\expandafter\gdef\csname num@d:es:novel_5:dacis-D:P:abs\endcsname{0.5}
\expandafter\gdef\csname num@d:es:novel_5:dacis-D:P:wins\endcsname{2}
\expandafter\gdef\csname num@d:es:novel_5:dacis-D:P:n\endcsname{5}
\expandafter\gdef\csname num@d:es:novel_5:dacis-D:P:worst\endcsname{1.6}
\expandafter\gdef\csname num@d:es:novel_5:dacis-G:PMP\endcsname{+0.7}
\expandafter\gdef\csname num@d:es:novel_5:dacis-G:PMP:lo\endcsname{\ensuremath{-}2.0}
\expandafter\gdef\csname num@d:es:novel_5:dacis-G:PMP:hi\endcsname{+3.4}
\expandafter\gdef\csname num@d:es:novel_5:dacis-G:PMP:abs\endcsname{0.7}
\expandafter\gdef\csname num@d:es:novel_5:dacis-G:PMP:wins\endcsname{3}
\expandafter\gdef\csname num@d:es:novel_5:dacis-G:PMP:n\endcsname{5}
\expandafter\gdef\csname num@d:es:novel_5:dacis-G:PMP:worst\endcsname{3.4}
\expandafter\gdef\csname num@d:es:novel_5:dacis-V:PMP\endcsname{+0.6}
\expandafter\gdef\csname num@d:es:novel_5:dacis-V:PMP:lo\endcsname{\ensuremath{-}2.4}
\expandafter\gdef\csname num@d:es:novel_5:dacis-V:PMP:hi\endcsname{+3.6}
\expandafter\gdef\csname num@d:es:novel_5:dacis-V:PMP:abs\endcsname{0.6}
\expandafter\gdef\csname num@d:es:novel_5:dacis-V:PMP:wins\endcsname{2}
\expandafter\gdef\csname num@d:es:novel_5:dacis-V:PMP:n\endcsname{5}
\expandafter\gdef\csname num@d:es:novel_5:dacis-V:PMP:worst\endcsname{3.6}
\expandafter\gdef\csname num@d:es:novel_5:l1-random:PMP\endcsname{\ensuremath{-}0.6}
\expandafter\gdef\csname num@d:es:novel_5:l1-random:PMP:lo\endcsname{\ensuremath{-}1.4}
\expandafter\gdef\csname num@d:es:novel_5:l1-random:PMP:hi\endcsname{+0.3}
\expandafter\gdef\csname num@d:es:novel_5:l1-random:PMP:abs\endcsname{0.6}
\expandafter\gdef\csname num@d:es:novel_5:l1-random:PMP:wins\endcsname{1}
\expandafter\gdef\csname num@d:es:novel_5:l1-random:PMP:n\endcsname{5}
\expandafter\gdef\csname num@d:es:novel_5:l1-random:PMP:worst\endcsname{1.4}
\expandafter\gdef\csname num@d:es:novel_5:l1-random:P\endcsname{+0.5}
\expandafter\gdef\csname num@d:es:novel_5:l1-random:P:lo\endcsname{\ensuremath{-}3.9}
\expandafter\gdef\csname num@d:es:novel_5:l1-random:P:hi\endcsname{+4.9}
\expandafter\gdef\csname num@d:es:novel_5:l1-random:P:abs\endcsname{0.5}
\expandafter\gdef\csname num@d:es:novel_5:l1-random:P:wins\endcsname{3}
\expandafter\gdef\csname num@d:es:novel_5:l1-random:P:n\endcsname{5}
\expandafter\gdef\csname num@d:es:novel_5:l1-random:P:worst\endcsname{4.9}
\expandafter\gdef\csname num@d:es:novel_5:taylor-random:PMP\endcsname{\ensuremath{-}0.6}
\expandafter\gdef\csname num@d:es:novel_5:taylor-random:PMP:lo\endcsname{\ensuremath{-}3.3}
\expandafter\gdef\csname num@d:es:novel_5:taylor-random:PMP:hi\endcsname{+2.2}
\expandafter\gdef\csname num@d:es:novel_5:taylor-random:PMP:abs\endcsname{0.6}
\expandafter\gdef\csname num@d:es:novel_5:taylor-random:PMP:wins\endcsname{3}
\expandafter\gdef\csname num@d:es:novel_5:taylor-random:PMP:n\endcsname{5}
\expandafter\gdef\csname num@d:es:novel_5:taylor-random:PMP:worst\endcsname{3.3}
\expandafter\gdef\csname num@d:es:novel_5:taylor-random:P\endcsname{+2.5}
\expandafter\gdef\csname num@d:es:novel_5:taylor-random:P:lo\endcsname{\ensuremath{-}0.9}
\expandafter\gdef\csname num@d:es:novel_5:taylor-random:P:hi\endcsname{+5.9}
\expandafter\gdef\csname num@d:es:novel_5:taylor-random:P:abs\endcsname{2.5}
\expandafter\gdef\csname num@d:es:novel_5:taylor-random:P:wins\endcsname{4}
\expandafter\gdef\csname num@d:es:novel_5:taylor-random:P:n\endcsname{5}
\expandafter\gdef\csname num@d:es:novel_5:taylor-random:P:worst\endcsname{5.9}
\expandafter\gdef\csname num@d:es:novel_5:bn_gamma-random:PMP\endcsname{\ensuremath{-}2.6}
\expandafter\gdef\csname num@d:es:novel_5:bn_gamma-random:PMP:lo\endcsname{\ensuremath{-}5.0}
\expandafter\gdef\csname num@d:es:novel_5:bn_gamma-random:PMP:hi\endcsname{\ensuremath{-}0.2}
\expandafter\gdef\csname num@d:es:novel_5:bn_gamma-random:PMP:abs\endcsname{2.6}
\expandafter\gdef\csname num@d:es:novel_5:bn_gamma-random:PMP:wins\endcsname{0}
\expandafter\gdef\csname num@d:es:novel_5:bn_gamma-random:PMP:n\endcsname{5}
\expandafter\gdef\csname num@d:es:novel_5:bn_gamma-random:PMP:worst\endcsname{5.0}
\expandafter\gdef\csname num@d:es:novel_5:D-random:PMP\endcsname{\ensuremath{-}1.5}
\expandafter\gdef\csname num@d:es:novel_5:D-random:PMP:lo\endcsname{\ensuremath{-}3.9}
\expandafter\gdef\csname num@d:es:novel_5:D-random:PMP:hi\endcsname{+0.8}
\expandafter\gdef\csname num@d:es:novel_5:D-random:PMP:abs\endcsname{1.5}
\expandafter\gdef\csname num@d:es:novel_5:D-random:PMP:wins\endcsname{2}
\expandafter\gdef\csname num@d:es:novel_5:D-random:PMP:n\endcsname{5}
\expandafter\gdef\csname num@d:es:novel_5:D-random:PMP:worst\endcsname{3.9}
\expandafter\gdef\csname num@d:es:novel_5:D-random:P\endcsname{+0.2}
\expandafter\gdef\csname num@d:es:novel_5:D-random:P:lo\endcsname{\ensuremath{-}2.5}
\expandafter\gdef\csname num@d:es:novel_5:D-random:P:hi\endcsname{+2.9}
\expandafter\gdef\csname num@d:es:novel_5:D-random:P:abs\endcsname{0.2}
\expandafter\gdef\csname num@d:es:novel_5:D-random:P:wins\endcsname{2}
\expandafter\gdef\csname num@d:es:novel_5:D-random:P:n\endcsname{5}
\expandafter\gdef\csname num@d:es:novel_5:D-random:P:worst\endcsname{2.9}
\expandafter\gdef\csname num@d:es:novel_5:G-random:PMP\endcsname{\ensuremath{-}1.1}
\expandafter\gdef\csname num@d:es:novel_5:G-random:PMP:lo\endcsname{\ensuremath{-}3.1}
\expandafter\gdef\csname num@d:es:novel_5:G-random:PMP:hi\endcsname{+0.8}
\expandafter\gdef\csname num@d:es:novel_5:G-random:PMP:abs\endcsname{1.1}
\expandafter\gdef\csname num@d:es:novel_5:G-random:PMP:wins\endcsname{1}
\expandafter\gdef\csname num@d:es:novel_5:G-random:PMP:n\endcsname{5}
\expandafter\gdef\csname num@d:es:novel_5:G-random:PMP:worst\endcsname{3.1}
\expandafter\gdef\csname num@d:es:novel_5:V-random:PMP\endcsname{\ensuremath{-}1.0}
\expandafter\gdef\csname num@d:es:novel_5:V-random:PMP:lo\endcsname{\ensuremath{-}3.2}
\expandafter\gdef\csname num@d:es:novel_5:V-random:PMP:hi\endcsname{+1.2}
\expandafter\gdef\csname num@d:es:novel_5:V-random:PMP:abs\endcsname{1.0}
\expandafter\gdef\csname num@d:es:novel_5:V-random:PMP:wins\endcsname{1}
\expandafter\gdef\csname num@d:es:novel_5:V-random:PMP:n\endcsname{5}
\expandafter\gdef\csname num@d:es:novel_5:V-random:PMP:worst\endcsname{3.2}
\expandafter\gdef\csname num@d:es:novel_5:PM-P\endcsname{\ensuremath{-}0.5}
\expandafter\gdef\csname num@d:es:novel_5:PM-P:lo\endcsname{\ensuremath{-}2.7}
\expandafter\gdef\csname num@d:es:novel_5:PM-P:hi\endcsname{+1.8}
\expandafter\gdef\csname num@d:es:novel_5:PM-P:abs\endcsname{0.5}
\expandafter\gdef\csname num@d:es:novel_5:PM-P:wins\endcsname{3}
\expandafter\gdef\csname num@d:es:novel_5:PM-P:n\endcsname{5}
\expandafter\gdef\csname num@d:es:novel_5:PM-P:worst\endcsname{2.7}
\expandafter\gdef\csname num@d:es:novel_5:MP-P\endcsname{+0.2}
\expandafter\gdef\csname num@d:es:novel_5:MP-P:lo\endcsname{\ensuremath{-}1.2}
\expandafter\gdef\csname num@d:es:novel_5:MP-P:hi\endcsname{+1.6}
\expandafter\gdef\csname num@d:es:novel_5:MP-P:abs\endcsname{0.2}
\expandafter\gdef\csname num@d:es:novel_5:MP-P:wins\endcsname{3}
\expandafter\gdef\csname num@d:es:novel_5:MP-P:n\endcsname{5}
\expandafter\gdef\csname num@d:es:novel_5:MP-P:worst\endcsname{1.6}
\expandafter\gdef\csname num@d:es:novel_5:PGP-P\endcsname{+0.3}
\expandafter\gdef\csname num@d:es:novel_5:PGP-P:lo\endcsname{\ensuremath{-}1.5}
\expandafter\gdef\csname num@d:es:novel_5:PGP-P:hi\endcsname{+2.1}
\expandafter\gdef\csname num@d:es:novel_5:PGP-P:abs\endcsname{0.3}
\expandafter\gdef\csname num@d:es:novel_5:PGP-P:wins\endcsname{3}
\expandafter\gdef\csname num@d:es:novel_5:PGP-P:n\endcsname{5}
\expandafter\gdef\csname num@d:es:novel_5:PGP-P:worst\endcsname{2.1}
\expandafter\gdef\csname num@d:es:novel_5:PMP-P\endcsname{+0.8}
\expandafter\gdef\csname num@d:es:novel_5:PMP-P:lo\endcsname{\ensuremath{-}1.9}
\expandafter\gdef\csname num@d:es:novel_5:PMP-P:hi\endcsname{+3.6}
\expandafter\gdef\csname num@d:es:novel_5:PMP-P:abs\endcsname{0.8}
\expandafter\gdef\csname num@d:es:novel_5:PMP-P:wins\endcsname{4}
\expandafter\gdef\csname num@d:es:novel_5:PMP-P:n\endcsname{5}
\expandafter\gdef\csname num@d:es:novel_5:PMP-P:worst\endcsname{3.6}
\expandafter\gdef\csname num@d:es:novel_5:densem-dense\endcsname{\ensuremath{-}0.8}
\expandafter\gdef\csname num@d:es:novel_5:densem-dense:lo\endcsname{\ensuremath{-}2.7}
\expandafter\gdef\csname num@d:es:novel_5:densem-dense:hi\endcsname{+1.0}
\expandafter\gdef\csname num@d:es:novel_5:densem-dense:abs\endcsname{0.8}
\expandafter\gdef\csname num@d:es:novel_5:densem-dense:wins\endcsname{1}
\expandafter\gdef\csname num@d:es:novel_5:densem-dense:n\endcsname{5}
\expandafter\gdef\csname num@d:es:novel_5:densem-dense:worst\endcsname{2.7}
\expandafter\gdef\csname num@d:es:novel_5:densermr-dense\endcsname{\ensuremath{-}0.5}
\expandafter\gdef\csname num@d:es:novel_5:densermr-dense:lo\endcsname{\ensuremath{-}1.1}
\expandafter\gdef\csname num@d:es:novel_5:densermr-dense:hi\endcsname{+0.1}
\expandafter\gdef\csname num@d:es:novel_5:densermr-dense:abs\endcsname{0.5}
\expandafter\gdef\csname num@d:es:novel_5:densermr-dense:wins\endcsname{1}
\expandafter\gdef\csname num@d:es:novel_5:densermr-dense:n\endcsname{5}
\expandafter\gdef\csname num@d:es:novel_5:densermr-dense:worst\endcsname{1.1}
\expandafter\gdef\csname num@d:es:novel_5:dense-dacisP\endcsname{+5.8}
\expandafter\gdef\csname num@d:es:novel_5:dense-dacisP:lo\endcsname{+3.8}
\expandafter\gdef\csname num@d:es:novel_5:dense-dacisP:hi\endcsname{+7.9}
\expandafter\gdef\csname num@d:es:novel_5:dense-dacisP:abs\endcsname{5.8}
\expandafter\gdef\csname num@d:es:novel_5:dense-dacisP:wins\endcsname{5}
\expandafter\gdef\csname num@d:es:novel_5:dense-dacisP:n\endcsname{5}
\expandafter\gdef\csname num@d:es:novel_5:dense-dacisP:worst\endcsname{7.9}
\expandafter\gdef\csname num@d:es:novel_5:dense-dacisPMP\endcsname{+5.0}
\expandafter\gdef\csname num@d:es:novel_5:dense-dacisPMP:lo\endcsname{+3.7}
\expandafter\gdef\csname num@d:es:novel_5:dense-dacisPMP:hi\endcsname{+6.4}
\expandafter\gdef\csname num@d:es:novel_5:dense-dacisPMP:abs\endcsname{5.0}
\expandafter\gdef\csname num@d:es:novel_5:dense-dacisPMP:wins\endcsname{5}
\expandafter\gdef\csname num@d:es:novel_5:dense-dacisPMP:n\endcsname{5}
\expandafter\gdef\csname num@d:es:novel_5:dense-dacisPMP:worst\endcsname{6.4}
\expandafter\gdef\csname num@d:es:novel_5:dense-randomP\endcsname{+6.5}
\expandafter\gdef\csname num@d:es:novel_5:dense-randomP:lo\endcsname{+4.2}
\expandafter\gdef\csname num@d:es:novel_5:dense-randomP:hi\endcsname{+8.8}
\expandafter\gdef\csname num@d:es:novel_5:dense-randomP:abs\endcsname{6.5}
\expandafter\gdef\csname num@d:es:novel_5:dense-randomP:wins\endcsname{5}
\expandafter\gdef\csname num@d:es:novel_5:dense-randomP:n\endcsname{5}
\expandafter\gdef\csname num@d:es:novel_5:dense-randomP:worst\endcsname{8.8}
\expandafter\gdef\csname num@d:es:novel_5:dense-randomPMP\endcsname{+4.6}
\expandafter\gdef\csname num@d:es:novel_5:dense-randomPMP:lo\endcsname{+2.6}
\expandafter\gdef\csname num@d:es:novel_5:dense-randomPMP:hi\endcsname{+6.5}
\expandafter\gdef\csname num@d:es:novel_5:dense-randomPMP:abs\endcsname{4.6}
\expandafter\gdef\csname num@d:es:novel_5:dense-randomPMP:wins\endcsname{5}
\expandafter\gdef\csname num@d:es:novel_5:dense-randomPMP:n\endcsname{5}
\expandafter\gdef\csname num@d:es:novel_5:dense-randomPMP:worst\endcsname{6.5}
\expandafter\gdef\csname num@d:es:novel_5:mnv3-dacisPMP\endcsname{+15.3}
\expandafter\gdef\csname num@d:es:novel_5:mnv3-dacisPMP:lo\endcsname{+9.4}
\expandafter\gdef\csname num@d:es:novel_5:mnv3-dacisPMP:hi\endcsname{+21.2}
\expandafter\gdef\csname num@d:es:novel_5:mnv3-dacisPMP:abs\endcsname{15.3}
\expandafter\gdef\csname num@d:es:novel_5:mnv3-dacisPMP:wins\endcsname{5}
\expandafter\gdef\csname num@d:es:novel_5:mnv3-dacisPMP:n\endcsname{5}
\expandafter\gdef\csname num@d:es:novel_5:mnv3-dacisPMP:worst\endcsname{21.2}
\expandafter\gdef\csname num@d:es:novel_5:effb0-dacisPMP\endcsname{+12.8}
\expandafter\gdef\csname num@d:es:novel_5:effb0-dacisPMP:lo\endcsname{+6.1}
\expandafter\gdef\csname num@d:es:novel_5:effb0-dacisPMP:hi\endcsname{+19.4}
\expandafter\gdef\csname num@d:es:novel_5:effb0-dacisPMP:abs\endcsname{12.8}
\expandafter\gdef\csname num@d:es:novel_5:effb0-dacisPMP:wins\endcsname{5}
\expandafter\gdef\csname num@d:es:novel_5:effb0-dacisPMP:n\endcsname{5}
\expandafter\gdef\csname num@d:es:novel_5:effb0-dacisPMP:worst\endcsname{19.4}
\expandafter\gdef\csname num@d:es:novel_5:mnv2-dacisPMP\endcsname{+8.3}
\expandafter\gdef\csname num@d:es:novel_5:mnv2-dacisPMP:lo\endcsname{+3.9}
\expandafter\gdef\csname num@d:es:novel_5:mnv2-dacisPMP:hi\endcsname{+12.7}
\expandafter\gdef\csname num@d:es:novel_5:mnv2-dacisPMP:abs\endcsname{8.3}
\expandafter\gdef\csname num@d:es:novel_5:mnv2-dacisPMP:wins\endcsname{5}
\expandafter\gdef\csname num@d:es:novel_5:mnv2-dacisPMP:n\endcsname{5}
\expandafter\gdef\csname num@d:es:novel_5:mnv2-dacisPMP:worst\endcsname{12.7}
\expandafter\gdef\csname num@d:es:novel_5:dacis-l1:PMP50\endcsname{\ensuremath{-}1.4}
\expandafter\gdef\csname num@d:es:novel_5:dacis-l1:PMP50:lo\endcsname{\ensuremath{-}4.2}
\expandafter\gdef\csname num@d:es:novel_5:dacis-l1:PMP50:hi\endcsname{+1.4}
\expandafter\gdef\csname num@d:es:novel_5:dacis-l1:PMP50:abs\endcsname{1.4}
\expandafter\gdef\csname num@d:es:novel_5:dacis-l1:PMP50:wins\endcsname{2}
\expandafter\gdef\csname num@d:es:novel_5:dacis-l1:PMP50:n\endcsname{5}
\expandafter\gdef\csname num@d:es:novel_5:dacis-l1:PMP50:worst\endcsname{4.2}
\expandafter\gdef\csname num@d:es:novel_5:dacis-random:PMP50\endcsname{\ensuremath{-}0.7}
\expandafter\gdef\csname num@d:es:novel_5:dacis-random:PMP50:lo\endcsname{\ensuremath{-}1.5}
\expandafter\gdef\csname num@d:es:novel_5:dacis-random:PMP50:hi\endcsname{+0.1}
\expandafter\gdef\csname num@d:es:novel_5:dacis-random:PMP50:abs\endcsname{0.7}
\expandafter\gdef\csname num@d:es:novel_5:dacis-random:PMP50:wins\endcsname{0}
\expandafter\gdef\csname num@d:es:novel_5:dacis-random:PMP50:n\endcsname{5}
\expandafter\gdef\csname num@d:es:novel_5:dacis-random:PMP50:worst\endcsname{1.5}
\expandafter\gdef\csname num@d:es:novel_5:dense-dacisPMP50\endcsname{+3.5}
\expandafter\gdef\csname num@d:es:novel_5:dense-dacisPMP50:lo\endcsname{+2.3}
\expandafter\gdef\csname num@d:es:novel_5:dense-dacisPMP50:hi\endcsname{+4.8}
\expandafter\gdef\csname num@d:es:novel_5:dense-dacisPMP50:abs\endcsname{3.5}
\expandafter\gdef\csname num@d:es:novel_5:dense-dacisPMP50:wins\endcsname{5}
\expandafter\gdef\csname num@d:es:novel_5:dense-dacisPMP50:n\endcsname{5}
\expandafter\gdef\csname num@d:es:novel_5:dense-dacisPMP50:worst\endcsname{4.8}
\expandafter\gdef\csname num@d:es:novel_5:dacis-l1:PMP20\endcsname{\ensuremath{-}0.0}
\expandafter\gdef\csname num@d:es:novel_5:dacis-l1:PMP20:lo\endcsname{\ensuremath{-}3.0}
\expandafter\gdef\csname num@d:es:novel_5:dacis-l1:PMP20:hi\endcsname{+2.9}
\expandafter\gdef\csname num@d:es:novel_5:dacis-l1:PMP20:abs\endcsname{0.0}
\expandafter\gdef\csname num@d:es:novel_5:dacis-l1:PMP20:wins\endcsname{2}
\expandafter\gdef\csname num@d:es:novel_5:dacis-l1:PMP20:n\endcsname{5}
\expandafter\gdef\csname num@d:es:novel_5:dacis-l1:PMP20:worst\endcsname{3.0}
\expandafter\gdef\csname num@d:es:novel_5:dacis-random:PMP20\endcsname{+0.0}
\expandafter\gdef\csname num@d:es:novel_5:dacis-random:PMP20:lo\endcsname{\ensuremath{-}2.3}
\expandafter\gdef\csname num@d:es:novel_5:dacis-random:PMP20:hi\endcsname{+2.4}
\expandafter\gdef\csname num@d:es:novel_5:dacis-random:PMP20:abs\endcsname{0.0}
\expandafter\gdef\csname num@d:es:novel_5:dacis-random:PMP20:wins\endcsname{2}
\expandafter\gdef\csname num@d:es:novel_5:dacis-random:PMP20:n\endcsname{5}
\expandafter\gdef\csname num@d:es:novel_5:dacis-random:PMP20:worst\endcsname{2.4}
\expandafter\gdef\csname num@d:es:novel_5:dense-dacisPMP20\endcsname{+7.6}
\expandafter\gdef\csname num@d:es:novel_5:dense-dacisPMP20:lo\endcsname{+3.9}
\expandafter\gdef\csname num@d:es:novel_5:dense-dacisPMP20:hi\endcsname{+11.3}
\expandafter\gdef\csname num@d:es:novel_5:dense-dacisPMP20:abs\endcsname{7.6}
\expandafter\gdef\csname num@d:es:novel_5:dense-dacisPMP20:wins\endcsname{5}
\expandafter\gdef\csname num@d:es:novel_5:dense-dacisPMP20:n\endcsname{5}
\expandafter\gdef\csname num@d:es:novel_5:dense-dacisPMP20:worst\endcsname{11.3}
\expandafter\gdef\csname num@d:es:novel_5:dacis-l1:PMP10\endcsname{+0.6}
\expandafter\gdef\csname num@d:es:novel_5:dacis-l1:PMP10:lo\endcsname{\ensuremath{-}3.5}
\expandafter\gdef\csname num@d:es:novel_5:dacis-l1:PMP10:hi\endcsname{+4.6}
\expandafter\gdef\csname num@d:es:novel_5:dacis-l1:PMP10:abs\endcsname{0.6}
\expandafter\gdef\csname num@d:es:novel_5:dacis-l1:PMP10:wins\endcsname{2}
\expandafter\gdef\csname num@d:es:novel_5:dacis-l1:PMP10:n\endcsname{5}
\expandafter\gdef\csname num@d:es:novel_5:dacis-l1:PMP10:worst\endcsname{4.6}
\expandafter\gdef\csname num@d:es:novel_5:dacis-random:PMP10\endcsname{+2.3}
\expandafter\gdef\csname num@d:es:novel_5:dacis-random:PMP10:lo\endcsname{\ensuremath{-}2.1}
\expandafter\gdef\csname num@d:es:novel_5:dacis-random:PMP10:hi\endcsname{+6.7}
\expandafter\gdef\csname num@d:es:novel_5:dacis-random:PMP10:abs\endcsname{2.3}
\expandafter\gdef\csname num@d:es:novel_5:dacis-random:PMP10:wins\endcsname{4}
\expandafter\gdef\csname num@d:es:novel_5:dacis-random:PMP10:n\endcsname{5}
\expandafter\gdef\csname num@d:es:novel_5:dacis-random:PMP10:worst\endcsname{6.7}
\expandafter\gdef\csname num@d:es:novel_5:dense-dacisPMP10\endcsname{+10.7}
\expandafter\gdef\csname num@d:es:novel_5:dense-dacisPMP10:lo\endcsname{+6.4}
\expandafter\gdef\csname num@d:es:novel_5:dense-dacisPMP10:hi\endcsname{+15.0}
\expandafter\gdef\csname num@d:es:novel_5:dense-dacisPMP10:abs\endcsname{10.7}
\expandafter\gdef\csname num@d:es:novel_5:dense-dacisPMP10:wins\endcsname{5}
\expandafter\gdef\csname num@d:es:novel_5:dense-dacisPMP10:n\endcsname{5}
\expandafter\gdef\csname num@d:es:novel_5:dense-dacisPMP10:worst\endcsname{15.0}
\expandafter\gdef\csname num@d:pv:novel_5:dacis-random:PMP\endcsname{\ensuremath{-}0.3}
\expandafter\gdef\csname num@d:pv:novel_5:dacis-random:PMP:lo\endcsname{\ensuremath{-}2.3}
\expandafter\gdef\csname num@d:pv:novel_5:dacis-random:PMP:hi\endcsname{+1.8}
\expandafter\gdef\csname num@d:pv:novel_5:dacis-random:PMP:abs\endcsname{0.3}
\expandafter\gdef\csname num@d:pv:novel_5:dacis-random:PMP:wins\endcsname{2}
\expandafter\gdef\csname num@d:pv:novel_5:dacis-random:PMP:n\endcsname{5}
\expandafter\gdef\csname num@d:pv:novel_5:dacis-random:PMP:worst\endcsname{2.3}
\expandafter\gdef\csname num@d:pv:novel_5:dacis-random:P\endcsname{+0.8}
\expandafter\gdef\csname num@d:pv:novel_5:dacis-random:P:lo\endcsname{\ensuremath{-}1.3}
\expandafter\gdef\csname num@d:pv:novel_5:dacis-random:P:hi\endcsname{+3.0}
\expandafter\gdef\csname num@d:pv:novel_5:dacis-random:P:abs\endcsname{0.8}
\expandafter\gdef\csname num@d:pv:novel_5:dacis-random:P:wins\endcsname{4}
\expandafter\gdef\csname num@d:pv:novel_5:dacis-random:P:n\endcsname{5}
\expandafter\gdef\csname num@d:pv:novel_5:dacis-random:P:worst\endcsname{3.0}
\expandafter\gdef\csname num@d:pv:novel_5:dacis-l1:PMP\endcsname{+0.2}
\expandafter\gdef\csname num@d:pv:novel_5:dacis-l1:PMP:lo\endcsname{\ensuremath{-}1.3}
\expandafter\gdef\csname num@d:pv:novel_5:dacis-l1:PMP:hi\endcsname{+1.7}
\expandafter\gdef\csname num@d:pv:novel_5:dacis-l1:PMP:abs\endcsname{0.2}
\expandafter\gdef\csname num@d:pv:novel_5:dacis-l1:PMP:wins\endcsname{2}
\expandafter\gdef\csname num@d:pv:novel_5:dacis-l1:PMP:n\endcsname{5}
\expandafter\gdef\csname num@d:pv:novel_5:dacis-l1:PMP:worst\endcsname{1.7}
\expandafter\gdef\csname num@d:pv:novel_5:dacis-l1:P\endcsname{+0.6}
\expandafter\gdef\csname num@d:pv:novel_5:dacis-l1:P:lo\endcsname{\ensuremath{-}1.9}
\expandafter\gdef\csname num@d:pv:novel_5:dacis-l1:P:hi\endcsname{+3.1}
\expandafter\gdef\csname num@d:pv:novel_5:dacis-l1:P:abs\endcsname{0.6}
\expandafter\gdef\csname num@d:pv:novel_5:dacis-l1:P:wins\endcsname{2}
\expandafter\gdef\csname num@d:pv:novel_5:dacis-l1:P:n\endcsname{5}
\expandafter\gdef\csname num@d:pv:novel_5:dacis-l1:P:worst\endcsname{3.1}
\expandafter\gdef\csname num@d:pv:novel_5:dacis-taylor:PMP\endcsname{\ensuremath{-}0.6}
\expandafter\gdef\csname num@d:pv:novel_5:dacis-taylor:PMP:lo\endcsname{\ensuremath{-}1.9}
\expandafter\gdef\csname num@d:pv:novel_5:dacis-taylor:PMP:hi\endcsname{+0.8}
\expandafter\gdef\csname num@d:pv:novel_5:dacis-taylor:PMP:abs\endcsname{0.6}
\expandafter\gdef\csname num@d:pv:novel_5:dacis-taylor:PMP:wins\endcsname{2}
\expandafter\gdef\csname num@d:pv:novel_5:dacis-taylor:PMP:n\endcsname{5}
\expandafter\gdef\csname num@d:pv:novel_5:dacis-taylor:PMP:worst\endcsname{1.9}
\expandafter\gdef\csname num@d:pv:novel_5:dacis-taylor:P\endcsname{\ensuremath{-}1.9}
\expandafter\gdef\csname num@d:pv:novel_5:dacis-taylor:P:lo\endcsname{\ensuremath{-}2.9}
\expandafter\gdef\csname num@d:pv:novel_5:dacis-taylor:P:hi\endcsname{\ensuremath{-}0.9}
\expandafter\gdef\csname num@d:pv:novel_5:dacis-taylor:P:abs\endcsname{1.9}
\expandafter\gdef\csname num@d:pv:novel_5:dacis-taylor:P:wins\endcsname{0}
\expandafter\gdef\csname num@d:pv:novel_5:dacis-taylor:P:n\endcsname{5}
\expandafter\gdef\csname num@d:pv:novel_5:dacis-taylor:P:worst\endcsname{2.9}
\expandafter\gdef\csname num@d:pv:novel_5:dacis-bn_gamma:PMP\endcsname{+1.5}
\expandafter\gdef\csname num@d:pv:novel_5:dacis-bn_gamma:PMP:lo\endcsname{\ensuremath{-}0.8}
\expandafter\gdef\csname num@d:pv:novel_5:dacis-bn_gamma:PMP:hi\endcsname{+3.8}
\expandafter\gdef\csname num@d:pv:novel_5:dacis-bn_gamma:PMP:abs\endcsname{1.5}
\expandafter\gdef\csname num@d:pv:novel_5:dacis-bn_gamma:PMP:wins\endcsname{4}
\expandafter\gdef\csname num@d:pv:novel_5:dacis-bn_gamma:PMP:n\endcsname{5}
\expandafter\gdef\csname num@d:pv:novel_5:dacis-bn_gamma:PMP:worst\endcsname{3.8}
\expandafter\gdef\csname num@d:pv:novel_5:dacis-bn_gamma:P\endcsname{+1.4}
\expandafter\gdef\csname num@d:pv:novel_5:dacis-bn_gamma:P:lo\endcsname{\ensuremath{-}2.5}
\expandafter\gdef\csname num@d:pv:novel_5:dacis-bn_gamma:P:hi\endcsname{+5.3}
\expandafter\gdef\csname num@d:pv:novel_5:dacis-bn_gamma:P:abs\endcsname{1.4}
\expandafter\gdef\csname num@d:pv:novel_5:dacis-bn_gamma:P:wins\endcsname{2}
\expandafter\gdef\csname num@d:pv:novel_5:dacis-bn_gamma:P:n\endcsname{5}
\expandafter\gdef\csname num@d:pv:novel_5:dacis-bn_gamma:P:worst\endcsname{5.3}
\expandafter\gdef\csname num@d:pv:novel_5:dacis-D:PMP\endcsname{+1.3}
\expandafter\gdef\csname num@d:pv:novel_5:dacis-D:PMP:lo\endcsname{\ensuremath{-}0.3}
\expandafter\gdef\csname num@d:pv:novel_5:dacis-D:PMP:hi\endcsname{+2.9}
\expandafter\gdef\csname num@d:pv:novel_5:dacis-D:PMP:abs\endcsname{1.3}
\expandafter\gdef\csname num@d:pv:novel_5:dacis-D:PMP:wins\endcsname{3}
\expandafter\gdef\csname num@d:pv:novel_5:dacis-D:PMP:n\endcsname{5}
\expandafter\gdef\csname num@d:pv:novel_5:dacis-D:PMP:worst\endcsname{2.9}
\expandafter\gdef\csname num@d:pv:novel_5:dacis-D:P\endcsname{+1.3}
\expandafter\gdef\csname num@d:pv:novel_5:dacis-D:P:lo\endcsname{\ensuremath{-}0.4}
\expandafter\gdef\csname num@d:pv:novel_5:dacis-D:P:hi\endcsname{+3.0}
\expandafter\gdef\csname num@d:pv:novel_5:dacis-D:P:abs\endcsname{1.3}
\expandafter\gdef\csname num@d:pv:novel_5:dacis-D:P:wins\endcsname{3}
\expandafter\gdef\csname num@d:pv:novel_5:dacis-D:P:n\endcsname{5}
\expandafter\gdef\csname num@d:pv:novel_5:dacis-D:P:worst\endcsname{3.0}
\expandafter\gdef\csname num@d:pv:novel_5:dacis-G:PMP\endcsname{\ensuremath{-}0.2}
\expandafter\gdef\csname num@d:pv:novel_5:dacis-G:PMP:lo\endcsname{\ensuremath{-}0.7}
\expandafter\gdef\csname num@d:pv:novel_5:dacis-G:PMP:hi\endcsname{+0.4}
\expandafter\gdef\csname num@d:pv:novel_5:dacis-G:PMP:abs\endcsname{0.2}
\expandafter\gdef\csname num@d:pv:novel_5:dacis-G:PMP:wins\endcsname{2}
\expandafter\gdef\csname num@d:pv:novel_5:dacis-G:PMP:n\endcsname{5}
\expandafter\gdef\csname num@d:pv:novel_5:dacis-G:PMP:worst\endcsname{0.7}
\expandafter\gdef\csname num@d:pv:novel_5:dacis-G:P\endcsname{\ensuremath{-}0.9}
\expandafter\gdef\csname num@d:pv:novel_5:dacis-G:P:lo\endcsname{\ensuremath{-}2.0}
\expandafter\gdef\csname num@d:pv:novel_5:dacis-G:P:hi\endcsname{+0.2}
\expandafter\gdef\csname num@d:pv:novel_5:dacis-G:P:abs\endcsname{0.9}
\expandafter\gdef\csname num@d:pv:novel_5:dacis-G:P:wins\endcsname{1}
\expandafter\gdef\csname num@d:pv:novel_5:dacis-G:P:n\endcsname{5}
\expandafter\gdef\csname num@d:pv:novel_5:dacis-G:P:worst\endcsname{2.0}
\expandafter\gdef\csname num@d:pv:novel_5:dacis-V:PMP\endcsname{+0.1}
\expandafter\gdef\csname num@d:pv:novel_5:dacis-V:PMP:lo\endcsname{\ensuremath{-}0.7}
\expandafter\gdef\csname num@d:pv:novel_5:dacis-V:PMP:hi\endcsname{+0.8}
\expandafter\gdef\csname num@d:pv:novel_5:dacis-V:PMP:abs\endcsname{0.1}
\expandafter\gdef\csname num@d:pv:novel_5:dacis-V:PMP:wins\endcsname{3}
\expandafter\gdef\csname num@d:pv:novel_5:dacis-V:PMP:n\endcsname{5}
\expandafter\gdef\csname num@d:pv:novel_5:dacis-V:PMP:worst\endcsname{0.8}
\expandafter\gdef\csname num@d:pv:novel_5:dacis-V:P\endcsname{\ensuremath{-}0.0}
\expandafter\gdef\csname num@d:pv:novel_5:dacis-V:P:lo\endcsname{\ensuremath{-}1.8}
\expandafter\gdef\csname num@d:pv:novel_5:dacis-V:P:hi\endcsname{+1.8}
\expandafter\gdef\csname num@d:pv:novel_5:dacis-V:P:abs\endcsname{0.0}
\expandafter\gdef\csname num@d:pv:novel_5:dacis-V:P:wins\endcsname{2}
\expandafter\gdef\csname num@d:pv:novel_5:dacis-V:P:n\endcsname{5}
\expandafter\gdef\csname num@d:pv:novel_5:dacis-V:P:worst\endcsname{1.8}
\expandafter\gdef\csname num@d:pv:novel_5:l1-random:PMP\endcsname{\ensuremath{-}0.5}
\expandafter\gdef\csname num@d:pv:novel_5:l1-random:PMP:lo\endcsname{\ensuremath{-}2.5}
\expandafter\gdef\csname num@d:pv:novel_5:l1-random:PMP:hi\endcsname{+1.6}
\expandafter\gdef\csname num@d:pv:novel_5:l1-random:PMP:abs\endcsname{0.5}
\expandafter\gdef\csname num@d:pv:novel_5:l1-random:PMP:wins\endcsname{2}
\expandafter\gdef\csname num@d:pv:novel_5:l1-random:PMP:n\endcsname{5}
\expandafter\gdef\csname num@d:pv:novel_5:l1-random:PMP:worst\endcsname{2.5}
\expandafter\gdef\csname num@d:pv:novel_5:l1-random:P\endcsname{+0.2}
\expandafter\gdef\csname num@d:pv:novel_5:l1-random:P:lo\endcsname{\ensuremath{-}1.9}
\expandafter\gdef\csname num@d:pv:novel_5:l1-random:P:hi\endcsname{+2.4}
\expandafter\gdef\csname num@d:pv:novel_5:l1-random:P:abs\endcsname{0.2}
\expandafter\gdef\csname num@d:pv:novel_5:l1-random:P:wins\endcsname{3}
\expandafter\gdef\csname num@d:pv:novel_5:l1-random:P:n\endcsname{5}
\expandafter\gdef\csname num@d:pv:novel_5:l1-random:P:worst\endcsname{2.4}
\expandafter\gdef\csname num@d:pv:novel_5:taylor-random:PMP\endcsname{+0.3}
\expandafter\gdef\csname num@d:pv:novel_5:taylor-random:PMP:lo\endcsname{\ensuremath{-}1.6}
\expandafter\gdef\csname num@d:pv:novel_5:taylor-random:PMP:hi\endcsname{+2.2}
\expandafter\gdef\csname num@d:pv:novel_5:taylor-random:PMP:abs\endcsname{0.3}
\expandafter\gdef\csname num@d:pv:novel_5:taylor-random:PMP:wins\endcsname{2}
\expandafter\gdef\csname num@d:pv:novel_5:taylor-random:PMP:n\endcsname{5}
\expandafter\gdef\csname num@d:pv:novel_5:taylor-random:PMP:worst\endcsname{2.2}
\expandafter\gdef\csname num@d:pv:novel_5:taylor-random:P\endcsname{+2.7}
\expandafter\gdef\csname num@d:pv:novel_5:taylor-random:P:lo\endcsname{+0.6}
\expandafter\gdef\csname num@d:pv:novel_5:taylor-random:P:hi\endcsname{+4.9}
\expandafter\gdef\csname num@d:pv:novel_5:taylor-random:P:abs\endcsname{2.7}
\expandafter\gdef\csname num@d:pv:novel_5:taylor-random:P:wins\endcsname{5}
\expandafter\gdef\csname num@d:pv:novel_5:taylor-random:P:n\endcsname{5}
\expandafter\gdef\csname num@d:pv:novel_5:taylor-random:P:worst\endcsname{4.9}
\expandafter\gdef\csname num@d:pv:novel_5:bn_gamma-random:PMP\endcsname{\ensuremath{-}1.7}
\expandafter\gdef\csname num@d:pv:novel_5:bn_gamma-random:PMP:lo\endcsname{\ensuremath{-}3.5}
\expandafter\gdef\csname num@d:pv:novel_5:bn_gamma-random:PMP:hi\endcsname{+0.0}
\expandafter\gdef\csname num@d:pv:novel_5:bn_gamma-random:PMP:abs\endcsname{1.7}
\expandafter\gdef\csname num@d:pv:novel_5:bn_gamma-random:PMP:wins\endcsname{1}
\expandafter\gdef\csname num@d:pv:novel_5:bn_gamma-random:PMP:n\endcsname{5}
\expandafter\gdef\csname num@d:pv:novel_5:bn_gamma-random:PMP:worst\endcsname{3.5}
\expandafter\gdef\csname num@d:pv:novel_5:bn_gamma-random:P\endcsname{\ensuremath{-}0.5}
\expandafter\gdef\csname num@d:pv:novel_5:bn_gamma-random:P:lo\endcsname{\ensuremath{-}4.6}
\expandafter\gdef\csname num@d:pv:novel_5:bn_gamma-random:P:hi\endcsname{+3.5}
\expandafter\gdef\csname num@d:pv:novel_5:bn_gamma-random:P:abs\endcsname{0.5}
\expandafter\gdef\csname num@d:pv:novel_5:bn_gamma-random:P:wins\endcsname{3}
\expandafter\gdef\csname num@d:pv:novel_5:bn_gamma-random:P:n\endcsname{5}
\expandafter\gdef\csname num@d:pv:novel_5:bn_gamma-random:P:worst\endcsname{4.6}
\expandafter\gdef\csname num@d:pv:novel_5:D-random:PMP\endcsname{\ensuremath{-}1.6}
\expandafter\gdef\csname num@d:pv:novel_5:D-random:PMP:lo\endcsname{\ensuremath{-}4.8}
\expandafter\gdef\csname num@d:pv:novel_5:D-random:PMP:hi\endcsname{+1.7}
\expandafter\gdef\csname num@d:pv:novel_5:D-random:PMP:abs\endcsname{1.6}
\expandafter\gdef\csname num@d:pv:novel_5:D-random:PMP:wins\endcsname{1}
\expandafter\gdef\csname num@d:pv:novel_5:D-random:PMP:n\endcsname{5}
\expandafter\gdef\csname num@d:pv:novel_5:D-random:PMP:worst\endcsname{4.8}
\expandafter\gdef\csname num@d:pv:novel_5:D-random:P\endcsname{\ensuremath{-}0.4}
\expandafter\gdef\csname num@d:pv:novel_5:D-random:P:lo\endcsname{\ensuremath{-}2.6}
\expandafter\gdef\csname num@d:pv:novel_5:D-random:P:hi\endcsname{+1.7}
\expandafter\gdef\csname num@d:pv:novel_5:D-random:P:abs\endcsname{0.4}
\expandafter\gdef\csname num@d:pv:novel_5:D-random:P:wins\endcsname{2}
\expandafter\gdef\csname num@d:pv:novel_5:D-random:P:n\endcsname{5}
\expandafter\gdef\csname num@d:pv:novel_5:D-random:P:worst\endcsname{2.6}
\expandafter\gdef\csname num@d:pv:novel_5:G-random:PMP\endcsname{\ensuremath{-}0.1}
\expandafter\gdef\csname num@d:pv:novel_5:G-random:PMP:lo\endcsname{\ensuremath{-}2.2}
\expandafter\gdef\csname num@d:pv:novel_5:G-random:PMP:hi\endcsname{+2.0}
\expandafter\gdef\csname num@d:pv:novel_5:G-random:PMP:abs\endcsname{0.1}
\expandafter\gdef\csname num@d:pv:novel_5:G-random:PMP:wins\endcsname{3}
\expandafter\gdef\csname num@d:pv:novel_5:G-random:PMP:n\endcsname{5}
\expandafter\gdef\csname num@d:pv:novel_5:G-random:PMP:worst\endcsname{2.2}
\expandafter\gdef\csname num@d:pv:novel_5:G-random:P\endcsname{+1.7}
\expandafter\gdef\csname num@d:pv:novel_5:G-random:P:lo\endcsname{\ensuremath{-}0.1}
\expandafter\gdef\csname num@d:pv:novel_5:G-random:P:hi\endcsname{+3.5}
\expandafter\gdef\csname num@d:pv:novel_5:G-random:P:abs\endcsname{1.7}
\expandafter\gdef\csname num@d:pv:novel_5:G-random:P:wins\endcsname{5}
\expandafter\gdef\csname num@d:pv:novel_5:G-random:P:n\endcsname{5}
\expandafter\gdef\csname num@d:pv:novel_5:G-random:P:worst\endcsname{3.5}
\expandafter\gdef\csname num@d:pv:novel_5:V-random:PMP\endcsname{\ensuremath{-}0.3}
\expandafter\gdef\csname num@d:pv:novel_5:V-random:PMP:lo\endcsname{\ensuremath{-}2.3}
\expandafter\gdef\csname num@d:pv:novel_5:V-random:PMP:hi\endcsname{+1.6}
\expandafter\gdef\csname num@d:pv:novel_5:V-random:PMP:abs\endcsname{0.3}
\expandafter\gdef\csname num@d:pv:novel_5:V-random:PMP:wins\endcsname{1}
\expandafter\gdef\csname num@d:pv:novel_5:V-random:PMP:n\endcsname{5}
\expandafter\gdef\csname num@d:pv:novel_5:V-random:PMP:worst\endcsname{2.3}
\expandafter\gdef\csname num@d:pv:novel_5:V-random:P\endcsname{+0.8}
\expandafter\gdef\csname num@d:pv:novel_5:V-random:P:lo\endcsname{\ensuremath{-}1.1}
\expandafter\gdef\csname num@d:pv:novel_5:V-random:P:hi\endcsname{+2.8}
\expandafter\gdef\csname num@d:pv:novel_5:V-random:P:abs\endcsname{0.8}
\expandafter\gdef\csname num@d:pv:novel_5:V-random:P:wins\endcsname{3}
\expandafter\gdef\csname num@d:pv:novel_5:V-random:P:n\endcsname{5}
\expandafter\gdef\csname num@d:pv:novel_5:V-random:P:worst\endcsname{2.8}
\expandafter\gdef\csname num@d:pv:novel_5:PM-P\endcsname{\ensuremath{-}0.7}
\expandafter\gdef\csname num@d:pv:novel_5:PM-P:lo\endcsname{\ensuremath{-}2.4}
\expandafter\gdef\csname num@d:pv:novel_5:PM-P:hi\endcsname{+1.0}
\expandafter\gdef\csname num@d:pv:novel_5:PM-P:abs\endcsname{0.7}
\expandafter\gdef\csname num@d:pv:novel_5:PM-P:wins\endcsname{2}
\expandafter\gdef\csname num@d:pv:novel_5:PM-P:n\endcsname{5}
\expandafter\gdef\csname num@d:pv:novel_5:PM-P:worst\endcsname{2.4}
\expandafter\gdef\csname num@d:pv:novel_5:MP-P\endcsname{+0.3}
\expandafter\gdef\csname num@d:pv:novel_5:MP-P:lo\endcsname{\ensuremath{-}1.2}
\expandafter\gdef\csname num@d:pv:novel_5:MP-P:hi\endcsname{+1.7}
\expandafter\gdef\csname num@d:pv:novel_5:MP-P:abs\endcsname{0.3}
\expandafter\gdef\csname num@d:pv:novel_5:MP-P:wins\endcsname{2}
\expandafter\gdef\csname num@d:pv:novel_5:MP-P:n\endcsname{5}
\expandafter\gdef\csname num@d:pv:novel_5:MP-P:worst\endcsname{1.7}
\expandafter\gdef\csname num@d:pv:novel_5:PGP-P\endcsname{+0.3}
\expandafter\gdef\csname num@d:pv:novel_5:PGP-P:lo\endcsname{\ensuremath{-}0.4}
\expandafter\gdef\csname num@d:pv:novel_5:PGP-P:hi\endcsname{+1.0}
\expandafter\gdef\csname num@d:pv:novel_5:PGP-P:abs\endcsname{0.3}
\expandafter\gdef\csname num@d:pv:novel_5:PGP-P:wins\endcsname{3}
\expandafter\gdef\csname num@d:pv:novel_5:PGP-P:n\endcsname{5}
\expandafter\gdef\csname num@d:pv:novel_5:PGP-P:worst\endcsname{1.0}
\expandafter\gdef\csname num@d:pv:novel_5:PMPM-P\endcsname{+0.1}
\expandafter\gdef\csname num@d:pv:novel_5:PMPM-P:lo\endcsname{\ensuremath{-}1.6}
\expandafter\gdef\csname num@d:pv:novel_5:PMPM-P:hi\endcsname{+1.7}
\expandafter\gdef\csname num@d:pv:novel_5:PMPM-P:abs\endcsname{0.1}
\expandafter\gdef\csname num@d:pv:novel_5:PMPM-P:wins\endcsname{2}
\expandafter\gdef\csname num@d:pv:novel_5:PMPM-P:n\endcsname{5}
\expandafter\gdef\csname num@d:pv:novel_5:PMPM-P:worst\endcsname{1.7}
\expandafter\gdef\csname num@d:pv:novel_5:PM2x-P\endcsname{\ensuremath{-}0.8}
\expandafter\gdef\csname num@d:pv:novel_5:PM2x-P:lo\endcsname{\ensuremath{-}2.3}
\expandafter\gdef\csname num@d:pv:novel_5:PM2x-P:hi\endcsname{+0.6}
\expandafter\gdef\csname num@d:pv:novel_5:PM2x-P:abs\endcsname{0.8}
\expandafter\gdef\csname num@d:pv:novel_5:PM2x-P:wins\endcsname{1}
\expandafter\gdef\csname num@d:pv:novel_5:PM2x-P:n\endcsname{5}
\expandafter\gdef\csname num@d:pv:novel_5:PM2x-P:worst\endcsname{2.3}
\expandafter\gdef\csname num@d:pv:novel_5:PMP-P\endcsname{+0.1}
\expandafter\gdef\csname num@d:pv:novel_5:PMP-P:lo\endcsname{\ensuremath{-}0.7}
\expandafter\gdef\csname num@d:pv:novel_5:PMP-P:hi\endcsname{+0.9}
\expandafter\gdef\csname num@d:pv:novel_5:PMP-P:abs\endcsname{0.1}
\expandafter\gdef\csname num@d:pv:novel_5:PMP-P:wins\endcsname{2}
\expandafter\gdef\csname num@d:pv:novel_5:PMP-P:n\endcsname{5}
\expandafter\gdef\csname num@d:pv:novel_5:PMP-P:worst\endcsname{0.9}
\expandafter\gdef\csname num@d:pv:novel_5:densem-dense\endcsname{\ensuremath{-}0.3}
\expandafter\gdef\csname num@d:pv:novel_5:densem-dense:lo\endcsname{\ensuremath{-}0.6}
\expandafter\gdef\csname num@d:pv:novel_5:densem-dense:hi\endcsname{+0.1}
\expandafter\gdef\csname num@d:pv:novel_5:densem-dense:abs\endcsname{0.3}
\expandafter\gdef\csname num@d:pv:novel_5:densem-dense:wins\endcsname{1}
\expandafter\gdef\csname num@d:pv:novel_5:densem-dense:n\endcsname{5}
\expandafter\gdef\csname num@d:pv:novel_5:densem-dense:worst\endcsname{0.6}
\expandafter\gdef\csname num@d:pv:novel_5:densermr-dense\endcsname{\ensuremath{-}0.2}
\expandafter\gdef\csname num@d:pv:novel_5:densermr-dense:lo\endcsname{\ensuremath{-}0.3}
\expandafter\gdef\csname num@d:pv:novel_5:densermr-dense:hi\endcsname{\ensuremath{-}0.1}
\expandafter\gdef\csname num@d:pv:novel_5:densermr-dense:abs\endcsname{0.2}
\expandafter\gdef\csname num@d:pv:novel_5:densermr-dense:wins\endcsname{0}
\expandafter\gdef\csname num@d:pv:novel_5:densermr-dense:n\endcsname{5}
\expandafter\gdef\csname num@d:pv:novel_5:densermr-dense:worst\endcsname{0.3}
\expandafter\gdef\csname num@d:pv:novel_5:dense-dacisP\endcsname{+5.0}
\expandafter\gdef\csname num@d:pv:novel_5:dense-dacisP:lo\endcsname{+2.4}
\expandafter\gdef\csname num@d:pv:novel_5:dense-dacisP:hi\endcsname{+7.5}
\expandafter\gdef\csname num@d:pv:novel_5:dense-dacisP:abs\endcsname{5.0}
\expandafter\gdef\csname num@d:pv:novel_5:dense-dacisP:wins\endcsname{5}
\expandafter\gdef\csname num@d:pv:novel_5:dense-dacisP:n\endcsname{5}
\expandafter\gdef\csname num@d:pv:novel_5:dense-dacisP:worst\endcsname{7.5}
\expandafter\gdef\csname num@d:pv:novel_5:dense-dacisPMP\endcsname{+4.9}
\expandafter\gdef\csname num@d:pv:novel_5:dense-dacisPMP:lo\endcsname{+2.4}
\expandafter\gdef\csname num@d:pv:novel_5:dense-dacisPMP:hi\endcsname{+7.3}
\expandafter\gdef\csname num@d:pv:novel_5:dense-dacisPMP:abs\endcsname{4.9}
\expandafter\gdef\csname num@d:pv:novel_5:dense-dacisPMP:wins\endcsname{5}
\expandafter\gdef\csname num@d:pv:novel_5:dense-dacisPMP:n\endcsname{5}
\expandafter\gdef\csname num@d:pv:novel_5:dense-dacisPMP:worst\endcsname{7.3}
\expandafter\gdef\csname num@d:pv:novel_5:dense-randomP\endcsname{+5.8}
\expandafter\gdef\csname num@d:pv:novel_5:dense-randomP:lo\endcsname{+2.3}
\expandafter\gdef\csname num@d:pv:novel_5:dense-randomP:hi\endcsname{+9.3}
\expandafter\gdef\csname num@d:pv:novel_5:dense-randomP:abs\endcsname{5.8}
\expandafter\gdef\csname num@d:pv:novel_5:dense-randomP:wins\endcsname{5}
\expandafter\gdef\csname num@d:pv:novel_5:dense-randomP:n\endcsname{5}
\expandafter\gdef\csname num@d:pv:novel_5:dense-randomP:worst\endcsname{9.3}
\expandafter\gdef\csname num@d:pv:novel_5:dense-randomPMP\endcsname{+4.6}
\expandafter\gdef\csname num@d:pv:novel_5:dense-randomPMP:lo\endcsname{+1.2}
\expandafter\gdef\csname num@d:pv:novel_5:dense-randomPMP:hi\endcsname{+7.9}
\expandafter\gdef\csname num@d:pv:novel_5:dense-randomPMP:abs\endcsname{4.6}
\expandafter\gdef\csname num@d:pv:novel_5:dense-randomPMP:wins\endcsname{5}
\expandafter\gdef\csname num@d:pv:novel_5:dense-randomPMP:n\endcsname{5}
\expandafter\gdef\csname num@d:pv:novel_5:dense-randomPMP:worst\endcsname{7.9}
\expandafter\gdef\csname num@d:pv:novel_5:norefine-full\endcsname{\ensuremath{-}0.9}
\expandafter\gdef\csname num@d:pv:novel_5:norefine-full:lo\endcsname{\ensuremath{-}1.7}
\expandafter\gdef\csname num@d:pv:novel_5:norefine-full:hi\endcsname{\ensuremath{-}0.1}
\expandafter\gdef\csname num@d:pv:novel_5:norefine-full:abs\endcsname{0.9}
\expandafter\gdef\csname num@d:pv:novel_5:norefine-full:wins\endcsname{1}
\expandafter\gdef\csname num@d:pv:novel_5:norefine-full:n\endcsname{5}
\expandafter\gdef\csname num@d:pv:novel_5:norefine-full:worst\endcsname{1.7}
\expandafter\gdef\csname num@d:pv:novel_5:mnv3-dacisPMP\endcsname{+4.5}
\expandafter\gdef\csname num@d:pv:novel_5:mnv3-dacisPMP:lo\endcsname{+1.9}
\expandafter\gdef\csname num@d:pv:novel_5:mnv3-dacisPMP:hi\endcsname{+7.1}
\expandafter\gdef\csname num@d:pv:novel_5:mnv3-dacisPMP:abs\endcsname{4.5}
\expandafter\gdef\csname num@d:pv:novel_5:mnv3-dacisPMP:wins\endcsname{5}
\expandafter\gdef\csname num@d:pv:novel_5:mnv3-dacisPMP:n\endcsname{5}
\expandafter\gdef\csname num@d:pv:novel_5:mnv3-dacisPMP:worst\endcsname{7.1}
\expandafter\gdef\csname num@d:pv:novel_5:effb0-dacisPMP\endcsname{+7.2}
\expandafter\gdef\csname num@d:pv:novel_5:effb0-dacisPMP:lo\endcsname{+4.1}
\expandafter\gdef\csname num@d:pv:novel_5:effb0-dacisPMP:hi\endcsname{+10.3}
\expandafter\gdef\csname num@d:pv:novel_5:effb0-dacisPMP:abs\endcsname{7.2}
\expandafter\gdef\csname num@d:pv:novel_5:effb0-dacisPMP:wins\endcsname{5}
\expandafter\gdef\csname num@d:pv:novel_5:effb0-dacisPMP:n\endcsname{5}
\expandafter\gdef\csname num@d:pv:novel_5:effb0-dacisPMP:worst\endcsname{10.3}
\expandafter\gdef\csname num@d:pv:novel_5:mnv2-dacisPMP\endcsname{+4.6}
\expandafter\gdef\csname num@d:pv:novel_5:mnv2-dacisPMP:lo\endcsname{+2.7}
\expandafter\gdef\csname num@d:pv:novel_5:mnv2-dacisPMP:hi\endcsname{+6.5}
\expandafter\gdef\csname num@d:pv:novel_5:mnv2-dacisPMP:abs\endcsname{4.6}
\expandafter\gdef\csname num@d:pv:novel_5:mnv2-dacisPMP:wins\endcsname{5}
\expandafter\gdef\csname num@d:pv:novel_5:mnv2-dacisPMP:n\endcsname{5}
\expandafter\gdef\csname num@d:pv:novel_5:mnv2-dacisPMP:worst\endcsname{6.5}
\expandafter\gdef\csname num@d:pv:novel_5:dacis-l1:PMP50\endcsname{\ensuremath{-}0.4}
\expandafter\gdef\csname num@d:pv:novel_5:dacis-l1:PMP50:lo\endcsname{\ensuremath{-}2.2}
\expandafter\gdef\csname num@d:pv:novel_5:dacis-l1:PMP50:hi\endcsname{+1.5}
\expandafter\gdef\csname num@d:pv:novel_5:dacis-l1:PMP50:abs\endcsname{0.4}
\expandafter\gdef\csname num@d:pv:novel_5:dacis-l1:PMP50:wins\endcsname{2}
\expandafter\gdef\csname num@d:pv:novel_5:dacis-l1:PMP50:n\endcsname{5}
\expandafter\gdef\csname num@d:pv:novel_5:dacis-l1:PMP50:worst\endcsname{2.2}
\expandafter\gdef\csname num@d:pv:novel_5:dense-dacisPMP50\endcsname{+2.8}
\expandafter\gdef\csname num@d:pv:novel_5:dense-dacisPMP50:lo\endcsname{+0.8}
\expandafter\gdef\csname num@d:pv:novel_5:dense-dacisPMP50:hi\endcsname{+4.8}
\expandafter\gdef\csname num@d:pv:novel_5:dense-dacisPMP50:abs\endcsname{2.8}
\expandafter\gdef\csname num@d:pv:novel_5:dense-dacisPMP50:wins\endcsname{5}
\expandafter\gdef\csname num@d:pv:novel_5:dense-dacisPMP50:n\endcsname{5}
\expandafter\gdef\csname num@d:pv:novel_5:dense-dacisPMP50:worst\endcsname{4.8}
\expandafter\gdef\csname num@d:pv:novel_5:dacis-l1:PMP20\endcsname{\ensuremath{-}0.4}
\expandafter\gdef\csname num@d:pv:novel_5:dacis-l1:PMP20:lo\endcsname{\ensuremath{-}2.1}
\expandafter\gdef\csname num@d:pv:novel_5:dacis-l1:PMP20:hi\endcsname{+1.3}
\expandafter\gdef\csname num@d:pv:novel_5:dacis-l1:PMP20:abs\endcsname{0.4}
\expandafter\gdef\csname num@d:pv:novel_5:dacis-l1:PMP20:wins\endcsname{2}
\expandafter\gdef\csname num@d:pv:novel_5:dacis-l1:PMP20:n\endcsname{5}
\expandafter\gdef\csname num@d:pv:novel_5:dacis-l1:PMP20:worst\endcsname{2.1}
\expandafter\gdef\csname num@d:pv:novel_5:dense-dacisPMP20\endcsname{+7.6}
\expandafter\gdef\csname num@d:pv:novel_5:dense-dacisPMP20:lo\endcsname{+4.3}
\expandafter\gdef\csname num@d:pv:novel_5:dense-dacisPMP20:hi\endcsname{+10.9}
\expandafter\gdef\csname num@d:pv:novel_5:dense-dacisPMP20:abs\endcsname{7.6}
\expandafter\gdef\csname num@d:pv:novel_5:dense-dacisPMP20:wins\endcsname{5}
\expandafter\gdef\csname num@d:pv:novel_5:dense-dacisPMP20:n\endcsname{5}
\expandafter\gdef\csname num@d:pv:novel_5:dense-dacisPMP20:worst\endcsname{10.9}
\expandafter\gdef\csname num@d:pv:novel_5:dacis-l1:PMP10\endcsname{\ensuremath{-}1.1}
\expandafter\gdef\csname num@d:pv:novel_5:dacis-l1:PMP10:lo\endcsname{\ensuremath{-}4.7}
\expandafter\gdef\csname num@d:pv:novel_5:dacis-l1:PMP10:hi\endcsname{+2.4}
\expandafter\gdef\csname num@d:pv:novel_5:dacis-l1:PMP10:abs\endcsname{1.1}
\expandafter\gdef\csname num@d:pv:novel_5:dacis-l1:PMP10:wins\endcsname{3}
\expandafter\gdef\csname num@d:pv:novel_5:dacis-l1:PMP10:n\endcsname{5}
\expandafter\gdef\csname num@d:pv:novel_5:dacis-l1:PMP10:worst\endcsname{4.7}
\expandafter\gdef\csname num@d:pv:novel_5:dense-dacisPMP10\endcsname{+13.2}
\expandafter\gdef\csname num@d:pv:novel_5:dense-dacisPMP10:lo\endcsname{+7.3}
\expandafter\gdef\csname num@d:pv:novel_5:dense-dacisPMP10:hi\endcsname{+19.1}
\expandafter\gdef\csname num@d:pv:novel_5:dense-dacisPMP10:abs\endcsname{13.2}
\expandafter\gdef\csname num@d:pv:novel_5:dense-dacisPMP10:wins\endcsname{5}
\expandafter\gdef\csname num@d:pv:novel_5:dense-dacisPMP10:n\endcsname{5}
\expandafter\gdef\csname num@d:pv:novel_5:dense-dacisPMP10:worst\endcsname{19.1}
\expandafter\gdef\csname num@d:pd:plantdoc_5:dacis-random:PMP\endcsname{+0.8}
\expandafter\gdef\csname num@d:pd:plantdoc_5:dacis-random:PMP:lo\endcsname{\ensuremath{-}0.4}
\expandafter\gdef\csname num@d:pd:plantdoc_5:dacis-random:PMP:hi\endcsname{+2.1}
\expandafter\gdef\csname num@d:pd:plantdoc_5:dacis-random:PMP:abs\endcsname{0.8}
\expandafter\gdef\csname num@d:pd:plantdoc_5:dacis-random:PMP:wins\endcsname{4}
\expandafter\gdef\csname num@d:pd:plantdoc_5:dacis-random:PMP:n\endcsname{5}
\expandafter\gdef\csname num@d:pd:plantdoc_5:dacis-random:PMP:worst\endcsname{2.1}
\expandafter\gdef\csname num@d:pd:plantdoc_5:dacis-random:P\endcsname{+0.7}
\expandafter\gdef\csname num@d:pd:plantdoc_5:dacis-random:P:lo\endcsname{\ensuremath{-}1.2}
\expandafter\gdef\csname num@d:pd:plantdoc_5:dacis-random:P:hi\endcsname{+2.6}
\expandafter\gdef\csname num@d:pd:plantdoc_5:dacis-random:P:abs\endcsname{0.7}
\expandafter\gdef\csname num@d:pd:plantdoc_5:dacis-random:P:wins\endcsname{3}
\expandafter\gdef\csname num@d:pd:plantdoc_5:dacis-random:P:n\endcsname{5}
\expandafter\gdef\csname num@d:pd:plantdoc_5:dacis-random:P:worst\endcsname{2.6}
\expandafter\gdef\csname num@d:pd:plantdoc_5:dacis-l1:PMP\endcsname{\ensuremath{-}0.7}
\expandafter\gdef\csname num@d:pd:plantdoc_5:dacis-l1:PMP:lo\endcsname{\ensuremath{-}2.4}
\expandafter\gdef\csname num@d:pd:plantdoc_5:dacis-l1:PMP:hi\endcsname{+1.1}
\expandafter\gdef\csname num@d:pd:plantdoc_5:dacis-l1:PMP:abs\endcsname{0.7}
\expandafter\gdef\csname num@d:pd:plantdoc_5:dacis-l1:PMP:wins\endcsname{2}
\expandafter\gdef\csname num@d:pd:plantdoc_5:dacis-l1:PMP:n\endcsname{5}
\expandafter\gdef\csname num@d:pd:plantdoc_5:dacis-l1:PMP:worst\endcsname{2.4}
\expandafter\gdef\csname num@d:pd:plantdoc_5:dacis-l1:P\endcsname{\ensuremath{-}1.4}
\expandafter\gdef\csname num@d:pd:plantdoc_5:dacis-l1:P:lo\endcsname{\ensuremath{-}3.3}
\expandafter\gdef\csname num@d:pd:plantdoc_5:dacis-l1:P:hi\endcsname{+0.5}
\expandafter\gdef\csname num@d:pd:plantdoc_5:dacis-l1:P:abs\endcsname{1.4}
\expandafter\gdef\csname num@d:pd:plantdoc_5:dacis-l1:P:wins\endcsname{1}
\expandafter\gdef\csname num@d:pd:plantdoc_5:dacis-l1:P:n\endcsname{5}
\expandafter\gdef\csname num@d:pd:plantdoc_5:dacis-l1:P:worst\endcsname{3.3}
\expandafter\gdef\csname num@d:pd:plantdoc_5:dacis-taylor:PMP\endcsname{\ensuremath{-}1.1}
\expandafter\gdef\csname num@d:pd:plantdoc_5:dacis-taylor:PMP:lo\endcsname{\ensuremath{-}2.4}
\expandafter\gdef\csname num@d:pd:plantdoc_5:dacis-taylor:PMP:hi\endcsname{+0.3}
\expandafter\gdef\csname num@d:pd:plantdoc_5:dacis-taylor:PMP:abs\endcsname{1.1}
\expandafter\gdef\csname num@d:pd:plantdoc_5:dacis-taylor:PMP:wins\endcsname{1}
\expandafter\gdef\csname num@d:pd:plantdoc_5:dacis-taylor:PMP:n\endcsname{5}
\expandafter\gdef\csname num@d:pd:plantdoc_5:dacis-taylor:PMP:worst\endcsname{2.4}
\expandafter\gdef\csname num@d:pd:plantdoc_5:dacis-taylor:P\endcsname{\ensuremath{-}0.6}
\expandafter\gdef\csname num@d:pd:plantdoc_5:dacis-taylor:P:lo\endcsname{\ensuremath{-}1.6}
\expandafter\gdef\csname num@d:pd:plantdoc_5:dacis-taylor:P:hi\endcsname{+0.4}
\expandafter\gdef\csname num@d:pd:plantdoc_5:dacis-taylor:P:abs\endcsname{0.6}
\expandafter\gdef\csname num@d:pd:plantdoc_5:dacis-taylor:P:wins\endcsname{1}
\expandafter\gdef\csname num@d:pd:plantdoc_5:dacis-taylor:P:n\endcsname{5}
\expandafter\gdef\csname num@d:pd:plantdoc_5:dacis-taylor:P:worst\endcsname{1.6}
\expandafter\gdef\csname num@d:pd:plantdoc_5:dacis-bn_gamma:PMP\endcsname{+1.9}
\expandafter\gdef\csname num@d:pd:plantdoc_5:dacis-bn_gamma:PMP:lo\endcsname{+0.7}
\expandafter\gdef\csname num@d:pd:plantdoc_5:dacis-bn_gamma:PMP:hi\endcsname{+3.1}
\expandafter\gdef\csname num@d:pd:plantdoc_5:dacis-bn_gamma:PMP:abs\endcsname{1.9}
\expandafter\gdef\csname num@d:pd:plantdoc_5:dacis-bn_gamma:PMP:wins\endcsname{5}
\expandafter\gdef\csname num@d:pd:plantdoc_5:dacis-bn_gamma:PMP:n\endcsname{5}
\expandafter\gdef\csname num@d:pd:plantdoc_5:dacis-bn_gamma:PMP:worst\endcsname{3.1}
\expandafter\gdef\csname num@d:pd:plantdoc_5:dacis-bn_gamma:P\endcsname{+2.9}
\expandafter\gdef\csname num@d:pd:plantdoc_5:dacis-bn_gamma:P:lo\endcsname{+1.2}
\expandafter\gdef\csname num@d:pd:plantdoc_5:dacis-bn_gamma:P:hi\endcsname{+4.6}
\expandafter\gdef\csname num@d:pd:plantdoc_5:dacis-bn_gamma:P:abs\endcsname{2.9}
\expandafter\gdef\csname num@d:pd:plantdoc_5:dacis-bn_gamma:P:wins\endcsname{5}
\expandafter\gdef\csname num@d:pd:plantdoc_5:dacis-bn_gamma:P:n\endcsname{5}
\expandafter\gdef\csname num@d:pd:plantdoc_5:dacis-bn_gamma:P:worst\endcsname{4.6}
\expandafter\gdef\csname num@d:pd:plantdoc_5:dacis-D:PMP\endcsname{+1.1}
\expandafter\gdef\csname num@d:pd:plantdoc_5:dacis-D:PMP:lo\endcsname{\ensuremath{-}0.6}
\expandafter\gdef\csname num@d:pd:plantdoc_5:dacis-D:PMP:hi\endcsname{+2.9}
\expandafter\gdef\csname num@d:pd:plantdoc_5:dacis-D:PMP:abs\endcsname{1.1}
\expandafter\gdef\csname num@d:pd:plantdoc_5:dacis-D:PMP:wins\endcsname{3}
\expandafter\gdef\csname num@d:pd:plantdoc_5:dacis-D:PMP:n\endcsname{5}
\expandafter\gdef\csname num@d:pd:plantdoc_5:dacis-D:PMP:worst\endcsname{2.9}
\expandafter\gdef\csname num@d:pd:plantdoc_5:dacis-D:P\endcsname{+0.5}
\expandafter\gdef\csname num@d:pd:plantdoc_5:dacis-D:P:lo\endcsname{\ensuremath{-}0.4}
\expandafter\gdef\csname num@d:pd:plantdoc_5:dacis-D:P:hi\endcsname{+1.5}
\expandafter\gdef\csname num@d:pd:plantdoc_5:dacis-D:P:abs\endcsname{0.5}
\expandafter\gdef\csname num@d:pd:plantdoc_5:dacis-D:P:wins\endcsname{4}
\expandafter\gdef\csname num@d:pd:plantdoc_5:dacis-D:P:n\endcsname{5}
\expandafter\gdef\csname num@d:pd:plantdoc_5:dacis-D:P:worst\endcsname{1.5}
\expandafter\gdef\csname num@d:pd:plantdoc_5:dacis-G:PMP\endcsname{+0.6}
\expandafter\gdef\csname num@d:pd:plantdoc_5:dacis-G:PMP:lo\endcsname{\ensuremath{-}0.8}
\expandafter\gdef\csname num@d:pd:plantdoc_5:dacis-G:PMP:hi\endcsname{+2.0}
\expandafter\gdef\csname num@d:pd:plantdoc_5:dacis-G:PMP:abs\endcsname{0.6}
\expandafter\gdef\csname num@d:pd:plantdoc_5:dacis-G:PMP:wins\endcsname{4}
\expandafter\gdef\csname num@d:pd:plantdoc_5:dacis-G:PMP:n\endcsname{5}
\expandafter\gdef\csname num@d:pd:plantdoc_5:dacis-G:PMP:worst\endcsname{2.0}
\expandafter\gdef\csname num@d:pd:plantdoc_5:dacis-G:P\endcsname{\ensuremath{-}0.3}
\expandafter\gdef\csname num@d:pd:plantdoc_5:dacis-G:P:lo\endcsname{\ensuremath{-}0.7}
\expandafter\gdef\csname num@d:pd:plantdoc_5:dacis-G:P:hi\endcsname{+0.2}
\expandafter\gdef\csname num@d:pd:plantdoc_5:dacis-G:P:abs\endcsname{0.3}
\expandafter\gdef\csname num@d:pd:plantdoc_5:dacis-G:P:wins\endcsname{2}
\expandafter\gdef\csname num@d:pd:plantdoc_5:dacis-G:P:n\endcsname{5}
\expandafter\gdef\csname num@d:pd:plantdoc_5:dacis-G:P:worst\endcsname{0.7}
\expandafter\gdef\csname num@d:pd:plantdoc_5:dacis-V:PMP\endcsname{+0.2}
\expandafter\gdef\csname num@d:pd:plantdoc_5:dacis-V:PMP:lo\endcsname{\ensuremath{-}0.4}
\expandafter\gdef\csname num@d:pd:plantdoc_5:dacis-V:PMP:hi\endcsname{+0.7}
\expandafter\gdef\csname num@d:pd:plantdoc_5:dacis-V:PMP:abs\endcsname{0.2}
\expandafter\gdef\csname num@d:pd:plantdoc_5:dacis-V:PMP:wins\endcsname{2}
\expandafter\gdef\csname num@d:pd:plantdoc_5:dacis-V:PMP:n\endcsname{5}
\expandafter\gdef\csname num@d:pd:plantdoc_5:dacis-V:PMP:worst\endcsname{0.7}
\expandafter\gdef\csname num@d:pd:plantdoc_5:dacis-V:P\endcsname{+0.2}
\expandafter\gdef\csname num@d:pd:plantdoc_5:dacis-V:P:lo\endcsname{\ensuremath{-}1.4}
\expandafter\gdef\csname num@d:pd:plantdoc_5:dacis-V:P:hi\endcsname{+1.7}
\expandafter\gdef\csname num@d:pd:plantdoc_5:dacis-V:P:abs\endcsname{0.2}
\expandafter\gdef\csname num@d:pd:plantdoc_5:dacis-V:P:wins\endcsname{3}
\expandafter\gdef\csname num@d:pd:plantdoc_5:dacis-V:P:n\endcsname{5}
\expandafter\gdef\csname num@d:pd:plantdoc_5:dacis-V:P:worst\endcsname{1.7}
\expandafter\gdef\csname num@d:pd:plantdoc_5:l1-random:PMP\endcsname{+1.5}
\expandafter\gdef\csname num@d:pd:plantdoc_5:l1-random:PMP:lo\endcsname{\ensuremath{-}0.3}
\expandafter\gdef\csname num@d:pd:plantdoc_5:l1-random:PMP:hi\endcsname{+3.2}
\expandafter\gdef\csname num@d:pd:plantdoc_5:l1-random:PMP:abs\endcsname{1.5}
\expandafter\gdef\csname num@d:pd:plantdoc_5:l1-random:PMP:wins\endcsname{5}
\expandafter\gdef\csname num@d:pd:plantdoc_5:l1-random:PMP:n\endcsname{5}
\expandafter\gdef\csname num@d:pd:plantdoc_5:l1-random:PMP:worst\endcsname{3.2}
\expandafter\gdef\csname num@d:pd:plantdoc_5:l1-random:P\endcsname{+2.1}
\expandafter\gdef\csname num@d:pd:plantdoc_5:l1-random:P:lo\endcsname{+1.1}
\expandafter\gdef\csname num@d:pd:plantdoc_5:l1-random:P:hi\endcsname{+3.1}
\expandafter\gdef\csname num@d:pd:plantdoc_5:l1-random:P:abs\endcsname{2.1}
\expandafter\gdef\csname num@d:pd:plantdoc_5:l1-random:P:wins\endcsname{5}
\expandafter\gdef\csname num@d:pd:plantdoc_5:l1-random:P:n\endcsname{5}
\expandafter\gdef\csname num@d:pd:plantdoc_5:l1-random:P:worst\endcsname{3.1}
\expandafter\gdef\csname num@d:pd:plantdoc_5:taylor-random:PMP\endcsname{+1.9}
\expandafter\gdef\csname num@d:pd:plantdoc_5:taylor-random:PMP:lo\endcsname{+0.8}
\expandafter\gdef\csname num@d:pd:plantdoc_5:taylor-random:PMP:hi\endcsname{+3.0}
\expandafter\gdef\csname num@d:pd:plantdoc_5:taylor-random:PMP:abs\endcsname{1.9}
\expandafter\gdef\csname num@d:pd:plantdoc_5:taylor-random:PMP:wins\endcsname{5}
\expandafter\gdef\csname num@d:pd:plantdoc_5:taylor-random:PMP:n\endcsname{5}
\expandafter\gdef\csname num@d:pd:plantdoc_5:taylor-random:PMP:worst\endcsname{3.0}
\expandafter\gdef\csname num@d:pd:plantdoc_5:taylor-random:P\endcsname{+1.3}
\expandafter\gdef\csname num@d:pd:plantdoc_5:taylor-random:P:lo\endcsname{\ensuremath{-}0.4}
\expandafter\gdef\csname num@d:pd:plantdoc_5:taylor-random:P:hi\endcsname{+2.9}
\expandafter\gdef\csname num@d:pd:plantdoc_5:taylor-random:P:abs\endcsname{1.3}
\expandafter\gdef\csname num@d:pd:plantdoc_5:taylor-random:P:wins\endcsname{3}
\expandafter\gdef\csname num@d:pd:plantdoc_5:taylor-random:P:n\endcsname{5}
\expandafter\gdef\csname num@d:pd:plantdoc_5:taylor-random:P:worst\endcsname{2.9}
\expandafter\gdef\csname num@d:pd:plantdoc_5:bn_gamma-random:PMP\endcsname{\ensuremath{-}1.1}
\expandafter\gdef\csname num@d:pd:plantdoc_5:bn_gamma-random:PMP:lo\endcsname{\ensuremath{-}2.7}
\expandafter\gdef\csname num@d:pd:plantdoc_5:bn_gamma-random:PMP:hi\endcsname{+0.5}
\expandafter\gdef\csname num@d:pd:plantdoc_5:bn_gamma-random:PMP:abs\endcsname{1.1}
\expandafter\gdef\csname num@d:pd:plantdoc_5:bn_gamma-random:PMP:wins\endcsname{1}
\expandafter\gdef\csname num@d:pd:plantdoc_5:bn_gamma-random:PMP:n\endcsname{5}
\expandafter\gdef\csname num@d:pd:plantdoc_5:bn_gamma-random:PMP:worst\endcsname{2.7}
\expandafter\gdef\csname num@d:pd:plantdoc_5:bn_gamma-random:P\endcsname{\ensuremath{-}2.2}
\expandafter\gdef\csname num@d:pd:plantdoc_5:bn_gamma-random:P:lo\endcsname{\ensuremath{-}3.1}
\expandafter\gdef\csname num@d:pd:plantdoc_5:bn_gamma-random:P:hi\endcsname{\ensuremath{-}1.4}
\expandafter\gdef\csname num@d:pd:plantdoc_5:bn_gamma-random:P:abs\endcsname{2.2}
\expandafter\gdef\csname num@d:pd:plantdoc_5:bn_gamma-random:P:wins\endcsname{0}
\expandafter\gdef\csname num@d:pd:plantdoc_5:bn_gamma-random:P:n\endcsname{5}
\expandafter\gdef\csname num@d:pd:plantdoc_5:bn_gamma-random:P:worst\endcsname{3.1}
\expandafter\gdef\csname num@d:pd:plantdoc_5:D-random:PMP\endcsname{\ensuremath{-}0.3}
\expandafter\gdef\csname num@d:pd:plantdoc_5:D-random:PMP:lo\endcsname{\ensuremath{-}2.3}
\expandafter\gdef\csname num@d:pd:plantdoc_5:D-random:PMP:hi\endcsname{+1.7}
\expandafter\gdef\csname num@d:pd:plantdoc_5:D-random:PMP:abs\endcsname{0.3}
\expandafter\gdef\csname num@d:pd:plantdoc_5:D-random:PMP:wins\endcsname{2}
\expandafter\gdef\csname num@d:pd:plantdoc_5:D-random:PMP:n\endcsname{5}
\expandafter\gdef\csname num@d:pd:plantdoc_5:D-random:PMP:worst\endcsname{2.3}
\expandafter\gdef\csname num@d:pd:plantdoc_5:D-random:P\endcsname{+0.2}
\expandafter\gdef\csname num@d:pd:plantdoc_5:D-random:P:lo\endcsname{\ensuremath{-}1.3}
\expandafter\gdef\csname num@d:pd:plantdoc_5:D-random:P:hi\endcsname{+1.7}
\expandafter\gdef\csname num@d:pd:plantdoc_5:D-random:P:abs\endcsname{0.2}
\expandafter\gdef\csname num@d:pd:plantdoc_5:D-random:P:wins\endcsname{3}
\expandafter\gdef\csname num@d:pd:plantdoc_5:D-random:P:n\endcsname{5}
\expandafter\gdef\csname num@d:pd:plantdoc_5:D-random:P:worst\endcsname{1.7}
\expandafter\gdef\csname num@d:pd:plantdoc_5:G-random:PMP\endcsname{+0.2}
\expandafter\gdef\csname num@d:pd:plantdoc_5:G-random:PMP:lo\endcsname{\ensuremath{-}1.8}
\expandafter\gdef\csname num@d:pd:plantdoc_5:G-random:PMP:hi\endcsname{+2.2}
\expandafter\gdef\csname num@d:pd:plantdoc_5:G-random:PMP:abs\endcsname{0.2}
\expandafter\gdef\csname num@d:pd:plantdoc_5:G-random:PMP:wins\endcsname{2}
\expandafter\gdef\csname num@d:pd:plantdoc_5:G-random:PMP:n\endcsname{5}
\expandafter\gdef\csname num@d:pd:plantdoc_5:G-random:PMP:worst\endcsname{2.2}
\expandafter\gdef\csname num@d:pd:plantdoc_5:G-random:P\endcsname{+1.0}
\expandafter\gdef\csname num@d:pd:plantdoc_5:G-random:P:lo\endcsname{\ensuremath{-}0.9}
\expandafter\gdef\csname num@d:pd:plantdoc_5:G-random:P:hi\endcsname{+2.8}
\expandafter\gdef\csname num@d:pd:plantdoc_5:G-random:P:abs\endcsname{1.0}
\expandafter\gdef\csname num@d:pd:plantdoc_5:G-random:P:wins\endcsname{4}
\expandafter\gdef\csname num@d:pd:plantdoc_5:G-random:P:n\endcsname{5}
\expandafter\gdef\csname num@d:pd:plantdoc_5:G-random:P:worst\endcsname{2.8}
\expandafter\gdef\csname num@d:pd:plantdoc_5:V-random:PMP\endcsname{+0.6}
\expandafter\gdef\csname num@d:pd:plantdoc_5:V-random:PMP:lo\endcsname{\ensuremath{-}0.3}
\expandafter\gdef\csname num@d:pd:plantdoc_5:V-random:PMP:hi\endcsname{+1.6}
\expandafter\gdef\csname num@d:pd:plantdoc_5:V-random:PMP:abs\endcsname{0.6}
\expandafter\gdef\csname num@d:pd:plantdoc_5:V-random:PMP:wins\endcsname{4}
\expandafter\gdef\csname num@d:pd:plantdoc_5:V-random:PMP:n\endcsname{5}
\expandafter\gdef\csname num@d:pd:plantdoc_5:V-random:PMP:worst\endcsname{1.6}
\expandafter\gdef\csname num@d:pd:plantdoc_5:V-random:P\endcsname{+0.5}
\expandafter\gdef\csname num@d:pd:plantdoc_5:V-random:P:lo\endcsname{\ensuremath{-}0.0}
\expandafter\gdef\csname num@d:pd:plantdoc_5:V-random:P:hi\endcsname{+1.1}
\expandafter\gdef\csname num@d:pd:plantdoc_5:V-random:P:abs\endcsname{0.5}
\expandafter\gdef\csname num@d:pd:plantdoc_5:V-random:P:wins\endcsname{5}
\expandafter\gdef\csname num@d:pd:plantdoc_5:V-random:P:n\endcsname{5}
\expandafter\gdef\csname num@d:pd:plantdoc_5:V-random:P:worst\endcsname{1.1}
\expandafter\gdef\csname num@d:pd:plantdoc_5:PM-P\endcsname{\ensuremath{-}0.8}
\expandafter\gdef\csname num@d:pd:plantdoc_5:PM-P:lo\endcsname{\ensuremath{-}1.6}
\expandafter\gdef\csname num@d:pd:plantdoc_5:PM-P:hi\endcsname{\ensuremath{-}0.0}
\expandafter\gdef\csname num@d:pd:plantdoc_5:PM-P:abs\endcsname{0.8}
\expandafter\gdef\csname num@d:pd:plantdoc_5:PM-P:wins\endcsname{1}
\expandafter\gdef\csname num@d:pd:plantdoc_5:PM-P:n\endcsname{5}
\expandafter\gdef\csname num@d:pd:plantdoc_5:PM-P:worst\endcsname{1.6}
\expandafter\gdef\csname num@d:pd:plantdoc_5:MP-P\endcsname{\ensuremath{-}0.7}
\expandafter\gdef\csname num@d:pd:plantdoc_5:MP-P:lo\endcsname{\ensuremath{-}1.5}
\expandafter\gdef\csname num@d:pd:plantdoc_5:MP-P:hi\endcsname{+0.2}
\expandafter\gdef\csname num@d:pd:plantdoc_5:MP-P:abs\endcsname{0.7}
\expandafter\gdef\csname num@d:pd:plantdoc_5:MP-P:wins\endcsname{1}
\expandafter\gdef\csname num@d:pd:plantdoc_5:MP-P:n\endcsname{5}
\expandafter\gdef\csname num@d:pd:plantdoc_5:MP-P:worst\endcsname{1.5}
\expandafter\gdef\csname num@d:pd:plantdoc_5:PGP-P\endcsname{\ensuremath{-}0.2}
\expandafter\gdef\csname num@d:pd:plantdoc_5:PGP-P:lo\endcsname{\ensuremath{-}0.8}
\expandafter\gdef\csname num@d:pd:plantdoc_5:PGP-P:hi\endcsname{+0.5}
\expandafter\gdef\csname num@d:pd:plantdoc_5:PGP-P:abs\endcsname{0.2}
\expandafter\gdef\csname num@d:pd:plantdoc_5:PGP-P:wins\endcsname{1}
\expandafter\gdef\csname num@d:pd:plantdoc_5:PGP-P:n\endcsname{5}
\expandafter\gdef\csname num@d:pd:plantdoc_5:PGP-P:worst\endcsname{0.8}
\expandafter\gdef\csname num@d:pd:plantdoc_5:PMPM-P\endcsname{+0.2}
\expandafter\gdef\csname num@d:pd:plantdoc_5:PMPM-P:lo\endcsname{\ensuremath{-}0.5}
\expandafter\gdef\csname num@d:pd:plantdoc_5:PMPM-P:hi\endcsname{+1.0}
\expandafter\gdef\csname num@d:pd:plantdoc_5:PMPM-P:abs\endcsname{0.2}
\expandafter\gdef\csname num@d:pd:plantdoc_5:PMPM-P:wins\endcsname{3}
\expandafter\gdef\csname num@d:pd:plantdoc_5:PMPM-P:n\endcsname{5}
\expandafter\gdef\csname num@d:pd:plantdoc_5:PMPM-P:worst\endcsname{1.0}
\expandafter\gdef\csname num@d:pd:plantdoc_5:PM2x-P\endcsname{\ensuremath{-}0.9}
\expandafter\gdef\csname num@d:pd:plantdoc_5:PM2x-P:lo\endcsname{\ensuremath{-}1.4}
\expandafter\gdef\csname num@d:pd:plantdoc_5:PM2x-P:hi\endcsname{\ensuremath{-}0.4}
\expandafter\gdef\csname num@d:pd:plantdoc_5:PM2x-P:abs\endcsname{0.9}
\expandafter\gdef\csname num@d:pd:plantdoc_5:PM2x-P:wins\endcsname{0}
\expandafter\gdef\csname num@d:pd:plantdoc_5:PM2x-P:n\endcsname{5}
\expandafter\gdef\csname num@d:pd:plantdoc_5:PM2x-P:worst\endcsname{1.4}
\expandafter\gdef\csname num@d:pd:plantdoc_5:PMP-P\endcsname{+0.1}
\expandafter\gdef\csname num@d:pd:plantdoc_5:PMP-P:lo\endcsname{\ensuremath{-}0.8}
\expandafter\gdef\csname num@d:pd:plantdoc_5:PMP-P:hi\endcsname{+1.1}
\expandafter\gdef\csname num@d:pd:plantdoc_5:PMP-P:abs\endcsname{0.1}
\expandafter\gdef\csname num@d:pd:plantdoc_5:PMP-P:wins\endcsname{2}
\expandafter\gdef\csname num@d:pd:plantdoc_5:PMP-P:n\endcsname{5}
\expandafter\gdef\csname num@d:pd:plantdoc_5:PMP-P:worst\endcsname{1.1}
\expandafter\gdef\csname num@d:pd:plantdoc_5:densem-dense\endcsname{\ensuremath{-}0.0}
\expandafter\gdef\csname num@d:pd:plantdoc_5:densem-dense:lo\endcsname{\ensuremath{-}0.4}
\expandafter\gdef\csname num@d:pd:plantdoc_5:densem-dense:hi\endcsname{+0.3}
\expandafter\gdef\csname num@d:pd:plantdoc_5:densem-dense:abs\endcsname{0.0}
\expandafter\gdef\csname num@d:pd:plantdoc_5:densem-dense:wins\endcsname{2}
\expandafter\gdef\csname num@d:pd:plantdoc_5:densem-dense:n\endcsname{5}
\expandafter\gdef\csname num@d:pd:plantdoc_5:densem-dense:worst\endcsname{0.4}
\expandafter\gdef\csname num@d:pd:plantdoc_5:densermr-dense\endcsname{+0.1}
\expandafter\gdef\csname num@d:pd:plantdoc_5:densermr-dense:lo\endcsname{\ensuremath{-}0.3}
\expandafter\gdef\csname num@d:pd:plantdoc_5:densermr-dense:hi\endcsname{+0.5}
\expandafter\gdef\csname num@d:pd:plantdoc_5:densermr-dense:abs\endcsname{0.1}
\expandafter\gdef\csname num@d:pd:plantdoc_5:densermr-dense:wins\endcsname{4}
\expandafter\gdef\csname num@d:pd:plantdoc_5:densermr-dense:n\endcsname{5}
\expandafter\gdef\csname num@d:pd:plantdoc_5:densermr-dense:worst\endcsname{0.5}
\expandafter\gdef\csname num@d:pd:plantdoc_5:dense-dacisP\endcsname{+8.1}
\expandafter\gdef\csname num@d:pd:plantdoc_5:dense-dacisP:lo\endcsname{+7.2}
\expandafter\gdef\csname num@d:pd:plantdoc_5:dense-dacisP:hi\endcsname{+9.1}
\expandafter\gdef\csname num@d:pd:plantdoc_5:dense-dacisP:abs\endcsname{8.1}
\expandafter\gdef\csname num@d:pd:plantdoc_5:dense-dacisP:wins\endcsname{5}
\expandafter\gdef\csname num@d:pd:plantdoc_5:dense-dacisP:n\endcsname{5}
\expandafter\gdef\csname num@d:pd:plantdoc_5:dense-dacisP:worst\endcsname{9.1}
\expandafter\gdef\csname num@d:pd:plantdoc_5:dense-dacisPMP\endcsname{+8.0}
\expandafter\gdef\csname num@d:pd:plantdoc_5:dense-dacisPMP:lo\endcsname{+6.6}
\expandafter\gdef\csname num@d:pd:plantdoc_5:dense-dacisPMP:hi\endcsname{+9.3}
\expandafter\gdef\csname num@d:pd:plantdoc_5:dense-dacisPMP:abs\endcsname{8.0}
\expandafter\gdef\csname num@d:pd:plantdoc_5:dense-dacisPMP:wins\endcsname{5}
\expandafter\gdef\csname num@d:pd:plantdoc_5:dense-dacisPMP:n\endcsname{5}
\expandafter\gdef\csname num@d:pd:plantdoc_5:dense-dacisPMP:worst\endcsname{9.3}
\expandafter\gdef\csname num@d:pd:plantdoc_5:dense-randomP\endcsname{+8.8}
\expandafter\gdef\csname num@d:pd:plantdoc_5:dense-randomP:lo\endcsname{+6.8}
\expandafter\gdef\csname num@d:pd:plantdoc_5:dense-randomP:hi\endcsname{+10.8}
\expandafter\gdef\csname num@d:pd:plantdoc_5:dense-randomP:abs\endcsname{8.8}
\expandafter\gdef\csname num@d:pd:plantdoc_5:dense-randomP:wins\endcsname{5}
\expandafter\gdef\csname num@d:pd:plantdoc_5:dense-randomP:n\endcsname{5}
\expandafter\gdef\csname num@d:pd:plantdoc_5:dense-randomP:worst\endcsname{10.8}
\expandafter\gdef\csname num@d:pd:plantdoc_5:dense-randomPMP\endcsname{+8.8}
\expandafter\gdef\csname num@d:pd:plantdoc_5:dense-randomPMP:lo\endcsname{+6.5}
\expandafter\gdef\csname num@d:pd:plantdoc_5:dense-randomPMP:hi\endcsname{+11.0}
\expandafter\gdef\csname num@d:pd:plantdoc_5:dense-randomPMP:abs\endcsname{8.8}
\expandafter\gdef\csname num@d:pd:plantdoc_5:dense-randomPMP:wins\endcsname{5}
\expandafter\gdef\csname num@d:pd:plantdoc_5:dense-randomPMP:n\endcsname{5}
\expandafter\gdef\csname num@d:pd:plantdoc_5:dense-randomPMP:worst\endcsname{11.0}
\expandafter\gdef\csname num@d:pd:plantdoc_5:norefine-full\endcsname{\ensuremath{-}0.5}
\expandafter\gdef\csname num@d:pd:plantdoc_5:norefine-full:lo\endcsname{\ensuremath{-}1.6}
\expandafter\gdef\csname num@d:pd:plantdoc_5:norefine-full:hi\endcsname{+0.6}
\expandafter\gdef\csname num@d:pd:plantdoc_5:norefine-full:abs\endcsname{0.5}
\expandafter\gdef\csname num@d:pd:plantdoc_5:norefine-full:wins\endcsname{1}
\expandafter\gdef\csname num@d:pd:plantdoc_5:norefine-full:n\endcsname{5}
\expandafter\gdef\csname num@d:pd:plantdoc_5:norefine-full:worst\endcsname{1.6}
\expandafter\gdef\csname num@d:pd:plantdoc_5:mnv3-dacisPMP\endcsname{+2.5}
\expandafter\gdef\csname num@d:pd:plantdoc_5:mnv3-dacisPMP:lo\endcsname{+0.2}
\expandafter\gdef\csname num@d:pd:plantdoc_5:mnv3-dacisPMP:hi\endcsname{+4.9}
\expandafter\gdef\csname num@d:pd:plantdoc_5:mnv3-dacisPMP:abs\endcsname{2.5}
\expandafter\gdef\csname num@d:pd:plantdoc_5:mnv3-dacisPMP:wins\endcsname{5}
\expandafter\gdef\csname num@d:pd:plantdoc_5:mnv3-dacisPMP:n\endcsname{5}
\expandafter\gdef\csname num@d:pd:plantdoc_5:mnv3-dacisPMP:worst\endcsname{4.9}
\expandafter\gdef\csname num@d:pd:plantdoc_5:effb0-dacisPMP\endcsname{+11.1}
\expandafter\gdef\csname num@d:pd:plantdoc_5:effb0-dacisPMP:lo\endcsname{+8.8}
\expandafter\gdef\csname num@d:pd:plantdoc_5:effb0-dacisPMP:hi\endcsname{+13.3}
\expandafter\gdef\csname num@d:pd:plantdoc_5:effb0-dacisPMP:abs\endcsname{11.1}
\expandafter\gdef\csname num@d:pd:plantdoc_5:effb0-dacisPMP:wins\endcsname{5}
\expandafter\gdef\csname num@d:pd:plantdoc_5:effb0-dacisPMP:n\endcsname{5}
\expandafter\gdef\csname num@d:pd:plantdoc_5:effb0-dacisPMP:worst\endcsname{13.3}
\expandafter\gdef\csname num@d:pd:plantdoc_5:mnv2-dacisPMP\endcsname{+7.7}
\expandafter\gdef\csname num@d:pd:plantdoc_5:mnv2-dacisPMP:lo\endcsname{+6.1}
\expandafter\gdef\csname num@d:pd:plantdoc_5:mnv2-dacisPMP:hi\endcsname{+9.2}
\expandafter\gdef\csname num@d:pd:plantdoc_5:mnv2-dacisPMP:abs\endcsname{7.7}
\expandafter\gdef\csname num@d:pd:plantdoc_5:mnv2-dacisPMP:wins\endcsname{5}
\expandafter\gdef\csname num@d:pd:plantdoc_5:mnv2-dacisPMP:n\endcsname{5}
\expandafter\gdef\csname num@d:pd:plantdoc_5:mnv2-dacisPMP:worst\endcsname{9.2}
\expandafter\gdef\csname num@d:pd:plantdoc_5:dacis-l1:PMP50\endcsname{+0.3}
\expandafter\gdef\csname num@d:pd:plantdoc_5:dacis-l1:PMP50:lo\endcsname{\ensuremath{-}1.1}
\expandafter\gdef\csname num@d:pd:plantdoc_5:dacis-l1:PMP50:hi\endcsname{+1.7}
\expandafter\gdef\csname num@d:pd:plantdoc_5:dacis-l1:PMP50:abs\endcsname{0.3}
\expandafter\gdef\csname num@d:pd:plantdoc_5:dacis-l1:PMP50:wins\endcsname{4}
\expandafter\gdef\csname num@d:pd:plantdoc_5:dacis-l1:PMP50:n\endcsname{5}
\expandafter\gdef\csname num@d:pd:plantdoc_5:dacis-l1:PMP50:worst\endcsname{1.7}
\expandafter\gdef\csname num@d:pd:plantdoc_5:dense-dacisPMP50\endcsname{+3.9}
\expandafter\gdef\csname num@d:pd:plantdoc_5:dense-dacisPMP50:lo\endcsname{+3.5}
\expandafter\gdef\csname num@d:pd:plantdoc_5:dense-dacisPMP50:hi\endcsname{+4.2}
\expandafter\gdef\csname num@d:pd:plantdoc_5:dense-dacisPMP50:abs\endcsname{3.9}
\expandafter\gdef\csname num@d:pd:plantdoc_5:dense-dacisPMP50:wins\endcsname{5}
\expandafter\gdef\csname num@d:pd:plantdoc_5:dense-dacisPMP50:n\endcsname{5}
\expandafter\gdef\csname num@d:pd:plantdoc_5:dense-dacisPMP50:worst\endcsname{4.2}
\expandafter\gdef\csname num@d:pd:plantdoc_5:dacis-l1:PMP20\endcsname{\ensuremath{-}1.3}
\expandafter\gdef\csname num@d:pd:plantdoc_5:dacis-l1:PMP20:lo\endcsname{\ensuremath{-}2.4}
\expandafter\gdef\csname num@d:pd:plantdoc_5:dacis-l1:PMP20:hi\endcsname{\ensuremath{-}0.1}
\expandafter\gdef\csname num@d:pd:plantdoc_5:dacis-l1:PMP20:abs\endcsname{1.3}
\expandafter\gdef\csname num@d:pd:plantdoc_5:dacis-l1:PMP20:wins\endcsname{1}
\expandafter\gdef\csname num@d:pd:plantdoc_5:dacis-l1:PMP20:n\endcsname{5}
\expandafter\gdef\csname num@d:pd:plantdoc_5:dacis-l1:PMP20:worst\endcsname{2.4}
\expandafter\gdef\csname num@d:pd:plantdoc_5:dense-dacisPMP20\endcsname{+11.4}
\expandafter\gdef\csname num@d:pd:plantdoc_5:dense-dacisPMP20:lo\endcsname{+9.7}
\expandafter\gdef\csname num@d:pd:plantdoc_5:dense-dacisPMP20:hi\endcsname{+13.0}
\expandafter\gdef\csname num@d:pd:plantdoc_5:dense-dacisPMP20:abs\endcsname{11.4}
\expandafter\gdef\csname num@d:pd:plantdoc_5:dense-dacisPMP20:wins\endcsname{5}
\expandafter\gdef\csname num@d:pd:plantdoc_5:dense-dacisPMP20:n\endcsname{5}
\expandafter\gdef\csname num@d:pd:plantdoc_5:dense-dacisPMP20:worst\endcsname{13.0}
\expandafter\gdef\csname num@d:pd:plantdoc_5:dacis-l1:PMP10\endcsname{+1.4}
\expandafter\gdef\csname num@d:pd:plantdoc_5:dacis-l1:PMP10:lo\endcsname{\ensuremath{-}0.3}
\expandafter\gdef\csname num@d:pd:plantdoc_5:dacis-l1:PMP10:hi\endcsname{+3.2}
\expandafter\gdef\csname num@d:pd:plantdoc_5:dacis-l1:PMP10:abs\endcsname{1.4}
\expandafter\gdef\csname num@d:pd:plantdoc_5:dacis-l1:PMP10:wins\endcsname{4}
\expandafter\gdef\csname num@d:pd:plantdoc_5:dacis-l1:PMP10:n\endcsname{5}
\expandafter\gdef\csname num@d:pd:plantdoc_5:dacis-l1:PMP10:worst\endcsname{3.2}
\expandafter\gdef\csname num@d:pd:plantdoc_5:dense-dacisPMP10\endcsname{+14.9}
\expandafter\gdef\csname num@d:pd:plantdoc_5:dense-dacisPMP10:lo\endcsname{+13.9}
\expandafter\gdef\csname num@d:pd:plantdoc_5:dense-dacisPMP10:hi\endcsname{+16.0}
\expandafter\gdef\csname num@d:pd:plantdoc_5:dense-dacisPMP10:abs\endcsname{14.9}
\expandafter\gdef\csname num@d:pd:plantdoc_5:dense-dacisPMP10:wins\endcsname{5}
\expandafter\gdef\csname num@d:pd:plantdoc_5:dense-dacisPMP10:n\endcsname{5}
\expandafter\gdef\csname num@d:pd:plantdoc_5:dense-dacisPMP10:worst\endcsname{16.0}
\expandafter\gdef\csname num@d:pv:novel_1:dacis-random:PMP\endcsname{\ensuremath{-}0.9}
\expandafter\gdef\csname num@d:pv:novel_1:dacis-random:PMP:lo\endcsname{\ensuremath{-}2.7}
\expandafter\gdef\csname num@d:pv:novel_1:dacis-random:PMP:hi\endcsname{+0.9}
\expandafter\gdef\csname num@d:pv:novel_1:dacis-random:PMP:abs\endcsname{0.9}
\expandafter\gdef\csname num@d:pv:novel_1:dacis-random:PMP:wins\endcsname{1}
\expandafter\gdef\csname num@d:pv:novel_1:dacis-random:PMP:n\endcsname{5}
\expandafter\gdef\csname num@d:pv:novel_1:dacis-random:PMP:worst\endcsname{2.7}
\expandafter\gdef\csname num@d:pv:novel_1:dacis-random:P\endcsname{+1.0}
\expandafter\gdef\csname num@d:pv:novel_1:dacis-random:P:lo\endcsname{\ensuremath{-}2.3}
\expandafter\gdef\csname num@d:pv:novel_1:dacis-random:P:hi\endcsname{+4.4}
\expandafter\gdef\csname num@d:pv:novel_1:dacis-random:P:abs\endcsname{1.0}
\expandafter\gdef\csname num@d:pv:novel_1:dacis-random:P:wins\endcsname{2}
\expandafter\gdef\csname num@d:pv:novel_1:dacis-random:P:n\endcsname{5}
\expandafter\gdef\csname num@d:pv:novel_1:dacis-random:P:worst\endcsname{4.4}
\expandafter\gdef\csname num@d:pv:novel_1:dacis-l1:PMP\endcsname{\ensuremath{-}0.8}
\expandafter\gdef\csname num@d:pv:novel_1:dacis-l1:PMP:lo\endcsname{\ensuremath{-}2.9}
\expandafter\gdef\csname num@d:pv:novel_1:dacis-l1:PMP:hi\endcsname{+1.2}
\expandafter\gdef\csname num@d:pv:novel_1:dacis-l1:PMP:abs\endcsname{0.8}
\expandafter\gdef\csname num@d:pv:novel_1:dacis-l1:PMP:wins\endcsname{1}
\expandafter\gdef\csname num@d:pv:novel_1:dacis-l1:PMP:n\endcsname{5}
\expandafter\gdef\csname num@d:pv:novel_1:dacis-l1:PMP:worst\endcsname{2.9}
\expandafter\gdef\csname num@d:pv:novel_1:dacis-l1:P\endcsname{\ensuremath{-}0.2}
\expandafter\gdef\csname num@d:pv:novel_1:dacis-l1:P:lo\endcsname{\ensuremath{-}3.7}
\expandafter\gdef\csname num@d:pv:novel_1:dacis-l1:P:hi\endcsname{+3.4}
\expandafter\gdef\csname num@d:pv:novel_1:dacis-l1:P:abs\endcsname{0.2}
\expandafter\gdef\csname num@d:pv:novel_1:dacis-l1:P:wins\endcsname{2}
\expandafter\gdef\csname num@d:pv:novel_1:dacis-l1:P:n\endcsname{5}
\expandafter\gdef\csname num@d:pv:novel_1:dacis-l1:P:worst\endcsname{3.7}
\expandafter\gdef\csname num@d:pv:novel_1:dacis-taylor:PMP\endcsname{\ensuremath{-}0.2}
\expandafter\gdef\csname num@d:pv:novel_1:dacis-taylor:PMP:lo\endcsname{\ensuremath{-}2.0}
\expandafter\gdef\csname num@d:pv:novel_1:dacis-taylor:PMP:hi\endcsname{+1.6}
\expandafter\gdef\csname num@d:pv:novel_1:dacis-taylor:PMP:abs\endcsname{0.2}
\expandafter\gdef\csname num@d:pv:novel_1:dacis-taylor:PMP:wins\endcsname{3}
\expandafter\gdef\csname num@d:pv:novel_1:dacis-taylor:PMP:n\endcsname{5}
\expandafter\gdef\csname num@d:pv:novel_1:dacis-taylor:PMP:worst\endcsname{2.0}
\expandafter\gdef\csname num@d:pv:novel_1:dacis-taylor:P\endcsname{\ensuremath{-}2.3}
\expandafter\gdef\csname num@d:pv:novel_1:dacis-taylor:P:lo\endcsname{\ensuremath{-}3.9}
\expandafter\gdef\csname num@d:pv:novel_1:dacis-taylor:P:hi\endcsname{\ensuremath{-}0.8}
\expandafter\gdef\csname num@d:pv:novel_1:dacis-taylor:P:abs\endcsname{2.3}
\expandafter\gdef\csname num@d:pv:novel_1:dacis-taylor:P:wins\endcsname{0}
\expandafter\gdef\csname num@d:pv:novel_1:dacis-taylor:P:n\endcsname{5}
\expandafter\gdef\csname num@d:pv:novel_1:dacis-taylor:P:worst\endcsname{3.9}
\expandafter\gdef\csname num@d:pv:novel_1:dacis-bn_gamma:PMP\endcsname{+1.1}
\expandafter\gdef\csname num@d:pv:novel_1:dacis-bn_gamma:PMP:lo\endcsname{\ensuremath{-}1.1}
\expandafter\gdef\csname num@d:pv:novel_1:dacis-bn_gamma:PMP:hi\endcsname{+3.3}
\expandafter\gdef\csname num@d:pv:novel_1:dacis-bn_gamma:PMP:abs\endcsname{1.1}
\expandafter\gdef\csname num@d:pv:novel_1:dacis-bn_gamma:PMP:wins\endcsname{4}
\expandafter\gdef\csname num@d:pv:novel_1:dacis-bn_gamma:PMP:n\endcsname{5}
\expandafter\gdef\csname num@d:pv:novel_1:dacis-bn_gamma:PMP:worst\endcsname{3.3}
\expandafter\gdef\csname num@d:pv:novel_1:dacis-bn_gamma:P\endcsname{+1.1}
\expandafter\gdef\csname num@d:pv:novel_1:dacis-bn_gamma:P:lo\endcsname{\ensuremath{-}3.0}
\expandafter\gdef\csname num@d:pv:novel_1:dacis-bn_gamma:P:hi\endcsname{+5.1}
\expandafter\gdef\csname num@d:pv:novel_1:dacis-bn_gamma:P:abs\endcsname{1.1}
\expandafter\gdef\csname num@d:pv:novel_1:dacis-bn_gamma:P:wins\endcsname{2}
\expandafter\gdef\csname num@d:pv:novel_1:dacis-bn_gamma:P:n\endcsname{5}
\expandafter\gdef\csname num@d:pv:novel_1:dacis-bn_gamma:P:worst\endcsname{5.1}
\expandafter\gdef\csname num@d:pv:novel_1:dacis-D:PMP\endcsname{+1.0}
\expandafter\gdef\csname num@d:pv:novel_1:dacis-D:PMP:lo\endcsname{\ensuremath{-}1.2}
\expandafter\gdef\csname num@d:pv:novel_1:dacis-D:PMP:hi\endcsname{+3.3}
\expandafter\gdef\csname num@d:pv:novel_1:dacis-D:PMP:abs\endcsname{1.0}
\expandafter\gdef\csname num@d:pv:novel_1:dacis-D:PMP:wins\endcsname{3}
\expandafter\gdef\csname num@d:pv:novel_1:dacis-D:PMP:n\endcsname{5}
\expandafter\gdef\csname num@d:pv:novel_1:dacis-D:PMP:worst\endcsname{3.3}
\expandafter\gdef\csname num@d:pv:novel_1:dacis-D:P\endcsname{+0.8}
\expandafter\gdef\csname num@d:pv:novel_1:dacis-D:P:lo\endcsname{\ensuremath{-}0.4}
\expandafter\gdef\csname num@d:pv:novel_1:dacis-D:P:hi\endcsname{+2.0}
\expandafter\gdef\csname num@d:pv:novel_1:dacis-D:P:abs\endcsname{0.8}
\expandafter\gdef\csname num@d:pv:novel_1:dacis-D:P:wins\endcsname{4}
\expandafter\gdef\csname num@d:pv:novel_1:dacis-D:P:n\endcsname{5}
\expandafter\gdef\csname num@d:pv:novel_1:dacis-D:P:worst\endcsname{2.0}
\expandafter\gdef\csname num@d:pv:novel_1:dacis-G:PMP\endcsname{\ensuremath{-}0.1}
\expandafter\gdef\csname num@d:pv:novel_1:dacis-G:PMP:lo\endcsname{\ensuremath{-}1.3}
\expandafter\gdef\csname num@d:pv:novel_1:dacis-G:PMP:hi\endcsname{+1.0}
\expandafter\gdef\csname num@d:pv:novel_1:dacis-G:PMP:abs\endcsname{0.1}
\expandafter\gdef\csname num@d:pv:novel_1:dacis-G:PMP:wins\endcsname{3}
\expandafter\gdef\csname num@d:pv:novel_1:dacis-G:PMP:n\endcsname{5}
\expandafter\gdef\csname num@d:pv:novel_1:dacis-G:PMP:worst\endcsname{1.3}
\expandafter\gdef\csname num@d:pv:novel_1:dacis-G:P\endcsname{\ensuremath{-}1.3}
\expandafter\gdef\csname num@d:pv:novel_1:dacis-G:P:lo\endcsname{\ensuremath{-}2.9}
\expandafter\gdef\csname num@d:pv:novel_1:dacis-G:P:hi\endcsname{+0.3}
\expandafter\gdef\csname num@d:pv:novel_1:dacis-G:P:abs\endcsname{1.3}
\expandafter\gdef\csname num@d:pv:novel_1:dacis-G:P:wins\endcsname{1}
\expandafter\gdef\csname num@d:pv:novel_1:dacis-G:P:n\endcsname{5}
\expandafter\gdef\csname num@d:pv:novel_1:dacis-G:P:worst\endcsname{2.9}
\expandafter\gdef\csname num@d:pv:novel_1:dacis-V:PMP\endcsname{+0.2}
\expandafter\gdef\csname num@d:pv:novel_1:dacis-V:PMP:lo\endcsname{\ensuremath{-}1.3}
\expandafter\gdef\csname num@d:pv:novel_1:dacis-V:PMP:hi\endcsname{+1.6}
\expandafter\gdef\csname num@d:pv:novel_1:dacis-V:PMP:abs\endcsname{0.2}
\expandafter\gdef\csname num@d:pv:novel_1:dacis-V:PMP:wins\endcsname{2}
\expandafter\gdef\csname num@d:pv:novel_1:dacis-V:PMP:n\endcsname{5}
\expandafter\gdef\csname num@d:pv:novel_1:dacis-V:PMP:worst\endcsname{1.6}
\expandafter\gdef\csname num@d:pv:novel_1:dacis-V:P\endcsname{\ensuremath{-}0.1}
\expandafter\gdef\csname num@d:pv:novel_1:dacis-V:P:lo\endcsname{\ensuremath{-}3.2}
\expandafter\gdef\csname num@d:pv:novel_1:dacis-V:P:hi\endcsname{+2.9}
\expandafter\gdef\csname num@d:pv:novel_1:dacis-V:P:abs\endcsname{0.1}
\expandafter\gdef\csname num@d:pv:novel_1:dacis-V:P:wins\endcsname{2}
\expandafter\gdef\csname num@d:pv:novel_1:dacis-V:P:n\endcsname{5}
\expandafter\gdef\csname num@d:pv:novel_1:dacis-V:P:worst\endcsname{3.2}
\expandafter\gdef\csname num@d:pv:novel_1:l1-random:PMP\endcsname{\ensuremath{-}0.1}
\expandafter\gdef\csname num@d:pv:novel_1:l1-random:PMP:lo\endcsname{\ensuremath{-}3.1}
\expandafter\gdef\csname num@d:pv:novel_1:l1-random:PMP:hi\endcsname{+2.9}
\expandafter\gdef\csname num@d:pv:novel_1:l1-random:PMP:abs\endcsname{0.1}
\expandafter\gdef\csname num@d:pv:novel_1:l1-random:PMP:wins\endcsname{3}
\expandafter\gdef\csname num@d:pv:novel_1:l1-random:PMP:n\endcsname{5}
\expandafter\gdef\csname num@d:pv:novel_1:l1-random:PMP:worst\endcsname{3.1}
\expandafter\gdef\csname num@d:pv:novel_1:l1-random:P\endcsname{+1.2}
\expandafter\gdef\csname num@d:pv:novel_1:l1-random:P:lo\endcsname{\ensuremath{-}1.5}
\expandafter\gdef\csname num@d:pv:novel_1:l1-random:P:hi\endcsname{+3.9}
\expandafter\gdef\csname num@d:pv:novel_1:l1-random:P:abs\endcsname{1.2}
\expandafter\gdef\csname num@d:pv:novel_1:l1-random:P:wins\endcsname{4}
\expandafter\gdef\csname num@d:pv:novel_1:l1-random:P:n\endcsname{5}
\expandafter\gdef\csname num@d:pv:novel_1:l1-random:P:worst\endcsname{3.9}
\expandafter\gdef\csname num@d:pv:novel_1:taylor-random:PMP\endcsname{\ensuremath{-}0.7}
\expandafter\gdef\csname num@d:pv:novel_1:taylor-random:PMP:lo\endcsname{\ensuremath{-}2.5}
\expandafter\gdef\csname num@d:pv:novel_1:taylor-random:PMP:hi\endcsname{+1.1}
\expandafter\gdef\csname num@d:pv:novel_1:taylor-random:PMP:abs\endcsname{0.7}
\expandafter\gdef\csname num@d:pv:novel_1:taylor-random:PMP:wins\endcsname{2}
\expandafter\gdef\csname num@d:pv:novel_1:taylor-random:PMP:n\endcsname{5}
\expandafter\gdef\csname num@d:pv:novel_1:taylor-random:PMP:worst\endcsname{2.5}
\expandafter\gdef\csname num@d:pv:novel_1:taylor-random:P\endcsname{+3.3}
\expandafter\gdef\csname num@d:pv:novel_1:taylor-random:P:lo\endcsname{+0.4}
\expandafter\gdef\csname num@d:pv:novel_1:taylor-random:P:hi\endcsname{+6.3}
\expandafter\gdef\csname num@d:pv:novel_1:taylor-random:P:abs\endcsname{3.3}
\expandafter\gdef\csname num@d:pv:novel_1:taylor-random:P:wins\endcsname{5}
\expandafter\gdef\csname num@d:pv:novel_1:taylor-random:P:n\endcsname{5}
\expandafter\gdef\csname num@d:pv:novel_1:taylor-random:P:worst\endcsname{6.3}
\expandafter\gdef\csname num@d:pv:novel_1:bn_gamma-random:PMP\endcsname{\ensuremath{-}2.0}
\expandafter\gdef\csname num@d:pv:novel_1:bn_gamma-random:PMP:lo\endcsname{\ensuremath{-}3.3}
\expandafter\gdef\csname num@d:pv:novel_1:bn_gamma-random:PMP:hi\endcsname{\ensuremath{-}0.6}
\expandafter\gdef\csname num@d:pv:novel_1:bn_gamma-random:PMP:abs\endcsname{2.0}
\expandafter\gdef\csname num@d:pv:novel_1:bn_gamma-random:PMP:wins\endcsname{0}
\expandafter\gdef\csname num@d:pv:novel_1:bn_gamma-random:PMP:n\endcsname{5}
\expandafter\gdef\csname num@d:pv:novel_1:bn_gamma-random:PMP:worst\endcsname{3.3}
\expandafter\gdef\csname num@d:pv:novel_1:bn_gamma-random:P\endcsname{\ensuremath{-}0.0}
\expandafter\gdef\csname num@d:pv:novel_1:bn_gamma-random:P:lo\endcsname{\ensuremath{-}5.3}
\expandafter\gdef\csname num@d:pv:novel_1:bn_gamma-random:P:hi\endcsname{+5.3}
\expandafter\gdef\csname num@d:pv:novel_1:bn_gamma-random:P:abs\endcsname{0.0}
\expandafter\gdef\csname num@d:pv:novel_1:bn_gamma-random:P:wins\endcsname{3}
\expandafter\gdef\csname num@d:pv:novel_1:bn_gamma-random:P:n\endcsname{5}
\expandafter\gdef\csname num@d:pv:novel_1:bn_gamma-random:P:worst\endcsname{5.3}
\expandafter\gdef\csname num@d:pv:novel_1:D-random:PMP\endcsname{\ensuremath{-}1.9}
\expandafter\gdef\csname num@d:pv:novel_1:D-random:PMP:lo\endcsname{\ensuremath{-}5.5}
\expandafter\gdef\csname num@d:pv:novel_1:D-random:PMP:hi\endcsname{+1.6}
\expandafter\gdef\csname num@d:pv:novel_1:D-random:PMP:abs\endcsname{1.9}
\expandafter\gdef\csname num@d:pv:novel_1:D-random:PMP:wins\endcsname{1}
\expandafter\gdef\csname num@d:pv:novel_1:D-random:PMP:n\endcsname{5}
\expandafter\gdef\csname num@d:pv:novel_1:D-random:PMP:worst\endcsname{5.5}
\expandafter\gdef\csname num@d:pv:novel_1:D-random:P\endcsname{+0.2}
\expandafter\gdef\csname num@d:pv:novel_1:D-random:P:lo\endcsname{\ensuremath{-}3.1}
\expandafter\gdef\csname num@d:pv:novel_1:D-random:P:hi\endcsname{+3.6}
\expandafter\gdef\csname num@d:pv:novel_1:D-random:P:abs\endcsname{0.2}
\expandafter\gdef\csname num@d:pv:novel_1:D-random:P:wins\endcsname{2}
\expandafter\gdef\csname num@d:pv:novel_1:D-random:P:n\endcsname{5}
\expandafter\gdef\csname num@d:pv:novel_1:D-random:P:worst\endcsname{3.6}
\expandafter\gdef\csname num@d:pv:novel_1:G-random:PMP\endcsname{\ensuremath{-}0.8}
\expandafter\gdef\csname num@d:pv:novel_1:G-random:PMP:lo\endcsname{\ensuremath{-}2.4}
\expandafter\gdef\csname num@d:pv:novel_1:G-random:PMP:hi\endcsname{+0.8}
\expandafter\gdef\csname num@d:pv:novel_1:G-random:PMP:abs\endcsname{0.8}
\expandafter\gdef\csname num@d:pv:novel_1:G-random:PMP:wins\endcsname{1}
\expandafter\gdef\csname num@d:pv:novel_1:G-random:PMP:n\endcsname{5}
\expandafter\gdef\csname num@d:pv:novel_1:G-random:PMP:worst\endcsname{2.4}
\expandafter\gdef\csname num@d:pv:novel_1:G-random:P\endcsname{+2.4}
\expandafter\gdef\csname num@d:pv:novel_1:G-random:P:lo\endcsname{+0.5}
\expandafter\gdef\csname num@d:pv:novel_1:G-random:P:hi\endcsname{+4.2}
\expandafter\gdef\csname num@d:pv:novel_1:G-random:P:abs\endcsname{2.4}
\expandafter\gdef\csname num@d:pv:novel_1:G-random:P:wins\endcsname{5}
\expandafter\gdef\csname num@d:pv:novel_1:G-random:P:n\endcsname{5}
\expandafter\gdef\csname num@d:pv:novel_1:G-random:P:worst\endcsname{4.2}
\expandafter\gdef\csname num@d:pv:novel_1:V-random:PMP\endcsname{\ensuremath{-}1.1}
\expandafter\gdef\csname num@d:pv:novel_1:V-random:PMP:lo\endcsname{\ensuremath{-}3.6}
\expandafter\gdef\csname num@d:pv:novel_1:V-random:PMP:hi\endcsname{+1.4}
\expandafter\gdef\csname num@d:pv:novel_1:V-random:PMP:abs\endcsname{1.1}
\expandafter\gdef\csname num@d:pv:novel_1:V-random:PMP:wins\endcsname{1}
\expandafter\gdef\csname num@d:pv:novel_1:V-random:PMP:n\endcsname{5}
\expandafter\gdef\csname num@d:pv:novel_1:V-random:PMP:worst\endcsname{3.6}
\expandafter\gdef\csname num@d:pv:novel_1:V-random:P\endcsname{+1.2}
\expandafter\gdef\csname num@d:pv:novel_1:V-random:P:lo\endcsname{\ensuremath{-}1.3}
\expandafter\gdef\csname num@d:pv:novel_1:V-random:P:hi\endcsname{+3.6}
\expandafter\gdef\csname num@d:pv:novel_1:V-random:P:abs\endcsname{1.2}
\expandafter\gdef\csname num@d:pv:novel_1:V-random:P:wins\endcsname{4}
\expandafter\gdef\csname num@d:pv:novel_1:V-random:P:n\endcsname{5}
\expandafter\gdef\csname num@d:pv:novel_1:V-random:P:worst\endcsname{3.6}
\expandafter\gdef\csname num@d:pv:novel_1:PM-P\endcsname{\ensuremath{-}0.2}
\expandafter\gdef\csname num@d:pv:novel_1:PM-P:lo\endcsname{\ensuremath{-}1.8}
\expandafter\gdef\csname num@d:pv:novel_1:PM-P:hi\endcsname{+1.4}
\expandafter\gdef\csname num@d:pv:novel_1:PM-P:abs\endcsname{0.2}
\expandafter\gdef\csname num@d:pv:novel_1:PM-P:wins\endcsname{2}
\expandafter\gdef\csname num@d:pv:novel_1:PM-P:n\endcsname{5}
\expandafter\gdef\csname num@d:pv:novel_1:PM-P:worst\endcsname{1.8}
\expandafter\gdef\csname num@d:pv:novel_1:MP-P\endcsname{+0.6}
\expandafter\gdef\csname num@d:pv:novel_1:MP-P:lo\endcsname{\ensuremath{-}0.8}
\expandafter\gdef\csname num@d:pv:novel_1:MP-P:hi\endcsname{+2.1}
\expandafter\gdef\csname num@d:pv:novel_1:MP-P:abs\endcsname{0.6}
\expandafter\gdef\csname num@d:pv:novel_1:MP-P:wins\endcsname{3}
\expandafter\gdef\csname num@d:pv:novel_1:MP-P:n\endcsname{5}
\expandafter\gdef\csname num@d:pv:novel_1:MP-P:worst\endcsname{2.1}
\expandafter\gdef\csname num@d:pv:novel_1:PGP-P\endcsname{+0.3}
\expandafter\gdef\csname num@d:pv:novel_1:PGP-P:lo\endcsname{\ensuremath{-}0.7}
\expandafter\gdef\csname num@d:pv:novel_1:PGP-P:hi\endcsname{+1.3}
\expandafter\gdef\csname num@d:pv:novel_1:PGP-P:abs\endcsname{0.3}
\expandafter\gdef\csname num@d:pv:novel_1:PGP-P:wins\endcsname{4}
\expandafter\gdef\csname num@d:pv:novel_1:PGP-P:n\endcsname{5}
\expandafter\gdef\csname num@d:pv:novel_1:PGP-P:worst\endcsname{1.3}
\expandafter\gdef\csname num@d:pv:novel_1:PMPM-P\endcsname{+0.5}
\expandafter\gdef\csname num@d:pv:novel_1:PMPM-P:lo\endcsname{\ensuremath{-}1.0}
\expandafter\gdef\csname num@d:pv:novel_1:PMPM-P:hi\endcsname{+1.9}
\expandafter\gdef\csname num@d:pv:novel_1:PMPM-P:abs\endcsname{0.5}
\expandafter\gdef\csname num@d:pv:novel_1:PMPM-P:wins\endcsname{3}
\expandafter\gdef\csname num@d:pv:novel_1:PMPM-P:n\endcsname{5}
\expandafter\gdef\csname num@d:pv:novel_1:PMPM-P:worst\endcsname{1.9}
\expandafter\gdef\csname num@d:pv:novel_1:PM2x-P\endcsname{\ensuremath{-}0.3}
\expandafter\gdef\csname num@d:pv:novel_1:PM2x-P:lo\endcsname{\ensuremath{-}1.4}
\expandafter\gdef\csname num@d:pv:novel_1:PM2x-P:hi\endcsname{+0.9}
\expandafter\gdef\csname num@d:pv:novel_1:PM2x-P:abs\endcsname{0.3}
\expandafter\gdef\csname num@d:pv:novel_1:PM2x-P:wins\endcsname{3}
\expandafter\gdef\csname num@d:pv:novel_1:PM2x-P:n\endcsname{5}
\expandafter\gdef\csname num@d:pv:novel_1:PM2x-P:worst\endcsname{1.4}
\expandafter\gdef\csname num@d:pv:novel_1:PMP-P\endcsname{\ensuremath{-}0.0}
\expandafter\gdef\csname num@d:pv:novel_1:PMP-P:lo\endcsname{\ensuremath{-}1.6}
\expandafter\gdef\csname num@d:pv:novel_1:PMP-P:hi\endcsname{+1.6}
\expandafter\gdef\csname num@d:pv:novel_1:PMP-P:abs\endcsname{0.0}
\expandafter\gdef\csname num@d:pv:novel_1:PMP-P:wins\endcsname{3}
\expandafter\gdef\csname num@d:pv:novel_1:PMP-P:n\endcsname{5}
\expandafter\gdef\csname num@d:pv:novel_1:PMP-P:worst\endcsname{1.6}
\expandafter\gdef\csname num@d:pv:novel_1:densem-dense\endcsname{\ensuremath{-}0.0}
\expandafter\gdef\csname num@d:pv:novel_1:densem-dense:lo\endcsname{\ensuremath{-}0.8}
\expandafter\gdef\csname num@d:pv:novel_1:densem-dense:hi\endcsname{+0.8}
\expandafter\gdef\csname num@d:pv:novel_1:densem-dense:abs\endcsname{0.0}
\expandafter\gdef\csname num@d:pv:novel_1:densem-dense:wins\endcsname{2}
\expandafter\gdef\csname num@d:pv:novel_1:densem-dense:n\endcsname{5}
\expandafter\gdef\csname num@d:pv:novel_1:densem-dense:worst\endcsname{0.8}
\expandafter\gdef\csname num@d:pv:novel_1:densermr-dense\endcsname{\ensuremath{-}0.3}
\expandafter\gdef\csname num@d:pv:novel_1:densermr-dense:lo\endcsname{\ensuremath{-}0.9}
\expandafter\gdef\csname num@d:pv:novel_1:densermr-dense:hi\endcsname{+0.3}
\expandafter\gdef\csname num@d:pv:novel_1:densermr-dense:abs\endcsname{0.3}
\expandafter\gdef\csname num@d:pv:novel_1:densermr-dense:wins\endcsname{1}
\expandafter\gdef\csname num@d:pv:novel_1:densermr-dense:n\endcsname{5}
\expandafter\gdef\csname num@d:pv:novel_1:densermr-dense:worst\endcsname{0.9}
\expandafter\gdef\csname num@d:pv:novel_1:dense-dacisP\endcsname{+6.8}
\expandafter\gdef\csname num@d:pv:novel_1:dense-dacisP:lo\endcsname{+4.0}
\expandafter\gdef\csname num@d:pv:novel_1:dense-dacisP:hi\endcsname{+9.6}
\expandafter\gdef\csname num@d:pv:novel_1:dense-dacisP:abs\endcsname{6.8}
\expandafter\gdef\csname num@d:pv:novel_1:dense-dacisP:wins\endcsname{5}
\expandafter\gdef\csname num@d:pv:novel_1:dense-dacisP:n\endcsname{5}
\expandafter\gdef\csname num@d:pv:novel_1:dense-dacisP:worst\endcsname{9.6}
\expandafter\gdef\csname num@d:pv:novel_1:dense-dacisPMP\endcsname{+6.8}
\expandafter\gdef\csname num@d:pv:novel_1:dense-dacisPMP:lo\endcsname{+4.2}
\expandafter\gdef\csname num@d:pv:novel_1:dense-dacisPMP:hi\endcsname{+9.4}
\expandafter\gdef\csname num@d:pv:novel_1:dense-dacisPMP:abs\endcsname{6.8}
\expandafter\gdef\csname num@d:pv:novel_1:dense-dacisPMP:wins\endcsname{5}
\expandafter\gdef\csname num@d:pv:novel_1:dense-dacisPMP:n\endcsname{5}
\expandafter\gdef\csname num@d:pv:novel_1:dense-dacisPMP:worst\endcsname{9.4}
\expandafter\gdef\csname num@d:pv:novel_1:dense-randomP\endcsname{+7.8}
\expandafter\gdef\csname num@d:pv:novel_1:dense-randomP:lo\endcsname{+4.2}
\expandafter\gdef\csname num@d:pv:novel_1:dense-randomP:hi\endcsname{+11.4}
\expandafter\gdef\csname num@d:pv:novel_1:dense-randomP:abs\endcsname{7.8}
\expandafter\gdef\csname num@d:pv:novel_1:dense-randomP:wins\endcsname{5}
\expandafter\gdef\csname num@d:pv:novel_1:dense-randomP:n\endcsname{5}
\expandafter\gdef\csname num@d:pv:novel_1:dense-randomP:worst\endcsname{11.4}
\expandafter\gdef\csname num@d:pv:novel_1:dense-randomPMP\endcsname{+5.9}
\expandafter\gdef\csname num@d:pv:novel_1:dense-randomPMP:lo\endcsname{+2.4}
\expandafter\gdef\csname num@d:pv:novel_1:dense-randomPMP:hi\endcsname{+9.4}
\expandafter\gdef\csname num@d:pv:novel_1:dense-randomPMP:abs\endcsname{5.9}
\expandafter\gdef\csname num@d:pv:novel_1:dense-randomPMP:wins\endcsname{5}
\expandafter\gdef\csname num@d:pv:novel_1:dense-randomPMP:n\endcsname{5}
\expandafter\gdef\csname num@d:pv:novel_1:dense-randomPMP:worst\endcsname{9.4}
\expandafter\gdef\csname num@d:pv:novel_1:norefine-full\endcsname{\ensuremath{-}0.9}
\expandafter\gdef\csname num@d:pv:novel_1:norefine-full:lo\endcsname{\ensuremath{-}2.3}
\expandafter\gdef\csname num@d:pv:novel_1:norefine-full:hi\endcsname{+0.5}
\expandafter\gdef\csname num@d:pv:novel_1:norefine-full:abs\endcsname{0.9}
\expandafter\gdef\csname num@d:pv:novel_1:norefine-full:wins\endcsname{0}
\expandafter\gdef\csname num@d:pv:novel_1:norefine-full:n\endcsname{5}
\expandafter\gdef\csname num@d:pv:novel_1:norefine-full:worst\endcsname{2.3}
\expandafter\gdef\csname num@d:pv:novel_1:mnv3-dacisPMP\endcsname{+5.2}
\expandafter\gdef\csname num@d:pv:novel_1:mnv3-dacisPMP:lo\endcsname{+2.7}
\expandafter\gdef\csname num@d:pv:novel_1:mnv3-dacisPMP:hi\endcsname{+7.8}
\expandafter\gdef\csname num@d:pv:novel_1:mnv3-dacisPMP:abs\endcsname{5.2}
\expandafter\gdef\csname num@d:pv:novel_1:mnv3-dacisPMP:wins\endcsname{5}
\expandafter\gdef\csname num@d:pv:novel_1:mnv3-dacisPMP:n\endcsname{5}
\expandafter\gdef\csname num@d:pv:novel_1:mnv3-dacisPMP:worst\endcsname{7.8}
\expandafter\gdef\csname num@d:pv:novel_1:effb0-dacisPMP\endcsname{+10.0}
\expandafter\gdef\csname num@d:pv:novel_1:effb0-dacisPMP:lo\endcsname{+8.0}
\expandafter\gdef\csname num@d:pv:novel_1:effb0-dacisPMP:hi\endcsname{+12.0}
\expandafter\gdef\csname num@d:pv:novel_1:effb0-dacisPMP:abs\endcsname{10.0}
\expandafter\gdef\csname num@d:pv:novel_1:effb0-dacisPMP:wins\endcsname{5}
\expandafter\gdef\csname num@d:pv:novel_1:effb0-dacisPMP:n\endcsname{5}
\expandafter\gdef\csname num@d:pv:novel_1:effb0-dacisPMP:worst\endcsname{12.0}
\expandafter\gdef\csname num@d:pv:novel_1:mnv2-dacisPMP\endcsname{+6.5}
\expandafter\gdef\csname num@d:pv:novel_1:mnv2-dacisPMP:lo\endcsname{+3.8}
\expandafter\gdef\csname num@d:pv:novel_1:mnv2-dacisPMP:hi\endcsname{+9.2}
\expandafter\gdef\csname num@d:pv:novel_1:mnv2-dacisPMP:abs\endcsname{6.5}
\expandafter\gdef\csname num@d:pv:novel_1:mnv2-dacisPMP:wins\endcsname{5}
\expandafter\gdef\csname num@d:pv:novel_1:mnv2-dacisPMP:n\endcsname{5}
\expandafter\gdef\csname num@d:pv:novel_1:mnv2-dacisPMP:worst\endcsname{9.2}
\expandafter\gdef\csname num@d:pv:novel_1:dacis-l1:PMP50\endcsname{\ensuremath{-}0.4}
\expandafter\gdef\csname num@d:pv:novel_1:dacis-l1:PMP50:lo\endcsname{\ensuremath{-}3.9}
\expandafter\gdef\csname num@d:pv:novel_1:dacis-l1:PMP50:hi\endcsname{+3.1}
\expandafter\gdef\csname num@d:pv:novel_1:dacis-l1:PMP50:abs\endcsname{0.4}
\expandafter\gdef\csname num@d:pv:novel_1:dacis-l1:PMP50:wins\endcsname{2}
\expandafter\gdef\csname num@d:pv:novel_1:dacis-l1:PMP50:n\endcsname{5}
\expandafter\gdef\csname num@d:pv:novel_1:dacis-l1:PMP50:worst\endcsname{3.9}
\expandafter\gdef\csname num@d:pv:novel_1:dense-dacisPMP50\endcsname{+4.0}
\expandafter\gdef\csname num@d:pv:novel_1:dense-dacisPMP50:lo\endcsname{+2.1}
\expandafter\gdef\csname num@d:pv:novel_1:dense-dacisPMP50:hi\endcsname{+5.9}
\expandafter\gdef\csname num@d:pv:novel_1:dense-dacisPMP50:abs\endcsname{4.0}
\expandafter\gdef\csname num@d:pv:novel_1:dense-dacisPMP50:wins\endcsname{5}
\expandafter\gdef\csname num@d:pv:novel_1:dense-dacisPMP50:n\endcsname{5}
\expandafter\gdef\csname num@d:pv:novel_1:dense-dacisPMP50:worst\endcsname{5.9}
\expandafter\gdef\csname num@d:pv:novel_1:dacis-l1:PMP20\endcsname{\ensuremath{-}1.6}
\expandafter\gdef\csname num@d:pv:novel_1:dacis-l1:PMP20:lo\endcsname{\ensuremath{-}3.4}
\expandafter\gdef\csname num@d:pv:novel_1:dacis-l1:PMP20:hi\endcsname{+0.2}
\expandafter\gdef\csname num@d:pv:novel_1:dacis-l1:PMP20:abs\endcsname{1.6}
\expandafter\gdef\csname num@d:pv:novel_1:dacis-l1:PMP20:wins\endcsname{1}
\expandafter\gdef\csname num@d:pv:novel_1:dacis-l1:PMP20:n\endcsname{5}
\expandafter\gdef\csname num@d:pv:novel_1:dacis-l1:PMP20:worst\endcsname{3.4}
\expandafter\gdef\csname num@d:pv:novel_1:dense-dacisPMP20\endcsname{+10.0}
\expandafter\gdef\csname num@d:pv:novel_1:dense-dacisPMP20:lo\endcsname{+7.5}
\expandafter\gdef\csname num@d:pv:novel_1:dense-dacisPMP20:hi\endcsname{+12.5}
\expandafter\gdef\csname num@d:pv:novel_1:dense-dacisPMP20:abs\endcsname{10.0}
\expandafter\gdef\csname num@d:pv:novel_1:dense-dacisPMP20:wins\endcsname{5}
\expandafter\gdef\csname num@d:pv:novel_1:dense-dacisPMP20:n\endcsname{5}
\expandafter\gdef\csname num@d:pv:novel_1:dense-dacisPMP20:worst\endcsname{12.5}
\expandafter\gdef\csname num@d:pv:novel_1:dacis-l1:PMP10\endcsname{\ensuremath{-}2.1}
\expandafter\gdef\csname num@d:pv:novel_1:dacis-l1:PMP10:lo\endcsname{\ensuremath{-}7.2}
\expandafter\gdef\csname num@d:pv:novel_1:dacis-l1:PMP10:hi\endcsname{+3.0}
\expandafter\gdef\csname num@d:pv:novel_1:dacis-l1:PMP10:abs\endcsname{2.1}
\expandafter\gdef\csname num@d:pv:novel_1:dacis-l1:PMP10:wins\endcsname{2}
\expandafter\gdef\csname num@d:pv:novel_1:dacis-l1:PMP10:n\endcsname{5}
\expandafter\gdef\csname num@d:pv:novel_1:dacis-l1:PMP10:worst\endcsname{7.2}
\expandafter\gdef\csname num@d:pv:novel_1:dense-dacisPMP10\endcsname{+15.8}
\expandafter\gdef\csname num@d:pv:novel_1:dense-dacisPMP10:lo\endcsname{+12.8}
\expandafter\gdef\csname num@d:pv:novel_1:dense-dacisPMP10:hi\endcsname{+18.8}
\expandafter\gdef\csname num@d:pv:novel_1:dense-dacisPMP10:abs\endcsname{15.8}
\expandafter\gdef\csname num@d:pv:novel_1:dense-dacisPMP10:wins\endcsname{5}
\expandafter\gdef\csname num@d:pv:novel_1:dense-dacisPMP10:n\endcsname{5}
\expandafter\gdef\csname num@d:pv:novel_1:dense-dacisPMP10:worst\endcsname{18.8}
\expandafter\gdef\csname num@d:es:novel_1:dacis-random:PMP\endcsname{\ensuremath{-}0.8}
\expandafter\gdef\csname num@d:es:novel_1:dacis-random:PMP:lo\endcsname{\ensuremath{-}3.3}
\expandafter\gdef\csname num@d:es:novel_1:dacis-random:PMP:hi\endcsname{+1.7}
\expandafter\gdef\csname num@d:es:novel_1:dacis-random:PMP:abs\endcsname{0.8}
\expandafter\gdef\csname num@d:es:novel_1:dacis-random:PMP:wins\endcsname{2}
\expandafter\gdef\csname num@d:es:novel_1:dacis-random:PMP:n\endcsname{5}
\expandafter\gdef\csname num@d:es:novel_1:dacis-random:PMP:worst\endcsname{3.3}
\expandafter\gdef\csname num@d:es:novel_1:dacis-random:P\endcsname{+1.4}
\expandafter\gdef\csname num@d:es:novel_1:dacis-random:P:lo\endcsname{\ensuremath{-}1.2}
\expandafter\gdef\csname num@d:es:novel_1:dacis-random:P:hi\endcsname{+4.0}
\expandafter\gdef\csname num@d:es:novel_1:dacis-random:P:abs\endcsname{1.4}
\expandafter\gdef\csname num@d:es:novel_1:dacis-random:P:wins\endcsname{4}
\expandafter\gdef\csname num@d:es:novel_1:dacis-random:P:n\endcsname{5}
\expandafter\gdef\csname num@d:es:novel_1:dacis-random:P:worst\endcsname{4.0}
\expandafter\gdef\csname num@d:es:novel_1:dacis-l1:PMP\endcsname{\ensuremath{-}0.1}
\expandafter\gdef\csname num@d:es:novel_1:dacis-l1:PMP:lo\endcsname{\ensuremath{-}3.5}
\expandafter\gdef\csname num@d:es:novel_1:dacis-l1:PMP:hi\endcsname{+3.4}
\expandafter\gdef\csname num@d:es:novel_1:dacis-l1:PMP:abs\endcsname{0.1}
\expandafter\gdef\csname num@d:es:novel_1:dacis-l1:PMP:wins\endcsname{3}
\expandafter\gdef\csname num@d:es:novel_1:dacis-l1:PMP:n\endcsname{5}
\expandafter\gdef\csname num@d:es:novel_1:dacis-l1:PMP:worst\endcsname{3.5}
\expandafter\gdef\csname num@d:es:novel_1:dacis-l1:P\endcsname{+0.3}
\expandafter\gdef\csname num@d:es:novel_1:dacis-l1:P:lo\endcsname{\ensuremath{-}2.5}
\expandafter\gdef\csname num@d:es:novel_1:dacis-l1:P:hi\endcsname{+3.1}
\expandafter\gdef\csname num@d:es:novel_1:dacis-l1:P:abs\endcsname{0.3}
\expandafter\gdef\csname num@d:es:novel_1:dacis-l1:P:wins\endcsname{2}
\expandafter\gdef\csname num@d:es:novel_1:dacis-l1:P:n\endcsname{5}
\expandafter\gdef\csname num@d:es:novel_1:dacis-l1:P:worst\endcsname{3.1}
\expandafter\gdef\csname num@d:es:novel_1:dacis-taylor:PMP\endcsname{+0.1}
\expandafter\gdef\csname num@d:es:novel_1:dacis-taylor:PMP:lo\endcsname{\ensuremath{-}4.5}
\expandafter\gdef\csname num@d:es:novel_1:dacis-taylor:PMP:hi\endcsname{+4.7}
\expandafter\gdef\csname num@d:es:novel_1:dacis-taylor:PMP:abs\endcsname{0.1}
\expandafter\gdef\csname num@d:es:novel_1:dacis-taylor:PMP:wins\endcsname{3}
\expandafter\gdef\csname num@d:es:novel_1:dacis-taylor:PMP:n\endcsname{5}
\expandafter\gdef\csname num@d:es:novel_1:dacis-taylor:PMP:worst\endcsname{4.7}
\expandafter\gdef\csname num@d:es:novel_1:dacis-taylor:P\endcsname{\ensuremath{-}1.3}
\expandafter\gdef\csname num@d:es:novel_1:dacis-taylor:P:lo\endcsname{\ensuremath{-}4.6}
\expandafter\gdef\csname num@d:es:novel_1:dacis-taylor:P:hi\endcsname{+1.9}
\expandafter\gdef\csname num@d:es:novel_1:dacis-taylor:P:abs\endcsname{1.3}
\expandafter\gdef\csname num@d:es:novel_1:dacis-taylor:P:wins\endcsname{1}
\expandafter\gdef\csname num@d:es:novel_1:dacis-taylor:P:n\endcsname{5}
\expandafter\gdef\csname num@d:es:novel_1:dacis-taylor:P:worst\endcsname{4.6}
\expandafter\gdef\csname num@d:es:novel_1:dacis-bn_gamma:PMP\endcsname{+2.0}
\expandafter\gdef\csname num@d:es:novel_1:dacis-bn_gamma:PMP:lo\endcsname{+1.0}
\expandafter\gdef\csname num@d:es:novel_1:dacis-bn_gamma:PMP:hi\endcsname{+2.9}
\expandafter\gdef\csname num@d:es:novel_1:dacis-bn_gamma:PMP:abs\endcsname{2.0}
\expandafter\gdef\csname num@d:es:novel_1:dacis-bn_gamma:PMP:wins\endcsname{5}
\expandafter\gdef\csname num@d:es:novel_1:dacis-bn_gamma:PMP:n\endcsname{5}
\expandafter\gdef\csname num@d:es:novel_1:dacis-bn_gamma:PMP:worst\endcsname{2.9}
\expandafter\gdef\csname num@d:es:novel_1:dacis-D:PMP\endcsname{+0.7}
\expandafter\gdef\csname num@d:es:novel_1:dacis-D:PMP:lo\endcsname{\ensuremath{-}3.0}
\expandafter\gdef\csname num@d:es:novel_1:dacis-D:PMP:hi\endcsname{+4.4}
\expandafter\gdef\csname num@d:es:novel_1:dacis-D:PMP:abs\endcsname{0.7}
\expandafter\gdef\csname num@d:es:novel_1:dacis-D:PMP:wins\endcsname{4}
\expandafter\gdef\csname num@d:es:novel_1:dacis-D:PMP:n\endcsname{5}
\expandafter\gdef\csname num@d:es:novel_1:dacis-D:PMP:worst\endcsname{4.4}
\expandafter\gdef\csname num@d:es:novel_1:dacis-D:P\endcsname{+0.1}
\expandafter\gdef\csname num@d:es:novel_1:dacis-D:P:lo\endcsname{\ensuremath{-}0.4}
\expandafter\gdef\csname num@d:es:novel_1:dacis-D:P:hi\endcsname{+0.6}
\expandafter\gdef\csname num@d:es:novel_1:dacis-D:P:abs\endcsname{0.1}
\expandafter\gdef\csname num@d:es:novel_1:dacis-D:P:wins\endcsname{3}
\expandafter\gdef\csname num@d:es:novel_1:dacis-D:P:n\endcsname{5}
\expandafter\gdef\csname num@d:es:novel_1:dacis-D:P:worst\endcsname{0.6}
\expandafter\gdef\csname num@d:es:novel_1:dacis-G:PMP\endcsname{+0.8}
\expandafter\gdef\csname num@d:es:novel_1:dacis-G:PMP:lo\endcsname{\ensuremath{-}1.0}
\expandafter\gdef\csname num@d:es:novel_1:dacis-G:PMP:hi\endcsname{+2.7}
\expandafter\gdef\csname num@d:es:novel_1:dacis-G:PMP:abs\endcsname{0.8}
\expandafter\gdef\csname num@d:es:novel_1:dacis-G:PMP:wins\endcsname{4}
\expandafter\gdef\csname num@d:es:novel_1:dacis-G:PMP:n\endcsname{5}
\expandafter\gdef\csname num@d:es:novel_1:dacis-G:PMP:worst\endcsname{2.7}
\expandafter\gdef\csname num@d:es:novel_1:dacis-V:PMP\endcsname{+0.5}
\expandafter\gdef\csname num@d:es:novel_1:dacis-V:PMP:lo\endcsname{\ensuremath{-}2.9}
\expandafter\gdef\csname num@d:es:novel_1:dacis-V:PMP:hi\endcsname{+4.0}
\expandafter\gdef\csname num@d:es:novel_1:dacis-V:PMP:abs\endcsname{0.5}
\expandafter\gdef\csname num@d:es:novel_1:dacis-V:PMP:wins\endcsname{2}
\expandafter\gdef\csname num@d:es:novel_1:dacis-V:PMP:n\endcsname{5}
\expandafter\gdef\csname num@d:es:novel_1:dacis-V:PMP:worst\endcsname{4.0}
\expandafter\gdef\csname num@d:es:novel_1:l1-random:PMP\endcsname{\ensuremath{-}0.7}
\expandafter\gdef\csname num@d:es:novel_1:l1-random:PMP:lo\endcsname{\ensuremath{-}2.1}
\expandafter\gdef\csname num@d:es:novel_1:l1-random:PMP:hi\endcsname{+0.6}
\expandafter\gdef\csname num@d:es:novel_1:l1-random:PMP:abs\endcsname{0.7}
\expandafter\gdef\csname num@d:es:novel_1:l1-random:PMP:wins\endcsname{1}
\expandafter\gdef\csname num@d:es:novel_1:l1-random:PMP:n\endcsname{5}
\expandafter\gdef\csname num@d:es:novel_1:l1-random:PMP:worst\endcsname{2.1}
\expandafter\gdef\csname num@d:es:novel_1:l1-random:P\endcsname{+1.1}
\expandafter\gdef\csname num@d:es:novel_1:l1-random:P:lo\endcsname{\ensuremath{-}2.2}
\expandafter\gdef\csname num@d:es:novel_1:l1-random:P:hi\endcsname{+4.4}
\expandafter\gdef\csname num@d:es:novel_1:l1-random:P:abs\endcsname{1.1}
\expandafter\gdef\csname num@d:es:novel_1:l1-random:P:wins\endcsname{2}
\expandafter\gdef\csname num@d:es:novel_1:l1-random:P:n\endcsname{5}
\expandafter\gdef\csname num@d:es:novel_1:l1-random:P:worst\endcsname{4.4}
\expandafter\gdef\csname num@d:es:novel_1:taylor-random:PMP\endcsname{\ensuremath{-}0.9}
\expandafter\gdef\csname num@d:es:novel_1:taylor-random:PMP:lo\endcsname{\ensuremath{-}3.5}
\expandafter\gdef\csname num@d:es:novel_1:taylor-random:PMP:hi\endcsname{+1.7}
\expandafter\gdef\csname num@d:es:novel_1:taylor-random:PMP:abs\endcsname{0.9}
\expandafter\gdef\csname num@d:es:novel_1:taylor-random:PMP:wins\endcsname{2}
\expandafter\gdef\csname num@d:es:novel_1:taylor-random:PMP:n\endcsname{5}
\expandafter\gdef\csname num@d:es:novel_1:taylor-random:PMP:worst\endcsname{3.5}
\expandafter\gdef\csname num@d:es:novel_1:taylor-random:P\endcsname{+2.8}
\expandafter\gdef\csname num@d:es:novel_1:taylor-random:P:lo\endcsname{+0.0}
\expandafter\gdef\csname num@d:es:novel_1:taylor-random:P:hi\endcsname{+5.5}
\expandafter\gdef\csname num@d:es:novel_1:taylor-random:P:abs\endcsname{2.8}
\expandafter\gdef\csname num@d:es:novel_1:taylor-random:P:wins\endcsname{5}
\expandafter\gdef\csname num@d:es:novel_1:taylor-random:P:n\endcsname{5}
\expandafter\gdef\csname num@d:es:novel_1:taylor-random:P:worst\endcsname{5.5}
\expandafter\gdef\csname num@d:es:novel_1:bn_gamma-random:PMP\endcsname{\ensuremath{-}2.8}
\expandafter\gdef\csname num@d:es:novel_1:bn_gamma-random:PMP:lo\endcsname{\ensuremath{-}5.7}
\expandafter\gdef\csname num@d:es:novel_1:bn_gamma-random:PMP:hi\endcsname{+0.1}
\expandafter\gdef\csname num@d:es:novel_1:bn_gamma-random:PMP:abs\endcsname{2.8}
\expandafter\gdef\csname num@d:es:novel_1:bn_gamma-random:PMP:wins\endcsname{1}
\expandafter\gdef\csname num@d:es:novel_1:bn_gamma-random:PMP:n\endcsname{5}
\expandafter\gdef\csname num@d:es:novel_1:bn_gamma-random:PMP:worst\endcsname{5.7}
\expandafter\gdef\csname num@d:es:novel_1:D-random:PMP\endcsname{\ensuremath{-}1.5}
\expandafter\gdef\csname num@d:es:novel_1:D-random:PMP:lo\endcsname{\ensuremath{-}3.5}
\expandafter\gdef\csname num@d:es:novel_1:D-random:PMP:hi\endcsname{+0.6}
\expandafter\gdef\csname num@d:es:novel_1:D-random:PMP:abs\endcsname{1.5}
\expandafter\gdef\csname num@d:es:novel_1:D-random:PMP:wins\endcsname{1}
\expandafter\gdef\csname num@d:es:novel_1:D-random:PMP:n\endcsname{5}
\expandafter\gdef\csname num@d:es:novel_1:D-random:PMP:worst\endcsname{3.5}
\expandafter\gdef\csname num@d:es:novel_1:D-random:P\endcsname{+1.3}
\expandafter\gdef\csname num@d:es:novel_1:D-random:P:lo\endcsname{\ensuremath{-}1.6}
\expandafter\gdef\csname num@d:es:novel_1:D-random:P:hi\endcsname{+4.3}
\expandafter\gdef\csname num@d:es:novel_1:D-random:P:abs\endcsname{1.3}
\expandafter\gdef\csname num@d:es:novel_1:D-random:P:wins\endcsname{4}
\expandafter\gdef\csname num@d:es:novel_1:D-random:P:n\endcsname{5}
\expandafter\gdef\csname num@d:es:novel_1:D-random:P:worst\endcsname{4.3}
\expandafter\gdef\csname num@d:es:novel_1:G-random:PMP\endcsname{\ensuremath{-}1.6}
\expandafter\gdef\csname num@d:es:novel_1:G-random:PMP:lo\endcsname{\ensuremath{-}3.6}
\expandafter\gdef\csname num@d:es:novel_1:G-random:PMP:hi\endcsname{+0.3}
\expandafter\gdef\csname num@d:es:novel_1:G-random:PMP:abs\endcsname{1.6}
\expandafter\gdef\csname num@d:es:novel_1:G-random:PMP:wins\endcsname{1}
\expandafter\gdef\csname num@d:es:novel_1:G-random:PMP:n\endcsname{5}
\expandafter\gdef\csname num@d:es:novel_1:G-random:PMP:worst\endcsname{3.6}
\expandafter\gdef\csname num@d:es:novel_1:V-random:PMP\endcsname{\ensuremath{-}1.3}
\expandafter\gdef\csname num@d:es:novel_1:V-random:PMP:lo\endcsname{\ensuremath{-}2.8}
\expandafter\gdef\csname num@d:es:novel_1:V-random:PMP:hi\endcsname{+0.1}
\expandafter\gdef\csname num@d:es:novel_1:V-random:PMP:abs\endcsname{1.3}
\expandafter\gdef\csname num@d:es:novel_1:V-random:PMP:wins\endcsname{1}
\expandafter\gdef\csname num@d:es:novel_1:V-random:PMP:n\endcsname{5}
\expandafter\gdef\csname num@d:es:novel_1:V-random:PMP:worst\endcsname{2.8}
\expandafter\gdef\csname num@d:es:novel_1:PM-P\endcsname{\ensuremath{-}0.2}
\expandafter\gdef\csname num@d:es:novel_1:PM-P:lo\endcsname{\ensuremath{-}2.4}
\expandafter\gdef\csname num@d:es:novel_1:PM-P:hi\endcsname{+1.9}
\expandafter\gdef\csname num@d:es:novel_1:PM-P:abs\endcsname{0.2}
\expandafter\gdef\csname num@d:es:novel_1:PM-P:wins\endcsname{3}
\expandafter\gdef\csname num@d:es:novel_1:PM-P:n\endcsname{5}
\expandafter\gdef\csname num@d:es:novel_1:PM-P:worst\endcsname{2.4}
\expandafter\gdef\csname num@d:es:novel_1:MP-P\endcsname{+0.3}
\expandafter\gdef\csname num@d:es:novel_1:MP-P:lo\endcsname{\ensuremath{-}1.0}
\expandafter\gdef\csname num@d:es:novel_1:MP-P:hi\endcsname{+1.7}
\expandafter\gdef\csname num@d:es:novel_1:MP-P:abs\endcsname{0.3}
\expandafter\gdef\csname num@d:es:novel_1:MP-P:wins\endcsname{2}
\expandafter\gdef\csname num@d:es:novel_1:MP-P:n\endcsname{5}
\expandafter\gdef\csname num@d:es:novel_1:MP-P:worst\endcsname{1.7}
\expandafter\gdef\csname num@d:es:novel_1:PGP-P\endcsname{+0.5}
\expandafter\gdef\csname num@d:es:novel_1:PGP-P:lo\endcsname{\ensuremath{-}1.6}
\expandafter\gdef\csname num@d:es:novel_1:PGP-P:hi\endcsname{+2.6}
\expandafter\gdef\csname num@d:es:novel_1:PGP-P:abs\endcsname{0.5}
\expandafter\gdef\csname num@d:es:novel_1:PGP-P:wins\endcsname{3}
\expandafter\gdef\csname num@d:es:novel_1:PGP-P:n\endcsname{5}
\expandafter\gdef\csname num@d:es:novel_1:PGP-P:worst\endcsname{2.6}
\expandafter\gdef\csname num@d:es:novel_1:PMP-P\endcsname{+0.6}
\expandafter\gdef\csname num@d:es:novel_1:PMP-P:lo\endcsname{\ensuremath{-}1.3}
\expandafter\gdef\csname num@d:es:novel_1:PMP-P:hi\endcsname{+2.4}
\expandafter\gdef\csname num@d:es:novel_1:PMP-P:abs\endcsname{0.6}
\expandafter\gdef\csname num@d:es:novel_1:PMP-P:wins\endcsname{4}
\expandafter\gdef\csname num@d:es:novel_1:PMP-P:n\endcsname{5}
\expandafter\gdef\csname num@d:es:novel_1:PMP-P:worst\endcsname{2.4}
\expandafter\gdef\csname num@d:es:novel_1:densem-dense\endcsname{\ensuremath{-}0.5}
\expandafter\gdef\csname num@d:es:novel_1:densem-dense:lo\endcsname{\ensuremath{-}2.5}
\expandafter\gdef\csname num@d:es:novel_1:densem-dense:hi\endcsname{+1.4}
\expandafter\gdef\csname num@d:es:novel_1:densem-dense:abs\endcsname{0.5}
\expandafter\gdef\csname num@d:es:novel_1:densem-dense:wins\endcsname{3}
\expandafter\gdef\csname num@d:es:novel_1:densem-dense:n\endcsname{5}
\expandafter\gdef\csname num@d:es:novel_1:densem-dense:worst\endcsname{2.5}
\expandafter\gdef\csname num@d:es:novel_1:densermr-dense\endcsname{\ensuremath{-}0.5}
\expandafter\gdef\csname num@d:es:novel_1:densermr-dense:lo\endcsname{\ensuremath{-}1.6}
\expandafter\gdef\csname num@d:es:novel_1:densermr-dense:hi\endcsname{+0.6}
\expandafter\gdef\csname num@d:es:novel_1:densermr-dense:abs\endcsname{0.5}
\expandafter\gdef\csname num@d:es:novel_1:densermr-dense:wins\endcsname{1}
\expandafter\gdef\csname num@d:es:novel_1:densermr-dense:n\endcsname{5}
\expandafter\gdef\csname num@d:es:novel_1:densermr-dense:worst\endcsname{1.6}
\expandafter\gdef\csname num@d:es:novel_1:dense-dacisP\endcsname{+4.0}
\expandafter\gdef\csname num@d:es:novel_1:dense-dacisP:lo\endcsname{+2.2}
\expandafter\gdef\csname num@d:es:novel_1:dense-dacisP:hi\endcsname{+5.7}
\expandafter\gdef\csname num@d:es:novel_1:dense-dacisP:abs\endcsname{4.0}
\expandafter\gdef\csname num@d:es:novel_1:dense-dacisP:wins\endcsname{5}
\expandafter\gdef\csname num@d:es:novel_1:dense-dacisP:n\endcsname{5}
\expandafter\gdef\csname num@d:es:novel_1:dense-dacisP:worst\endcsname{5.7}
\expandafter\gdef\csname num@d:es:novel_1:dense-dacisPMP\endcsname{+3.4}
\expandafter\gdef\csname num@d:es:novel_1:dense-dacisPMP:lo\endcsname{+2.1}
\expandafter\gdef\csname num@d:es:novel_1:dense-dacisPMP:hi\endcsname{+4.8}
\expandafter\gdef\csname num@d:es:novel_1:dense-dacisPMP:abs\endcsname{3.4}
\expandafter\gdef\csname num@d:es:novel_1:dense-dacisPMP:wins\endcsname{5}
\expandafter\gdef\csname num@d:es:novel_1:dense-dacisPMP:n\endcsname{5}
\expandafter\gdef\csname num@d:es:novel_1:dense-dacisPMP:worst\endcsname{4.8}
\expandafter\gdef\csname num@d:es:novel_1:dense-randomP\endcsname{+5.4}
\expandafter\gdef\csname num@d:es:novel_1:dense-randomP:lo\endcsname{+3.7}
\expandafter\gdef\csname num@d:es:novel_1:dense-randomP:hi\endcsname{+7.1}
\expandafter\gdef\csname num@d:es:novel_1:dense-randomP:abs\endcsname{5.4}
\expandafter\gdef\csname num@d:es:novel_1:dense-randomP:wins\endcsname{5}
\expandafter\gdef\csname num@d:es:novel_1:dense-randomP:n\endcsname{5}
\expandafter\gdef\csname num@d:es:novel_1:dense-randomP:worst\endcsname{7.1}
\expandafter\gdef\csname num@d:es:novel_1:dense-randomPMP\endcsname{+2.6}
\expandafter\gdef\csname num@d:es:novel_1:dense-randomPMP:lo\endcsname{+0.7}
\expandafter\gdef\csname num@d:es:novel_1:dense-randomPMP:hi\endcsname{+4.6}
\expandafter\gdef\csname num@d:es:novel_1:dense-randomPMP:abs\endcsname{2.6}
\expandafter\gdef\csname num@d:es:novel_1:dense-randomPMP:wins\endcsname{5}
\expandafter\gdef\csname num@d:es:novel_1:dense-randomPMP:n\endcsname{5}
\expandafter\gdef\csname num@d:es:novel_1:dense-randomPMP:worst\endcsname{4.6}
\expandafter\gdef\csname num@d:es:novel_1:mnv3-dacisPMP\endcsname{+11.7}
\expandafter\gdef\csname num@d:es:novel_1:mnv3-dacisPMP:lo\endcsname{+6.0}
\expandafter\gdef\csname num@d:es:novel_1:mnv3-dacisPMP:hi\endcsname{+17.5}
\expandafter\gdef\csname num@d:es:novel_1:mnv3-dacisPMP:abs\endcsname{11.7}
\expandafter\gdef\csname num@d:es:novel_1:mnv3-dacisPMP:wins\endcsname{5}
\expandafter\gdef\csname num@d:es:novel_1:mnv3-dacisPMP:n\endcsname{5}
\expandafter\gdef\csname num@d:es:novel_1:mnv3-dacisPMP:worst\endcsname{17.5}
\expandafter\gdef\csname num@d:es:novel_1:effb0-dacisPMP\endcsname{+8.4}
\expandafter\gdef\csname num@d:es:novel_1:effb0-dacisPMP:lo\endcsname{+3.2}
\expandafter\gdef\csname num@d:es:novel_1:effb0-dacisPMP:hi\endcsname{+13.6}
\expandafter\gdef\csname num@d:es:novel_1:effb0-dacisPMP:abs\endcsname{8.4}
\expandafter\gdef\csname num@d:es:novel_1:effb0-dacisPMP:wins\endcsname{5}
\expandafter\gdef\csname num@d:es:novel_1:effb0-dacisPMP:n\endcsname{5}
\expandafter\gdef\csname num@d:es:novel_1:effb0-dacisPMP:worst\endcsname{13.6}
\expandafter\gdef\csname num@d:es:novel_1:mnv2-dacisPMP\endcsname{+7.6}
\expandafter\gdef\csname num@d:es:novel_1:mnv2-dacisPMP:lo\endcsname{+4.3}
\expandafter\gdef\csname num@d:es:novel_1:mnv2-dacisPMP:hi\endcsname{+10.9}
\expandafter\gdef\csname num@d:es:novel_1:mnv2-dacisPMP:abs\endcsname{7.6}
\expandafter\gdef\csname num@d:es:novel_1:mnv2-dacisPMP:wins\endcsname{5}
\expandafter\gdef\csname num@d:es:novel_1:mnv2-dacisPMP:n\endcsname{5}
\expandafter\gdef\csname num@d:es:novel_1:mnv2-dacisPMP:worst\endcsname{10.9}
\expandafter\gdef\csname num@d:es:novel_1:dacis-l1:PMP50\endcsname{\ensuremath{-}0.9}
\expandafter\gdef\csname num@d:es:novel_1:dacis-l1:PMP50:lo\endcsname{\ensuremath{-}3.8}
\expandafter\gdef\csname num@d:es:novel_1:dacis-l1:PMP50:hi\endcsname{+2.1}
\expandafter\gdef\csname num@d:es:novel_1:dacis-l1:PMP50:abs\endcsname{0.9}
\expandafter\gdef\csname num@d:es:novel_1:dacis-l1:PMP50:wins\endcsname{2}
\expandafter\gdef\csname num@d:es:novel_1:dacis-l1:PMP50:n\endcsname{5}
\expandafter\gdef\csname num@d:es:novel_1:dacis-l1:PMP50:worst\endcsname{3.8}
\expandafter\gdef\csname num@d:es:novel_1:dacis-random:PMP50\endcsname{+0.0}
\expandafter\gdef\csname num@d:es:novel_1:dacis-random:PMP50:lo\endcsname{\ensuremath{-}1.4}
\expandafter\gdef\csname num@d:es:novel_1:dacis-random:PMP50:hi\endcsname{+1.4}
\expandafter\gdef\csname num@d:es:novel_1:dacis-random:PMP50:abs\endcsname{0.0}
\expandafter\gdef\csname num@d:es:novel_1:dacis-random:PMP50:wins\endcsname{2}
\expandafter\gdef\csname num@d:es:novel_1:dacis-random:PMP50:n\endcsname{5}
\expandafter\gdef\csname num@d:es:novel_1:dacis-random:PMP50:worst\endcsname{1.4}
\expandafter\gdef\csname num@d:es:novel_1:dense-dacisPMP50\endcsname{+2.4}
\expandafter\gdef\csname num@d:es:novel_1:dense-dacisPMP50:lo\endcsname{+1.2}
\expandafter\gdef\csname num@d:es:novel_1:dense-dacisPMP50:hi\endcsname{+3.7}
\expandafter\gdef\csname num@d:es:novel_1:dense-dacisPMP50:abs\endcsname{2.4}
\expandafter\gdef\csname num@d:es:novel_1:dense-dacisPMP50:wins\endcsname{5}
\expandafter\gdef\csname num@d:es:novel_1:dense-dacisPMP50:n\endcsname{5}
\expandafter\gdef\csname num@d:es:novel_1:dense-dacisPMP50:worst\endcsname{3.7}
\expandafter\gdef\csname num@d:es:novel_1:dacis-l1:PMP20\endcsname{+0.7}
\expandafter\gdef\csname num@d:es:novel_1:dacis-l1:PMP20:lo\endcsname{\ensuremath{-}2.4}
\expandafter\gdef\csname num@d:es:novel_1:dacis-l1:PMP20:hi\endcsname{+3.8}
\expandafter\gdef\csname num@d:es:novel_1:dacis-l1:PMP20:abs\endcsname{0.7}
\expandafter\gdef\csname num@d:es:novel_1:dacis-l1:PMP20:wins\endcsname{3}
\expandafter\gdef\csname num@d:es:novel_1:dacis-l1:PMP20:n\endcsname{5}
\expandafter\gdef\csname num@d:es:novel_1:dacis-l1:PMP20:worst\endcsname{3.8}
\expandafter\gdef\csname num@d:es:novel_1:dacis-random:PMP20\endcsname{+0.8}
\expandafter\gdef\csname num@d:es:novel_1:dacis-random:PMP20:lo\endcsname{\ensuremath{-}1.8}
\expandafter\gdef\csname num@d:es:novel_1:dacis-random:PMP20:hi\endcsname{+3.4}
\expandafter\gdef\csname num@d:es:novel_1:dacis-random:PMP20:abs\endcsname{0.8}
\expandafter\gdef\csname num@d:es:novel_1:dacis-random:PMP20:wins\endcsname{2}
\expandafter\gdef\csname num@d:es:novel_1:dacis-random:PMP20:n\endcsname{5}
\expandafter\gdef\csname num@d:es:novel_1:dacis-random:PMP20:worst\endcsname{3.4}
\expandafter\gdef\csname num@d:es:novel_1:dense-dacisPMP20\endcsname{+4.7}
\expandafter\gdef\csname num@d:es:novel_1:dense-dacisPMP20:lo\endcsname{+2.0}
\expandafter\gdef\csname num@d:es:novel_1:dense-dacisPMP20:hi\endcsname{+7.4}
\expandafter\gdef\csname num@d:es:novel_1:dense-dacisPMP20:abs\endcsname{4.7}
\expandafter\gdef\csname num@d:es:novel_1:dense-dacisPMP20:wins\endcsname{5}
\expandafter\gdef\csname num@d:es:novel_1:dense-dacisPMP20:n\endcsname{5}
\expandafter\gdef\csname num@d:es:novel_1:dense-dacisPMP20:worst\endcsname{7.4}
\expandafter\gdef\csname num@d:es:novel_1:dacis-l1:PMP10\endcsname{\ensuremath{-}1.0}
\expandafter\gdef\csname num@d:es:novel_1:dacis-l1:PMP10:lo\endcsname{\ensuremath{-}4.5}
\expandafter\gdef\csname num@d:es:novel_1:dacis-l1:PMP10:hi\endcsname{+2.5}
\expandafter\gdef\csname num@d:es:novel_1:dacis-l1:PMP10:abs\endcsname{1.0}
\expandafter\gdef\csname num@d:es:novel_1:dacis-l1:PMP10:wins\endcsname{1}
\expandafter\gdef\csname num@d:es:novel_1:dacis-l1:PMP10:n\endcsname{5}
\expandafter\gdef\csname num@d:es:novel_1:dacis-l1:PMP10:worst\endcsname{4.5}
\expandafter\gdef\csname num@d:es:novel_1:dacis-random:PMP10\endcsname{+1.6}
\expandafter\gdef\csname num@d:es:novel_1:dacis-random:PMP10:lo\endcsname{\ensuremath{-}1.1}
\expandafter\gdef\csname num@d:es:novel_1:dacis-random:PMP10:hi\endcsname{+4.3}
\expandafter\gdef\csname num@d:es:novel_1:dacis-random:PMP10:abs\endcsname{1.6}
\expandafter\gdef\csname num@d:es:novel_1:dacis-random:PMP10:wins\endcsname{3}
\expandafter\gdef\csname num@d:es:novel_1:dacis-random:PMP10:n\endcsname{5}
\expandafter\gdef\csname num@d:es:novel_1:dacis-random:PMP10:worst\endcsname{4.3}
\expandafter\gdef\csname num@d:es:novel_1:dense-dacisPMP10\endcsname{+7.7}
\expandafter\gdef\csname num@d:es:novel_1:dense-dacisPMP10:lo\endcsname{+4.0}
\expandafter\gdef\csname num@d:es:novel_1:dense-dacisPMP10:hi\endcsname{+11.3}
\expandafter\gdef\csname num@d:es:novel_1:dense-dacisPMP10:abs\endcsname{7.7}
\expandafter\gdef\csname num@d:es:novel_1:dense-dacisPMP10:wins\endcsname{5}
\expandafter\gdef\csname num@d:es:novel_1:dense-dacisPMP10:n\endcsname{5}
\expandafter\gdef\csname num@d:es:novel_1:dense-dacisPMP10:worst\endcsname{11.3}
\expandafter\gdef\csname num@d:rs:novel_1:dacis-random:PMP\endcsname{\ensuremath{-}1.5}
\expandafter\gdef\csname num@d:rs:novel_1:dacis-random:PMP:lo\endcsname{\ensuremath{-}2.6}
\expandafter\gdef\csname num@d:rs:novel_1:dacis-random:PMP:hi\endcsname{\ensuremath{-}0.5}
\expandafter\gdef\csname num@d:rs:novel_1:dacis-random:PMP:abs\endcsname{1.5}
\expandafter\gdef\csname num@d:rs:novel_1:dacis-random:PMP:wins\endcsname{0}
\expandafter\gdef\csname num@d:rs:novel_1:dacis-random:PMP:n\endcsname{5}
\expandafter\gdef\csname num@d:rs:novel_1:dacis-random:PMP:worst\endcsname{2.6}
\expandafter\gdef\csname num@d:rs:novel_1:dacis-random:P\endcsname{\ensuremath{-}0.6}
\expandafter\gdef\csname num@d:rs:novel_1:dacis-random:P:lo\endcsname{\ensuremath{-}2.4}
\expandafter\gdef\csname num@d:rs:novel_1:dacis-random:P:hi\endcsname{+1.2}
\expandafter\gdef\csname num@d:rs:novel_1:dacis-random:P:abs\endcsname{0.6}
\expandafter\gdef\csname num@d:rs:novel_1:dacis-random:P:wins\endcsname{1}
\expandafter\gdef\csname num@d:rs:novel_1:dacis-random:P:n\endcsname{5}
\expandafter\gdef\csname num@d:rs:novel_1:dacis-random:P:worst\endcsname{2.4}
\expandafter\gdef\csname num@d:rs:novel_1:dacis-l1:PMP\endcsname{\ensuremath{-}2.2}
\expandafter\gdef\csname num@d:rs:novel_1:dacis-l1:PMP:lo\endcsname{\ensuremath{-}2.7}
\expandafter\gdef\csname num@d:rs:novel_1:dacis-l1:PMP:hi\endcsname{\ensuremath{-}1.7}
\expandafter\gdef\csname num@d:rs:novel_1:dacis-l1:PMP:abs\endcsname{2.2}
\expandafter\gdef\csname num@d:rs:novel_1:dacis-l1:PMP:wins\endcsname{0}
\expandafter\gdef\csname num@d:rs:novel_1:dacis-l1:PMP:n\endcsname{5}
\expandafter\gdef\csname num@d:rs:novel_1:dacis-l1:PMP:worst\endcsname{2.7}
\expandafter\gdef\csname num@d:rs:novel_1:dacis-l1:P\endcsname{\ensuremath{-}2.0}
\expandafter\gdef\csname num@d:rs:novel_1:dacis-l1:P:lo\endcsname{\ensuremath{-}3.4}
\expandafter\gdef\csname num@d:rs:novel_1:dacis-l1:P:hi\endcsname{\ensuremath{-}0.7}
\expandafter\gdef\csname num@d:rs:novel_1:dacis-l1:P:abs\endcsname{2.0}
\expandafter\gdef\csname num@d:rs:novel_1:dacis-l1:P:wins\endcsname{0}
\expandafter\gdef\csname num@d:rs:novel_1:dacis-l1:P:n\endcsname{5}
\expandafter\gdef\csname num@d:rs:novel_1:dacis-l1:P:worst\endcsname{3.4}
\expandafter\gdef\csname num@d:rs:novel_1:dacis-taylor:PMP\endcsname{\ensuremath{-}1.6}
\expandafter\gdef\csname num@d:rs:novel_1:dacis-taylor:PMP:lo\endcsname{\ensuremath{-}2.2}
\expandafter\gdef\csname num@d:rs:novel_1:dacis-taylor:PMP:hi\endcsname{\ensuremath{-}1.1}
\expandafter\gdef\csname num@d:rs:novel_1:dacis-taylor:PMP:abs\endcsname{1.6}
\expandafter\gdef\csname num@d:rs:novel_1:dacis-taylor:PMP:wins\endcsname{0}
\expandafter\gdef\csname num@d:rs:novel_1:dacis-taylor:PMP:n\endcsname{5}
\expandafter\gdef\csname num@d:rs:novel_1:dacis-taylor:PMP:worst\endcsname{2.2}
\expandafter\gdef\csname num@d:rs:novel_1:dacis-taylor:P\endcsname{\ensuremath{-}2.3}
\expandafter\gdef\csname num@d:rs:novel_1:dacis-taylor:P:lo\endcsname{\ensuremath{-}3.9}
\expandafter\gdef\csname num@d:rs:novel_1:dacis-taylor:P:hi\endcsname{\ensuremath{-}0.6}
\expandafter\gdef\csname num@d:rs:novel_1:dacis-taylor:P:abs\endcsname{2.3}
\expandafter\gdef\csname num@d:rs:novel_1:dacis-taylor:P:wins\endcsname{0}
\expandafter\gdef\csname num@d:rs:novel_1:dacis-taylor:P:n\endcsname{5}
\expandafter\gdef\csname num@d:rs:novel_1:dacis-taylor:P:worst\endcsname{3.9}
\expandafter\gdef\csname num@d:rs:novel_1:dacis-bn_gamma:PMP\endcsname{\ensuremath{-}0.1}
\expandafter\gdef\csname num@d:rs:novel_1:dacis-bn_gamma:PMP:lo\endcsname{\ensuremath{-}0.8}
\expandafter\gdef\csname num@d:rs:novel_1:dacis-bn_gamma:PMP:hi\endcsname{+0.5}
\expandafter\gdef\csname num@d:rs:novel_1:dacis-bn_gamma:PMP:abs\endcsname{0.1}
\expandafter\gdef\csname num@d:rs:novel_1:dacis-bn_gamma:PMP:wins\endcsname{2}
\expandafter\gdef\csname num@d:rs:novel_1:dacis-bn_gamma:PMP:n\endcsname{5}
\expandafter\gdef\csname num@d:rs:novel_1:dacis-bn_gamma:PMP:worst\endcsname{0.8}
\expandafter\gdef\csname num@d:rs:novel_1:dacis-D:PMP\endcsname{+0.7}
\expandafter\gdef\csname num@d:rs:novel_1:dacis-D:PMP:lo\endcsname{\ensuremath{-}0.5}
\expandafter\gdef\csname num@d:rs:novel_1:dacis-D:PMP:hi\endcsname{+1.8}
\expandafter\gdef\csname num@d:rs:novel_1:dacis-D:PMP:abs\endcsname{0.7}
\expandafter\gdef\csname num@d:rs:novel_1:dacis-D:PMP:wins\endcsname{4}
\expandafter\gdef\csname num@d:rs:novel_1:dacis-D:PMP:n\endcsname{5}
\expandafter\gdef\csname num@d:rs:novel_1:dacis-D:PMP:worst\endcsname{1.8}
\expandafter\gdef\csname num@d:rs:novel_1:dacis-D:P\endcsname{+0.1}
\expandafter\gdef\csname num@d:rs:novel_1:dacis-D:P:lo\endcsname{\ensuremath{-}1.0}
\expandafter\gdef\csname num@d:rs:novel_1:dacis-D:P:hi\endcsname{+1.3}
\expandafter\gdef\csname num@d:rs:novel_1:dacis-D:P:abs\endcsname{0.1}
\expandafter\gdef\csname num@d:rs:novel_1:dacis-D:P:wins\endcsname{2}
\expandafter\gdef\csname num@d:rs:novel_1:dacis-D:P:n\endcsname{5}
\expandafter\gdef\csname num@d:rs:novel_1:dacis-D:P:worst\endcsname{1.3}
\expandafter\gdef\csname num@d:rs:novel_1:dacis-G:PMP\endcsname{\ensuremath{-}0.3}
\expandafter\gdef\csname num@d:rs:novel_1:dacis-G:PMP:lo\endcsname{\ensuremath{-}1.6}
\expandafter\gdef\csname num@d:rs:novel_1:dacis-G:PMP:hi\endcsname{+1.0}
\expandafter\gdef\csname num@d:rs:novel_1:dacis-G:PMP:abs\endcsname{0.3}
\expandafter\gdef\csname num@d:rs:novel_1:dacis-G:PMP:wins\endcsname{3}
\expandafter\gdef\csname num@d:rs:novel_1:dacis-G:PMP:n\endcsname{5}
\expandafter\gdef\csname num@d:rs:novel_1:dacis-G:PMP:worst\endcsname{1.6}
\expandafter\gdef\csname num@d:rs:novel_1:dacis-V:PMP\endcsname{\ensuremath{-}0.6}
\expandafter\gdef\csname num@d:rs:novel_1:dacis-V:PMP:lo\endcsname{\ensuremath{-}1.5}
\expandafter\gdef\csname num@d:rs:novel_1:dacis-V:PMP:hi\endcsname{+0.3}
\expandafter\gdef\csname num@d:rs:novel_1:dacis-V:PMP:abs\endcsname{0.6}
\expandafter\gdef\csname num@d:rs:novel_1:dacis-V:PMP:wins\endcsname{1}
\expandafter\gdef\csname num@d:rs:novel_1:dacis-V:PMP:n\endcsname{5}
\expandafter\gdef\csname num@d:rs:novel_1:dacis-V:PMP:worst\endcsname{1.5}
\expandafter\gdef\csname num@d:rs:novel_1:l1-random:PMP\endcsname{+0.6}
\expandafter\gdef\csname num@d:rs:novel_1:l1-random:PMP:lo\endcsname{\ensuremath{-}0.3}
\expandafter\gdef\csname num@d:rs:novel_1:l1-random:PMP:hi\endcsname{+1.6}
\expandafter\gdef\csname num@d:rs:novel_1:l1-random:PMP:abs\endcsname{0.6}
\expandafter\gdef\csname num@d:rs:novel_1:l1-random:PMP:wins\endcsname{4}
\expandafter\gdef\csname num@d:rs:novel_1:l1-random:PMP:n\endcsname{5}
\expandafter\gdef\csname num@d:rs:novel_1:l1-random:PMP:worst\endcsname{1.6}
\expandafter\gdef\csname num@d:rs:novel_1:l1-random:P\endcsname{+1.5}
\expandafter\gdef\csname num@d:rs:novel_1:l1-random:P:lo\endcsname{+0.1}
\expandafter\gdef\csname num@d:rs:novel_1:l1-random:P:hi\endcsname{+2.8}
\expandafter\gdef\csname num@d:rs:novel_1:l1-random:P:abs\endcsname{1.5}
\expandafter\gdef\csname num@d:rs:novel_1:l1-random:P:wins\endcsname{5}
\expandafter\gdef\csname num@d:rs:novel_1:l1-random:P:n\endcsname{5}
\expandafter\gdef\csname num@d:rs:novel_1:l1-random:P:worst\endcsname{2.8}
\expandafter\gdef\csname num@d:rs:novel_1:taylor-random:PMP\endcsname{+0.1}
\expandafter\gdef\csname num@d:rs:novel_1:taylor-random:PMP:lo\endcsname{\ensuremath{-}0.8}
\expandafter\gdef\csname num@d:rs:novel_1:taylor-random:PMP:hi\endcsname{+1.1}
\expandafter\gdef\csname num@d:rs:novel_1:taylor-random:PMP:abs\endcsname{0.1}
\expandafter\gdef\csname num@d:rs:novel_1:taylor-random:PMP:wins\endcsname{4}
\expandafter\gdef\csname num@d:rs:novel_1:taylor-random:PMP:n\endcsname{5}
\expandafter\gdef\csname num@d:rs:novel_1:taylor-random:PMP:worst\endcsname{1.1}
\expandafter\gdef\csname num@d:rs:novel_1:taylor-random:P\endcsname{+1.7}
\expandafter\gdef\csname num@d:rs:novel_1:taylor-random:P:lo\endcsname{+0.8}
\expandafter\gdef\csname num@d:rs:novel_1:taylor-random:P:hi\endcsname{+2.7}
\expandafter\gdef\csname num@d:rs:novel_1:taylor-random:P:abs\endcsname{1.7}
\expandafter\gdef\csname num@d:rs:novel_1:taylor-random:P:wins\endcsname{5}
\expandafter\gdef\csname num@d:rs:novel_1:taylor-random:P:n\endcsname{5}
\expandafter\gdef\csname num@d:rs:novel_1:taylor-random:P:worst\endcsname{2.7}
\expandafter\gdef\csname num@d:rs:novel_1:bn_gamma-random:PMP\endcsname{\ensuremath{-}1.4}
\expandafter\gdef\csname num@d:rs:novel_1:bn_gamma-random:PMP:lo\endcsname{\ensuremath{-}2.6}
\expandafter\gdef\csname num@d:rs:novel_1:bn_gamma-random:PMP:hi\endcsname{\ensuremath{-}0.2}
\expandafter\gdef\csname num@d:rs:novel_1:bn_gamma-random:PMP:abs\endcsname{1.4}
\expandafter\gdef\csname num@d:rs:novel_1:bn_gamma-random:PMP:wins\endcsname{1}
\expandafter\gdef\csname num@d:rs:novel_1:bn_gamma-random:PMP:n\endcsname{5}
\expandafter\gdef\csname num@d:rs:novel_1:bn_gamma-random:PMP:worst\endcsname{2.6}
\expandafter\gdef\csname num@d:rs:novel_1:D-random:PMP\endcsname{\ensuremath{-}2.2}
\expandafter\gdef\csname num@d:rs:novel_1:D-random:PMP:lo\endcsname{\ensuremath{-}4.1}
\expandafter\gdef\csname num@d:rs:novel_1:D-random:PMP:hi\endcsname{\ensuremath{-}0.2}
\expandafter\gdef\csname num@d:rs:novel_1:D-random:PMP:abs\endcsname{2.2}
\expandafter\gdef\csname num@d:rs:novel_1:D-random:PMP:wins\endcsname{1}
\expandafter\gdef\csname num@d:rs:novel_1:D-random:PMP:n\endcsname{5}
\expandafter\gdef\csname num@d:rs:novel_1:D-random:PMP:worst\endcsname{4.1}
\expandafter\gdef\csname num@d:rs:novel_1:D-random:P\endcsname{\ensuremath{-}0.7}
\expandafter\gdef\csname num@d:rs:novel_1:D-random:P:lo\endcsname{\ensuremath{-}2.1}
\expandafter\gdef\csname num@d:rs:novel_1:D-random:P:hi\endcsname{+0.7}
\expandafter\gdef\csname num@d:rs:novel_1:D-random:P:abs\endcsname{0.7}
\expandafter\gdef\csname num@d:rs:novel_1:D-random:P:wins\endcsname{2}
\expandafter\gdef\csname num@d:rs:novel_1:D-random:P:n\endcsname{5}
\expandafter\gdef\csname num@d:rs:novel_1:D-random:P:worst\endcsname{2.1}
\expandafter\gdef\csname num@d:rs:novel_1:G-random:PMP\endcsname{\ensuremath{-}1.2}
\expandafter\gdef\csname num@d:rs:novel_1:G-random:PMP:lo\endcsname{\ensuremath{-}3.2}
\expandafter\gdef\csname num@d:rs:novel_1:G-random:PMP:hi\endcsname{+0.8}
\expandafter\gdef\csname num@d:rs:novel_1:G-random:PMP:abs\endcsname{1.2}
\expandafter\gdef\csname num@d:rs:novel_1:G-random:PMP:wins\endcsname{1}
\expandafter\gdef\csname num@d:rs:novel_1:G-random:PMP:n\endcsname{5}
\expandafter\gdef\csname num@d:rs:novel_1:G-random:PMP:worst\endcsname{3.2}
\expandafter\gdef\csname num@d:rs:novel_1:V-random:PMP\endcsname{\ensuremath{-}0.9}
\expandafter\gdef\csname num@d:rs:novel_1:V-random:PMP:lo\endcsname{\ensuremath{-}1.5}
\expandafter\gdef\csname num@d:rs:novel_1:V-random:PMP:hi\endcsname{\ensuremath{-}0.3}
\expandafter\gdef\csname num@d:rs:novel_1:V-random:PMP:abs\endcsname{0.9}
\expandafter\gdef\csname num@d:rs:novel_1:V-random:PMP:wins\endcsname{0}
\expandafter\gdef\csname num@d:rs:novel_1:V-random:PMP:n\endcsname{5}
\expandafter\gdef\csname num@d:rs:novel_1:V-random:PMP:worst\endcsname{1.5}
\expandafter\gdef\csname num@d:rs:novel_1:PM-P\endcsname{+1.2}
\expandafter\gdef\csname num@d:rs:novel_1:PM-P:lo\endcsname{+0.6}
\expandafter\gdef\csname num@d:rs:novel_1:PM-P:hi\endcsname{+1.8}
\expandafter\gdef\csname num@d:rs:novel_1:PM-P:abs\endcsname{1.2}
\expandafter\gdef\csname num@d:rs:novel_1:PM-P:wins\endcsname{5}
\expandafter\gdef\csname num@d:rs:novel_1:PM-P:n\endcsname{5}
\expandafter\gdef\csname num@d:rs:novel_1:PM-P:worst\endcsname{1.8}
\expandafter\gdef\csname num@d:rs:novel_1:MP-P\endcsname{+0.4}
\expandafter\gdef\csname num@d:rs:novel_1:MP-P:lo\endcsname{\ensuremath{-}0.7}
\expandafter\gdef\csname num@d:rs:novel_1:MP-P:hi\endcsname{+1.5}
\expandafter\gdef\csname num@d:rs:novel_1:MP-P:abs\endcsname{0.4}
\expandafter\gdef\csname num@d:rs:novel_1:MP-P:wins\endcsname{3}
\expandafter\gdef\csname num@d:rs:novel_1:MP-P:n\endcsname{5}
\expandafter\gdef\csname num@d:rs:novel_1:MP-P:worst\endcsname{1.5}
\expandafter\gdef\csname num@d:rs:novel_1:PGP-P\endcsname{+0.9}
\expandafter\gdef\csname num@d:rs:novel_1:PGP-P:lo\endcsname{\ensuremath{-}0.1}
\expandafter\gdef\csname num@d:rs:novel_1:PGP-P:hi\endcsname{+1.9}
\expandafter\gdef\csname num@d:rs:novel_1:PGP-P:abs\endcsname{0.9}
\expandafter\gdef\csname num@d:rs:novel_1:PGP-P:wins\endcsname{4}
\expandafter\gdef\csname num@d:rs:novel_1:PGP-P:n\endcsname{5}
\expandafter\gdef\csname num@d:rs:novel_1:PGP-P:worst\endcsname{1.9}
\expandafter\gdef\csname num@d:rs:novel_1:PMP-P\endcsname{+0.4}
\expandafter\gdef\csname num@d:rs:novel_1:PMP-P:lo\endcsname{\ensuremath{-}0.6}
\expandafter\gdef\csname num@d:rs:novel_1:PMP-P:hi\endcsname{+1.5}
\expandafter\gdef\csname num@d:rs:novel_1:PMP-P:abs\endcsname{0.4}
\expandafter\gdef\csname num@d:rs:novel_1:PMP-P:wins\endcsname{4}
\expandafter\gdef\csname num@d:rs:novel_1:PMP-P:n\endcsname{5}
\expandafter\gdef\csname num@d:rs:novel_1:PMP-P:worst\endcsname{1.5}
\expandafter\gdef\csname num@d:rs:novel_1:densem-dense\endcsname{+0.5}
\expandafter\gdef\csname num@d:rs:novel_1:densem-dense:lo\endcsname{\ensuremath{-}0.1}
\expandafter\gdef\csname num@d:rs:novel_1:densem-dense:hi\endcsname{+1.0}
\expandafter\gdef\csname num@d:rs:novel_1:densem-dense:abs\endcsname{0.5}
\expandafter\gdef\csname num@d:rs:novel_1:densem-dense:wins\endcsname{5}
\expandafter\gdef\csname num@d:rs:novel_1:densem-dense:n\endcsname{5}
\expandafter\gdef\csname num@d:rs:novel_1:densem-dense:worst\endcsname{1.0}
\expandafter\gdef\csname num@d:rs:novel_1:densermr-dense\endcsname{\ensuremath{-}0.3}
\expandafter\gdef\csname num@d:rs:novel_1:densermr-dense:lo\endcsname{\ensuremath{-}0.8}
\expandafter\gdef\csname num@d:rs:novel_1:densermr-dense:hi\endcsname{+0.1}
\expandafter\gdef\csname num@d:rs:novel_1:densermr-dense:abs\endcsname{0.3}
\expandafter\gdef\csname num@d:rs:novel_1:densermr-dense:wins\endcsname{2}
\expandafter\gdef\csname num@d:rs:novel_1:densermr-dense:n\endcsname{5}
\expandafter\gdef\csname num@d:rs:novel_1:densermr-dense:worst\endcsname{0.8}
\expandafter\gdef\csname num@d:rs:novel_1:dense-dacisP\endcsname{+8.3}
\expandafter\gdef\csname num@d:rs:novel_1:dense-dacisP:lo\endcsname{+6.6}
\expandafter\gdef\csname num@d:rs:novel_1:dense-dacisP:hi\endcsname{+10.0}
\expandafter\gdef\csname num@d:rs:novel_1:dense-dacisP:abs\endcsname{8.3}
\expandafter\gdef\csname num@d:rs:novel_1:dense-dacisP:wins\endcsname{5}
\expandafter\gdef\csname num@d:rs:novel_1:dense-dacisP:n\endcsname{5}
\expandafter\gdef\csname num@d:rs:novel_1:dense-dacisP:worst\endcsname{10.0}
\expandafter\gdef\csname num@d:rs:novel_1:dense-dacisPMP\endcsname{+7.9}
\expandafter\gdef\csname num@d:rs:novel_1:dense-dacisPMP:lo\endcsname{+6.5}
\expandafter\gdef\csname num@d:rs:novel_1:dense-dacisPMP:hi\endcsname{+9.2}
\expandafter\gdef\csname num@d:rs:novel_1:dense-dacisPMP:abs\endcsname{7.9}
\expandafter\gdef\csname num@d:rs:novel_1:dense-dacisPMP:wins\endcsname{5}
\expandafter\gdef\csname num@d:rs:novel_1:dense-dacisPMP:n\endcsname{5}
\expandafter\gdef\csname num@d:rs:novel_1:dense-dacisPMP:worst\endcsname{9.2}
\expandafter\gdef\csname num@d:rs:novel_1:dense-randomP\endcsname{+7.8}
\expandafter\gdef\csname num@d:rs:novel_1:dense-randomP:lo\endcsname{+5.6}
\expandafter\gdef\csname num@d:rs:novel_1:dense-randomP:hi\endcsname{+9.9}
\expandafter\gdef\csname num@d:rs:novel_1:dense-randomP:abs\endcsname{7.8}
\expandafter\gdef\csname num@d:rs:novel_1:dense-randomP:wins\endcsname{5}
\expandafter\gdef\csname num@d:rs:novel_1:dense-randomP:n\endcsname{5}
\expandafter\gdef\csname num@d:rs:novel_1:dense-randomP:worst\endcsname{9.9}
\expandafter\gdef\csname num@d:rs:novel_1:dense-randomPMP\endcsname{+6.3}
\expandafter\gdef\csname num@d:rs:novel_1:dense-randomPMP:lo\endcsname{+4.5}
\expandafter\gdef\csname num@d:rs:novel_1:dense-randomPMP:hi\endcsname{+8.2}
\expandafter\gdef\csname num@d:rs:novel_1:dense-randomPMP:abs\endcsname{6.3}
\expandafter\gdef\csname num@d:rs:novel_1:dense-randomPMP:wins\endcsname{5}
\expandafter\gdef\csname num@d:rs:novel_1:dense-randomPMP:n\endcsname{5}
\expandafter\gdef\csname num@d:rs:novel_1:dense-randomPMP:worst\endcsname{8.2}
\expandafter\gdef\csname num@d:rs:novel_1:mnv3-dacisPMP\endcsname{+10.4}
\expandafter\gdef\csname num@d:rs:novel_1:mnv3-dacisPMP:lo\endcsname{+9.1}
\expandafter\gdef\csname num@d:rs:novel_1:mnv3-dacisPMP:hi\endcsname{+11.7}
\expandafter\gdef\csname num@d:rs:novel_1:mnv3-dacisPMP:abs\endcsname{10.4}
\expandafter\gdef\csname num@d:rs:novel_1:mnv3-dacisPMP:wins\endcsname{5}
\expandafter\gdef\csname num@d:rs:novel_1:mnv3-dacisPMP:n\endcsname{5}
\expandafter\gdef\csname num@d:rs:novel_1:mnv3-dacisPMP:worst\endcsname{11.7}
\expandafter\gdef\csname num@d:rs:novel_1:effb0-dacisPMP\endcsname{+13.9}
\expandafter\gdef\csname num@d:rs:novel_1:effb0-dacisPMP:lo\endcsname{+11.7}
\expandafter\gdef\csname num@d:rs:novel_1:effb0-dacisPMP:hi\endcsname{+16.0}
\expandafter\gdef\csname num@d:rs:novel_1:effb0-dacisPMP:abs\endcsname{13.9}
\expandafter\gdef\csname num@d:rs:novel_1:effb0-dacisPMP:wins\endcsname{5}
\expandafter\gdef\csname num@d:rs:novel_1:effb0-dacisPMP:n\endcsname{5}
\expandafter\gdef\csname num@d:rs:novel_1:effb0-dacisPMP:worst\endcsname{16.0}
\expandafter\gdef\csname num@d:rs:novel_1:mnv2-dacisPMP\endcsname{+8.4}
\expandafter\gdef\csname num@d:rs:novel_1:mnv2-dacisPMP:lo\endcsname{+5.6}
\expandafter\gdef\csname num@d:rs:novel_1:mnv2-dacisPMP:hi\endcsname{+11.1}
\expandafter\gdef\csname num@d:rs:novel_1:mnv2-dacisPMP:abs\endcsname{8.4}
\expandafter\gdef\csname num@d:rs:novel_1:mnv2-dacisPMP:wins\endcsname{5}
\expandafter\gdef\csname num@d:rs:novel_1:mnv2-dacisPMP:n\endcsname{5}
\expandafter\gdef\csname num@d:rs:novel_1:mnv2-dacisPMP:worst\endcsname{11.1}
\expandafter\gdef\csname num@d:rs:novel_1:dacis-l1:PMP50\endcsname{\ensuremath{-}1.3}
\expandafter\gdef\csname num@d:rs:novel_1:dacis-l1:PMP50:lo\endcsname{\ensuremath{-}2.1}
\expandafter\gdef\csname num@d:rs:novel_1:dacis-l1:PMP50:hi\endcsname{\ensuremath{-}0.5}
\expandafter\gdef\csname num@d:rs:novel_1:dacis-l1:PMP50:abs\endcsname{1.3}
\expandafter\gdef\csname num@d:rs:novel_1:dacis-l1:PMP50:wins\endcsname{0}
\expandafter\gdef\csname num@d:rs:novel_1:dacis-l1:PMP50:n\endcsname{5}
\expandafter\gdef\csname num@d:rs:novel_1:dacis-l1:PMP50:worst\endcsname{2.1}
\expandafter\gdef\csname num@d:rs:novel_1:dacis-random:PMP50\endcsname{\ensuremath{-}0.4}
\expandafter\gdef\csname num@d:rs:novel_1:dacis-random:PMP50:lo\endcsname{\ensuremath{-}1.9}
\expandafter\gdef\csname num@d:rs:novel_1:dacis-random:PMP50:hi\endcsname{+1.1}
\expandafter\gdef\csname num@d:rs:novel_1:dacis-random:PMP50:abs\endcsname{0.4}
\expandafter\gdef\csname num@d:rs:novel_1:dacis-random:PMP50:wins\endcsname{2}
\expandafter\gdef\csname num@d:rs:novel_1:dacis-random:PMP50:n\endcsname{5}
\expandafter\gdef\csname num@d:rs:novel_1:dacis-random:PMP50:worst\endcsname{1.9}
\expandafter\gdef\csname num@d:rs:novel_1:dense-dacisPMP50\endcsname{+4.0}
\expandafter\gdef\csname num@d:rs:novel_1:dense-dacisPMP50:lo\endcsname{+3.2}
\expandafter\gdef\csname num@d:rs:novel_1:dense-dacisPMP50:hi\endcsname{+4.7}
\expandafter\gdef\csname num@d:rs:novel_1:dense-dacisPMP50:abs\endcsname{4.0}
\expandafter\gdef\csname num@d:rs:novel_1:dense-dacisPMP50:wins\endcsname{5}
\expandafter\gdef\csname num@d:rs:novel_1:dense-dacisPMP50:n\endcsname{5}
\expandafter\gdef\csname num@d:rs:novel_1:dense-dacisPMP50:worst\endcsname{4.7}
\expandafter\gdef\csname num@d:rs:novel_1:dacis-l1:PMP20\endcsname{\ensuremath{-}4.2}
\expandafter\gdef\csname num@d:rs:novel_1:dacis-l1:PMP20:lo\endcsname{\ensuremath{-}6.2}
\expandafter\gdef\csname num@d:rs:novel_1:dacis-l1:PMP20:hi\endcsname{\ensuremath{-}2.2}
\expandafter\gdef\csname num@d:rs:novel_1:dacis-l1:PMP20:abs\endcsname{4.2}
\expandafter\gdef\csname num@d:rs:novel_1:dacis-l1:PMP20:wins\endcsname{0}
\expandafter\gdef\csname num@d:rs:novel_1:dacis-l1:PMP20:n\endcsname{5}
\expandafter\gdef\csname num@d:rs:novel_1:dacis-l1:PMP20:worst\endcsname{6.2}
\expandafter\gdef\csname num@d:rs:novel_1:dacis-random:PMP20\endcsname{\ensuremath{-}0.6}
\expandafter\gdef\csname num@d:rs:novel_1:dacis-random:PMP20:lo\endcsname{\ensuremath{-}1.9}
\expandafter\gdef\csname num@d:rs:novel_1:dacis-random:PMP20:hi\endcsname{+0.7}
\expandafter\gdef\csname num@d:rs:novel_1:dacis-random:PMP20:abs\endcsname{0.6}
\expandafter\gdef\csname num@d:rs:novel_1:dacis-random:PMP20:wins\endcsname{2}
\expandafter\gdef\csname num@d:rs:novel_1:dacis-random:PMP20:n\endcsname{5}
\expandafter\gdef\csname num@d:rs:novel_1:dacis-random:PMP20:worst\endcsname{1.9}
\expandafter\gdef\csname num@d:rs:novel_1:dense-dacisPMP20\endcsname{+10.9}
\expandafter\gdef\csname num@d:rs:novel_1:dense-dacisPMP20:lo\endcsname{+7.5}
\expandafter\gdef\csname num@d:rs:novel_1:dense-dacisPMP20:hi\endcsname{+14.2}
\expandafter\gdef\csname num@d:rs:novel_1:dense-dacisPMP20:abs\endcsname{10.9}
\expandafter\gdef\csname num@d:rs:novel_1:dense-dacisPMP20:wins\endcsname{5}
\expandafter\gdef\csname num@d:rs:novel_1:dense-dacisPMP20:n\endcsname{5}
\expandafter\gdef\csname num@d:rs:novel_1:dense-dacisPMP20:worst\endcsname{14.2}
\expandafter\gdef\csname num@d:rs:novel_1:dacis-l1:PMP10\endcsname{\ensuremath{-}0.8}
\expandafter\gdef\csname num@d:rs:novel_1:dacis-l1:PMP10:lo\endcsname{\ensuremath{-}1.8}
\expandafter\gdef\csname num@d:rs:novel_1:dacis-l1:PMP10:hi\endcsname{+0.1}
\expandafter\gdef\csname num@d:rs:novel_1:dacis-l1:PMP10:abs\endcsname{0.8}
\expandafter\gdef\csname num@d:rs:novel_1:dacis-l1:PMP10:wins\endcsname{1}
\expandafter\gdef\csname num@d:rs:novel_1:dacis-l1:PMP10:n\endcsname{5}
\expandafter\gdef\csname num@d:rs:novel_1:dacis-l1:PMP10:worst\endcsname{1.8}
\expandafter\gdef\csname num@d:rs:novel_1:dacis-random:PMP10\endcsname{+0.9}
\expandafter\gdef\csname num@d:rs:novel_1:dacis-random:PMP10:lo\endcsname{\ensuremath{-}1.1}
\expandafter\gdef\csname num@d:rs:novel_1:dacis-random:PMP10:hi\endcsname{+2.9}
\expandafter\gdef\csname num@d:rs:novel_1:dacis-random:PMP10:abs\endcsname{0.9}
\expandafter\gdef\csname num@d:rs:novel_1:dacis-random:PMP10:wins\endcsname{4}
\expandafter\gdef\csname num@d:rs:novel_1:dacis-random:PMP10:n\endcsname{5}
\expandafter\gdef\csname num@d:rs:novel_1:dacis-random:PMP10:worst\endcsname{2.9}
\expandafter\gdef\csname num@d:rs:novel_1:dense-dacisPMP10\endcsname{+15.8}
\expandafter\gdef\csname num@d:rs:novel_1:dense-dacisPMP10:lo\endcsname{+14.0}
\expandafter\gdef\csname num@d:rs:novel_1:dense-dacisPMP10:hi\endcsname{+17.7}
\expandafter\gdef\csname num@d:rs:novel_1:dense-dacisPMP10:abs\endcsname{15.8}
\expandafter\gdef\csname num@d:rs:novel_1:dense-dacisPMP10:wins\endcsname{5}
\expandafter\gdef\csname num@d:rs:novel_1:dense-dacisPMP10:n\endcsname{5}
\expandafter\gdef\csname num@d:rs:novel_1:dense-dacisPMP10:worst\endcsname{17.7}
\expandafter\gdef\csname num@d:e3:novel_5:dacis-l1:PMP\endcsname{\ensuremath{-}0.3}
\expandafter\gdef\csname num@d:e3:novel_5:dacis-l1:PMP:lo\endcsname{\ensuremath{-}3.7}
\expandafter\gdef\csname num@d:e3:novel_5:dacis-l1:PMP:hi\endcsname{+3.1}
\expandafter\gdef\csname num@d:e3:novel_5:dacis-l1:PMP:abs\endcsname{0.3}
\expandafter\gdef\csname num@d:e3:novel_5:dacis-l1:PMP:wins\endcsname{2}
\expandafter\gdef\csname num@d:e3:novel_5:dacis-l1:PMP:n\endcsname{5}
\expandafter\gdef\csname num@d:e3:novel_5:dacis-l1:PMP:worst\endcsname{3.7}
\expandafter\gdef\csname num@d:e3:novel_5:densem-dense\endcsname{\ensuremath{-}0.4}
\expandafter\gdef\csname num@d:e3:novel_5:densem-dense:lo\endcsname{\ensuremath{-}1.6}
\expandafter\gdef\csname num@d:e3:novel_5:densem-dense:hi\endcsname{+0.8}
\expandafter\gdef\csname num@d:e3:novel_5:densem-dense:abs\endcsname{0.4}
\expandafter\gdef\csname num@d:e3:novel_5:densem-dense:wins\endcsname{1}
\expandafter\gdef\csname num@d:e3:novel_5:densem-dense:n\endcsname{5}
\expandafter\gdef\csname num@d:e3:novel_5:densem-dense:worst\endcsname{1.6}
\expandafter\gdef\csname num@d:e3:novel_5:densermr-dense\endcsname{\ensuremath{-}0.6}
\expandafter\gdef\csname num@d:e3:novel_5:densermr-dense:lo\endcsname{\ensuremath{-}1.5}
\expandafter\gdef\csname num@d:e3:novel_5:densermr-dense:hi\endcsname{+0.2}
\expandafter\gdef\csname num@d:e3:novel_5:densermr-dense:abs\endcsname{0.6}
\expandafter\gdef\csname num@d:e3:novel_5:densermr-dense:wins\endcsname{1}
\expandafter\gdef\csname num@d:e3:novel_5:densermr-dense:n\endcsname{5}
\expandafter\gdef\csname num@d:e3:novel_5:densermr-dense:worst\endcsname{1.5}
\expandafter\gdef\csname num@d:e3:novel_5:dense-dacisPMP\endcsname{+5.2}
\expandafter\gdef\csname num@d:e3:novel_5:dense-dacisPMP:lo\endcsname{+0.5}
\expandafter\gdef\csname num@d:e3:novel_5:dense-dacisPMP:hi\endcsname{+9.8}
\expandafter\gdef\csname num@d:e3:novel_5:dense-dacisPMP:abs\endcsname{5.2}
\expandafter\gdef\csname num@d:e3:novel_5:dense-dacisPMP:wins\endcsname{5}
\expandafter\gdef\csname num@d:e3:novel_5:dense-dacisPMP:n\endcsname{5}
\expandafter\gdef\csname num@d:e3:novel_5:dense-dacisPMP:worst\endcsname{9.8}
\expandafter\gdef\csname num@d:e3:novel_5:mnv3-dacisPMP\endcsname{+11.2}
\expandafter\gdef\csname num@d:e3:novel_5:mnv3-dacisPMP:lo\endcsname{+4.3}
\expandafter\gdef\csname num@d:e3:novel_5:mnv3-dacisPMP:hi\endcsname{+18.0}
\expandafter\gdef\csname num@d:e3:novel_5:mnv3-dacisPMP:abs\endcsname{11.2}
\expandafter\gdef\csname num@d:e3:novel_5:mnv3-dacisPMP:wins\endcsname{5}
\expandafter\gdef\csname num@d:e3:novel_5:mnv3-dacisPMP:n\endcsname{5}
\expandafter\gdef\csname num@d:e3:novel_5:mnv3-dacisPMP:worst\endcsname{18.0}
\expandafter\gdef\csname num@d:pv:novel_10:dacis-random:PMP\endcsname{\ensuremath{-}0.1}
\expandafter\gdef\csname num@d:pv:novel_10:dacis-random:PMP:lo\endcsname{\ensuremath{-}1.8}
\expandafter\gdef\csname num@d:pv:novel_10:dacis-random:PMP:hi\endcsname{+1.7}
\expandafter\gdef\csname num@d:pv:novel_10:dacis-random:PMP:abs\endcsname{0.1}
\expandafter\gdef\csname num@d:pv:novel_10:dacis-random:PMP:wins\endcsname{3}
\expandafter\gdef\csname num@d:pv:novel_10:dacis-random:PMP:n\endcsname{5}
\expandafter\gdef\csname num@d:pv:novel_10:dacis-random:PMP:worst\endcsname{1.8}
\expandafter\gdef\csname num@d:pv:novel_10:dacis-random:P\endcsname{+0.5}
\expandafter\gdef\csname num@d:pv:novel_10:dacis-random:P:lo\endcsname{\ensuremath{-}1.2}
\expandafter\gdef\csname num@d:pv:novel_10:dacis-random:P:hi\endcsname{+2.2}
\expandafter\gdef\csname num@d:pv:novel_10:dacis-random:P:abs\endcsname{0.5}
\expandafter\gdef\csname num@d:pv:novel_10:dacis-random:P:wins\endcsname{3}
\expandafter\gdef\csname num@d:pv:novel_10:dacis-random:P:n\endcsname{5}
\expandafter\gdef\csname num@d:pv:novel_10:dacis-random:P:worst\endcsname{2.2}
\expandafter\gdef\csname num@d:pv:novel_10:dacis-l1:PMP\endcsname{+0.1}
\expandafter\gdef\csname num@d:pv:novel_10:dacis-l1:PMP:lo\endcsname{\ensuremath{-}1.1}
\expandafter\gdef\csname num@d:pv:novel_10:dacis-l1:PMP:hi\endcsname{+1.4}
\expandafter\gdef\csname num@d:pv:novel_10:dacis-l1:PMP:abs\endcsname{0.1}
\expandafter\gdef\csname num@d:pv:novel_10:dacis-l1:PMP:wins\endcsname{2}
\expandafter\gdef\csname num@d:pv:novel_10:dacis-l1:PMP:n\endcsname{5}
\expandafter\gdef\csname num@d:pv:novel_10:dacis-l1:PMP:worst\endcsname{1.4}
\expandafter\gdef\csname num@d:pv:novel_10:dacis-l1:P\endcsname{+0.4}
\expandafter\gdef\csname num@d:pv:novel_10:dacis-l1:P:lo\endcsname{\ensuremath{-}1.5}
\expandafter\gdef\csname num@d:pv:novel_10:dacis-l1:P:hi\endcsname{+2.4}
\expandafter\gdef\csname num@d:pv:novel_10:dacis-l1:P:abs\endcsname{0.4}
\expandafter\gdef\csname num@d:pv:novel_10:dacis-l1:P:wins\endcsname{2}
\expandafter\gdef\csname num@d:pv:novel_10:dacis-l1:P:n\endcsname{5}
\expandafter\gdef\csname num@d:pv:novel_10:dacis-l1:P:worst\endcsname{2.4}
\expandafter\gdef\csname num@d:pv:novel_10:dacis-taylor:PMP\endcsname{\ensuremath{-}0.5}
\expandafter\gdef\csname num@d:pv:novel_10:dacis-taylor:PMP:lo\endcsname{\ensuremath{-}1.7}
\expandafter\gdef\csname num@d:pv:novel_10:dacis-taylor:PMP:hi\endcsname{+0.6}
\expandafter\gdef\csname num@d:pv:novel_10:dacis-taylor:PMP:abs\endcsname{0.5}
\expandafter\gdef\csname num@d:pv:novel_10:dacis-taylor:PMP:wins\endcsname{2}
\expandafter\gdef\csname num@d:pv:novel_10:dacis-taylor:PMP:n\endcsname{5}
\expandafter\gdef\csname num@d:pv:novel_10:dacis-taylor:PMP:worst\endcsname{1.7}
\expandafter\gdef\csname num@d:pv:novel_10:dacis-taylor:P\endcsname{\ensuremath{-}1.7}
\expandafter\gdef\csname num@d:pv:novel_10:dacis-taylor:P:lo\endcsname{\ensuremath{-}2.4}
\expandafter\gdef\csname num@d:pv:novel_10:dacis-taylor:P:hi\endcsname{\ensuremath{-}0.9}
\expandafter\gdef\csname num@d:pv:novel_10:dacis-taylor:P:abs\endcsname{1.7}
\expandafter\gdef\csname num@d:pv:novel_10:dacis-taylor:P:wins\endcsname{0}
\expandafter\gdef\csname num@d:pv:novel_10:dacis-taylor:P:n\endcsname{5}
\expandafter\gdef\csname num@d:pv:novel_10:dacis-taylor:P:worst\endcsname{2.4}
\expandafter\gdef\csname num@d:pv:novel_10:dacis-bn_gamma:PMP\endcsname{+1.4}
\expandafter\gdef\csname num@d:pv:novel_10:dacis-bn_gamma:PMP:lo\endcsname{\ensuremath{-}0.4}
\expandafter\gdef\csname num@d:pv:novel_10:dacis-bn_gamma:PMP:hi\endcsname{+3.2}
\expandafter\gdef\csname num@d:pv:novel_10:dacis-bn_gamma:PMP:abs\endcsname{1.4}
\expandafter\gdef\csname num@d:pv:novel_10:dacis-bn_gamma:PMP:wins\endcsname{4}
\expandafter\gdef\csname num@d:pv:novel_10:dacis-bn_gamma:PMP:n\endcsname{5}
\expandafter\gdef\csname num@d:pv:novel_10:dacis-bn_gamma:PMP:worst\endcsname{3.2}
\expandafter\gdef\csname num@d:pv:novel_10:dacis-bn_gamma:P\endcsname{+1.1}
\expandafter\gdef\csname num@d:pv:novel_10:dacis-bn_gamma:P:lo\endcsname{\ensuremath{-}1.7}
\expandafter\gdef\csname num@d:pv:novel_10:dacis-bn_gamma:P:hi\endcsname{+4.0}
\expandafter\gdef\csname num@d:pv:novel_10:dacis-bn_gamma:P:abs\endcsname{1.1}
\expandafter\gdef\csname num@d:pv:novel_10:dacis-bn_gamma:P:wins\endcsname{3}
\expandafter\gdef\csname num@d:pv:novel_10:dacis-bn_gamma:P:n\endcsname{5}
\expandafter\gdef\csname num@d:pv:novel_10:dacis-bn_gamma:P:worst\endcsname{4.0}
\expandafter\gdef\csname num@d:pv:novel_10:dacis-D:PMP\endcsname{+1.2}
\expandafter\gdef\csname num@d:pv:novel_10:dacis-D:PMP:lo\endcsname{\ensuremath{-}0.1}
\expandafter\gdef\csname num@d:pv:novel_10:dacis-D:PMP:hi\endcsname{+2.5}
\expandafter\gdef\csname num@d:pv:novel_10:dacis-D:PMP:abs\endcsname{1.2}
\expandafter\gdef\csname num@d:pv:novel_10:dacis-D:PMP:wins\endcsname{5}
\expandafter\gdef\csname num@d:pv:novel_10:dacis-D:PMP:n\endcsname{5}
\expandafter\gdef\csname num@d:pv:novel_10:dacis-D:PMP:worst\endcsname{2.5}
\expandafter\gdef\csname num@d:pv:novel_10:dacis-D:P\endcsname{+0.9}
\expandafter\gdef\csname num@d:pv:novel_10:dacis-D:P:lo\endcsname{\ensuremath{-}0.5}
\expandafter\gdef\csname num@d:pv:novel_10:dacis-D:P:hi\endcsname{+2.4}
\expandafter\gdef\csname num@d:pv:novel_10:dacis-D:P:abs\endcsname{0.9}
\expandafter\gdef\csname num@d:pv:novel_10:dacis-D:P:wins\endcsname{3}
\expandafter\gdef\csname num@d:pv:novel_10:dacis-D:P:n\endcsname{5}
\expandafter\gdef\csname num@d:pv:novel_10:dacis-D:P:worst\endcsname{2.4}
\expandafter\gdef\csname num@d:pv:novel_10:dacis-G:PMP\endcsname{\ensuremath{-}0.1}
\expandafter\gdef\csname num@d:pv:novel_10:dacis-G:PMP:lo\endcsname{\ensuremath{-}0.3}
\expandafter\gdef\csname num@d:pv:novel_10:dacis-G:PMP:hi\endcsname{+0.2}
\expandafter\gdef\csname num@d:pv:novel_10:dacis-G:PMP:abs\endcsname{0.1}
\expandafter\gdef\csname num@d:pv:novel_10:dacis-G:PMP:wins\endcsname{2}
\expandafter\gdef\csname num@d:pv:novel_10:dacis-G:PMP:n\endcsname{5}
\expandafter\gdef\csname num@d:pv:novel_10:dacis-G:PMP:worst\endcsname{0.3}
\expandafter\gdef\csname num@d:pv:novel_10:dacis-G:P\endcsname{\ensuremath{-}0.8}
\expandafter\gdef\csname num@d:pv:novel_10:dacis-G:P:lo\endcsname{\ensuremath{-}1.6}
\expandafter\gdef\csname num@d:pv:novel_10:dacis-G:P:hi\endcsname{\ensuremath{-}0.1}
\expandafter\gdef\csname num@d:pv:novel_10:dacis-G:P:abs\endcsname{0.8}
\expandafter\gdef\csname num@d:pv:novel_10:dacis-G:P:wins\endcsname{0}
\expandafter\gdef\csname num@d:pv:novel_10:dacis-G:P:n\endcsname{5}
\expandafter\gdef\csname num@d:pv:novel_10:dacis-G:P:worst\endcsname{1.6}
\expandafter\gdef\csname num@d:pv:novel_10:dacis-V:PMP\endcsname{+0.0}
\expandafter\gdef\csname num@d:pv:novel_10:dacis-V:PMP:lo\endcsname{\ensuremath{-}0.6}
\expandafter\gdef\csname num@d:pv:novel_10:dacis-V:PMP:hi\endcsname{+0.7}
\expandafter\gdef\csname num@d:pv:novel_10:dacis-V:PMP:abs\endcsname{0.0}
\expandafter\gdef\csname num@d:pv:novel_10:dacis-V:PMP:wins\endcsname{3}
\expandafter\gdef\csname num@d:pv:novel_10:dacis-V:PMP:n\endcsname{5}
\expandafter\gdef\csname num@d:pv:novel_10:dacis-V:PMP:worst\endcsname{0.7}
\expandafter\gdef\csname num@d:pv:novel_10:dacis-V:P\endcsname{\ensuremath{-}0.1}
\expandafter\gdef\csname num@d:pv:novel_10:dacis-V:P:lo\endcsname{\ensuremath{-}1.6}
\expandafter\gdef\csname num@d:pv:novel_10:dacis-V:P:hi\endcsname{+1.3}
\expandafter\gdef\csname num@d:pv:novel_10:dacis-V:P:abs\endcsname{0.1}
\expandafter\gdef\csname num@d:pv:novel_10:dacis-V:P:wins\endcsname{2}
\expandafter\gdef\csname num@d:pv:novel_10:dacis-V:P:n\endcsname{5}
\expandafter\gdef\csname num@d:pv:novel_10:dacis-V:P:worst\endcsname{1.6}
\expandafter\gdef\csname num@d:pv:novel_10:l1-random:PMP\endcsname{\ensuremath{-}0.2}
\expandafter\gdef\csname num@d:pv:novel_10:l1-random:PMP:lo\endcsname{\ensuremath{-}1.7}
\expandafter\gdef\csname num@d:pv:novel_10:l1-random:PMP:hi\endcsname{+1.3}
\expandafter\gdef\csname num@d:pv:novel_10:l1-random:PMP:abs\endcsname{0.2}
\expandafter\gdef\csname num@d:pv:novel_10:l1-random:PMP:wins\endcsname{2}
\expandafter\gdef\csname num@d:pv:novel_10:l1-random:PMP:n\endcsname{5}
\expandafter\gdef\csname num@d:pv:novel_10:l1-random:PMP:worst\endcsname{1.7}
\expandafter\gdef\csname num@d:pv:novel_10:l1-random:P\endcsname{+0.1}
\expandafter\gdef\csname num@d:pv:novel_10:l1-random:P:lo\endcsname{\ensuremath{-}1.8}
\expandafter\gdef\csname num@d:pv:novel_10:l1-random:P:hi\endcsname{+2.0}
\expandafter\gdef\csname num@d:pv:novel_10:l1-random:P:abs\endcsname{0.1}
\expandafter\gdef\csname num@d:pv:novel_10:l1-random:P:wins\endcsname{3}
\expandafter\gdef\csname num@d:pv:novel_10:l1-random:P:n\endcsname{5}
\expandafter\gdef\csname num@d:pv:novel_10:l1-random:P:worst\endcsname{2.0}
\expandafter\gdef\csname num@d:pv:novel_10:taylor-random:PMP\endcsname{+0.5}
\expandafter\gdef\csname num@d:pv:novel_10:taylor-random:PMP:lo\endcsname{\ensuremath{-}1.4}
\expandafter\gdef\csname num@d:pv:novel_10:taylor-random:PMP:hi\endcsname{+2.3}
\expandafter\gdef\csname num@d:pv:novel_10:taylor-random:PMP:abs\endcsname{0.5}
\expandafter\gdef\csname num@d:pv:novel_10:taylor-random:PMP:wins\endcsname{2}
\expandafter\gdef\csname num@d:pv:novel_10:taylor-random:PMP:n\endcsname{5}
\expandafter\gdef\csname num@d:pv:novel_10:taylor-random:PMP:worst\endcsname{2.3}
\expandafter\gdef\csname num@d:pv:novel_10:taylor-random:P\endcsname{+2.2}
\expandafter\gdef\csname num@d:pv:novel_10:taylor-random:P:lo\endcsname{+0.3}
\expandafter\gdef\csname num@d:pv:novel_10:taylor-random:P:hi\endcsname{+4.1}
\expandafter\gdef\csname num@d:pv:novel_10:taylor-random:P:abs\endcsname{2.2}
\expandafter\gdef\csname num@d:pv:novel_10:taylor-random:P:wins\endcsname{5}
\expandafter\gdef\csname num@d:pv:novel_10:taylor-random:P:n\endcsname{5}
\expandafter\gdef\csname num@d:pv:novel_10:taylor-random:P:worst\endcsname{4.1}
\expandafter\gdef\csname num@d:pv:novel_10:bn_gamma-random:PMP\endcsname{\ensuremath{-}1.5}
\expandafter\gdef\csname num@d:pv:novel_10:bn_gamma-random:PMP:lo\endcsname{\ensuremath{-}3.1}
\expandafter\gdef\csname num@d:pv:novel_10:bn_gamma-random:PMP:hi\endcsname{+0.2}
\expandafter\gdef\csname num@d:pv:novel_10:bn_gamma-random:PMP:abs\endcsname{1.5}
\expandafter\gdef\csname num@d:pv:novel_10:bn_gamma-random:PMP:wins\endcsname{1}
\expandafter\gdef\csname num@d:pv:novel_10:bn_gamma-random:PMP:n\endcsname{5}
\expandafter\gdef\csname num@d:pv:novel_10:bn_gamma-random:PMP:worst\endcsname{3.1}
\expandafter\gdef\csname num@d:pv:novel_10:bn_gamma-random:P\endcsname{\ensuremath{-}0.6}
\expandafter\gdef\csname num@d:pv:novel_10:bn_gamma-random:P:lo\endcsname{\ensuremath{-}3.9}
\expandafter\gdef\csname num@d:pv:novel_10:bn_gamma-random:P:hi\endcsname{+2.7}
\expandafter\gdef\csname num@d:pv:novel_10:bn_gamma-random:P:abs\endcsname{0.6}
\expandafter\gdef\csname num@d:pv:novel_10:bn_gamma-random:P:wins\endcsname{2}
\expandafter\gdef\csname num@d:pv:novel_10:bn_gamma-random:P:n\endcsname{5}
\expandafter\gdef\csname num@d:pv:novel_10:bn_gamma-random:P:worst\endcsname{3.9}
\expandafter\gdef\csname num@d:pv:novel_10:D-random:PMP\endcsname{\ensuremath{-}1.2}
\expandafter\gdef\csname num@d:pv:novel_10:D-random:PMP:lo\endcsname{\ensuremath{-}3.8}
\expandafter\gdef\csname num@d:pv:novel_10:D-random:PMP:hi\endcsname{+1.3}
\expandafter\gdef\csname num@d:pv:novel_10:D-random:PMP:abs\endcsname{1.2}
\expandafter\gdef\csname num@d:pv:novel_10:D-random:PMP:wins\endcsname{1}
\expandafter\gdef\csname num@d:pv:novel_10:D-random:PMP:n\endcsname{5}
\expandafter\gdef\csname num@d:pv:novel_10:D-random:PMP:worst\endcsname{3.8}
\expandafter\gdef\csname num@d:pv:novel_10:D-random:P\endcsname{\ensuremath{-}0.4}
\expandafter\gdef\csname num@d:pv:novel_10:D-random:P:lo\endcsname{\ensuremath{-}2.1}
\expandafter\gdef\csname num@d:pv:novel_10:D-random:P:hi\endcsname{+1.3}
\expandafter\gdef\csname num@d:pv:novel_10:D-random:P:abs\endcsname{0.4}
\expandafter\gdef\csname num@d:pv:novel_10:D-random:P:wins\endcsname{1}
\expandafter\gdef\csname num@d:pv:novel_10:D-random:P:n\endcsname{5}
\expandafter\gdef\csname num@d:pv:novel_10:D-random:P:worst\endcsname{2.1}
\expandafter\gdef\csname num@d:pv:novel_10:G-random:PMP\endcsname{\ensuremath{-}0.0}
\expandafter\gdef\csname num@d:pv:novel_10:G-random:PMP:lo\endcsname{\ensuremath{-}1.8}
\expandafter\gdef\csname num@d:pv:novel_10:G-random:PMP:hi\endcsname{+1.8}
\expandafter\gdef\csname num@d:pv:novel_10:G-random:PMP:abs\endcsname{0.0}
\expandafter\gdef\csname num@d:pv:novel_10:G-random:PMP:wins\endcsname{3}
\expandafter\gdef\csname num@d:pv:novel_10:G-random:PMP:n\endcsname{5}
\expandafter\gdef\csname num@d:pv:novel_10:G-random:PMP:worst\endcsname{1.8}
\expandafter\gdef\csname num@d:pv:novel_10:G-random:P\endcsname{+1.4}
\expandafter\gdef\csname num@d:pv:novel_10:G-random:P:lo\endcsname{\ensuremath{-}0.5}
\expandafter\gdef\csname num@d:pv:novel_10:G-random:P:hi\endcsname{+3.2}
\expandafter\gdef\csname num@d:pv:novel_10:G-random:P:abs\endcsname{1.4}
\expandafter\gdef\csname num@d:pv:novel_10:G-random:P:wins\endcsname{4}
\expandafter\gdef\csname num@d:pv:novel_10:G-random:P:n\endcsname{5}
\expandafter\gdef\csname num@d:pv:novel_10:G-random:P:worst\endcsname{3.2}
\expandafter\gdef\csname num@d:pv:novel_10:V-random:PMP\endcsname{\ensuremath{-}0.1}
\expandafter\gdef\csname num@d:pv:novel_10:V-random:PMP:lo\endcsname{\ensuremath{-}1.5}
\expandafter\gdef\csname num@d:pv:novel_10:V-random:PMP:hi\endcsname{+1.3}
\expandafter\gdef\csname num@d:pv:novel_10:V-random:PMP:abs\endcsname{0.1}
\expandafter\gdef\csname num@d:pv:novel_10:V-random:PMP:wins\endcsname{1}
\expandafter\gdef\csname num@d:pv:novel_10:V-random:PMP:n\endcsname{5}
\expandafter\gdef\csname num@d:pv:novel_10:V-random:PMP:worst\endcsname{1.5}
\expandafter\gdef\csname num@d:pv:novel_10:V-random:P\endcsname{+0.6}
\expandafter\gdef\csname num@d:pv:novel_10:V-random:P:lo\endcsname{\ensuremath{-}1.0}
\expandafter\gdef\csname num@d:pv:novel_10:V-random:P:hi\endcsname{+2.3}
\expandafter\gdef\csname num@d:pv:novel_10:V-random:P:abs\endcsname{0.6}
\expandafter\gdef\csname num@d:pv:novel_10:V-random:P:wins\endcsname{3}
\expandafter\gdef\csname num@d:pv:novel_10:V-random:P:n\endcsname{5}
\expandafter\gdef\csname num@d:pv:novel_10:V-random:P:worst\endcsname{2.3}
\expandafter\gdef\csname num@d:pv:novel_10:PM-P\endcsname{\ensuremath{-}0.8}
\expandafter\gdef\csname num@d:pv:novel_10:PM-P:lo\endcsname{\ensuremath{-}2.2}
\expandafter\gdef\csname num@d:pv:novel_10:PM-P:hi\endcsname{+0.7}
\expandafter\gdef\csname num@d:pv:novel_10:PM-P:abs\endcsname{0.8}
\expandafter\gdef\csname num@d:pv:novel_10:PM-P:wins\endcsname{1}
\expandafter\gdef\csname num@d:pv:novel_10:PM-P:n\endcsname{5}
\expandafter\gdef\csname num@d:pv:novel_10:PM-P:worst\endcsname{2.2}
\expandafter\gdef\csname num@d:pv:novel_10:MP-P\endcsname{+0.2}
\expandafter\gdef\csname num@d:pv:novel_10:MP-P:lo\endcsname{\ensuremath{-}1.3}
\expandafter\gdef\csname num@d:pv:novel_10:MP-P:hi\endcsname{+1.6}
\expandafter\gdef\csname num@d:pv:novel_10:MP-P:abs\endcsname{0.2}
\expandafter\gdef\csname num@d:pv:novel_10:MP-P:wins\endcsname{1}
\expandafter\gdef\csname num@d:pv:novel_10:MP-P:n\endcsname{5}
\expandafter\gdef\csname num@d:pv:novel_10:MP-P:worst\endcsname{1.6}
\expandafter\gdef\csname num@d:pv:novel_10:PGP-P\endcsname{+0.4}
\expandafter\gdef\csname num@d:pv:novel_10:PGP-P:lo\endcsname{\ensuremath{-}0.4}
\expandafter\gdef\csname num@d:pv:novel_10:PGP-P:hi\endcsname{+1.2}
\expandafter\gdef\csname num@d:pv:novel_10:PGP-P:abs\endcsname{0.4}
\expandafter\gdef\csname num@d:pv:novel_10:PGP-P:wins\endcsname{3}
\expandafter\gdef\csname num@d:pv:novel_10:PGP-P:n\endcsname{5}
\expandafter\gdef\csname num@d:pv:novel_10:PGP-P:worst\endcsname{1.2}
\expandafter\gdef\csname num@d:pv:novel_10:PMPM-P\endcsname{\ensuremath{-}0.0}
\expandafter\gdef\csname num@d:pv:novel_10:PMPM-P:lo\endcsname{\ensuremath{-}1.5}
\expandafter\gdef\csname num@d:pv:novel_10:PMPM-P:hi\endcsname{+1.5}
\expandafter\gdef\csname num@d:pv:novel_10:PMPM-P:abs\endcsname{0.0}
\expandafter\gdef\csname num@d:pv:novel_10:PMPM-P:wins\endcsname{2}
\expandafter\gdef\csname num@d:pv:novel_10:PMPM-P:n\endcsname{5}
\expandafter\gdef\csname num@d:pv:novel_10:PMPM-P:worst\endcsname{1.5}
\expandafter\gdef\csname num@d:pv:novel_10:PM2x-P\endcsname{\ensuremath{-}0.8}
\expandafter\gdef\csname num@d:pv:novel_10:PM2x-P:lo\endcsname{\ensuremath{-}2.1}
\expandafter\gdef\csname num@d:pv:novel_10:PM2x-P:hi\endcsname{+0.6}
\expandafter\gdef\csname num@d:pv:novel_10:PM2x-P:abs\endcsname{0.8}
\expandafter\gdef\csname num@d:pv:novel_10:PM2x-P:wins\endcsname{1}
\expandafter\gdef\csname num@d:pv:novel_10:PM2x-P:n\endcsname{5}
\expandafter\gdef\csname num@d:pv:novel_10:PM2x-P:worst\endcsname{2.1}
\expandafter\gdef\csname num@d:pv:novel_10:PMP-P\endcsname{+0.2}
\expandafter\gdef\csname num@d:pv:novel_10:PMP-P:lo\endcsname{\ensuremath{-}0.3}
\expandafter\gdef\csname num@d:pv:novel_10:PMP-P:hi\endcsname{+0.7}
\expandafter\gdef\csname num@d:pv:novel_10:PMP-P:abs\endcsname{0.2}
\expandafter\gdef\csname num@d:pv:novel_10:PMP-P:wins\endcsname{2}
\expandafter\gdef\csname num@d:pv:novel_10:PMP-P:n\endcsname{5}
\expandafter\gdef\csname num@d:pv:novel_10:PMP-P:worst\endcsname{0.7}
\expandafter\gdef\csname num@d:pv:novel_10:densem-dense\endcsname{\ensuremath{-}0.3}
\expandafter\gdef\csname num@d:pv:novel_10:densem-dense:lo\endcsname{\ensuremath{-}0.6}
\expandafter\gdef\csname num@d:pv:novel_10:densem-dense:hi\endcsname{\ensuremath{-}0.1}
\expandafter\gdef\csname num@d:pv:novel_10:densem-dense:abs\endcsname{0.3}
\expandafter\gdef\csname num@d:pv:novel_10:densem-dense:wins\endcsname{0}
\expandafter\gdef\csname num@d:pv:novel_10:densem-dense:n\endcsname{5}
\expandafter\gdef\csname num@d:pv:novel_10:densem-dense:worst\endcsname{0.6}
\expandafter\gdef\csname num@d:pv:novel_10:densermr-dense\endcsname{\ensuremath{-}0.1}
\expandafter\gdef\csname num@d:pv:novel_10:densermr-dense:lo\endcsname{\ensuremath{-}0.2}
\expandafter\gdef\csname num@d:pv:novel_10:densermr-dense:hi\endcsname{+0.0}
\expandafter\gdef\csname num@d:pv:novel_10:densermr-dense:abs\endcsname{0.1}
\expandafter\gdef\csname num@d:pv:novel_10:densermr-dense:wins\endcsname{1}
\expandafter\gdef\csname num@d:pv:novel_10:densermr-dense:n\endcsname{5}
\expandafter\gdef\csname num@d:pv:novel_10:densermr-dense:worst\endcsname{0.2}
\expandafter\gdef\csname num@d:pv:novel_10:dense-dacisP\endcsname{+4.5}
\expandafter\gdef\csname num@d:pv:novel_10:dense-dacisP:lo\endcsname{+2.1}
\expandafter\gdef\csname num@d:pv:novel_10:dense-dacisP:hi\endcsname{+6.9}
\expandafter\gdef\csname num@d:pv:novel_10:dense-dacisP:abs\endcsname{4.5}
\expandafter\gdef\csname num@d:pv:novel_10:dense-dacisP:wins\endcsname{5}
\expandafter\gdef\csname num@d:pv:novel_10:dense-dacisP:n\endcsname{5}
\expandafter\gdef\csname num@d:pv:novel_10:dense-dacisP:worst\endcsname{6.9}
\expandafter\gdef\csname num@d:pv:novel_10:dense-dacisPMP\endcsname{+4.2}
\expandafter\gdef\csname num@d:pv:novel_10:dense-dacisPMP:lo\endcsname{+1.9}
\expandafter\gdef\csname num@d:pv:novel_10:dense-dacisPMP:hi\endcsname{+6.6}
\expandafter\gdef\csname num@d:pv:novel_10:dense-dacisPMP:abs\endcsname{4.2}
\expandafter\gdef\csname num@d:pv:novel_10:dense-dacisPMP:wins\endcsname{5}
\expandafter\gdef\csname num@d:pv:novel_10:dense-dacisPMP:n\endcsname{5}
\expandafter\gdef\csname num@d:pv:novel_10:dense-dacisPMP:worst\endcsname{6.6}
\expandafter\gdef\csname num@d:pv:novel_10:dense-randomP\endcsname{+5.0}
\expandafter\gdef\csname num@d:pv:novel_10:dense-randomP:lo\endcsname{+1.8}
\expandafter\gdef\csname num@d:pv:novel_10:dense-randomP:hi\endcsname{+8.2}
\expandafter\gdef\csname num@d:pv:novel_10:dense-randomP:abs\endcsname{5.0}
\expandafter\gdef\csname num@d:pv:novel_10:dense-randomP:wins\endcsname{5}
\expandafter\gdef\csname num@d:pv:novel_10:dense-randomP:n\endcsname{5}
\expandafter\gdef\csname num@d:pv:novel_10:dense-randomP:worst\endcsname{8.2}
\expandafter\gdef\csname num@d:pv:novel_10:dense-randomPMP\endcsname{+4.2}
\expandafter\gdef\csname num@d:pv:novel_10:dense-randomPMP:lo\endcsname{+1.0}
\expandafter\gdef\csname num@d:pv:novel_10:dense-randomPMP:hi\endcsname{+7.4}
\expandafter\gdef\csname num@d:pv:novel_10:dense-randomPMP:abs\endcsname{4.2}
\expandafter\gdef\csname num@d:pv:novel_10:dense-randomPMP:wins\endcsname{5}
\expandafter\gdef\csname num@d:pv:novel_10:dense-randomPMP:n\endcsname{5}
\expandafter\gdef\csname num@d:pv:novel_10:dense-randomPMP:worst\endcsname{7.4}
\expandafter\gdef\csname num@d:pv:novel_10:norefine-full\endcsname{\ensuremath{-}0.9}
\expandafter\gdef\csname num@d:pv:novel_10:norefine-full:lo\endcsname{\ensuremath{-}1.8}
\expandafter\gdef\csname num@d:pv:novel_10:norefine-full:hi\endcsname{\ensuremath{-}0.1}
\expandafter\gdef\csname num@d:pv:novel_10:norefine-full:abs\endcsname{0.9}
\expandafter\gdef\csname num@d:pv:novel_10:norefine-full:wins\endcsname{0}
\expandafter\gdef\csname num@d:pv:novel_10:norefine-full:n\endcsname{5}
\expandafter\gdef\csname num@d:pv:novel_10:norefine-full:worst\endcsname{1.8}
\expandafter\gdef\csname num@d:pv:novel_10:mnv3-dacisPMP\endcsname{+4.2}
\expandafter\gdef\csname num@d:pv:novel_10:mnv3-dacisPMP:lo\endcsname{+1.8}
\expandafter\gdef\csname num@d:pv:novel_10:mnv3-dacisPMP:hi\endcsname{+6.7}
\expandafter\gdef\csname num@d:pv:novel_10:mnv3-dacisPMP:abs\endcsname{4.2}
\expandafter\gdef\csname num@d:pv:novel_10:mnv3-dacisPMP:wins\endcsname{5}
\expandafter\gdef\csname num@d:pv:novel_10:mnv3-dacisPMP:n\endcsname{5}
\expandafter\gdef\csname num@d:pv:novel_10:mnv3-dacisPMP:worst\endcsname{6.7}
\expandafter\gdef\csname num@d:pv:novel_10:effb0-dacisPMP\endcsname{+6.5}
\expandafter\gdef\csname num@d:pv:novel_10:effb0-dacisPMP:lo\endcsname{+3.5}
\expandafter\gdef\csname num@d:pv:novel_10:effb0-dacisPMP:hi\endcsname{+9.6}
\expandafter\gdef\csname num@d:pv:novel_10:effb0-dacisPMP:abs\endcsname{6.5}
\expandafter\gdef\csname num@d:pv:novel_10:effb0-dacisPMP:wins\endcsname{5}
\expandafter\gdef\csname num@d:pv:novel_10:effb0-dacisPMP:n\endcsname{5}
\expandafter\gdef\csname num@d:pv:novel_10:effb0-dacisPMP:worst\endcsname{9.6}
\expandafter\gdef\csname num@d:pv:novel_10:mnv2-dacisPMP\endcsname{+4.3}
\expandafter\gdef\csname num@d:pv:novel_10:mnv2-dacisPMP:lo\endcsname{+2.4}
\expandafter\gdef\csname num@d:pv:novel_10:mnv2-dacisPMP:hi\endcsname{+6.2}
\expandafter\gdef\csname num@d:pv:novel_10:mnv2-dacisPMP:abs\endcsname{4.3}
\expandafter\gdef\csname num@d:pv:novel_10:mnv2-dacisPMP:wins\endcsname{5}
\expandafter\gdef\csname num@d:pv:novel_10:mnv2-dacisPMP:n\endcsname{5}
\expandafter\gdef\csname num@d:pv:novel_10:mnv2-dacisPMP:worst\endcsname{6.2}
\expandafter\gdef\csname num@d:pv:novel_10:dacis-l1:PMP50\endcsname{\ensuremath{-}0.4}
\expandafter\gdef\csname num@d:pv:novel_10:dacis-l1:PMP50:lo\endcsname{\ensuremath{-}1.6}
\expandafter\gdef\csname num@d:pv:novel_10:dacis-l1:PMP50:hi\endcsname{+0.8}
\expandafter\gdef\csname num@d:pv:novel_10:dacis-l1:PMP50:abs\endcsname{0.4}
\expandafter\gdef\csname num@d:pv:novel_10:dacis-l1:PMP50:wins\endcsname{2}
\expandafter\gdef\csname num@d:pv:novel_10:dacis-l1:PMP50:n\endcsname{5}
\expandafter\gdef\csname num@d:pv:novel_10:dacis-l1:PMP50:worst\endcsname{1.6}
\expandafter\gdef\csname num@d:pv:novel_10:dense-dacisPMP50\endcsname{+2.5}
\expandafter\gdef\csname num@d:pv:novel_10:dense-dacisPMP50:lo\endcsname{+0.6}
\expandafter\gdef\csname num@d:pv:novel_10:dense-dacisPMP50:hi\endcsname{+4.4}
\expandafter\gdef\csname num@d:pv:novel_10:dense-dacisPMP50:abs\endcsname{2.5}
\expandafter\gdef\csname num@d:pv:novel_10:dense-dacisPMP50:wins\endcsname{5}
\expandafter\gdef\csname num@d:pv:novel_10:dense-dacisPMP50:n\endcsname{5}
\expandafter\gdef\csname num@d:pv:novel_10:dense-dacisPMP50:worst\endcsname{4.4}
\expandafter\gdef\csname num@d:pv:novel_10:dacis-l1:PMP20\endcsname{\ensuremath{-}0.4}
\expandafter\gdef\csname num@d:pv:novel_10:dacis-l1:PMP20:lo\endcsname{\ensuremath{-}1.8}
\expandafter\gdef\csname num@d:pv:novel_10:dacis-l1:PMP20:hi\endcsname{+1.0}
\expandafter\gdef\csname num@d:pv:novel_10:dacis-l1:PMP20:abs\endcsname{0.4}
\expandafter\gdef\csname num@d:pv:novel_10:dacis-l1:PMP20:wins\endcsname{2}
\expandafter\gdef\csname num@d:pv:novel_10:dacis-l1:PMP20:n\endcsname{5}
\expandafter\gdef\csname num@d:pv:novel_10:dacis-l1:PMP20:worst\endcsname{1.8}
\expandafter\gdef\csname num@d:pv:novel_10:dense-dacisPMP20\endcsname{+6.8}
\expandafter\gdef\csname num@d:pv:novel_10:dense-dacisPMP20:lo\endcsname{+3.5}
\expandafter\gdef\csname num@d:pv:novel_10:dense-dacisPMP20:hi\endcsname{+10.2}
\expandafter\gdef\csname num@d:pv:novel_10:dense-dacisPMP20:abs\endcsname{6.8}
\expandafter\gdef\csname num@d:pv:novel_10:dense-dacisPMP20:wins\endcsname{5}
\expandafter\gdef\csname num@d:pv:novel_10:dense-dacisPMP20:n\endcsname{5}
\expandafter\gdef\csname num@d:pv:novel_10:dense-dacisPMP20:worst\endcsname{10.2}
\expandafter\gdef\csname num@d:pv:novel_10:dacis-l1:PMP10\endcsname{\ensuremath{-}1.0}
\expandafter\gdef\csname num@d:pv:novel_10:dacis-l1:PMP10:lo\endcsname{\ensuremath{-}3.8}
\expandafter\gdef\csname num@d:pv:novel_10:dacis-l1:PMP10:hi\endcsname{+1.7}
\expandafter\gdef\csname num@d:pv:novel_10:dacis-l1:PMP10:abs\endcsname{1.0}
\expandafter\gdef\csname num@d:pv:novel_10:dacis-l1:PMP10:wins\endcsname{3}
\expandafter\gdef\csname num@d:pv:novel_10:dacis-l1:PMP10:n\endcsname{5}
\expandafter\gdef\csname num@d:pv:novel_10:dacis-l1:PMP10:worst\endcsname{3.8}
\expandafter\gdef\csname num@d:pv:novel_10:dense-dacisPMP10\endcsname{+12.1}
\expandafter\gdef\csname num@d:pv:novel_10:dense-dacisPMP10:lo\endcsname{+5.9}
\expandafter\gdef\csname num@d:pv:novel_10:dense-dacisPMP10:hi\endcsname{+18.3}
\expandafter\gdef\csname num@d:pv:novel_10:dense-dacisPMP10:abs\endcsname{12.1}
\expandafter\gdef\csname num@d:pv:novel_10:dense-dacisPMP10:wins\endcsname{5}
\expandafter\gdef\csname num@d:pv:novel_10:dense-dacisPMP10:n\endcsname{5}
\expandafter\gdef\csname num@d:pv:novel_10:dense-dacisPMP10:worst\endcsname{18.3}
\expandafter\gdef\csname num@d:es:novel_10:dacis-random:PMP\endcsname{\ensuremath{-}0.5}
\expandafter\gdef\csname num@d:es:novel_10:dacis-random:PMP:lo\endcsname{\ensuremath{-}4.1}
\expandafter\gdef\csname num@d:es:novel_10:dacis-random:PMP:hi\endcsname{+3.2}
\expandafter\gdef\csname num@d:es:novel_10:dacis-random:PMP:abs\endcsname{0.5}
\expandafter\gdef\csname num@d:es:novel_10:dacis-random:PMP:wins\endcsname{3}
\expandafter\gdef\csname num@d:es:novel_10:dacis-random:PMP:n\endcsname{5}
\expandafter\gdef\csname num@d:es:novel_10:dacis-random:PMP:worst\endcsname{4.1}
\expandafter\gdef\csname num@d:es:novel_10:dacis-random:P\endcsname{+0.4}
\expandafter\gdef\csname num@d:es:novel_10:dacis-random:P:lo\endcsname{\ensuremath{-}1.6}
\expandafter\gdef\csname num@d:es:novel_10:dacis-random:P:hi\endcsname{+2.3}
\expandafter\gdef\csname num@d:es:novel_10:dacis-random:P:abs\endcsname{0.4}
\expandafter\gdef\csname num@d:es:novel_10:dacis-random:P:wins\endcsname{3}
\expandafter\gdef\csname num@d:es:novel_10:dacis-random:P:n\endcsname{5}
\expandafter\gdef\csname num@d:es:novel_10:dacis-random:P:worst\endcsname{2.3}
\expandafter\gdef\csname num@d:es:novel_10:dacis-l1:PMP\endcsname{+0.1}
\expandafter\gdef\csname num@d:es:novel_10:dacis-l1:PMP:lo\endcsname{\ensuremath{-}2.9}
\expandafter\gdef\csname num@d:es:novel_10:dacis-l1:PMP:hi\endcsname{+3.1}
\expandafter\gdef\csname num@d:es:novel_10:dacis-l1:PMP:abs\endcsname{0.1}
\expandafter\gdef\csname num@d:es:novel_10:dacis-l1:PMP:wins\endcsname{3}
\expandafter\gdef\csname num@d:es:novel_10:dacis-l1:PMP:n\endcsname{5}
\expandafter\gdef\csname num@d:es:novel_10:dacis-l1:PMP:worst\endcsname{3.1}
\expandafter\gdef\csname num@d:es:novel_10:dacis-l1:P\endcsname{\ensuremath{-}0.1}
\expandafter\gdef\csname num@d:es:novel_10:dacis-l1:P:lo\endcsname{\ensuremath{-}3.9}
\expandafter\gdef\csname num@d:es:novel_10:dacis-l1:P:hi\endcsname{+3.7}
\expandafter\gdef\csname num@d:es:novel_10:dacis-l1:P:abs\endcsname{0.1}
\expandafter\gdef\csname num@d:es:novel_10:dacis-l1:P:wins\endcsname{2}
\expandafter\gdef\csname num@d:es:novel_10:dacis-l1:P:n\endcsname{5}
\expandafter\gdef\csname num@d:es:novel_10:dacis-l1:P:worst\endcsname{3.9}
\expandafter\gdef\csname num@d:es:novel_10:dacis-taylor:PMP\endcsname{\ensuremath{-}0.0}
\expandafter\gdef\csname num@d:es:novel_10:dacis-taylor:PMP:lo\endcsname{\ensuremath{-}5.4}
\expandafter\gdef\csname num@d:es:novel_10:dacis-taylor:PMP:hi\endcsname{+5.3}
\expandafter\gdef\csname num@d:es:novel_10:dacis-taylor:PMP:abs\endcsname{0.0}
\expandafter\gdef\csname num@d:es:novel_10:dacis-taylor:PMP:wins\endcsname{2}
\expandafter\gdef\csname num@d:es:novel_10:dacis-taylor:PMP:n\endcsname{5}
\expandafter\gdef\csname num@d:es:novel_10:dacis-taylor:PMP:worst\endcsname{5.4}
\expandafter\gdef\csname num@d:es:novel_10:dacis-taylor:P\endcsname{\ensuremath{-}1.8}
\expandafter\gdef\csname num@d:es:novel_10:dacis-taylor:P:lo\endcsname{\ensuremath{-}5.4}
\expandafter\gdef\csname num@d:es:novel_10:dacis-taylor:P:hi\endcsname{+1.8}
\expandafter\gdef\csname num@d:es:novel_10:dacis-taylor:P:abs\endcsname{1.8}
\expandafter\gdef\csname num@d:es:novel_10:dacis-taylor:P:wins\endcsname{1}
\expandafter\gdef\csname num@d:es:novel_10:dacis-taylor:P:n\endcsname{5}
\expandafter\gdef\csname num@d:es:novel_10:dacis-taylor:P:worst\endcsname{5.4}
\expandafter\gdef\csname num@d:es:novel_10:dacis-bn_gamma:PMP\endcsname{+2.2}
\expandafter\gdef\csname num@d:es:novel_10:dacis-bn_gamma:PMP:lo\endcsname{\ensuremath{-}0.2}
\expandafter\gdef\csname num@d:es:novel_10:dacis-bn_gamma:PMP:hi\endcsname{+4.6}
\expandafter\gdef\csname num@d:es:novel_10:dacis-bn_gamma:PMP:abs\endcsname{2.2}
\expandafter\gdef\csname num@d:es:novel_10:dacis-bn_gamma:PMP:wins\endcsname{5}
\expandafter\gdef\csname num@d:es:novel_10:dacis-bn_gamma:PMP:n\endcsname{5}
\expandafter\gdef\csname num@d:es:novel_10:dacis-bn_gamma:PMP:worst\endcsname{4.6}
\expandafter\gdef\csname num@d:es:novel_10:dacis-D:PMP\endcsname{+0.8}
\expandafter\gdef\csname num@d:es:novel_10:dacis-D:PMP:lo\endcsname{\ensuremath{-}3.3}
\expandafter\gdef\csname num@d:es:novel_10:dacis-D:PMP:hi\endcsname{+4.9}
\expandafter\gdef\csname num@d:es:novel_10:dacis-D:PMP:abs\endcsname{0.8}
\expandafter\gdef\csname num@d:es:novel_10:dacis-D:PMP:wins\endcsname{4}
\expandafter\gdef\csname num@d:es:novel_10:dacis-D:PMP:n\endcsname{5}
\expandafter\gdef\csname num@d:es:novel_10:dacis-D:PMP:worst\endcsname{4.9}
\expandafter\gdef\csname num@d:es:novel_10:dacis-D:P\endcsname{+0.4}
\expandafter\gdef\csname num@d:es:novel_10:dacis-D:P:lo\endcsname{\ensuremath{-}1.2}
\expandafter\gdef\csname num@d:es:novel_10:dacis-D:P:hi\endcsname{+1.9}
\expandafter\gdef\csname num@d:es:novel_10:dacis-D:P:abs\endcsname{0.4}
\expandafter\gdef\csname num@d:es:novel_10:dacis-D:P:wins\endcsname{2}
\expandafter\gdef\csname num@d:es:novel_10:dacis-D:P:n\endcsname{5}
\expandafter\gdef\csname num@d:es:novel_10:dacis-D:P:worst\endcsname{1.9}
\expandafter\gdef\csname num@d:es:novel_10:dacis-G:PMP\endcsname{+0.5}
\expandafter\gdef\csname num@d:es:novel_10:dacis-G:PMP:lo\endcsname{\ensuremath{-}2.5}
\expandafter\gdef\csname num@d:es:novel_10:dacis-G:PMP:hi\endcsname{+3.5}
\expandafter\gdef\csname num@d:es:novel_10:dacis-G:PMP:abs\endcsname{0.5}
\expandafter\gdef\csname num@d:es:novel_10:dacis-G:PMP:wins\endcsname{2}
\expandafter\gdef\csname num@d:es:novel_10:dacis-G:PMP:n\endcsname{5}
\expandafter\gdef\csname num@d:es:novel_10:dacis-G:PMP:worst\endcsname{3.5}
\expandafter\gdef\csname num@d:es:novel_10:dacis-V:PMP\endcsname{+0.5}
\expandafter\gdef\csname num@d:es:novel_10:dacis-V:PMP:lo\endcsname{\ensuremath{-}2.7}
\expandafter\gdef\csname num@d:es:novel_10:dacis-V:PMP:hi\endcsname{+3.7}
\expandafter\gdef\csname num@d:es:novel_10:dacis-V:PMP:abs\endcsname{0.5}
\expandafter\gdef\csname num@d:es:novel_10:dacis-V:PMP:wins\endcsname{3}
\expandafter\gdef\csname num@d:es:novel_10:dacis-V:PMP:n\endcsname{5}
\expandafter\gdef\csname num@d:es:novel_10:dacis-V:PMP:worst\endcsname{3.7}
\expandafter\gdef\csname num@d:es:novel_10:l1-random:PMP\endcsname{\ensuremath{-}0.6}
\expandafter\gdef\csname num@d:es:novel_10:l1-random:PMP:lo\endcsname{\ensuremath{-}2.0}
\expandafter\gdef\csname num@d:es:novel_10:l1-random:PMP:hi\endcsname{+0.9}
\expandafter\gdef\csname num@d:es:novel_10:l1-random:PMP:abs\endcsname{0.6}
\expandafter\gdef\csname num@d:es:novel_10:l1-random:PMP:wins\endcsname{2}
\expandafter\gdef\csname num@d:es:novel_10:l1-random:PMP:n\endcsname{5}
\expandafter\gdef\csname num@d:es:novel_10:l1-random:PMP:worst\endcsname{2.0}
\expandafter\gdef\csname num@d:es:novel_10:l1-random:P\endcsname{+0.4}
\expandafter\gdef\csname num@d:es:novel_10:l1-random:P:lo\endcsname{\ensuremath{-}3.9}
\expandafter\gdef\csname num@d:es:novel_10:l1-random:P:hi\endcsname{+4.8}
\expandafter\gdef\csname num@d:es:novel_10:l1-random:P:abs\endcsname{0.4}
\expandafter\gdef\csname num@d:es:novel_10:l1-random:P:wins\endcsname{3}
\expandafter\gdef\csname num@d:es:novel_10:l1-random:P:n\endcsname{5}
\expandafter\gdef\csname num@d:es:novel_10:l1-random:P:worst\endcsname{4.8}
\expandafter\gdef\csname num@d:es:novel_10:taylor-random:PMP\endcsname{\ensuremath{-}0.5}
\expandafter\gdef\csname num@d:es:novel_10:taylor-random:PMP:lo\endcsname{\ensuremath{-}2.8}
\expandafter\gdef\csname num@d:es:novel_10:taylor-random:PMP:hi\endcsname{+1.9}
\expandafter\gdef\csname num@d:es:novel_10:taylor-random:PMP:abs\endcsname{0.5}
\expandafter\gdef\csname num@d:es:novel_10:taylor-random:PMP:wins\endcsname{3}
\expandafter\gdef\csname num@d:es:novel_10:taylor-random:PMP:n\endcsname{5}
\expandafter\gdef\csname num@d:es:novel_10:taylor-random:PMP:worst\endcsname{2.8}
\expandafter\gdef\csname num@d:es:novel_10:taylor-random:P\endcsname{+2.2}
\expandafter\gdef\csname num@d:es:novel_10:taylor-random:P:lo\endcsname{\ensuremath{-}1.1}
\expandafter\gdef\csname num@d:es:novel_10:taylor-random:P:hi\endcsname{+5.4}
\expandafter\gdef\csname num@d:es:novel_10:taylor-random:P:abs\endcsname{2.2}
\expandafter\gdef\csname num@d:es:novel_10:taylor-random:P:wins\endcsname{4}
\expandafter\gdef\csname num@d:es:novel_10:taylor-random:P:n\endcsname{5}
\expandafter\gdef\csname num@d:es:novel_10:taylor-random:P:worst\endcsname{5.4}
\expandafter\gdef\csname num@d:es:novel_10:bn_gamma-random:PMP\endcsname{\ensuremath{-}2.7}
\expandafter\gdef\csname num@d:es:novel_10:bn_gamma-random:PMP:lo\endcsname{\ensuremath{-}5.3}
\expandafter\gdef\csname num@d:es:novel_10:bn_gamma-random:PMP:hi\endcsname{\ensuremath{-}0.1}
\expandafter\gdef\csname num@d:es:novel_10:bn_gamma-random:PMP:abs\endcsname{2.7}
\expandafter\gdef\csname num@d:es:novel_10:bn_gamma-random:PMP:wins\endcsname{0}
\expandafter\gdef\csname num@d:es:novel_10:bn_gamma-random:PMP:n\endcsname{5}
\expandafter\gdef\csname num@d:es:novel_10:bn_gamma-random:PMP:worst\endcsname{5.3}
\expandafter\gdef\csname num@d:es:novel_10:D-random:PMP\endcsname{\ensuremath{-}1.3}
\expandafter\gdef\csname num@d:es:novel_10:D-random:PMP:lo\endcsname{\ensuremath{-}3.8}
\expandafter\gdef\csname num@d:es:novel_10:D-random:PMP:hi\endcsname{+1.1}
\expandafter\gdef\csname num@d:es:novel_10:D-random:PMP:abs\endcsname{1.3}
\expandafter\gdef\csname num@d:es:novel_10:D-random:PMP:wins\endcsname{2}
\expandafter\gdef\csname num@d:es:novel_10:D-random:PMP:n\endcsname{5}
\expandafter\gdef\csname num@d:es:novel_10:D-random:PMP:worst\endcsname{3.8}
\expandafter\gdef\csname num@d:es:novel_10:D-random:P\endcsname{+0.0}
\expandafter\gdef\csname num@d:es:novel_10:D-random:P:lo\endcsname{\ensuremath{-}2.4}
\expandafter\gdef\csname num@d:es:novel_10:D-random:P:hi\endcsname{+2.4}
\expandafter\gdef\csname num@d:es:novel_10:D-random:P:abs\endcsname{0.0}
\expandafter\gdef\csname num@d:es:novel_10:D-random:P:wins\endcsname{2}
\expandafter\gdef\csname num@d:es:novel_10:D-random:P:n\endcsname{5}
\expandafter\gdef\csname num@d:es:novel_10:D-random:P:worst\endcsname{2.4}
\expandafter\gdef\csname num@d:es:novel_10:G-random:PMP\endcsname{\ensuremath{-}1.0}
\expandafter\gdef\csname num@d:es:novel_10:G-random:PMP:lo\endcsname{\ensuremath{-}3.1}
\expandafter\gdef\csname num@d:es:novel_10:G-random:PMP:hi\endcsname{+1.1}
\expandafter\gdef\csname num@d:es:novel_10:G-random:PMP:abs\endcsname{1.0}
\expandafter\gdef\csname num@d:es:novel_10:G-random:PMP:wins\endcsname{1}
\expandafter\gdef\csname num@d:es:novel_10:G-random:PMP:n\endcsname{5}
\expandafter\gdef\csname num@d:es:novel_10:G-random:PMP:worst\endcsname{3.1}
\expandafter\gdef\csname num@d:es:novel_10:V-random:PMP\endcsname{\ensuremath{-}1.0}
\expandafter\gdef\csname num@d:es:novel_10:V-random:PMP:lo\endcsname{\ensuremath{-}3.6}
\expandafter\gdef\csname num@d:es:novel_10:V-random:PMP:hi\endcsname{+1.7}
\expandafter\gdef\csname num@d:es:novel_10:V-random:PMP:abs\endcsname{1.0}
\expandafter\gdef\csname num@d:es:novel_10:V-random:PMP:wins\endcsname{1}
\expandafter\gdef\csname num@d:es:novel_10:V-random:PMP:n\endcsname{5}
\expandafter\gdef\csname num@d:es:novel_10:V-random:PMP:worst\endcsname{3.6}
\expandafter\gdef\csname num@d:es:novel_10:PM-P\endcsname{\ensuremath{-}0.5}
\expandafter\gdef\csname num@d:es:novel_10:PM-P:lo\endcsname{\ensuremath{-}2.7}
\expandafter\gdef\csname num@d:es:novel_10:PM-P:hi\endcsname{+1.8}
\expandafter\gdef\csname num@d:es:novel_10:PM-P:abs\endcsname{0.5}
\expandafter\gdef\csname num@d:es:novel_10:PM-P:wins\endcsname{2}
\expandafter\gdef\csname num@d:es:novel_10:PM-P:n\endcsname{5}
\expandafter\gdef\csname num@d:es:novel_10:PM-P:worst\endcsname{2.7}
\expandafter\gdef\csname num@d:es:novel_10:MP-P\endcsname{+0.2}
\expandafter\gdef\csname num@d:es:novel_10:MP-P:lo\endcsname{\ensuremath{-}1.4}
\expandafter\gdef\csname num@d:es:novel_10:MP-P:hi\endcsname{+1.8}
\expandafter\gdef\csname num@d:es:novel_10:MP-P:abs\endcsname{0.2}
\expandafter\gdef\csname num@d:es:novel_10:MP-P:wins\endcsname{3}
\expandafter\gdef\csname num@d:es:novel_10:MP-P:n\endcsname{5}
\expandafter\gdef\csname num@d:es:novel_10:MP-P:worst\endcsname{1.8}
\expandafter\gdef\csname num@d:es:novel_10:PGP-P\endcsname{+0.2}
\expandafter\gdef\csname num@d:es:novel_10:PGP-P:lo\endcsname{\ensuremath{-}1.9}
\expandafter\gdef\csname num@d:es:novel_10:PGP-P:hi\endcsname{+2.4}
\expandafter\gdef\csname num@d:es:novel_10:PGP-P:abs\endcsname{0.2}
\expandafter\gdef\csname num@d:es:novel_10:PGP-P:wins\endcsname{3}
\expandafter\gdef\csname num@d:es:novel_10:PGP-P:n\endcsname{5}
\expandafter\gdef\csname num@d:es:novel_10:PGP-P:worst\endcsname{2.4}
\expandafter\gdef\csname num@d:es:novel_10:PMP-P\endcsname{+0.8}
\expandafter\gdef\csname num@d:es:novel_10:PMP-P:lo\endcsname{\ensuremath{-}2.6}
\expandafter\gdef\csname num@d:es:novel_10:PMP-P:hi\endcsname{+4.2}
\expandafter\gdef\csname num@d:es:novel_10:PMP-P:abs\endcsname{0.8}
\expandafter\gdef\csname num@d:es:novel_10:PMP-P:wins\endcsname{3}
\expandafter\gdef\csname num@d:es:novel_10:PMP-P:n\endcsname{5}
\expandafter\gdef\csname num@d:es:novel_10:PMP-P:worst\endcsname{4.2}
\expandafter\gdef\csname num@d:es:novel_10:densem-dense\endcsname{\ensuremath{-}1.0}
\expandafter\gdef\csname num@d:es:novel_10:densem-dense:lo\endcsname{\ensuremath{-}2.8}
\expandafter\gdef\csname num@d:es:novel_10:densem-dense:hi\endcsname{+0.7}
\expandafter\gdef\csname num@d:es:novel_10:densem-dense:abs\endcsname{1.0}
\expandafter\gdef\csname num@d:es:novel_10:densem-dense:wins\endcsname{1}
\expandafter\gdef\csname num@d:es:novel_10:densem-dense:n\endcsname{5}
\expandafter\gdef\csname num@d:es:novel_10:densem-dense:worst\endcsname{2.8}
\expandafter\gdef\csname num@d:es:novel_10:densermr-dense\endcsname{\ensuremath{-}0.6}
\expandafter\gdef\csname num@d:es:novel_10:densermr-dense:lo\endcsname{\ensuremath{-}1.3}
\expandafter\gdef\csname num@d:es:novel_10:densermr-dense:hi\endcsname{+0.1}
\expandafter\gdef\csname num@d:es:novel_10:densermr-dense:abs\endcsname{0.6}
\expandafter\gdef\csname num@d:es:novel_10:densermr-dense:wins\endcsname{1}
\expandafter\gdef\csname num@d:es:novel_10:densermr-dense:n\endcsname{5}
\expandafter\gdef\csname num@d:es:novel_10:densermr-dense:worst\endcsname{1.3}
\expandafter\gdef\csname num@d:es:novel_10:dense-dacisP\endcsname{+6.3}
\expandafter\gdef\csname num@d:es:novel_10:dense-dacisP:lo\endcsname{+4.3}
\expandafter\gdef\csname num@d:es:novel_10:dense-dacisP:hi\endcsname{+8.3}
\expandafter\gdef\csname num@d:es:novel_10:dense-dacisP:abs\endcsname{6.3}
\expandafter\gdef\csname num@d:es:novel_10:dense-dacisP:wins\endcsname{5}
\expandafter\gdef\csname num@d:es:novel_10:dense-dacisP:n\endcsname{5}
\expandafter\gdef\csname num@d:es:novel_10:dense-dacisP:worst\endcsname{8.3}
\expandafter\gdef\csname num@d:es:novel_10:dense-dacisPMP\endcsname{+5.5}
\expandafter\gdef\csname num@d:es:novel_10:dense-dacisPMP:lo\endcsname{+3.3}
\expandafter\gdef\csname num@d:es:novel_10:dense-dacisPMP:hi\endcsname{+7.7}
\expandafter\gdef\csname num@d:es:novel_10:dense-dacisPMP:abs\endcsname{5.5}
\expandafter\gdef\csname num@d:es:novel_10:dense-dacisPMP:wins\endcsname{5}
\expandafter\gdef\csname num@d:es:novel_10:dense-dacisPMP:n\endcsname{5}
\expandafter\gdef\csname num@d:es:novel_10:dense-dacisPMP:worst\endcsname{7.7}
\expandafter\gdef\csname num@d:es:novel_10:dense-randomP\endcsname{+6.7}
\expandafter\gdef\csname num@d:es:novel_10:dense-randomP:lo\endcsname{+4.3}
\expandafter\gdef\csname num@d:es:novel_10:dense-randomP:hi\endcsname{+9.1}
\expandafter\gdef\csname num@d:es:novel_10:dense-randomP:abs\endcsname{6.7}
\expandafter\gdef\csname num@d:es:novel_10:dense-randomP:wins\endcsname{5}
\expandafter\gdef\csname num@d:es:novel_10:dense-randomP:n\endcsname{5}
\expandafter\gdef\csname num@d:es:novel_10:dense-randomP:worst\endcsname{9.1}
\expandafter\gdef\csname num@d:es:novel_10:dense-randomPMP\endcsname{+5.0}
\expandafter\gdef\csname num@d:es:novel_10:dense-randomPMP:lo\endcsname{+3.2}
\expandafter\gdef\csname num@d:es:novel_10:dense-randomPMP:hi\endcsname{+6.8}
\expandafter\gdef\csname num@d:es:novel_10:dense-randomPMP:abs\endcsname{5.0}
\expandafter\gdef\csname num@d:es:novel_10:dense-randomPMP:wins\endcsname{5}
\expandafter\gdef\csname num@d:es:novel_10:dense-randomPMP:n\endcsname{5}
\expandafter\gdef\csname num@d:es:novel_10:dense-randomPMP:worst\endcsname{6.8}
\expandafter\gdef\csname num@d:es:novel_10:mnv3-dacisPMP\endcsname{+16.3}
\expandafter\gdef\csname num@d:es:novel_10:mnv3-dacisPMP:lo\endcsname{+10.6}
\expandafter\gdef\csname num@d:es:novel_10:mnv3-dacisPMP:hi\endcsname{+22.0}
\expandafter\gdef\csname num@d:es:novel_10:mnv3-dacisPMP:abs\endcsname{16.3}
\expandafter\gdef\csname num@d:es:novel_10:mnv3-dacisPMP:wins\endcsname{5}
\expandafter\gdef\csname num@d:es:novel_10:mnv3-dacisPMP:n\endcsname{5}
\expandafter\gdef\csname num@d:es:novel_10:mnv3-dacisPMP:worst\endcsname{22.0}
\expandafter\gdef\csname num@d:es:novel_10:effb0-dacisPMP\endcsname{+14.7}
\expandafter\gdef\csname num@d:es:novel_10:effb0-dacisPMP:lo\endcsname{+7.8}
\expandafter\gdef\csname num@d:es:novel_10:effb0-dacisPMP:hi\endcsname{+21.6}
\expandafter\gdef\csname num@d:es:novel_10:effb0-dacisPMP:abs\endcsname{14.7}
\expandafter\gdef\csname num@d:es:novel_10:effb0-dacisPMP:wins\endcsname{5}
\expandafter\gdef\csname num@d:es:novel_10:effb0-dacisPMP:n\endcsname{5}
\expandafter\gdef\csname num@d:es:novel_10:effb0-dacisPMP:worst\endcsname{21.6}
\expandafter\gdef\csname num@d:es:novel_10:mnv2-dacisPMP\endcsname{+8.2}
\expandafter\gdef\csname num@d:es:novel_10:mnv2-dacisPMP:lo\endcsname{+3.8}
\expandafter\gdef\csname num@d:es:novel_10:mnv2-dacisPMP:hi\endcsname{+12.5}
\expandafter\gdef\csname num@d:es:novel_10:mnv2-dacisPMP:abs\endcsname{8.2}
\expandafter\gdef\csname num@d:es:novel_10:mnv2-dacisPMP:wins\endcsname{5}
\expandafter\gdef\csname num@d:es:novel_10:mnv2-dacisPMP:n\endcsname{5}
\expandafter\gdef\csname num@d:es:novel_10:mnv2-dacisPMP:worst\endcsname{12.5}
\expandafter\gdef\csname num@d:es:novel_10:dacis-l1:PMP50\endcsname{\ensuremath{-}1.5}
\expandafter\gdef\csname num@d:es:novel_10:dacis-l1:PMP50:lo\endcsname{\ensuremath{-}3.9}
\expandafter\gdef\csname num@d:es:novel_10:dacis-l1:PMP50:hi\endcsname{+0.9}
\expandafter\gdef\csname num@d:es:novel_10:dacis-l1:PMP50:abs\endcsname{1.5}
\expandafter\gdef\csname num@d:es:novel_10:dacis-l1:PMP50:wins\endcsname{1}
\expandafter\gdef\csname num@d:es:novel_10:dacis-l1:PMP50:n\endcsname{5}
\expandafter\gdef\csname num@d:es:novel_10:dacis-l1:PMP50:worst\endcsname{3.9}
\expandafter\gdef\csname num@d:es:novel_10:dacis-random:PMP50\endcsname{\ensuremath{-}0.6}
\expandafter\gdef\csname num@d:es:novel_10:dacis-random:PMP50:lo\endcsname{\ensuremath{-}1.6}
\expandafter\gdef\csname num@d:es:novel_10:dacis-random:PMP50:hi\endcsname{+0.3}
\expandafter\gdef\csname num@d:es:novel_10:dacis-random:PMP50:abs\endcsname{0.6}
\expandafter\gdef\csname num@d:es:novel_10:dacis-random:PMP50:wins\endcsname{1}
\expandafter\gdef\csname num@d:es:novel_10:dacis-random:PMP50:n\endcsname{5}
\expandafter\gdef\csname num@d:es:novel_10:dacis-random:PMP50:worst\endcsname{1.6}
\expandafter\gdef\csname num@d:es:novel_10:dense-dacisPMP50\endcsname{+3.7}
\expandafter\gdef\csname num@d:es:novel_10:dense-dacisPMP50:lo\endcsname{+2.6}
\expandafter\gdef\csname num@d:es:novel_10:dense-dacisPMP50:hi\endcsname{+4.9}
\expandafter\gdef\csname num@d:es:novel_10:dense-dacisPMP50:abs\endcsname{3.7}
\expandafter\gdef\csname num@d:es:novel_10:dense-dacisPMP50:wins\endcsname{5}
\expandafter\gdef\csname num@d:es:novel_10:dense-dacisPMP50:n\endcsname{5}
\expandafter\gdef\csname num@d:es:novel_10:dense-dacisPMP50:worst\endcsname{4.9}
\expandafter\gdef\csname num@d:es:novel_10:dacis-l1:PMP20\endcsname{\ensuremath{-}0.3}
\expandafter\gdef\csname num@d:es:novel_10:dacis-l1:PMP20:lo\endcsname{\ensuremath{-}3.6}
\expandafter\gdef\csname num@d:es:novel_10:dacis-l1:PMP20:hi\endcsname{+3.0}
\expandafter\gdef\csname num@d:es:novel_10:dacis-l1:PMP20:abs\endcsname{0.3}
\expandafter\gdef\csname num@d:es:novel_10:dacis-l1:PMP20:wins\endcsname{2}
\expandafter\gdef\csname num@d:es:novel_10:dacis-l1:PMP20:n\endcsname{5}
\expandafter\gdef\csname num@d:es:novel_10:dacis-l1:PMP20:worst\endcsname{3.6}
\expandafter\gdef\csname num@d:es:novel_10:dacis-random:PMP20\endcsname{\ensuremath{-}0.5}
\expandafter\gdef\csname num@d:es:novel_10:dacis-random:PMP20:lo\endcsname{\ensuremath{-}2.8}
\expandafter\gdef\csname num@d:es:novel_10:dacis-random:PMP20:hi\endcsname{+1.8}
\expandafter\gdef\csname num@d:es:novel_10:dacis-random:PMP20:abs\endcsname{0.5}
\expandafter\gdef\csname num@d:es:novel_10:dacis-random:PMP20:wins\endcsname{2}
\expandafter\gdef\csname num@d:es:novel_10:dacis-random:PMP20:n\endcsname{5}
\expandafter\gdef\csname num@d:es:novel_10:dacis-random:PMP20:worst\endcsname{2.8}
\expandafter\gdef\csname num@d:es:novel_10:dense-dacisPMP20\endcsname{+8.3}
\expandafter\gdef\csname num@d:es:novel_10:dense-dacisPMP20:lo\endcsname{+4.9}
\expandafter\gdef\csname num@d:es:novel_10:dense-dacisPMP20:hi\endcsname{+11.7}
\expandafter\gdef\csname num@d:es:novel_10:dense-dacisPMP20:abs\endcsname{8.3}
\expandafter\gdef\csname num@d:es:novel_10:dense-dacisPMP20:wins\endcsname{5}
\expandafter\gdef\csname num@d:es:novel_10:dense-dacisPMP20:n\endcsname{5}
\expandafter\gdef\csname num@d:es:novel_10:dense-dacisPMP20:worst\endcsname{11.7}
\expandafter\gdef\csname num@d:es:novel_10:dacis-l1:PMP10\endcsname{+0.7}
\expandafter\gdef\csname num@d:es:novel_10:dacis-l1:PMP10:lo\endcsname{\ensuremath{-}3.3}
\expandafter\gdef\csname num@d:es:novel_10:dacis-l1:PMP10:hi\endcsname{+4.8}
\expandafter\gdef\csname num@d:es:novel_10:dacis-l1:PMP10:abs\endcsname{0.7}
\expandafter\gdef\csname num@d:es:novel_10:dacis-l1:PMP10:wins\endcsname{2}
\expandafter\gdef\csname num@d:es:novel_10:dacis-l1:PMP10:n\endcsname{5}
\expandafter\gdef\csname num@d:es:novel_10:dacis-l1:PMP10:worst\endcsname{4.8}
\expandafter\gdef\csname num@d:es:novel_10:dacis-random:PMP10\endcsname{+2.2}
\expandafter\gdef\csname num@d:es:novel_10:dacis-random:PMP10:lo\endcsname{\ensuremath{-}2.0}
\expandafter\gdef\csname num@d:es:novel_10:dacis-random:PMP10:hi\endcsname{+6.5}
\expandafter\gdef\csname num@d:es:novel_10:dacis-random:PMP10:abs\endcsname{2.2}
\expandafter\gdef\csname num@d:es:novel_10:dacis-random:PMP10:wins\endcsname{4}
\expandafter\gdef\csname num@d:es:novel_10:dacis-random:PMP10:n\endcsname{5}
\expandafter\gdef\csname num@d:es:novel_10:dacis-random:PMP10:worst\endcsname{6.5}
\expandafter\gdef\csname num@d:es:novel_10:dense-dacisPMP10\endcsname{+11.3}
\expandafter\gdef\csname num@d:es:novel_10:dense-dacisPMP10:lo\endcsname{+7.0}
\expandafter\gdef\csname num@d:es:novel_10:dense-dacisPMP10:hi\endcsname{+15.6}
\expandafter\gdef\csname num@d:es:novel_10:dense-dacisPMP10:abs\endcsname{11.3}
\expandafter\gdef\csname num@d:es:novel_10:dense-dacisPMP10:wins\endcsname{5}
\expandafter\gdef\csname num@d:es:novel_10:dense-dacisPMP10:n\endcsname{5}
\expandafter\gdef\csname num@d:es:novel_10:dense-dacisPMP10:worst\endcsname{15.6}
\expandafter\gdef\csname num@d:rs:novel_10:dacis-random:PMP\endcsname{\ensuremath{-}0.4}
\expandafter\gdef\csname num@d:rs:novel_10:dacis-random:PMP:lo\endcsname{\ensuremath{-}1.0}
\expandafter\gdef\csname num@d:rs:novel_10:dacis-random:PMP:hi\endcsname{+0.1}
\expandafter\gdef\csname num@d:rs:novel_10:dacis-random:PMP:abs\endcsname{0.4}
\expandafter\gdef\csname num@d:rs:novel_10:dacis-random:PMP:wins\endcsname{1}
\expandafter\gdef\csname num@d:rs:novel_10:dacis-random:PMP:n\endcsname{5}
\expandafter\gdef\csname num@d:rs:novel_10:dacis-random:PMP:worst\endcsname{1.0}
\expandafter\gdef\csname num@d:rs:novel_10:dacis-random:P\endcsname{+0.1}
\expandafter\gdef\csname num@d:rs:novel_10:dacis-random:P:lo\endcsname{\ensuremath{-}1.1}
\expandafter\gdef\csname num@d:rs:novel_10:dacis-random:P:hi\endcsname{+1.4}
\expandafter\gdef\csname num@d:rs:novel_10:dacis-random:P:abs\endcsname{0.1}
\expandafter\gdef\csname num@d:rs:novel_10:dacis-random:P:wins\endcsname{4}
\expandafter\gdef\csname num@d:rs:novel_10:dacis-random:P:n\endcsname{5}
\expandafter\gdef\csname num@d:rs:novel_10:dacis-random:P:worst\endcsname{1.4}
\expandafter\gdef\csname num@d:rs:novel_10:dacis-l1:PMP\endcsname{\ensuremath{-}0.4}
\expandafter\gdef\csname num@d:rs:novel_10:dacis-l1:PMP:lo\endcsname{\ensuremath{-}1.0}
\expandafter\gdef\csname num@d:rs:novel_10:dacis-l1:PMP:hi\endcsname{+0.3}
\expandafter\gdef\csname num@d:rs:novel_10:dacis-l1:PMP:abs\endcsname{0.4}
\expandafter\gdef\csname num@d:rs:novel_10:dacis-l1:PMP:wins\endcsname{1}
\expandafter\gdef\csname num@d:rs:novel_10:dacis-l1:PMP:n\endcsname{5}
\expandafter\gdef\csname num@d:rs:novel_10:dacis-l1:PMP:worst\endcsname{1.0}
\expandafter\gdef\csname num@d:rs:novel_10:dacis-l1:P\endcsname{\ensuremath{-}0.6}
\expandafter\gdef\csname num@d:rs:novel_10:dacis-l1:P:lo\endcsname{\ensuremath{-}1.6}
\expandafter\gdef\csname num@d:rs:novel_10:dacis-l1:P:hi\endcsname{+0.4}
\expandafter\gdef\csname num@d:rs:novel_10:dacis-l1:P:abs\endcsname{0.6}
\expandafter\gdef\csname num@d:rs:novel_10:dacis-l1:P:wins\endcsname{2}
\expandafter\gdef\csname num@d:rs:novel_10:dacis-l1:P:n\endcsname{5}
\expandafter\gdef\csname num@d:rs:novel_10:dacis-l1:P:worst\endcsname{1.6}
\expandafter\gdef\csname num@d:rs:novel_10:dacis-taylor:PMP\endcsname{\ensuremath{-}1.4}
\expandafter\gdef\csname num@d:rs:novel_10:dacis-taylor:PMP:lo\endcsname{\ensuremath{-}1.9}
\expandafter\gdef\csname num@d:rs:novel_10:dacis-taylor:PMP:hi\endcsname{\ensuremath{-}1.0}
\expandafter\gdef\csname num@d:rs:novel_10:dacis-taylor:PMP:abs\endcsname{1.4}
\expandafter\gdef\csname num@d:rs:novel_10:dacis-taylor:PMP:wins\endcsname{0}
\expandafter\gdef\csname num@d:rs:novel_10:dacis-taylor:PMP:n\endcsname{5}
\expandafter\gdef\csname num@d:rs:novel_10:dacis-taylor:PMP:worst\endcsname{1.9}
\expandafter\gdef\csname num@d:rs:novel_10:dacis-taylor:P\endcsname{\ensuremath{-}2.2}
\expandafter\gdef\csname num@d:rs:novel_10:dacis-taylor:P:lo\endcsname{\ensuremath{-}2.7}
\expandafter\gdef\csname num@d:rs:novel_10:dacis-taylor:P:hi\endcsname{\ensuremath{-}1.8}
\expandafter\gdef\csname num@d:rs:novel_10:dacis-taylor:P:abs\endcsname{2.2}
\expandafter\gdef\csname num@d:rs:novel_10:dacis-taylor:P:wins\endcsname{0}
\expandafter\gdef\csname num@d:rs:novel_10:dacis-taylor:P:n\endcsname{5}
\expandafter\gdef\csname num@d:rs:novel_10:dacis-taylor:P:worst\endcsname{2.7}
\expandafter\gdef\csname num@d:rs:novel_10:dacis-bn_gamma:PMP\endcsname{+0.9}
\expandafter\gdef\csname num@d:rs:novel_10:dacis-bn_gamma:PMP:lo\endcsname{\ensuremath{-}0.1}
\expandafter\gdef\csname num@d:rs:novel_10:dacis-bn_gamma:PMP:hi\endcsname{+2.0}
\expandafter\gdef\csname num@d:rs:novel_10:dacis-bn_gamma:PMP:abs\endcsname{0.9}
\expandafter\gdef\csname num@d:rs:novel_10:dacis-bn_gamma:PMP:wins\endcsname{5}
\expandafter\gdef\csname num@d:rs:novel_10:dacis-bn_gamma:PMP:n\endcsname{5}
\expandafter\gdef\csname num@d:rs:novel_10:dacis-bn_gamma:PMP:worst\endcsname{2.0}
\expandafter\gdef\csname num@d:rs:novel_10:dacis-D:PMP\endcsname{+1.0}
\expandafter\gdef\csname num@d:rs:novel_10:dacis-D:PMP:lo\endcsname{+0.4}
\expandafter\gdef\csname num@d:rs:novel_10:dacis-D:PMP:hi\endcsname{+1.7}
\expandafter\gdef\csname num@d:rs:novel_10:dacis-D:PMP:abs\endcsname{1.0}
\expandafter\gdef\csname num@d:rs:novel_10:dacis-D:PMP:wins\endcsname{5}
\expandafter\gdef\csname num@d:rs:novel_10:dacis-D:PMP:n\endcsname{5}
\expandafter\gdef\csname num@d:rs:novel_10:dacis-D:PMP:worst\endcsname{1.7}
\expandafter\gdef\csname num@d:rs:novel_10:dacis-D:P\endcsname{+0.6}
\expandafter\gdef\csname num@d:rs:novel_10:dacis-D:P:lo\endcsname{\ensuremath{-}0.3}
\expandafter\gdef\csname num@d:rs:novel_10:dacis-D:P:hi\endcsname{+1.4}
\expandafter\gdef\csname num@d:rs:novel_10:dacis-D:P:abs\endcsname{0.6}
\expandafter\gdef\csname num@d:rs:novel_10:dacis-D:P:wins\endcsname{4}
\expandafter\gdef\csname num@d:rs:novel_10:dacis-D:P:n\endcsname{5}
\expandafter\gdef\csname num@d:rs:novel_10:dacis-D:P:worst\endcsname{1.4}
\expandafter\gdef\csname num@d:rs:novel_10:dacis-G:PMP\endcsname{\ensuremath{-}0.5}
\expandafter\gdef\csname num@d:rs:novel_10:dacis-G:PMP:lo\endcsname{\ensuremath{-}1.3}
\expandafter\gdef\csname num@d:rs:novel_10:dacis-G:PMP:hi\endcsname{+0.4}
\expandafter\gdef\csname num@d:rs:novel_10:dacis-G:PMP:abs\endcsname{0.5}
\expandafter\gdef\csname num@d:rs:novel_10:dacis-G:PMP:wins\endcsname{1}
\expandafter\gdef\csname num@d:rs:novel_10:dacis-G:PMP:n\endcsname{5}
\expandafter\gdef\csname num@d:rs:novel_10:dacis-G:PMP:worst\endcsname{1.3}
\expandafter\gdef\csname num@d:rs:novel_10:dacis-V:PMP\endcsname{\ensuremath{-}0.4}
\expandafter\gdef\csname num@d:rs:novel_10:dacis-V:PMP:lo\endcsname{\ensuremath{-}1.2}
\expandafter\gdef\csname num@d:rs:novel_10:dacis-V:PMP:hi\endcsname{+0.4}
\expandafter\gdef\csname num@d:rs:novel_10:dacis-V:PMP:abs\endcsname{0.4}
\expandafter\gdef\csname num@d:rs:novel_10:dacis-V:PMP:wins\endcsname{1}
\expandafter\gdef\csname num@d:rs:novel_10:dacis-V:PMP:n\endcsname{5}
\expandafter\gdef\csname num@d:rs:novel_10:dacis-V:PMP:worst\endcsname{1.2}
\expandafter\gdef\csname num@d:rs:novel_10:l1-random:PMP\endcsname{\ensuremath{-}0.1}
\expandafter\gdef\csname num@d:rs:novel_10:l1-random:PMP:lo\endcsname{\ensuremath{-}0.5}
\expandafter\gdef\csname num@d:rs:novel_10:l1-random:PMP:hi\endcsname{+0.3}
\expandafter\gdef\csname num@d:rs:novel_10:l1-random:PMP:abs\endcsname{0.1}
\expandafter\gdef\csname num@d:rs:novel_10:l1-random:PMP:wins\endcsname{2}
\expandafter\gdef\csname num@d:rs:novel_10:l1-random:PMP:n\endcsname{5}
\expandafter\gdef\csname num@d:rs:novel_10:l1-random:PMP:worst\endcsname{0.5}
\expandafter\gdef\csname num@d:rs:novel_10:l1-random:P\endcsname{+0.7}
\expandafter\gdef\csname num@d:rs:novel_10:l1-random:P:lo\endcsname{\ensuremath{-}0.2}
\expandafter\gdef\csname num@d:rs:novel_10:l1-random:P:hi\endcsname{+1.7}
\expandafter\gdef\csname num@d:rs:novel_10:l1-random:P:abs\endcsname{0.7}
\expandafter\gdef\csname num@d:rs:novel_10:l1-random:P:wins\endcsname{4}
\expandafter\gdef\csname num@d:rs:novel_10:l1-random:P:n\endcsname{5}
\expandafter\gdef\csname num@d:rs:novel_10:l1-random:P:worst\endcsname{1.7}
\expandafter\gdef\csname num@d:rs:novel_10:taylor-random:PMP\endcsname{+1.0}
\expandafter\gdef\csname num@d:rs:novel_10:taylor-random:PMP:lo\endcsname{+0.3}
\expandafter\gdef\csname num@d:rs:novel_10:taylor-random:PMP:hi\endcsname{+1.7}
\expandafter\gdef\csname num@d:rs:novel_10:taylor-random:PMP:abs\endcsname{1.0}
\expandafter\gdef\csname num@d:rs:novel_10:taylor-random:PMP:wins\endcsname{5}
\expandafter\gdef\csname num@d:rs:novel_10:taylor-random:PMP:n\endcsname{5}
\expandafter\gdef\csname num@d:rs:novel_10:taylor-random:PMP:worst\endcsname{1.7}
\expandafter\gdef\csname num@d:rs:novel_10:taylor-random:P\endcsname{+2.4}
\expandafter\gdef\csname num@d:rs:novel_10:taylor-random:P:lo\endcsname{+1.1}
\expandafter\gdef\csname num@d:rs:novel_10:taylor-random:P:hi\endcsname{+3.6}
\expandafter\gdef\csname num@d:rs:novel_10:taylor-random:P:abs\endcsname{2.4}
\expandafter\gdef\csname num@d:rs:novel_10:taylor-random:P:wins\endcsname{5}
\expandafter\gdef\csname num@d:rs:novel_10:taylor-random:P:n\endcsname{5}
\expandafter\gdef\csname num@d:rs:novel_10:taylor-random:P:worst\endcsname{3.6}
\expandafter\gdef\csname num@d:rs:novel_10:bn_gamma-random:PMP\endcsname{\ensuremath{-}1.4}
\expandafter\gdef\csname num@d:rs:novel_10:bn_gamma-random:PMP:lo\endcsname{\ensuremath{-}1.9}
\expandafter\gdef\csname num@d:rs:novel_10:bn_gamma-random:PMP:hi\endcsname{\ensuremath{-}0.8}
\expandafter\gdef\csname num@d:rs:novel_10:bn_gamma-random:PMP:abs\endcsname{1.4}
\expandafter\gdef\csname num@d:rs:novel_10:bn_gamma-random:PMP:wins\endcsname{0}
\expandafter\gdef\csname num@d:rs:novel_10:bn_gamma-random:PMP:n\endcsname{5}
\expandafter\gdef\csname num@d:rs:novel_10:bn_gamma-random:PMP:worst\endcsname{1.9}
\expandafter\gdef\csname num@d:rs:novel_10:D-random:PMP\endcsname{\ensuremath{-}1.5}
\expandafter\gdef\csname num@d:rs:novel_10:D-random:PMP:lo\endcsname{\ensuremath{-}2.4}
\expandafter\gdef\csname num@d:rs:novel_10:D-random:PMP:hi\endcsname{\ensuremath{-}0.5}
\expandafter\gdef\csname num@d:rs:novel_10:D-random:PMP:abs\endcsname{1.5}
\expandafter\gdef\csname num@d:rs:novel_10:D-random:PMP:wins\endcsname{0}
\expandafter\gdef\csname num@d:rs:novel_10:D-random:PMP:n\endcsname{5}
\expandafter\gdef\csname num@d:rs:novel_10:D-random:PMP:worst\endcsname{2.4}
\expandafter\gdef\csname num@d:rs:novel_10:D-random:P\endcsname{\ensuremath{-}0.4}
\expandafter\gdef\csname num@d:rs:novel_10:D-random:P:lo\endcsname{\ensuremath{-}1.2}
\expandafter\gdef\csname num@d:rs:novel_10:D-random:P:hi\endcsname{+0.4}
\expandafter\gdef\csname num@d:rs:novel_10:D-random:P:abs\endcsname{0.4}
\expandafter\gdef\csname num@d:rs:novel_10:D-random:P:wins\endcsname{2}
\expandafter\gdef\csname num@d:rs:novel_10:D-random:P:n\endcsname{5}
\expandafter\gdef\csname num@d:rs:novel_10:D-random:P:worst\endcsname{1.2}
\expandafter\gdef\csname num@d:rs:novel_10:G-random:PMP\endcsname{+0.0}
\expandafter\gdef\csname num@d:rs:novel_10:G-random:PMP:lo\endcsname{\ensuremath{-}0.7}
\expandafter\gdef\csname num@d:rs:novel_10:G-random:PMP:hi\endcsname{+0.7}
\expandafter\gdef\csname num@d:rs:novel_10:G-random:PMP:abs\endcsname{0.0}
\expandafter\gdef\csname num@d:rs:novel_10:G-random:PMP:wins\endcsname{3}
\expandafter\gdef\csname num@d:rs:novel_10:G-random:PMP:n\endcsname{5}
\expandafter\gdef\csname num@d:rs:novel_10:G-random:PMP:worst\endcsname{0.7}
\expandafter\gdef\csname num@d:rs:novel_10:V-random:PMP\endcsname{\ensuremath{-}0.0}
\expandafter\gdef\csname num@d:rs:novel_10:V-random:PMP:lo\endcsname{\ensuremath{-}0.7}
\expandafter\gdef\csname num@d:rs:novel_10:V-random:PMP:hi\endcsname{+0.7}
\expandafter\gdef\csname num@d:rs:novel_10:V-random:PMP:abs\endcsname{0.0}
\expandafter\gdef\csname num@d:rs:novel_10:V-random:PMP:wins\endcsname{2}
\expandafter\gdef\csname num@d:rs:novel_10:V-random:PMP:n\endcsname{5}
\expandafter\gdef\csname num@d:rs:novel_10:V-random:PMP:worst\endcsname{0.7}
\expandafter\gdef\csname num@d:rs:novel_10:PM-P\endcsname{\ensuremath{-}0.2}
\expandafter\gdef\csname num@d:rs:novel_10:PM-P:lo\endcsname{\ensuremath{-}0.9}
\expandafter\gdef\csname num@d:rs:novel_10:PM-P:hi\endcsname{+0.6}
\expandafter\gdef\csname num@d:rs:novel_10:PM-P:abs\endcsname{0.2}
\expandafter\gdef\csname num@d:rs:novel_10:PM-P:wins\endcsname{2}
\expandafter\gdef\csname num@d:rs:novel_10:PM-P:n\endcsname{5}
\expandafter\gdef\csname num@d:rs:novel_10:PM-P:worst\endcsname{0.9}
\expandafter\gdef\csname num@d:rs:novel_10:MP-P\endcsname{+0.6}
\expandafter\gdef\csname num@d:rs:novel_10:MP-P:lo\endcsname{+0.0}
\expandafter\gdef\csname num@d:rs:novel_10:MP-P:hi\endcsname{+1.2}
\expandafter\gdef\csname num@d:rs:novel_10:MP-P:abs\endcsname{0.6}
\expandafter\gdef\csname num@d:rs:novel_10:MP-P:wins\endcsname{5}
\expandafter\gdef\csname num@d:rs:novel_10:MP-P:n\endcsname{5}
\expandafter\gdef\csname num@d:rs:novel_10:MP-P:worst\endcsname{1.2}
\expandafter\gdef\csname num@d:rs:novel_10:PGP-P\endcsname{+0.9}
\expandafter\gdef\csname num@d:rs:novel_10:PGP-P:lo\endcsname{+0.6}
\expandafter\gdef\csname num@d:rs:novel_10:PGP-P:hi\endcsname{+1.2}
\expandafter\gdef\csname num@d:rs:novel_10:PGP-P:abs\endcsname{0.9}
\expandafter\gdef\csname num@d:rs:novel_10:PGP-P:wins\endcsname{5}
\expandafter\gdef\csname num@d:rs:novel_10:PGP-P:n\endcsname{5}
\expandafter\gdef\csname num@d:rs:novel_10:PGP-P:worst\endcsname{1.2}
\expandafter\gdef\csname num@d:rs:novel_10:PMP-P\endcsname{+0.7}
\expandafter\gdef\csname num@d:rs:novel_10:PMP-P:lo\endcsname{+0.2}
\expandafter\gdef\csname num@d:rs:novel_10:PMP-P:hi\endcsname{+1.3}
\expandafter\gdef\csname num@d:rs:novel_10:PMP-P:abs\endcsname{0.7}
\expandafter\gdef\csname num@d:rs:novel_10:PMP-P:wins\endcsname{5}
\expandafter\gdef\csname num@d:rs:novel_10:PMP-P:n\endcsname{5}
\expandafter\gdef\csname num@d:rs:novel_10:PMP-P:worst\endcsname{1.3}
\expandafter\gdef\csname num@d:rs:novel_10:densem-dense\endcsname{\ensuremath{-}0.4}
\expandafter\gdef\csname num@d:rs:novel_10:densem-dense:lo\endcsname{\ensuremath{-}0.5}
\expandafter\gdef\csname num@d:rs:novel_10:densem-dense:hi\endcsname{\ensuremath{-}0.2}
\expandafter\gdef\csname num@d:rs:novel_10:densem-dense:abs\endcsname{0.4}
\expandafter\gdef\csname num@d:rs:novel_10:densem-dense:wins\endcsname{0}
\expandafter\gdef\csname num@d:rs:novel_10:densem-dense:n\endcsname{5}
\expandafter\gdef\csname num@d:rs:novel_10:densem-dense:worst\endcsname{0.5}
\expandafter\gdef\csname num@d:rs:novel_10:densermr-dense\endcsname{\ensuremath{-}0.2}
\expandafter\gdef\csname num@d:rs:novel_10:densermr-dense:lo\endcsname{\ensuremath{-}0.3}
\expandafter\gdef\csname num@d:rs:novel_10:densermr-dense:hi\endcsname{\ensuremath{-}0.0}
\expandafter\gdef\csname num@d:rs:novel_10:densermr-dense:abs\endcsname{0.2}
\expandafter\gdef\csname num@d:rs:novel_10:densermr-dense:wins\endcsname{0}
\expandafter\gdef\csname num@d:rs:novel_10:densermr-dense:n\endcsname{5}
\expandafter\gdef\csname num@d:rs:novel_10:densermr-dense:worst\endcsname{0.3}
\expandafter\gdef\csname num@d:rs:novel_10:dense-dacisP\endcsname{+6.5}
\expandafter\gdef\csname num@d:rs:novel_10:dense-dacisP:lo\endcsname{+5.1}
\expandafter\gdef\csname num@d:rs:novel_10:dense-dacisP:hi\endcsname{+7.9}
\expandafter\gdef\csname num@d:rs:novel_10:dense-dacisP:abs\endcsname{6.5}
\expandafter\gdef\csname num@d:rs:novel_10:dense-dacisP:wins\endcsname{5}
\expandafter\gdef\csname num@d:rs:novel_10:dense-dacisP:n\endcsname{5}
\expandafter\gdef\csname num@d:rs:novel_10:dense-dacisP:worst\endcsname{7.9}
\expandafter\gdef\csname num@d:rs:novel_10:dense-dacisPMP\endcsname{+5.7}
\expandafter\gdef\csname num@d:rs:novel_10:dense-dacisPMP:lo\endcsname{+4.5}
\expandafter\gdef\csname num@d:rs:novel_10:dense-dacisPMP:hi\endcsname{+7.0}
\expandafter\gdef\csname num@d:rs:novel_10:dense-dacisPMP:abs\endcsname{5.7}
\expandafter\gdef\csname num@d:rs:novel_10:dense-dacisPMP:wins\endcsname{5}
\expandafter\gdef\csname num@d:rs:novel_10:dense-dacisPMP:n\endcsname{5}
\expandafter\gdef\csname num@d:rs:novel_10:dense-dacisPMP:worst\endcsname{7.0}
\expandafter\gdef\csname num@d:rs:novel_10:dense-randomP\endcsname{+6.6}
\expandafter\gdef\csname num@d:rs:novel_10:dense-randomP:lo\endcsname{+4.5}
\expandafter\gdef\csname num@d:rs:novel_10:dense-randomP:hi\endcsname{+8.8}
\expandafter\gdef\csname num@d:rs:novel_10:dense-randomP:abs\endcsname{6.6}
\expandafter\gdef\csname num@d:rs:novel_10:dense-randomP:wins\endcsname{5}
\expandafter\gdef\csname num@d:rs:novel_10:dense-randomP:n\endcsname{5}
\expandafter\gdef\csname num@d:rs:novel_10:dense-randomP:worst\endcsname{8.8}
\expandafter\gdef\csname num@d:rs:novel_10:dense-randomPMP\endcsname{+5.3}
\expandafter\gdef\csname num@d:rs:novel_10:dense-randomPMP:lo\endcsname{+3.9}
\expandafter\gdef\csname num@d:rs:novel_10:dense-randomPMP:hi\endcsname{+6.7}
\expandafter\gdef\csname num@d:rs:novel_10:dense-randomPMP:abs\endcsname{5.3}
\expandafter\gdef\csname num@d:rs:novel_10:dense-randomPMP:wins\endcsname{5}
\expandafter\gdef\csname num@d:rs:novel_10:dense-randomPMP:n\endcsname{5}
\expandafter\gdef\csname num@d:rs:novel_10:dense-randomPMP:worst\endcsname{6.7}
\expandafter\gdef\csname num@d:rs:novel_10:mnv3-dacisPMP\endcsname{+6.8}
\expandafter\gdef\csname num@d:rs:novel_10:mnv3-dacisPMP:lo\endcsname{+5.6}
\expandafter\gdef\csname num@d:rs:novel_10:mnv3-dacisPMP:hi\endcsname{+8.1}
\expandafter\gdef\csname num@d:rs:novel_10:mnv3-dacisPMP:abs\endcsname{6.8}
\expandafter\gdef\csname num@d:rs:novel_10:mnv3-dacisPMP:wins\endcsname{5}
\expandafter\gdef\csname num@d:rs:novel_10:mnv3-dacisPMP:n\endcsname{5}
\expandafter\gdef\csname num@d:rs:novel_10:mnv3-dacisPMP:worst\endcsname{8.1}
\expandafter\gdef\csname num@d:rs:novel_10:effb0-dacisPMP\endcsname{+9.0}
\expandafter\gdef\csname num@d:rs:novel_10:effb0-dacisPMP:lo\endcsname{+6.7}
\expandafter\gdef\csname num@d:rs:novel_10:effb0-dacisPMP:hi\endcsname{+11.2}
\expandafter\gdef\csname num@d:rs:novel_10:effb0-dacisPMP:abs\endcsname{9.0}
\expandafter\gdef\csname num@d:rs:novel_10:effb0-dacisPMP:wins\endcsname{5}
\expandafter\gdef\csname num@d:rs:novel_10:effb0-dacisPMP:n\endcsname{5}
\expandafter\gdef\csname num@d:rs:novel_10:effb0-dacisPMP:worst\endcsname{11.2}
\expandafter\gdef\csname num@d:rs:novel_10:mnv2-dacisPMP\endcsname{+5.5}
\expandafter\gdef\csname num@d:rs:novel_10:mnv2-dacisPMP:lo\endcsname{+4.0}
\expandafter\gdef\csname num@d:rs:novel_10:mnv2-dacisPMP:hi\endcsname{+7.0}
\expandafter\gdef\csname num@d:rs:novel_10:mnv2-dacisPMP:abs\endcsname{5.5}
\expandafter\gdef\csname num@d:rs:novel_10:mnv2-dacisPMP:wins\endcsname{5}
\expandafter\gdef\csname num@d:rs:novel_10:mnv2-dacisPMP:n\endcsname{5}
\expandafter\gdef\csname num@d:rs:novel_10:mnv2-dacisPMP:worst\endcsname{7.0}
\expandafter\gdef\csname num@d:rs:novel_10:dacis-l1:PMP50\endcsname{\ensuremath{-}0.1}
\expandafter\gdef\csname num@d:rs:novel_10:dacis-l1:PMP50:lo\endcsname{\ensuremath{-}0.9}
\expandafter\gdef\csname num@d:rs:novel_10:dacis-l1:PMP50:hi\endcsname{+0.6}
\expandafter\gdef\csname num@d:rs:novel_10:dacis-l1:PMP50:abs\endcsname{0.1}
\expandafter\gdef\csname num@d:rs:novel_10:dacis-l1:PMP50:wins\endcsname{1}
\expandafter\gdef\csname num@d:rs:novel_10:dacis-l1:PMP50:n\endcsname{5}
\expandafter\gdef\csname num@d:rs:novel_10:dacis-l1:PMP50:worst\endcsname{0.9}
\expandafter\gdef\csname num@d:rs:novel_10:dacis-random:PMP50\endcsname{\ensuremath{-}0.3}
\expandafter\gdef\csname num@d:rs:novel_10:dacis-random:PMP50:lo\endcsname{\ensuremath{-}1.4}
\expandafter\gdef\csname num@d:rs:novel_10:dacis-random:PMP50:hi\endcsname{+0.8}
\expandafter\gdef\csname num@d:rs:novel_10:dacis-random:PMP50:abs\endcsname{0.3}
\expandafter\gdef\csname num@d:rs:novel_10:dacis-random:PMP50:wins\endcsname{2}
\expandafter\gdef\csname num@d:rs:novel_10:dacis-random:PMP50:n\endcsname{5}
\expandafter\gdef\csname num@d:rs:novel_10:dacis-random:PMP50:worst\endcsname{1.4}
\expandafter\gdef\csname num@d:rs:novel_10:dense-dacisPMP50\endcsname{+2.6}
\expandafter\gdef\csname num@d:rs:novel_10:dense-dacisPMP50:lo\endcsname{+2.1}
\expandafter\gdef\csname num@d:rs:novel_10:dense-dacisPMP50:hi\endcsname{+3.2}
\expandafter\gdef\csname num@d:rs:novel_10:dense-dacisPMP50:abs\endcsname{2.6}
\expandafter\gdef\csname num@d:rs:novel_10:dense-dacisPMP50:wins\endcsname{5}
\expandafter\gdef\csname num@d:rs:novel_10:dense-dacisPMP50:n\endcsname{5}
\expandafter\gdef\csname num@d:rs:novel_10:dense-dacisPMP50:worst\endcsname{3.2}
\expandafter\gdef\csname num@d:rs:novel_10:dacis-l1:PMP20\endcsname{\ensuremath{-}1.5}
\expandafter\gdef\csname num@d:rs:novel_10:dacis-l1:PMP20:lo\endcsname{\ensuremath{-}3.3}
\expandafter\gdef\csname num@d:rs:novel_10:dacis-l1:PMP20:hi\endcsname{+0.3}
\expandafter\gdef\csname num@d:rs:novel_10:dacis-l1:PMP20:abs\endcsname{1.5}
\expandafter\gdef\csname num@d:rs:novel_10:dacis-l1:PMP20:wins\endcsname{1}
\expandafter\gdef\csname num@d:rs:novel_10:dacis-l1:PMP20:n\endcsname{5}
\expandafter\gdef\csname num@d:rs:novel_10:dacis-l1:PMP20:worst\endcsname{3.3}
\expandafter\gdef\csname num@d:rs:novel_10:dacis-random:PMP20\endcsname{+0.5}
\expandafter\gdef\csname num@d:rs:novel_10:dacis-random:PMP20:lo\endcsname{\ensuremath{-}0.7}
\expandafter\gdef\csname num@d:rs:novel_10:dacis-random:PMP20:hi\endcsname{+1.6}
\expandafter\gdef\csname num@d:rs:novel_10:dacis-random:PMP20:abs\endcsname{0.5}
\expandafter\gdef\csname num@d:rs:novel_10:dacis-random:PMP20:wins\endcsname{3}
\expandafter\gdef\csname num@d:rs:novel_10:dacis-random:PMP20:n\endcsname{5}
\expandafter\gdef\csname num@d:rs:novel_10:dacis-random:PMP20:worst\endcsname{1.6}
\expandafter\gdef\csname num@d:rs:novel_10:dense-dacisPMP20\endcsname{+8.6}
\expandafter\gdef\csname num@d:rs:novel_10:dense-dacisPMP20:lo\endcsname{+5.8}
\expandafter\gdef\csname num@d:rs:novel_10:dense-dacisPMP20:hi\endcsname{+11.4}
\expandafter\gdef\csname num@d:rs:novel_10:dense-dacisPMP20:abs\endcsname{8.6}
\expandafter\gdef\csname num@d:rs:novel_10:dense-dacisPMP20:wins\endcsname{5}
\expandafter\gdef\csname num@d:rs:novel_10:dense-dacisPMP20:n\endcsname{5}
\expandafter\gdef\csname num@d:rs:novel_10:dense-dacisPMP20:worst\endcsname{11.4}
\expandafter\gdef\csname num@d:rs:novel_10:dacis-l1:PMP10\endcsname{+1.2}
\expandafter\gdef\csname num@d:rs:novel_10:dacis-l1:PMP10:lo\endcsname{+0.2}
\expandafter\gdef\csname num@d:rs:novel_10:dacis-l1:PMP10:hi\endcsname{+2.2}
\expandafter\gdef\csname num@d:rs:novel_10:dacis-l1:PMP10:abs\endcsname{1.2}
\expandafter\gdef\csname num@d:rs:novel_10:dacis-l1:PMP10:wins\endcsname{5}
\expandafter\gdef\csname num@d:rs:novel_10:dacis-l1:PMP10:n\endcsname{5}
\expandafter\gdef\csname num@d:rs:novel_10:dacis-l1:PMP10:worst\endcsname{2.2}
\expandafter\gdef\csname num@d:rs:novel_10:dacis-random:PMP10\endcsname{+3.3}
\expandafter\gdef\csname num@d:rs:novel_10:dacis-random:PMP10:lo\endcsname{\ensuremath{-}0.7}
\expandafter\gdef\csname num@d:rs:novel_10:dacis-random:PMP10:hi\endcsname{+7.3}
\expandafter\gdef\csname num@d:rs:novel_10:dacis-random:PMP10:abs\endcsname{3.3}
\expandafter\gdef\csname num@d:rs:novel_10:dacis-random:PMP10:wins\endcsname{5}
\expandafter\gdef\csname num@d:rs:novel_10:dacis-random:PMP10:n\endcsname{5}
\expandafter\gdef\csname num@d:rs:novel_10:dacis-random:PMP10:worst\endcsname{7.3}
\expandafter\gdef\csname num@d:rs:novel_10:dense-dacisPMP10\endcsname{+14.1}
\expandafter\gdef\csname num@d:rs:novel_10:dense-dacisPMP10:lo\endcsname{+11.8}
\expandafter\gdef\csname num@d:rs:novel_10:dense-dacisPMP10:hi\endcsname{+16.4}
\expandafter\gdef\csname num@d:rs:novel_10:dense-dacisPMP10:abs\endcsname{14.1}
\expandafter\gdef\csname num@d:rs:novel_10:dense-dacisPMP10:wins\endcsname{5}
\expandafter\gdef\csname num@d:rs:novel_10:dense-dacisPMP10:n\endcsname{5}
\expandafter\gdef\csname num@d:rs:novel_10:dense-dacisPMP10:worst\endcsname{16.4}
\expandafter\gdef\csname num@spread:rs:novel_5:P\endcsname{2.9}
\expandafter\gdef\csname num@best:rs:novel_5:P\endcsname{Taylor}
\expandafter\gdef\csname num@worst:rs:novel_5:P\endcsname{Fisher}
\expandafter\gdef\csname num@bestv:rs:novel_5:P\endcsname{82.2}
\expandafter\gdef\csname num@worstv:rs:novel_5:P\endcsname{79.3}
\expandafter\gdef\csname num@nscores:rs:novel_5:P\endcsname{5}
\expandafter\gdef\csname num@partsd:rs:novel_5:P\endcsname{4.3}
\expandafter\gdef\csname num@kw:rs:novel_5:P\endcsname{0.74}
\expandafter\gdef\csname num@kwp:rs:novel_5:P\endcsname{0.01}
\expandafter\gdef\csname num@randrank:rs:novel_5:P\endcsname{4}
\expandafter\gdef\csname num@spread:rs:novel_5:PMP\endcsname{2.7}
\expandafter\gdef\csname num@best:rs:novel_5:PMP\endcsname{Taylor}
\expandafter\gdef\csname num@worst:rs:novel_5:PMP\endcsname{Fisher}
\expandafter\gdef\csname num@bestv:rs:novel_5:PMP\endcsname{82.2}
\expandafter\gdef\csname num@worstv:rs:novel_5:PMP\endcsname{79.5}
\expandafter\gdef\csname num@nscores:rs:novel_5:PMP\endcsname{8}
\expandafter\gdef\csname num@partsd:rs:novel_5:PMP\endcsname{4.1}
\expandafter\gdef\csname num@kw:rs:novel_5:PMP\endcsname{0.68}
\expandafter\gdef\csname num@kwp:rs:novel_5:PMP\endcsname{0.00}
\expandafter\gdef\csname num@randrank:rs:novel_5:PMP\endcsname{3}
\expandafter\gdef\csname num@spread:es:novel_5:P\endcsname{2.5}
\expandafter\gdef\csname num@best:es:novel_5:P\endcsname{Taylor}
\expandafter\gdef\csname num@worst:es:novel_5:P\endcsname{random}
\expandafter\gdef\csname num@bestv:es:novel_5:P\endcsname{58.1}
\expandafter\gdef\csname num@worstv:es:novel_5:P\endcsname{55.6}
\expandafter\gdef\csname num@nscores:es:novel_5:P\endcsname{5}
\expandafter\gdef\csname num@partsd:es:novel_5:P\endcsname{6.8}
\expandafter\gdef\csname num@kw:es:novel_5:P\endcsname{0.27}
\expandafter\gdef\csname num@kwp:es:novel_5:P\endcsname{0.25}
\expandafter\gdef\csname num@randrank:es:novel_5:P\endcsname{5}
\expandafter\gdef\csname num@spread:es:novel_5:PMP\endcsname{2.6}
\expandafter\gdef\csname num@best:es:novel_5:PMP\endcsname{random}
\expandafter\gdef\csname num@worst:es:novel_5:PMP\endcsname{BN scale}
\expandafter\gdef\csname num@bestv:es:novel_5:PMP\endcsname{57.6}
\expandafter\gdef\csname num@worstv:es:novel_5:PMP\endcsname{55.0}
\expandafter\gdef\csname num@nscores:es:novel_5:PMP\endcsname{8}
\expandafter\gdef\csname num@partsd:es:novel_5:PMP\endcsname{6.1}
\expandafter\gdef\csname num@kw:es:novel_5:PMP\endcsname{0.22}
\expandafter\gdef\csname num@kwp:es:novel_5:PMP\endcsname{0.36}
\expandafter\gdef\csname num@randrank:es:novel_5:PMP\endcsname{1}
\expandafter\gdef\csname num@spread:pv:novel_5:P\endcsname{3.3}
\expandafter\gdef\csname num@best:pv:novel_5:P\endcsname{Taylor}
\expandafter\gdef\csname num@worst:pv:novel_5:P\endcsname{BN scale}
\expandafter\gdef\csname num@bestv:pv:novel_5:P\endcsname{88.3}
\expandafter\gdef\csname num@worstv:pv:novel_5:P\endcsname{85.1}
\expandafter\gdef\csname num@nscores:pv:novel_5:P\endcsname{8}
\expandafter\gdef\csname num@partsd:pv:novel_5:P\endcsname{7.8}
\expandafter\gdef\csname num@kw:pv:novel_5:P\endcsname{0.52}
\expandafter\gdef\csname num@kwp:pv:novel_5:P\endcsname{0.01}
\expandafter\gdef\csname num@randrank:pv:novel_5:P\endcsname{6}
\expandafter\gdef\csname num@spread:pv:novel_5:PMP\endcsname{2.0}
\expandafter\gdef\csname num@best:pv:novel_5:PMP\endcsname{Taylor}
\expandafter\gdef\csname num@worst:pv:novel_5:PMP\endcsname{BN scale}
\expandafter\gdef\csname num@bestv:pv:novel_5:PMP\endcsname{87.1}
\expandafter\gdef\csname num@worstv:pv:novel_5:PMP\endcsname{85.1}
\expandafter\gdef\csname num@nscores:pv:novel_5:PMP\endcsname{8}
\expandafter\gdef\csname num@partsd:pv:novel_5:PMP\endcsname{7.7}
\expandafter\gdef\csname num@kw:pv:novel_5:PMP\endcsname{0.24}
\expandafter\gdef\csname num@kwp:pv:novel_5:PMP\endcsname{0.30}
\expandafter\gdef\csname num@randrank:pv:novel_5:PMP\endcsname{2}
\expandafter\gdef\csname num@spread:pd:plantdoc_5:P\endcsname{4.3}
\expandafter\gdef\csname num@best:pd:plantdoc_5:P\endcsname{filter L1}
\expandafter\gdef\csname num@worst:pd:plantdoc_5:P\endcsname{BN scale}
\expandafter\gdef\csname num@bestv:pd:plantdoc_5:P\endcsname{48.7}
\expandafter\gdef\csname num@worstv:pd:plantdoc_5:P\endcsname{44.4}
\expandafter\gdef\csname num@nscores:pd:plantdoc_5:P\endcsname{8}
\expandafter\gdef\csname num@partsd:pd:plantdoc_5:P\endcsname{1.1}
\expandafter\gdef\csname num@kw:pd:plantdoc_5:P\endcsname{0.58}
\expandafter\gdef\csname num@kwp:pd:plantdoc_5:P\endcsname{0.01}
\expandafter\gdef\csname num@randrank:pd:plantdoc_5:P\endcsname{7}
\expandafter\gdef\csname num@spread:pd:plantdoc_5:PMP\endcsname{3.0}
\expandafter\gdef\csname num@best:pd:plantdoc_5:PMP\endcsname{Taylor}
\expandafter\gdef\csname num@worst:pd:plantdoc_5:PMP\endcsname{BN scale}
\expandafter\gdef\csname num@bestv:pd:plantdoc_5:PMP\endcsname{48.5}
\expandafter\gdef\csname num@worstv:pd:plantdoc_5:PMP\endcsname{45.5}
\expandafter\gdef\csname num@nscores:pd:plantdoc_5:PMP\endcsname{8}
\expandafter\gdef\csname num@partsd:pd:plantdoc_5:PMP\endcsname{1.2}
\expandafter\gdef\csname num@kw:pd:plantdoc_5:PMP\endcsname{0.52}
\expandafter\gdef\csname num@kwp:pd:plantdoc_5:PMP\endcsname{0.01}
\expandafter\gdef\csname num@randrank:pd:plantdoc_5:PMP\endcsname{6}
\expandafter\gdef\csname num@abl:max\endcsname{1.3}
\expandafter\gdef\csname num@abl:maxname\endcsname{the Fisher ratio alone}
\expandafter\gdef\csname num@abl:maxsign\endcsname{below}
\expandafter\gdef\csname num@abl:count\endcsname{9}
\expandafter\gdef\csname num@abl:nsig\endcsname{1}
\expandafter\gdef\csname num@lam:val5:default\endcsname{88.5}
\expandafter\gdef\csname num@lam:val5:best\endcsname{88.9}
\expandafter\gdef\csname num@lam:val5:worst\endcsname{87.9}
\expandafter\gdef\csname num@lam:val5:count\endcsname{8}
\expandafter\gdef\csname num@epse:resisc45\endcsname{0.2}
\expandafter\gdef\csname num@epse:eurosat13\endcsname{0.1}
\expandafter\gdef\csname num@epse:eurosat3\endcsname{0.1}
\expandafter\gdef\csname num@epse:plantvillage\endcsname{0.1}
\expandafter\gdef\csname num@gain:rs:novel_5:PMP\endcsname{1.2}
\expandafter\gdef\csname num@gainhi:rs:novel_5:PMP\endcsname{2.0}
\expandafter\gdef\csname num@gainname:rs:novel_5:PMP\endcsname{Taylor}
\expandafter\gdef\csname num@gain:es:novel_5:PMP\endcsname{-0.4}
\expandafter\gdef\csname num@gainhi:es:novel_5:PMP\endcsname{2.7}
\expandafter\gdef\csname num@gainname:es:novel_5:PMP\endcsname{DACIS}
\expandafter\gdef\csname num@gain:pv:novel_5:PMP\endcsname{0.3}
\expandafter\gdef\csname num@gainhi:pv:novel_5:PMP\endcsname{2.2}
\expandafter\gdef\csname num@gainname:pv:novel_5:PMP\endcsname{Taylor}
\expandafter\gdef\csname num@gain:pd:plantdoc_5:PMP\endcsname{1.9}
\expandafter\gdef\csname num@gainhi:pd:plantdoc_5:PMP\endcsname{3.0}
\expandafter\gdef\csname num@gainname:pd:plantdoc_5:PMP\endcsname{Taylor}
\expandafter\gdef\csname num@cmp:total:PMP\endcsname{28}
\expandafter\gdef\csname num@cmp:above:PMP\endcsname{2}
\expandafter\gdef\csname num@cmp:below:PMP\endcsname{3}
\expandafter\gdef\csname num@cmp:abovelist:PMP\endcsname{Taylor on RESISC45; Taylor on PlantDoc}
\expandafter\gdef\csname num@cmp:belowlist:PMP\endcsname{BN scale on RESISC45; Fisher on RESISC45; BN scale on EuroSAT-MS}
\expandafter\gdef\csname num@cmp:chance:PMP\endcsname{1.4}
\expandafter\gdef\csname num@cmp:maxgain:PMP\endcsname{1.9}
\expandafter\gdef\csname num@cmp:bonf:PMP\endcsname{0.01}
\expandafter\gdef\csname num@cmp:bonfabove:PMP\endcsname{0.26}
\expandafter\gdef\csname num@cmp:halfw:PMP\endcsname{1.7}
\expandafter\gdef\csname num@cmp:halfwmax:PMP\endcsname{3.2}
\expandafter\gdef\csname num@gain:rs:novel_5:P\endcsname{2.5}
\expandafter\gdef\csname num@gainhi:rs:novel_5:P\endcsname{3.8}
\expandafter\gdef\csname num@gainname:rs:novel_5:P\endcsname{Taylor}
\expandafter\gdef\csname num@gain:es:novel_5:P\endcsname{2.5}
\expandafter\gdef\csname num@gainhi:es:novel_5:P\endcsname{5.9}
\expandafter\gdef\csname num@gainname:es:novel_5:P\endcsname{Taylor}
\expandafter\gdef\csname num@gain:pv:novel_5:P\endcsname{2.7}
\expandafter\gdef\csname num@gainhi:pv:novel_5:P\endcsname{4.9}
\expandafter\gdef\csname num@gainname:pv:novel_5:P\endcsname{Taylor}
\expandafter\gdef\csname num@gain:pd:plantdoc_5:P\endcsname{2.1}
\expandafter\gdef\csname num@gainhi:pd:plantdoc_5:P\endcsname{3.1}
\expandafter\gdef\csname num@gainname:pd:plantdoc_5:P\endcsname{filter L1}
\expandafter\gdef\csname num@cmp:total:P\endcsname{22}
\expandafter\gdef\csname num@cmp:above:P\endcsname{3}
\expandafter\gdef\csname num@cmp:below:P\endcsname{1}
\expandafter\gdef\csname num@cmp:abovelist:P\endcsname{Taylor on RESISC45; Taylor on PlantVillage; filter L1 on PlantDoc}
\expandafter\gdef\csname num@cmp:belowlist:P\endcsname{BN scale on PlantDoc}
\expandafter\gdef\csname num@cmp:chance:P\endcsname{1.1}
\expandafter\gdef\csname num@cmp:maxgain:P\endcsname{2.7}
\expandafter\gdef\csname num@cmp:bonf:P\endcsname{0.05}
\expandafter\gdef\csname num@cmp:bonfabove:P\endcsname{0.11}
\expandafter\gdef\csname num@cmp:halfw:P\endcsname{2.0}
\expandafter\gdef\csname num@cmp:halfwmax:P\endcsname{4.4}
\expandafter\gdef\csname num@kw:max\endcsname{0.68}
\expandafter\gdef\csname num@kw:min\endcsname{0.22}
\expandafter\gdef\csname num@cost:variant\endcsname{resisc45}
\expandafter\gdef\csname num@cost:macs-red:30\endcsname{58}
\expandafter\gdef\csname num@cost:cpu-red:30\endcsname{39}
\expandafter\gdef\csname num@cost:gpu-red:30\endcsname{11}
\expandafter\gdef\csname num@cost:macs-red:50\endcsname{33}
\expandafter\gdef\csname num@cost:cpu-red:50\endcsname{21}
\expandafter\gdef\csname num@cost:gpu-red:50\endcsname{4}
\expandafter\gdef\csname num@cost:macs-red:20\endcsname{70}
\expandafter\gdef\csname num@cost:cpu-red:20\endcsname{47}
\expandafter\gdef\csname num@cost:gpu-red:20\endcsname{14}
\expandafter\gdef\csname num@cost:macs-red:10\endcsname{82}
\expandafter\gdef\csname num@cost:cpu-red:10\endcsname{57}
\expandafter\gdef\csname num@cost:gpu-red:10\endcsname{26}
\expandafter\gdef\csname num@cost:mobilenet_v3_small:macs\endcsname{0.018}
\expandafter\gdef\csname num@cost:mobilenet_v3_small:cpu\endcsname{2.3}
\expandafter\gdef\csname num@cost:mobilenet_v3_small:params\endcsname{0.93}
\expandafter\gdef\csname num@cost:mobilenet_v2:macs\endcsname{0.098}
\expandafter\gdef\csname num@cost:mobilenet_v2:cpu\endcsname{4.3}
\expandafter\gdef\csname num@cost:mobilenet_v2:params\endcsname{2.22}
\expandafter\gdef\csname num@cost:efficientnet_b0:macs\endcsname{0.126}
\expandafter\gdef\csname num@cost:efficientnet_b0:cpu\endcsname{6.3}
\expandafter\gdef\csname num@cost:efficientnet_b0:params\endcsname{4.01}
\expandafter\gdef\csname num@tune:best\endcsname{0.0001}
\expandafter\gdef\csname num@tune:S0:1\endcsname{86.7}
\expandafter\gdef\csname num@tune:S0:5\endcsname{96.1}
\expandafter\gdef\csname num@tune:0.0001:1\endcsname{85.4}
\expandafter\gdef\csname num@tune:0.0001:5\endcsname{95.3}
\expandafter\gdef\csname num@tune:0.00025:1\endcsname{84.9}
\expandafter\gdef\csname num@tune:0.00025:5\endcsname{95.0}
\expandafter\gdef\csname num@tune:0.0005:1\endcsname{84.3}
\expandafter\gdef\csname num@tune:0.0005:5\endcsname{94.7}
\expandafter\gdef\csname num@tune:0.001:1\endcsname{83.6}
\expandafter\gdef\csname num@tune:0.001:5\endcsname{94.5}
\expandafter\gdef\csname num@tune:0.002:1\endcsname{82.7}
\expandafter\gdef\csname num@tune:0.002:5\endcsname{94.2}
\expandafter\gdef\csname num@tune-rs:best\endcsname{0.0001}
\expandafter\gdef\csname num@tune-rs:S0:1\endcsname{74.2}
\expandafter\gdef\csname num@tune-rs:S0:5\endcsname{90.1}
\expandafter\gdef\csname num@tune-rs:0.0001:1\endcsname{75.8}
\expandafter\gdef\csname num@tune-rs:0.0001:5\endcsname{90.0}
\expandafter\gdef\csname num@tune-rs:0.00025:1\endcsname{76.1}
\expandafter\gdef\csname num@tune-rs:0.00025:5\endcsname{89.8}
\expandafter\gdef\csname num@tune-rs:0.0005:1\endcsname{75.9}
\expandafter\gdef\csname num@tune-rs:0.0005:5\endcsname{89.5}
\expandafter\gdef\csname num@tune-rs:0.001:1\endcsname{75.2}
\expandafter\gdef\csname num@tune-rs:0.001:5\endcsname{89.0}
\expandafter\gdef\csname num@tune-rs:0.002:1\endcsname{75.0}
\expandafter\gdef\csname num@tune-rs:0.002:5\endcsname{89.1}
\expandafter\gdef\csname num@hw-gpu\endcsname{NVIDIA RTX PRO 6000 Blackwell Server Edition}
\expandafter\gdef\csname num@hw-cpu\endcsname{a 20-thread server CPU (model not reported by the host)}
\expandafter\gdef\csname num@hw-gpumem\endcsname{95.0}
}{}

% \ci{tag}: the 95% interval of the paired difference d:<tag>, written [lo,hi] and breakable after the comma
\newcommand{\ci}[1]{[\num{d:#1:lo},\allowbreak\num{d:#1:hi}]}

% small helper for pm values in running text
\newcommand{\pmv}[2]{$#1\pm#2$}
\newcommand{\key}[1]{\textbf{#1}}
